\documentclass[11pt]{article}

\usepackage[utf8]{inputenc}
% microtype 의 폰트 확장은 scalable 폰트를 요구한다. 기본 CM 이 비트맵(Type3)으로
% 잡히는 배포판에서는 lmodern 없이는 "auto expansion is only possible with
% scalable fonts" 로 컴파일이 죽는다.
\usepackage{lmodern}
\usepackage[T1]{fontenc}
\usepackage[margin=1in]{geometry}
\usepackage{amsmath,amssymb}
\usepackage{graphicx}
\usepackage{longtable,booktabs}
\usepackage{microtype}
% svgnames 를 줘야 teal 이 정의된다. 없으면 \cite 하나마다
% "Undefined color `teal'" 이 뜬다 (본편 산출물에서 732회).
\usepackage[svgnames]{xcolor}
% hyperref goes last; it makes \cite clickable through to the bibliography.
\usepackage[colorlinks=true,linkcolor=blue,citecolor=teal,urlcolor=magenta]{hyperref}

% pandoc emits \tightlist but expects the preamble to define it.
\providecommand{\tightlist}{%
  \setlength{\itemsep}{0pt}\setlength{\parskip}{0pt}}

% If the compile times out on Overleaf, comment out \tableofcontents below first.
\title{Retrieval-Augmented Generation for Large Language Models: A Comprehensive Survey of Paradigms, Techniques, and Future Directions}
\author{Generated by AutoSurvey}
\date{}

\begin{document}
\maketitle
\tableofcontents
\newpage

\section{Introduction and Foundations of Retrieval-Augmented Generation}

\subsection{Background: Limitations of Large Language Models and the Emergence of Retrieval-Augmented Generation}

Large Language Models (LLMs) have demonstrated remarkable capabilities across a wide spectrum of natural language processing tasks, from dialogue generation and complex reasoning to knowledge retention and program synthesis \cite{ref1,ref2}. Their proficiency stems from parametric memory---knowledge encoded directly into model weights during large-scale pretraining, enabling them to store vast amounts of factual information and exhibit emergent abilities \cite{ref3,ref4,ref5}. However, despite these impressive capabilities, LLMs face several fundamental limitations that constrain their reliability and applicability in real-world settings. These limitations include hallucination, reliance on static and outdated training data, opacity of reasoning processes, and deficiencies in knowledge-intensive tasks \cite{ref1,ref6}.

First, hallucination remains one of the most critical obstacles to the reliable deployment of LLMs. Parametric knowledge alone is insufficient for producing consistently factually correct outputs, and LLMs frequently generate plausible-sounding but unfaithful or factually incorrect content, particularly when confronted with long-tail, domain-specific, or rarely encountered queries \cite{ref7,ref8,ref9}. Even when models are equipped with correct and sufficient retrieved context, they can still produce hallucinations due to imbalances between how they utilize external context and their internal parametric knowledge \cite{ref10,ref11}.

Second, LLMs rely on static training corpora, making them inherently limited in addressing rapidly evolving information, real-time knowledge updates, and domain-specific queries that fall outside their training distribution \cite{ref12,ref8,ref13}. This static nature leads to temporal degradation, where the model's knowledge becomes increasingly outdated, and infrequent knowledge that appears rarely in training data remains poorly captured \cite{ref14}. Parametric memory pools and model editing approaches offer only partial solutions, as they are costly to update and do not scale effectively \cite{ref15}.

Third, the opacity of reasoning processes poses significant challenges for interpretability and trustworthiness. LLMs lack inherent mechanisms for tracing their outputs back to verifiable sources, making it difficult to understand, explain, or audit their decisions \cite{ref16}. The uninterpretable nature of parametric memorization complicates efforts to identify and prevent hallucinations, as there is no transparent mechanism to determine which internal knowledge was invoked during generation \cite{ref14}.

Fourth, LLMs exhibit deficiencies in knowledge-intensive tasks that require precise factual retrieval, multi-hop reasoning, and integration of specialized domain expertise. While scaling model parameters improves general capabilities, it does not eliminate errors on knowledge-intensive benchmarks, and performance disparities persist when compared to task-specific architectures \cite{ref17,ref18}.

To address these limitations, Retrieval-Augmented Generation (RAG) has emerged as a promising paradigm that synergistically merges the parametric memory of LLMs with non-parametric external knowledge bases \cite{ref1,ref19}. Formally, RAG is a framework that integrates a retriever module, which accesses external corpora or databases, with a generator LLM, allowing the model to ground its responses in up-to-date, verifiable, and contextually relevant knowledge \cite{ref12}. The foundational architecture of RAG consists of three interconnected components: (1) the retriever, which identifies and retrieves the most relevant documents or passages from external knowledge sources; (2) the generator, which produces the final output conditioned on both the query and the retrieved evidence; and (3) the augmentation mechanism, which defines how retrieved information is integrated into the generation process---whether through contextual concatenation, attention modulation, or parameter-level injection \cite{ref8,ref5}.

RAG systems operate by leveraging external knowledge to enhance the accuracy and credibility of generation, particularly for knowledge-intensive tasks, while enabling continuous knowledge updates and seamless integration of domain-specific information \cite{ref6}. By shifting knowledge integration from context-level augmentation to memory-level interaction, RAG frameworks can consistently outperform purely parametric approaches, with improved stability and accuracy in long-context and multi-hop reasoning scenarios \cite{ref15}. The overarching objectives driving the RAG field include: (a) improving factual accuracy and reducing hallucination through grounded generation; (b) enabling dynamic and up-to-date knowledge access without retraining; (c) enhancing interpretability by providing attributable evidence for model outputs; (d) supporting domain-specific adaptation through customizable external knowledge bases; and (e) improving robustness to noisy or conflicting information \cite{ref1,ref20}.

In summary, RAG represents a principled approach to augmenting LLMs with non-parametric external knowledge, addressing fundamental limitations of static parametric memory while preserving the generative strengths of large-scale pretrained models. The paradigm has evolved from simple retrieve-then-read pipelines to sophisticated frameworks encompassing dynamic retrieval, parametric knowledge injection, and agentic reasoning, establishing itself as a foundational technique for reliable and knowledge-grounded natural language generation \cite{ref19,ref21}.

\subsection{The Evolution of RAG Paradigms: From Naive to Agentic}

The evolution of Retrieval-Augmented Generation (RAG) is a story of progressive sophistication, moving from simple, static pipelines to complex, autonomous systems capable of dynamic reasoning and decision-making. This trajectory is not merely a sequence of incremental improvements but a fundamental shift in how we conceptualize the integration of external knowledge with large language models (LLMs). The journey begins with the foundational ``Naive RAG'' paradigm, progresses through the optimized ``Advanced RAG,'' is reshaped by the composable ``Modular RAG'' framework, and culminates in the emergence of ``Agentic RAG,'' where autonomous agents orchestrate the entire process \cite{ref22}. This foundational evolution directly informs the modern system landscape introduced in the previous subsection, as each evolutionary stage has given rise to distinct architectural choices that the current taxonomy---spanning Graph-based, Multimodal, Multilingual, Reasoning-Enhanced, and Domain-Specific RAG---seeks to organize.

\textbf{Naive RAG: The Foundational Retrieve-Read Paradigm}

The initial paradigm, often referred to as Naive RAG, established the core tripartite architecture that underpins all subsequent developments: indexing, retrieval, and generation \cite{ref23}. In this basic ``retrieve-then-generate'' process, a user query is first used to retrieve a set of relevant documents from a pre-indexed knowledge base. These documents are then concatenated with the original query and fed into an LLM to generate a grounded response. While this approach marked a significant advancement over the limitations of LLMs with static training data, it quickly revealed critical shortcomings. The system is rigid and static, with a one-pass retrieval that cannot adapt to the specific needs of complex queries. It is susceptible to retrieving irrelevant or noisy information, which can lead to hallucinations and degraded performance. Furthermore, the single-shot nature of the process fails on tasks that require multi-step reasoning or iterative information gathering. This foundational model, while elegant in its simplicity, is fundamentally limited in its ability to handle the complexity and dynamism of real-world information-seeking tasks \cite{ref24,ref23}. Indeed, the flat document stores characteristic of this paradigm are precisely what Graph-based RAG later emerged to overcome, as the absence of relational structure limits the system's ability to answer multi-hop queries.

\textbf{Advanced RAG: Enhancing the Pipeline}

To address these limitations, the field evolved to ``Advanced RAG,'' which focuses on optimizing each stage of the naive pipeline. These enhancements can be broadly categorized as pre-retrieval, post-retrieval, and generation optimizations.

\begin{itemize}
\item
  \textbf{Pre-retrieval optimizations} aim to improve the quality of the input before a search is conducted. This often involves query transformation techniques such as rewriting, expansion, and decomposition to better align the user's question with the format of the indexed documents. By clarifying the query and expanding it with related terms, the retrieval process becomes more likely to surface relevant evidence \cite{ref23}.
\item
  \textbf{Post-retrieval optimizations} target the results returned by the retriever. A primary technique here is reranking, which uses a more sophisticated model to re-order the initially retrieved documents and push the most relevant ones to the top. Another crucial strategy is context compression, which filters out redundant or irrelevant passages from the retrieved documents before they are passed to the LLM, thereby reducing noise and improving the signal-to-noise ratio \cite{ref23}.
\item
  \textbf{Generation optimizations} focus on how the LLM uses the retrieved context. This can involve techniques to align the model's attention to key evidence or to guide it through a chain-of-thought reasoning process that explicitly references the retrieved information \cite{ref21}. These latter techniques foreshadow the Reasoning-Enhanced RAG systems described in the taxonomy, where reasoning is not a post-hoc addition but an integrated component of the pipeline.
\end{itemize}

While Advanced RAG yielded substantial improvements in accuracy and reliability, these systems are still fundamentally linear and rigid. Each component in the pipeline, from retriever to reranker to generator, performs a fixed function in a fixed order. This lack of flexibility makes them difficult to adapt to new tasks or data modalities without significant manual re-engineering \cite{ref23}. This rigidity also limits the seamless integration of heterogeneous data types, a challenge that Multimodal RAG would later confront by requiring more flexible, cross-modal retrieval and fusion mechanisms.

\textbf{Modular RAG: From Rigid Pipelines to Composable Frameworks}

The next major evolution was the introduction of the ``Modular RAG'' framework, a conceptual shift inspired by LEGO-like building blocks \cite{ref23}. This paradigm moves beyond the traditional linear architecture by decomposing complex RAG systems into independent, reusable modules and specialized operators. These modules, such as search, verification, and refinement, can be arranged and recombined in various configurations---linear, conditional, branching, or looping---to create custom workflows tailored to specific tasks. This composability allows for a high degree of flexibility and reconfigurability. For instance, a system can be designed with a routing mechanism that directs simpler queries to a fast, direct path while sending complex, multi-hop queries to a more elaborate pipeline involving iterative search and reasoning. This modular design philosophy transcends the constraints of the ``retrieve-then-generate'' process, embracing a more advanced design that integrates routing, scheduling, and fusion mechanisms. It provides a more unified and practical framework for conceptualizing and deploying RAG technologies, as new functionalities can be added by plugging in new modules rather than redesigning the entire system \cite{ref23}. This modularity is a crucial enabler for the five system variants outlined in the taxonomy: Graph-based RAG can be realized as a specialized retrieval module, Multimodal RAG as a fusion module, and Domain-Specific RAG as a configuration of task-specific modules, all within the same composable architecture.

\textbf{Agentic RAG: The Shift to Autonomous Orchestration}

The culmination of this evolution is ``Agentic RAG,'' a paradigm that represents a fundamental change in the role of the LLM within the RAG workflow. Unlike its predecessors, Agentic RAG transcends static workflows by embedding autonomous AI agents into the RAG pipeline \cite{ref22,ref25}. Here, the LLM acts not just as a generator but as a central agent that orchestrates the entire process. It is empowered with agentic design patterns---including reflection, planning, tool use, and multi-agent collaboration---to dynamically manage retrieval strategies, iteratively refine contextual understanding, and adapt its workflow on the fly \cite{ref22,ref26}. For example, the agent can decide whether a query necessitates a retrieval step at all, choose which specific retrieval tool to use (e.g., keyword search, semantic search, or graph traversal), and judge whether the information retrieved in a first pass is sufficient to form an answer or if further searches are required \cite{ref27,ref28}. This decomposes the retrieval process from a single black-box operation into a series of controllable, agent-driven decisions, enabling the system to handle complex, multi-step reasoning tasks that are impossible for earlier paradigms \cite{ref29}. The system can reason about its own information needs, plan a series of actions, and reflect on the results to correct course, leading to higher-quality and more robust answers \cite{ref26,ref21}. This agentic paradigm directly embodies the Reasoning-Enhanced RAG category in the taxonomy, particularly the Agentic RAG sub-type, and provides the adaptive orchestration layer through which Graph-based, Multimodal, Multilingual, and Domain-Specific modules can be intelligently combined and deployed.

This agentic approach can be characterized by the autonomy and cardinality of the agents involved, ranging from a single agent that handles all tasks to complex multi-agent systems where specialized agents collaborate on planning, retrieval, and synthesis \cite{ref22}. Yet, this increased power comes with new challenges, including the need for sophisticated optimization and the potential for compounding errors and unexplainable behavior \cite{ref25,ref30,ref31}. Furthermore, the strategic-operational loop introduces complexities where decisions made at a higher planning level must be faithfully executed by lower-level components, and failures in this chain can lead to suboptimal performance \cite{ref30}. In essence, the evolution from Naive to Agentic RAG is a progression from static, linear tools to dynamic, intelligent systems that can plan, act, and learn, representing a significant leap toward more capable and reliable AI. This evolutionary trajectory provides the essential background for the detailed exploration of the five system variants that follows, as each variant represents a particular instantiation and specialization of the broader RAG design space that this evolution has made possible.

\subsection{A Unified Taxonomy and Scope of Modern RAG Systems}

The rapid proliferation of RAG methodologies has given rise to a sprawling and often fragmented design space, making it increasingly difficult for researchers and practitioners to navigate the landscape, compare approaches, and identify appropriate solutions for specific problems. Early efforts to structure this space, such as the foundational taxonomy of Naive, Advanced, and Modular RAG \cite{ref1}, provided a crucial starting point by framing the evolution from simple retrieve-then-read pipelines to more flexible, component-based systems. However, as RAG has matured, its scope has broadened dramatically, incorporating sophisticated architectural paradigms and specialized application contexts that defy simple categorization. Current research reflects this expansion, with efforts to conceptualize RAG applications through multi-dimensional taxonomies that capture the defining characteristics of the technology across different domains \cite{ref32}. This survey builds upon these foundations to present a unified taxonomy that organizes the modern RAG design space into five principal, yet interconnected, system variants: Graph-based RAG, Multimodal RAG, Multilingual RAG, Reasoning-Enhanced RAG, and Domain-Specific RAG. This structure is designed not to be a rigid classification but rather a lens through which to understand the key innovations, challenges, and synergies that define the current state of the art.

\textbf{Graph-based RAG} represents one of the most significant paradigmatic shifts in the field. Traditional RAG systems that operate over flat document stores often struggle with complex queries requiring multi-hop reasoning and a holistic understanding of interconnected information \cite{ref33,ref28,ref34}. In contrast, GraphRAG leverages the intrinsic relational structure of knowledge, encoding entities and their connections to enable more precise, context-aware, and explainable retrieval \cite{ref35,ref36,ref37}. The core innovation lies in the \emph{representation} of knowledge as a graph, which can be constructed through various techniques, including token co-occurrence, LLM-extracted triples, or entity-centric pipelines \cite{ref37}. This structured representation supports advanced retrieval strategies that go beyond simple vector similarity, such as graph traversal, Personalized PageRank, and spreading activation, which are critical for answering queries that demand an understanding of relationships and hierarchies \cite{ref37,ref38,ref39}. Systems like LightRAG and SubgraphRAG exemplify different approaches within this space, with the former integrating graph structures to capture complex inter-dependencies \cite{ref40} and the latter focusing on efficient subgraph retrieval to balance effectiveness and efficiency \cite{ref41}. Furthermore, the field is advancing toward more dynamic and self-evolving graph structures, as demonstrated by frameworks like EvoGraph-R1, which reconceptualize knowledge graphs as interactive environments that can be refined during the reasoning process \cite{ref42}. This inclusion of graph-based approaches is driven by the recognition that structured knowledge is pervasive and that leveraging its topology is essential for achieving sophisticated reasoning capabilities \cite{ref43}.

\textbf{Multimodal RAG} (MRAG) extends the RAG paradigm beyond the textual domain to address the reality that modern knowledge repositories are inherently multimodal \cite{ref44}. These systems aim to retrieve and reason over information from diverse modalities, including text, images, tables, audio, and video, to provide more comprehensive and grounded answers \cite{ref45}. The fundamental challenge in this variant is the cross-modal alignment and fusion of disparate information types \cite{ref45}. Advanced frameworks address this by constructing unified multimodal knowledge graphs, where textual entities and visual regions are fused into coherent semantic nodes \cite{ref46,ref47}. This allows systems to capture cross-modal relationships and support complex reasoning that spans different data formats, moving beyond document-level retrieval to more granular, fragment-level or region-level retrieval \cite{ref48}. The practical importance of this variant is underscored by the proliferation of multimodal benchmarks and specialized frameworks that seek to evaluate and improve performance on tasks like visual question answering, multimodal fact verification, and document-centric QA \cite{ref49,ref50}. The integration of graphs with multimodal data, as seen in systems like MegaRAG, is a particularly promising direction for enabling deep reasoning over complex, long-form documents \cite{ref51}.

\textbf{Multilingual RAG} (mRAG) tackles the challenges of knowledge retrieval and generation across different languages. A key issue in this domain is the presence of language preference biases, where models may favor information in a particular language, potentially at the expense of accuracy \cite{ref52}. Strategies to mitigate these challenges are varied and include translating the user question into the language of the knowledge base (tRAG), retrieving evidence in a multilingual space (MultiRAG), or translating retrieved documents into the query language (CrossRAG). More sophisticated approaches involve quality-aware translation tagging, as in QTT-RAG, to ensure only reliable translated content is used, and knowledge fusion frameworks like DKM-RAG that integrate information from multiple languages to generate a comprehensive answer. The evaluation of these systems requires large-scale, dedicated benchmarks that can effectively measure cross-lingual retrieval and generation capabilities, such as M4-RAG and XRAG. The development of robust mRAG systems is critical for ensuring that the benefits of retrieval augmentation are accessible across the globe and are not limited to high-resource languages.

\textbf{Reasoning-Enhanced RAG} represents a class of systems designed to synergize retrieval with advanced reasoning capabilities to tackle complex, multi-step problems. While basic RAG improves factuality by injecting external knowledge, it often falls short on tasks that demand multi-hop inference \cite{ref21}. Conversely, purely reasoning-oriented approaches may hallucinate or mis-ground facts \cite{ref21}. Reasoning-Enhanced RAG systems address this by integrating techniques like Chain-of-Thought (CoT), Chain-of-Noting, and structured planning directly into the RAG pipeline \cite{ref53,ref21}. This synergy is bidirectional: reasoning can optimize each stage of the RAG pipeline, such as improving query decomposition for retrieval, while retrieved knowledge can supply missing premises to support complex inference \cite{ref53}. This category also encompasses Agentic RAG, where autonomous agents are deployed to dynamically plan retrieval strategies, act on the results, and iteratively refine their approach based on self-reflection and tool use \cite{ref28,ref21}. Systems like EvoGraph-R1, which formulates retrieval as a Markov Decision Process, epitomize this trend toward dynamic, adaptive, and reason-driven retrieval \cite{ref42}, while frameworks like MIXRAG employ Mixture-of-Experts to dynamically select specialized retrievers based on query intent, exemplifying a move toward more intelligent and adaptive retrieval strategies \cite{ref54}.

Finally, \textbf{Domain-Specific RAG} encompasses systems that are tailored to meet the unique requirements and challenges of specialized fields. These systems often leverage a combination of the aforementioned variants to address the specific knowledge representation, retrieval, and reasoning needs of a particular domain. For instance, in the medical domain, systems like BioMol-MQA and the graph-based DSRAG integrate multimodal knowledge graphs to reason over complex interactions between molecules, clinical notes, and knowledge graphs \cite{ref55,ref56}. Similarly, in e-commerce, comparative analyses and frameworks like DA-RAG highlight the need for retrieval pipelines that can handle large-scale, structured knowledge bases \cite{ref57,ref39}. Even within general-purpose systems, the growing need to adapt to a diverse set of queries and data types is driving the development of hybrid frameworks that intelligently route queries to appropriate retrieval mechanisms \cite{ref58}, and the pursuit of ``best practices'' is a central theme for industrial deployments, as seen in the medical domain \cite{ref59}.

In summary, this taxonomy illustrates the dynamic and evolving nature of RAG. These five variants are not mutually exclusive; rather, they are often combined to build the sophisticated systems required for complex real-world applications. The following sections of this survey will delve into the core pipeline components, advanced architectures, and specific application contexts that define each of these categories, providing a comprehensive map of the modern RAG landscape.

\subsection{Key Challenges, Scope, and Contributions of this Survey}

The preceding taxonomy delineates the primary architectural paradigms that define the modern RAG design space, highlighting the diverse mechanisms through which retrieval can be augmented by graphs, modalities, languages, reasoning, and domain-specific constraints. However, a comprehensive understanding of the field requires not only a map of its solutions but also a clear articulation of the persistent problems that these solutions seek to address. Despite the significant advances in Retrieval-Augmented Generation (RAG) architectures and their demonstrated efficacy across the diverse applications described above, the field remains constrained by a set of persistent and interconnected challenges that impede the deployment of fully reliable, scalable, and secure systems. These challenges are not merely peripheral engineering concerns; they are fundamental to the effectiveness of the entire RAG paradigm, regardless of the architectural variant employed. This survey is motivated by the need to systematically articulate these challenges, delineate the current scope of research, and provide a structured foundation for future inquiry. In this subsection, we first identify the primary open challenges that define the frontier of RAG research, then clarify the scope and positioning of this survey relative to existing reviews, and finally enumerate the specific contributions this work makes to researchers and practitioners.

\textbf{Key Challenges Constraining RAG Effectiveness.} The first and perhaps most fundamental challenge, which cuts across all the variants discussed previously, is \textbf{retrieval precision}. The efficacy of a RAG system is inextricably linked to the quality of the documents retrieved; irrelevant or semantically misaligned documents directly compromise the accuracy of the final generated response. This problem is particularly acute in specialized domains where general-purpose dense retrievers struggle with nuanced language, while high-accuracy in-domain models are often computationally prohibitive \cite{ref60}. Furthermore, the retrieval stage is susceptible to a phenomenon termed Document-Level Retrieval Mismatch, where the retriever selects information from entirely incorrect source documents, a failure mode that is especially prevalent in large databases of structurally similar documents such as those found in the legal domain \cite{ref61}. Even when relevant documents are retrieved, the inclusion of noisy or irrelevant passages can degrade performance, and the system's ability to filter this noise is a measure of its robustness. The challenge is compounded by the finding that RAG deployment must be highly selective, with variable recall thresholds and failure modes affecting up to 12.6\% of samples even when perfect documents are provided \cite{ref62}.

A second major challenge is \textbf{scalability and efficiency}, which becomes increasingly pressing as the graph-based, multimodal, and domain-specific systems described in the previous section grow in complexity and data volume. As knowledge bases grow, the computational cost of indexing, retrieving, and processing large volumes of context increases substantially. The multi-stage pipeline of retrieval and generation introduces unique workload characteristics and performance bottlenecks, including significant temporal locality and highly skewed hot document access patterns that are not captured by generic LLM inference traces \cite{ref63}. The need to balance latency, throughput, and memory usage against accuracy is a persistent operational concern. Moreover, the requirement to retrieve and process multiple documents to answer complex queries amplifies these costs, demanding sophisticated caching, pruning, and acceleration techniques.

Third, \textbf{noise management and robustness against imperfect retrieval} remain critical hurdles, particularly relevant to the Reasoning-Enhanced RAG systems that aim to perform multi-hop inference over potentially noisy evidence. Real-world retrieval environments are inherently noisy, and recent research has shown that noise is not always detrimental; a taxonomy of seven distinct noise types reveals that some forms of noise can be beneficial to LLMs, enhancing certain aspects of model capability, while others are harmful \cite{ref64}. This nuanced understanding complicates the design of universally effective noise-filtering mechanisms. Furthermore, Retrieval-Augmented Generation systems exhibit a surprising fragility to linguistic variations in user queries. Systematic analyses of how variations in formality, readability, politeness, and grammatical correctness impact performance reveal that linguistic reformulations can lead to relative performance drops of up to 40\% in retrieval and generation stages, highlighting a critical vulnerability to error propagation in real-world user interactions \cite{ref65}. Similarly, adaptive RAG methods, which dynamically trigger retrieval, show a critical robustness gap where small surface-level changes in semantically identical queries dramatically alter retrieval behavior and accuracy \cite{ref66}.

Fourth, \textbf{security and privacy vulnerabilities} constitute a growing and severe threat to RAG deployment, a concern that is amplified by the increasing integration of external knowledge sources in the multimodal and domain-specific variants. The integration of external knowledge bases introduces a new attack surface distinct from traditional LLM risks. These threats range from data poisoning and indirect prompt injection to membership inference and knowledge base extraction \cite{ref67,ref68}. Research has demonstrated that the safety guardrails of LLMs themselves can be exploited as a denial-of-service attack vector by injecting minimalistic jailbreak texts into the knowledge base \cite{ref69}. Moreover, the RAG pipeline is vulnerable to end-to-end indirect prompt manipulations, where attacks optimized to boost the ranking of malicious documents can achieve a consistent success rate of around 40\%, which can rise to 60\% when ambiguous answers are considered, with limited configuration changes able to thwart these attacks \cite{ref70}. Privacy attacks, such as membership inference against the external database, can successfully breach the confidentiality of sensitive retrieval data \cite{ref71}. These security challenges are compounded by fairness concerns, where RAG can undermine fairness alignment without fine-tuning, even when external datasets are fully censored and supposedly unbiased \cite{ref72}.

Fifth, \textbf{the dynamic conflict between parametric and non-parametric knowledge} remains a fundamental technical challenge that underpins the efficacy of all RAG variants. The model must effectively reconcile its internal, parametric memory with newly retrieved, external evidence. Conflicts between these sources can lead to inaccurate or hallucinated outputs, and the system's ability to resolve these conflicts in a credibility-aware manner is crucial for ensuring factual fidelity \cite{ref73}. This challenge is intertwined with the broader goal of ensuring faithfulness, attribution, and reduced hallucination in generated responses, and it becomes particularly acute when integrating knowledge from heterogeneous sources, such as the structured graphs and multimodal data discussed earlier.

\textbf{Scope and Distinctions from Related Reviews.} Given the rapidly expanding body of literature on RAG, both algorithmic innovations and system-level analyses have proliferated. While several recent surveys have focused on systematic reviews and comparisons of new state-of-the-art techniques against their predecessors, there remains a gap in extensive experimental comparisons that dissect the interplay between various components \cite{ref74}. Existing surveys often prioritize taxonomy and algorithmic novelty, sometimes at the expense of a detailed, component-specific assessment of engineering trade-offs and security implications. Our survey distinguishes itself by providing a holistic, end-to-end perspective that integrates three critical lenses: (1) the \emph{architectural and algorithmic evolution} from Naive to Agentic RAG, as charted in the preceding taxonomy, (2) the \emph{engineering and operational challenges} of deployment, including efficiency, scalability, and robustness, and (3) the \emph{security, privacy, and trustworthiness} dimensions that are essential for real-world adoption. While existing work has addressed some of these aspects in isolation, such as dedicated security surveys \cite{ref75} or robustness benchmarks \cite{ref76}, we aim to synthesize these threads into a unified framework. Furthermore, we move beyond accuracy-centric evaluations to encompass a multi-dimensional assessment of trustworthiness, including factuality, robustness, fairness, transparency, accountability, and privacy, guided by frameworks like the Trust-RAG Compass \cite{ref77}. This survey also addresses the critical gap in systematic, human-centered evaluation of RAG outputs, which remains underexplored despite the increasing deployment of RAG in user-facing applications \cite{ref78}.

\textbf{Intended Contributions.} Building upon the taxonomy presented in the previous sections, this survey makes several targeted contributions to both researchers and practitioners. First, it provides a \textbf{comprehensive and structured taxonomy} of RAG systems, tracing the paradigm's evolution and categorizing the design space along axes of architecture, augmentation strategy, and training paradigm. This taxonomy aims to clarify the relationships between disparate approaches and facilitate a more principled comparison of methods. Second, it offers an \textbf{in-depth synthesis of core pipeline components}, from indexing and retrieval to generation enhancement and knowledge integration, while critically analyzing the trade-offs between precision, recall, scalability, and computational cost. Third, it delivers a \textbf{systematic review of security and privacy threats}, formalizing the attack surface and providing a taxonomy of defense mechanisms, thereby serving as a foundational reference for building trustworthy systems \cite{ref79}. Fourth, we consolidate and compare the landscape of \textbf{evaluation methodologies and benchmarks}, providing guidance on component-level versus end-to-end assessment and highlighting the shift toward interpretable, stage-decomposed evaluation. Fifth, we synthesize \textbf{practical lessons and engineering methodologies} from real-world deployments and case studies, offering evidence-based guidance for designing, building, and maintaining robust RAG systems in production environments. Finally, we articulate a \textbf{research roadmap}, identifying persistent open challenges and promising future directions that span next-generation adaptive retrieval, the integration of causal structures, and the development of human-centric, trustworthy RAG systems. By bridging the gap between algorithmic innovation, system engineering, and security considerations, this survey aims to serve as a definitive resource for the next generation of retrieval-augmented language modeling research, building directly upon the architectural foundations established in the preceding discussion.

\section{Core Pipeline Components and Retrieval Strategies}

\subsection{Indexing and Knowledge Representation Strategies}

The indexing stage constitutes the foundational layer of any Retrieval-Augmented Generation (RAG) system, determining the granularity, structure, and accessibility of the external knowledge that will later be retrieved and synthesized. The design decisions made during this phase---ranging from how documents are segmented and embedded to how the resulting representations are organized into searchable structures---exert a profound influence on downstream retrieval precision, generation faithfulness, and overall system efficiency. This subsection systematically examines the core techniques and strategic considerations that define modern indexing pipelines, focusing on chunking strategies, embedding methods, and advanced index architectures. Crucially, these choices are inseparable from the retrieval paradigms discussed in the preceding subsection: the effectiveness of sparse, dense, late-interaction, graph-based, and hybrid retrieval all presuppose an index that has been optimally segmented and encoded. At the same time, the indexing decisions directly shape the candidate space that the retrieval stage must navigate, as well as the granularity of evidence available to the generator downstream.

\textbf{Chunking Strategies: Segmentation as a First-Class Decision.} The segmentation of source documents into retrievable units, or chunks, is arguably the most consequential preprocessing step, as it defines the atomic unit of evidence that retrieval algorithms will rank and that generators will consume. A systematic large-scale evaluation comparing 36 segmentation methods across six knowledge domains found that content-aware chunking significantly outperforms naive fixed-length splitting, with the top-performing Paragraph Group Chunking achieving a mean nDCG@5 of 0.459, while simple fixed-size character chunking lagged substantially below 0.244 \cite{ref80}. This performance gap underscores that the choice of chunking strategy is not a trivial implementation detail but a critical performance lever. Within the design space, strategies vary from fixed-size token windows to structure-aware methods that respect document semantics. Unsurprisingly, the optimal strategy is task-dependent: simple structure-based methods such as sentence or paragraph segmentation excel in standard in-corpus retrieval, while LLM-guided methods like LumberChunker demonstrate superiority in in-document needle-in-a-haystack settings \cite{ref81}. This finding challenges the assumption of a universally superior chunking method, suggesting instead that the intended retrieval objective must dictate the segmentation approach.

The granularity of chunking introduces an inherent precision-recall trade-off that cannot be resolved by tuning chunk size alone. Fixed-size chunking with smaller chunks (64--128 tokens) proves optimal for datasets with concise, fact-based answers, whereas larger chunks (512--1024 tokens) improve performance on tasks requiring broader contextual understanding \cite{ref82}. This sensitivity is further compounded by the interaction between chunk size and the embedding model employed; different embedding backbones exhibit distinct chunking sensitivities, with some models leveraging larger contexts for long-range retrieval while others excel with smaller, entity-focused segments \cite{ref82}. Moreover, empirical studies indicate that overlap, a common strategy to mitigate boundary fragmentation, provides no measurable benefit to retrieval quality while increasing indexing cost \cite{ref83}. Instead, adaptive and query-aware approaches such as Query-Adaptive Semantic Chunking (QASC), which dynamically constructs chunks by integrating user queries into the segmentation process, have demonstrated F1-score improvements of 18--27\% over fixed chunking by enabling more relevant and coherent retrieval \cite{ref84}.

\textbf{Structure-Aware and Hierarchical Segmentation.} For domain-specific corpora, preserving document structure during segmentation proves critical---particularly because several retrieval paradigms discussed earlier, such as graph-based traversal and hybrid fusion, rely on well-defined structural units to operate effectively. In the legal domain, chunking strategies aligned with the inherent structure---particularly section and subsection-based retrieval---achieve the highest recall, with more complex approaches that override this structure performing worse and at greater computational cost \cite{ref85}. Similarly, structure-aware chunking for tabular data, which operates on row-level units and constructs a hierarchical Row Tree, reduces chunk count by up to 56\% compared to key-value baselines while improving MRR from 0.3576 to 0.5945, highlighting the importance of maintaining structural integrity in semi-structured data \cite{ref86}. In the financial domain, element-based chunking guided by document understanding models improves retrieval accuracy over paragraph-level approaches \cite{ref87}. Hierarchical frameworks extend this concept by organizing chunks into progressively larger semantic units. SproutRAG, for instance, constructs binary chunking trees via learned inter-sentence attention, enabling multi-granularity retrieval without additional LLM calls or lossy summarization, achieving a 6.1\% average improvement in information efficiency across benchmarks \cite{ref88}. The structure-preserving rationale extends to content-addressable representations, where approaches such as Self-Conditioned Positional HNSW modify embeddings to account for document positional priors, mitigating the failure mode where top-k retrieval returns near-adjacent chunks that repeat evidence and waste prompt budget \cite{ref89}. These structural considerations are particularly salient for the hybrid retrieval architectures described in the following subsection, where heterogeneous evidence units must be combined coherently.

\textbf{Metadata Enrichment and Cross-Document Synthesis.} Modern indexing increasingly leverages LLM-generated metadata to enhance retrieval, directly influencing the effectiveness of both dense and hybrid retrieval paradigms. Systematic studies demonstrate that metadata-enriched approaches consistently outperform content-only baselines, with recursive chunking paired with TF-IDF weighted embeddings achieving 82.5\% precision and prefix-fusion embeddings achieving NDCG of 0.813 in enterprise settings \cite{ref90}. Embedding metadata directly into the indexed representations increases intra-document cohesion, reduces inter-document confusion, and widens the separation between relevant and irrelevant chunks, thereby addressing the problem of cross-document confusion that arises in highly homogeneous corpora such as regulatory filings \cite{ref91}. At the corpus level, Cross-Document Topic-Aligned (CDTA) chunking goes further by reconstructing knowledge at the corpus level---identifying topics across documents, mapping segments to those topics, and synthesizing them into unified chunks---achieving 93\% faithfulness on multi-hop reasoning tasks, a 12\% improvement over contextual retrieval \cite{ref92}. This approach directly targets the knowledge fragmentation problem inherent to within-document segmentation, and it pairs naturally with the graph-based traversal methods described earlier, which also seek to consolidate evidence across document boundaries.

\textbf{Embedding Methods and Vector Compression.} The choice of embedding technique governs how chunks are projected into the semantic space, and thus determines the effectiveness of the dense, late-interaction, and hybrid retrieval paradigms that operate on these representations. Strategies range from content-only embeddings to dense, sparse, and metadata-fused variants. Late interaction and multi-vector retrieval, while powerful, suffer from storage and computation costs that grow linearly with document length; attention-guided clustering (AGC) has emerged as an effective query-agnostic compression method, identifying semantically salient regions as cluster centroids and consistently outperforming parameterized compression methods while maintaining competitive performance against uncompressed indexes across text, visual-document, and video retrieval tasks \cite{ref93}. Additionally, the document-side early compression problem---where a short but decisive span is diluted during document encoding---is addressed by training-free aggregation strategies such as DICE, which splits documents into chunks, encodes them independently, and aggregates them back into a single vector, improving long-document retrieval scores from 30.0 to 90.0 for Passkey \textgreater4k tokens \cite{ref94}. These compression strategies are particularly relevant in light of the scalability findings discussed in the following subsection, where corpus growth exposes efficiency bottlenecks in candidate generation.

\textbf{Advanced Index Structures and Filtering.} Beyond chunking and embedding, the organization of chunks within an index affects retrieval efficiency and accuracy, and it conditions the feasibility of the adaptive and progressive search strategies that will be examined in the following subsection. Entity-centric tables, as demonstrated by methods like SlimRAG, construct a compact entity-to-chunk table based on semantic embeddings, enabling entity-aware retrieval without graph traversal and achieving strong accuracy while drastically reducing index size (RITU of 16.31 vs.~56+ for graph-based baselines) \cite{ref95}. More complex architectures include multi-partite graphs and hierarchical organizations that support coarse-to-fine retrieval. HiKEY, for instance, reconstructs a logical heterogeneous graph via Document Hierarchical Parsing to explicitly encode parent-child relationships, thereby improving retrieval recall by up to 12.9\% via hierarchical routing \cite{ref96}. Chunk filtering strategies provide a complementary optimization: entity-based filtering reduces vector index size by approximately 25--36\% while preserving retrieval quality, effectively reducing redundancy introduced during chunking \cite{ref97}. Finally, the notion of separating the physical index entry from the evidence payload, as in M-RAG's semantic key-value indexing, uncouples retrieval-oriented compact records from generation-oriented evidence, offering competitive accuracy under tight token budgets and stronger robustness to expanding corpora \cite{ref98}.

In summary, effective indexing in RAG requires a holistic view that balances semantic coherence, structural fidelity, operational cost, and retrieval adaptability. The literature shows that no single strategy is universally optimal; rather, the design must be guided by the domain, the characteristics of the corpus, the target queries, and the computational constraints of the deployment environment. The decisions made at this stage directly determine the ceiling for the retrieval mechanisms that follow, and the next subsection will examine how these indexed representations are searched, merged, and ranked to deliver the evidence that grounds generation.

\subsection{Retrieval Mechanisms and Hybrid Search Architectures}

The retrieval stage is the critical gateway of a RAG system, determining which external evidence is available for the generator. The quality of this stage fundamentally limits the ceiling of the entire pipeline, as no amount of downstream generation enhancement can compensate for evidence that was never retrieved. Retrieval mechanisms have evolved along several distinct paradigms---sparse lexical methods, dense semantic models, late-interaction architectures, graph-based traversal, and a family of hybrid approaches that seek to combine their complementary strengths. This taxonomy is foundational to understanding the current design space, and its effectiveness is directly conditioned by the indexing decisions described in the preceding subsection: the chunking granularity, embedding method, and index architecture determine the candidate space that these retrieval mechanisms must navigate. At the same time, the output of this stage---a ranked set of candidate evidence---constitutes the raw material upon which the query processing and result refinement techniques of the following subsection will operate.

Sparse lexical retrieval, epitomized by BM25, remains a remarkably robust baseline even in the era of neural models. Its effectiveness is rooted in exact keyword matching, weighted by corpus-wide significance using term frequency and inverse document frequency signals. Several recent studies underscore its surprising competitiveness. A systematic benchmarking of retrieval strategies on financial documents containing mixed text and tables found that BM25 outperformed state-of-the-art dense retrieval methods on this domain, challenging the common assumption that semantic search universally dominates \cite{ref99}. This finding is echoed in analyses of scaling behavior, where lexical retrieval demonstrated a distinct advantage as corpus size grows. In a controlled study varying corpus size across 28 nested tiers, BM25 overtook agentic search methods at around 10 million corpus tokens and led at every larger shared tier, anchoring the low-cost end of the Pareto frontier without requiring LLM-based construction \cite{ref100}. The continued relevance of BM25 is also supported by evidence that it tends to be more robust than dense models to position bias, showing less degradation when relevant information appears later in a passage \cite{ref101}. However, BM25 is fundamentally blind to the semantic vocabulary gap: when a relevant document phrases an answer differently from the query, lexical methods fail to retrieve it, and no amount of downstream reranking can recover a document that was never in the candidate set \cite{ref102}. This limitation is particularly consequential in light of the indexing strategies discussed earlier: structure-aware chunking and metadata enrichment can partially mitigate this gap by improving the surface-level alignment between chunks and queries, but the fundamental lexical bottleneck remains.

Dense retrieval was developed to bridge this gap by encoding queries and documents into low-dimensional contextualized embedding spaces. Single-vector bi-encoders such as DPR map each passage to a single embedding, typically derived from a CLS token, allowing for efficient similarity search via approximate nearest neighbor algorithms. These models have shown strong performance in open-domain QA, with DPR achieving a top-1 accuracy of 50.17\% on the NQ dataset in a comparative study \cite{ref103}. However, the single-vector paradigm often over-compresses local relevance signals, a limitation that late-interaction models were designed to address. ColBERT introduces a late-interaction architecture that encodes queries and documents into token-level embeddings, computing relevance through the fine-grained MaxSim operation \cite{ref104}. This yields effectiveness approaching full cross-encoders while maintaining the efficiency of pre-computed document representations. Direct comparisons between the two families have shown that multiple representations are statistically more effective than single representations for MAP and MRR@10, particularly for queries that are hardest for BM25, including definitional queries and those with complex information needs \cite{ref105}. ColBERTv2 further improves the practical viability of late interaction by coupling aggressive residual compression with denoised supervision, reducing the space footprint by 6--10× while improving quality across a wide range of benchmarks \cite{ref106}. Subsequent optimizations have attempted to make late interaction affordable through various means: SLIM maps contextualized token vectors to a sparse lexical space, enabling compatibility with off-the-shelf lexical search libraries like Lucene \cite{ref107}. SPLATE learns an MLM adapter to map frozen token embeddings to a sparse vocabulary space, allowing candidates to be generated via traditional sparse retrieval that runs under 10ms on CPUs \cite{ref108}. A sparse coding approach using Sparse Autoencoders replaced K-means clustering entirely, achieving a 15× reduction in indexing time relative to ColBERTv2 while halving retrieval latency and improving performance \cite{ref109}. Token pruning techniques have shown that query embedding pruning based on collection frequency can reduce the number of documents retrieved by 70\% without statistically significant effectiveness losses, yielding a 2.65× speedup in end-to-end response time \cite{ref110}. Routing mechanisms have also been proposed, where a learned router dynamically directs queries to their most informative representations, reducing computational complexity by up to 30× while maintaining competitive accuracy \cite{ref111}. The emergence of non-Transformer backbones such as Mamba also provides promising alternatives, with linear time scaling in sequence length showing particular promise for long-text retrieval tasks \cite{ref112}. These efficiency-focused innovations are particularly relevant given the index compression and organization techniques described in the previous subsection, which similarly aim to reduce the computational burden of large-scale retrieval.

Knowledge graph-based retrieval represents a structurally distinct paradigm that leverages explicit relational semantics. Rather than relying on surface-level similarity or learned representations alone, graph-based methods construct indices that connect entities, documents, and concepts through typed relations, enabling traversal-based retrieval that can support multi-hop reasoning. These approaches often require substantial construction effort but can yield richer structural coverage. FlexStructRAG illustrates a flexible approach by jointly constructing a knowledge graph for binary relations, a knowledge hypergraph for n-ary relations, and structure-aware semantic clusters, supporting entity-, edge-, hyperedge-, and cluster-level retrieval that can be adaptively combined at inference time \cite{ref113}. The trade-off between construction cost and structural benefit is substantial: graph-based RAG approaches have been observed to encounter construction walls before reaching deployment scale, and their scalable variants may remain below BM25 accuracy at shared corpus tiers \cite{ref100}. This finding reinforces the importance of the indexing-stage trade-offs discussed earlier: the hierarchical and entity-centric index structures introduced there can offer some of the structural benefits of graph-based methods at a fraction of the construction cost.

The complementary strengths of these paradigms have driven the development of hybrid retrieval architectures that combine multiple retrieval mechanisms. Hybrid systems aim to maximize lexical and semantic matching simultaneously while minimizing the shortcomings of each approach \cite{ref114}. A key finding from financial document retrieval is that fusing BM25 with a compact dense encoder improves reference-level Hit@10 by roughly 28\% over either component alone on a corrected corpus with encoder-sized windows \cite{ref115}. A two-stage pipeline combining hybrid retrieval with cross-encoder neural reranking achieves Recall@5 of 0.816 and MRR@3 of 0.605 on heterogeneous text-and-table financial documents, outperforming all single-stage methods by a large margin \cite{ref99}. While simple weighted fusion is effective for general-purpose retrieval, the interaction between the fusion weight and query characteristics is complex. An oracle over the interpolation-weight grid showed 21.8\% headroom on financial filings, yet none of three lightweight adaptive routers---a score-confidence heuristic, a random forest over query features, and a ridge regressor over query embeddings---established statistically reliable improvements over the fixed blend \cite{ref115}. Nonetheless, dynamic alpha tuning approaches that leverage an LLM to evaluate the top-1 results from both dense and sparse branches and calibrate the fusion weight per query have demonstrated significant improvements over fixed-weighting baselines across multiple evaluation metrics \cite{ref116}. These hybrid architectures align naturally with the metadata-enriched indexing strategies from the previous subsection, which enhance both lexical and semantic matching simultaneously.

The design of hybrid systems involves critical choices about which components to combine and how to fuse their outputs. Evidence from an analysis of advanced hybrid search architectures reveals a ``weakest link'' phenomenon, where a weak retrieval path can substantially degrade overall accuracy, underscoring the need for path-wise quality assessment before fusion \cite{ref117}. The same study identifies data-dependent optimal configurations, precluding a one-size-fits-all solution, and introduces Tensor-based Re-ranking Fusion (TRF) as a high-efficacy alternative to mainstream fusion methods \cite{ref117}. A retrieval unit granularity choice interacts substantially with the retrieval mechanism: proposition-level indexing---where atomic factoid expressions serve as the retrieval unit---significantly outperforms traditional passage or sentence-based methods in dense retrieval \cite{ref118}. This granularity choice also affects downstream QA performance by delivering more condensed, question-relevant context. When evaluating the retrieval unit itself, a critical methodological caveat arises: if the retrieval unit is larger than the dense encoder's input window, the dense model never sees a large share of the labeled evidence, confounding comparisons against full-text sparse baselines \cite{ref115}. This observation motivates the use of encoder-sized windows for fair comparisons between sparse and dense methods, and it echoes the chunking-size trade-offs highlighted in the previous subsection, where the interaction between segmentation granularity and embedding capacity was shown to be a critical performance determinant.

Finally, progressive and adaptive search algorithms have been developed to balance computational cost with precision. Query performance prediction methods that estimate retrieval effectiveness for a given query have shown promise but display significant variability across collections and rankers, limiting their utility for selective query processing \cite{ref119}. Similarly, coherence-based predictors adapted for dense retrieval have demonstrated improved accuracy when using neural embedding representations \cite{ref120}. A classifier-based approach that selects a sparse, dense, or hybrid strategy per query based on the query alone demonstrates an improved range of efficiency/effectiveness trade-offs, allowing system operators to choose an appropriate strategy based on target latency and resource budget \cite{ref121}. Guided traversal techniques that leverage a traditional sparse model's signal to efficiently traverse the index allow more effective learned sparse models to execute fewer scoring operations, boosting efficiency by up to a factor of four without measurable effectiveness loss \cite{ref122}. These adaptive mechanisms, combined with advances in embedding and routing, enable retrieval systems to navigate the accuracy-cost frontier more effectively as corpus scales grow. They also foreshadow the adaptive and corrective mechanisms of the following subsection, which extend this principle of dynamic resource allocation from the retrieval stage into query processing and result refinement.

\subsection{Query Processing, Reranking, and Retrieval Optimization}

The efficacy of Retrieval-Augmented Generation (RAG) hinges critically on the quality of the evidence retrieved, which in turn is determined by the effectiveness of the query processing and result refinement stages. A poorly formulated query or a noisy candidate set can severely degrade the final generation quality, regardless of the generator's capability. This subsection dissects the techniques designed to enhance retrieval quality through query transformation and result refinement, covering a spectrum of strategies from initial query rewriting and expansion to sophisticated reranking and adaptive retrieval mechanisms. Together with the retrieval architectures discussed previously and the systematic pipeline optimization explored next, these techniques form the connective tissue that determines how effectively evidence flows from the index to the generator.

A foundational challenge in RAG is the lexical gap between a user's query and the language used in relevant documents. Traditional sparse retrieval methods like BM25 suffer from vocabulary mismatch, a problem that query expansion (QE) aims to mitigate. The advent of pre-trained and large language models (PLMs/LLMs) has fundamentally reshaped the QE design space, introducing new capabilities such as stronger contextualization, controllable generation, and instruction-following \cite{ref123}. This survey highlights that modern QE methods are organized along dimensions of where expansion is injected, how it is grounded in corpus evidence, and how structured knowledge is incorporated \cite{ref123}. For instance, generative pseudo-document generation, as opposed to simple synonym addition, can better bridge the gap between short queries and long documents, a finding particularly relevant in dense retrieval contexts \cite{ref124}. However, all expansion techniques are not equally beneficial. A key finding is that query expansion can improve the generalization of even strong cross-encoder rankers if the process involves generating high-quality keywords through a reasoning chain and aggregating the ranking results of each expanded query via fusion, a departure from the notion that expansion primarily benefits weaker first-stage retrievers \cite{ref125}. Classical pseudo-relevance feedback (PRF) methods also benefit from LLM integration, where an LLM acts as a filter to judge the relevance of documents in the initial top-k ranking before expansion terms are estimated, mitigating topic drift and improving robustness \cite{ref126}. Furthermore, the theoretical equivalence between query and key expansion has been demonstrated, enabling techniques like Evolving Retrieval Memory (ERM) to convert transient query-time gains into persistent index-side improvements without any inference-time overhead \cite{ref127}. This interplay between query-side processing and index-side construction echoes the hybrid retrieval architectures discussed in the previous subsection, where the boundary between online and offline computation is increasingly fluid.

For complex, multi-hop questions that require synthesizing information across multiple documents, simple query expansion is often insufficient as relevant facts are distributed across sources. This necessitates more structural query transformation, primarily through decomposition. A straightforward yet effective pipeline incorporates question decomposition, where an LLM decomposes the original query into sub-questions, retrieves passages for each, and merges the candidate pool for a final reranking step \cite{ref128}. This approach has been shown to effectively assemble complementary documents and improve retrieval and answer accuracy on benchmarks like MultiHop-RAG and HotpotQA \cite{ref128}. The efficacy of decomposition, however, is stage-dependent. Research indicates that decomposition during initial retrieval can harm performance due to semantic dilution, while it substantially improves reranking by enabling more fine-grained constraint verification \cite{ref129}. This insight has led to stage-aware frameworks that retain a monolithic query for retrieval but employ sub-queries exclusively during reranking \cite{ref129}. Other frameworks, such as DeCoR, advocate for a structured refinement approach that decomposes a query into explicit reasoning steps and compresses supporting evidence from retrieved documents, proving to be more efficient and effective than relying solely on the generative capabilities of larger LLMs \cite{ref130}. The strategic goal is often to control the search process adaptively; for instance, the Adaptive Complex Query Optimization (ACQO) framework uses reinforcement learning to determine \emph{when} to decompose a query and \emph{how} to fuse the scores from multiple sub-query results, preventing unstable training and reducing unnecessary search \cite{ref131}. Similarly, a strategy-guided approach for query rewriting, known as SAGE, demonstrates that guiding LLMs with a concise set of expert-crafted strategies---such as semantic expansion and entity disambiguation---can substantially improve retrieval effectiveness \cite{ref132}. This highlights a broader principle: the transformation of queries should not be a free-form generation task but a guided process aimed at eliciting decision-relevant evidence. For instance, Hypothesis-Conditioned Query Rewriting (HCQR) reorients retrieval from topic-oriented to evidence-oriented by generating queries that seek evidence to support, distinguish, and verify an initial hypothesis derived from the question and candidate answers \cite{ref133}.

Once a candidate set of documents is retrieved, the task shifts to result refinement to ensure the most relevant and non-redundant information is passed to the generator. Reranking is a critical step here, transitioning from fast, first-stage retrieval to more precise, often more computationally intensive, scoring. A strong cross-encoder reranker often accounts for the majority of an RAG pipeline's quality, and the benefits of other complex retrieval enhancements must be evaluated in its presence \cite{ref134}. The landscape of reranking has evolved from pointwise relevance scoring to more holistic, set-level, and listwise approaches. Traditional rerankers score passages independently, which can fail in multi-hop scenarios where evidence usefulness is conditional on its compatibility with other passages. To address this, set-level retrieval frameworks formulate multi-hop retrieval as a query-set compatibility scoring problem, training retrievers to rank complete and compatible evidence sets above incomplete or noisy alternatives \cite{ref135}. Another perspective is to shift from ranking to explicit set selection, where a model like SETR identifies the information requirements of a query through Chain-of-Thought reasoning and selects an optimal set of passages that collectively satisfy these requirements, outperforming traditional rerankers \cite{ref136}. The set-selection problem has also been cast as a Quadratic Unconstrained Binary Optimization (QUBO) problem, formulating an energy function that balances relevance, coverage, and redundancy to find an optimal evidence subset, effectively separating combinatorial selection from semantic answer generation \cite{ref137}. A unified ``Expand-then-Refine'' paradigm, exemplified by OptiSet, leverages set-level utility changes to derive preference labels and train a model for compact, high-gain evidence set selection \cite{ref138}. These approaches are complemented by architecture designs that model inter-document dependencies, such as a dual-view cascaded reranking framework that uses a local scorer for fine-grained query-document relevance and a global scorer to model inter-document dependencies, dynamically fusing the two views for multi-hop reranking \cite{ref139}. Furthermore, the allocation of reasoning budget in the reranking stage itself is a critical consideration, with studies showing that moderate listwise reranking depth can yield larger gains than increasing search-time reasoning, at a substantially lower token cost \cite{ref140}. This observation directly informs the cost-accuracy trade-offs that the pipeline optimization frameworks discussed in the next subsection aim to navigate systematically.

Given the computational cost of processing long retrieved contexts, query-aware compression is essential for maximizing information density within a token budget. MergeRAG rethinks augmentation as a dynamic optimization problem, using a scoring agent to restructure retrieved contexts through symmetric merging (to consolidate weak signals) and asymmetric merging (to reduce redundancy), thereby addressing the limitations of static top-k filtering \cite{ref141}. Similarly, BRIEF compresses retrieved documents into dense textual summaries via a query-aware process, effectively bridging retrieval and inference while lowering latency and costs \cite{ref142}.

Beyond these static and iterative pipelines, the frontier of RAG research includes corrective and adaptive retrieval mechanisms that optimize the retrieval process based on intermediate outcomes. More advanced methods, such as Self-Correcting RAG, formalize context selection as a multi-dimensional multiple-choice knapsack problem to maximize information density while removing redundancy, and further employ NLI-guided Monte Carlo Tree Search to explore reasoning trajectories and ensure faithfulness \cite{ref143}. Adaptive controllers are also being developed to escalate retrieval effort based on query complexity. The selection of an optimal strategy itself can be learned through reinforcement learning, where a strategy-adaptive engine learns to select the best query rewriting strategy, reducing unnecessary exploration and lowering inference costs \cite{ref132}. Furthermore, these mechanisms can be systematically designed by formulating RAG design as an architecture search problem, as demonstrated by the RAISE framework, which evaluates optimization methods for RAG pipelines under standardized search spaces \cite{ref144}. This connection foreshadows the central theme of the next subsection: the shift from ad-hoc, manual pipeline composition toward principled, automated search over the configuration space. In conclusion, the field is moving toward more intelligent and adaptive pipelines that not only retrieve but actively plan, decompose, refine, and verify the evidence that grounds the generation process, navigating a complex landscape of search strategies and scoring functions to balance precision, recall, and computational efficiency---a balance that increasingly demands the systematic optimization approaches explored in the remainder of this survey.

\subsection{Pipeline Composition, Hyperparameter Tuning, and Benchmarking}

The performance of a Retrieval-Augmented Generation (RAG) system is not determined by any single component in isolation but by the intricate and often non-obvious interactions between its many design choices. From query rewriting and chunking strategies to retrieval depth, reranking, and context compression, the configuration space is vast. In practice, these choices are frequently made based on heuristics or default settings, which can lead to suboptimal performance and hinder reproducibility. Recognizing this, a growing body of research treats RAG design not as a manual tuning exercise but as a formal architecture search problem, emphasizing the need for systematic and principled approaches to pipeline composition and hyperparameter optimization. This shift toward systematic optimization is a natural extension of the corrective and adaptive retrieval mechanisms discussed in the previous subsection, where the focus on \emph{when} and \emph{how} to intervene in the retrieval process is now elevated to the level of the entire pipeline architecture.

A foundational step in this direction is the development of comprehensive frameworks for architecture search. The RAISE framework formalizes the challenge of optimizing RAG pipelines as a hyperparameter optimization problem, providing a standardized benchmark with 13 search algorithms evaluated across seven public text and multimodal datasets \cite{ref144}. Its experiments reveal that optimization performance is highly task-dependent, cautioning against the assumption that a single search strategy is universally superior. This suggests that practitioners should select optimization algorithms based on their specific domain and requirements. Similarly, RAGSmith offers a modular framework that treats RAG design as an end-to-end architecture search over a massive space of 46,080 feasible pipeline configurations spanning nine technique families \cite{ref145}. Using a genetic search algorithm, RAGSmith consistently outperformed a naive RAG baseline by an average of +3.8\% across various domains, with gains of up to +12.5\% in retrieval and +7.5\% in generation. Notably, it discovered a robust backbone of vector retrieval combined with post-generation reflection, augmented by domain-dependent choices for expansion, reranking, and prompt reordering. These findings underscore the value of automated search in uncovering high-performing configurations that might not be intuitive. For practitioners with more constrained resources, AutoRAGTuner provides a declarative, configuration-driven framework that automates the entire RAG lifecycle, including construction, execution, evaluation, and optimization \cite{ref146}. Its modular architecture and integration of an adaptive Bayesian optimization engine allow it to consistently outperform default baselines across diverse pipelines, from vanilla to graph-based, while its declarative language can reduce code churn for architectural adjustments by up to 95\%.

Beyond these comprehensive frameworks, a significant body of work focuses on isolating the contribution of individual components through ablation studies and controlled comparisons. For instance, a controlled empirical study in the biomedical domain compared five retrieval strategies---Dense Vector Search, Hybrid BM25+Dense, Cross-Encoder Reranking, Multi-Query Expansion, and Maximal Marginal Relevance---while holding the generation model, vector store, and embeddings fixed \cite{ref147}. This setup ensured that observed differences were attributable solely to the retrieval strategy. The study found that Cross-Encoder Reranking achieved the best composite score, while the Dense baseline performed nearly as well, and surprisingly, Multi-Query Expansion had the weakest contextual precision, suggesting that naive query diversification can introduce noise. This highlights the critical role of controlled experimentation in making informed design decisions, and it directly connects to the earlier discussion of query expansion techniques, reinforcing that not all enhancements are universally beneficial. Extending this to heterogeneous documents, a benchmark study on text-and-table financial documents compared ten retrieval strategies and found that a two-stage pipeline combining hybrid retrieval with neural reranking significantly outperformed all single-stage methods, while also challenging the common assumption that dense retrieval universally dominates by showing BM25 outperforming state-of-the-art dense retrieval on financial documents \cite{ref99}. However, an important counterpoint emerges from a study on the HetDocQA benchmark, which found that once a strong cross-encoder reranker is present, many subsequent retrieval enhancements---such as hierarchical summarization, graph expansion, and rank fusion---provide no reliable gains, with only query expansion and a novel calibrated corrector (SSCC) yielding improvements \cite{ref134}. This suggests that the marginal utility of additional complex modules diminishes when a robust reranker is in place, echoing the earlier observation about rerankers accounting for the majority of pipeline quality, and providing a critical insight for cost-effective pipeline design.

The sensitivity of RAG systems to hyperparameters is a recurring theme, with several studies analyzing its impact on the speed-quality trade-off. An analysis of vector stores, chunking policies, and reranking found a clear speed-accuracy trade-off, where Chroma processed queries faster while Faiss yielded higher retrieval precision \cite{ref148}. It also revealed that naive fixed-length chunking outperformed semantic segmentation while being quicker, and that while reranking improves retrieval quality, it increases runtime by a factor of five. This type of analysis provides practical guidance for balancing computational cost and accuracy, complementing the earlier discussion of token-cost trade-offs in reranking depth. Furthermore, a comprehensive study of hyper-parameter optimization (HPO) methods found that RAG HPO can be done efficiently with greedy or random search and significantly boosts performance, but recommends optimizing model selection first rather than following the pipeline order \cite{ref149}. For a multi-objective perspective, researchers have introduced the first approach for optimizing cost, latency, safety, and alignment simultaneously across entire LLM and RAG systems, finding that Bayesian optimization methods significantly outperform baselines in achieving a superior Pareto front \cite{ref150}. In the retrieval stage, a framework called CAMI formalizes multi-index construction as a budgeted, multi-objective portfolio selection problem, showing that it can systematically identify high-recall portfolios while using up to 5x less budget than random search \cite{ref151}.

Finally, metadata-driven enhancements have emerged as a powerful tool for improving retrieval quality, bridging the gap between structured knowledge and unstructured text. A systematic framework for enterprise knowledge retrieval demonstrated that LLM-generated metadata consistently improves retrieval accuracy over content-only baselines, with a combination of recursive chunking and TF-IDF weighted embeddings achieving high precision \cite{ref90}. This approach was further validated in the financial domain, where embedding chunk metadata directly with text (``contextual chunks'') provided the most significant performance gains, even outperforming powerful rerankers \cite{ref152}. The concept of a ``Meta Knowledge Summary'' for metadata-based clusters has also been proposed, which augments user queries to significantly improve retrieval precision and recall \cite{ref153}. In summary, the field is moving toward a more rigorous, evidence-based approach to RAG pipeline design, where automated search, systematic ablation, and an understanding of hyperparameter sensitivity are essential for building high-performing, efficient, and cost-effective systems. This principled optimization of the entire pipeline, encompassing the query transformation and result refinement techniques detailed earlier, constitutes the final layer of sophistication in RAG systems, ensuring that every component works in concert to deliver reliable, grounded, and efficient generation.

\section{Generation Enhancement and Knowledge Integration}

\subsection{Augmentation Strategies and Contextual Compression}

The integration of retrieved information into the generation process constitutes a pivotal determinant of RAG system performance. While effective retrieval ensures that relevant evidence reaches the generator, the manner in which this evidence is processed, filtered, and presented can profoundly influence output quality, faithfulness, and computational efficiency. The retrieved context is rarely pristine; it often contains redundant passages, irrelevant noise, and conflicting information that can dilute the model's attention and exacerbate issues such as the ``lost-in-the-middle'' phenomenon. Consequently, augmentation strategies and contextual compression have emerged as essential components of the RAG pipeline, spanning pre-retrieval, mid-retrieval, and post-retrieval processing stages. This stage of context curation acts as the critical interface between the retrieval module and the reasoning frameworks discussed in the following section, establishing the quality and structure of the evidence that reasoning and conflict resolution mechanisms will subsequently operate upon.

Early augmentation approaches focused on refining the query and the retrieval process itself. Query rewriting, expansion, and decomposition are designed to improve the alignment between user intent and the knowledge base. These pre-retrieval techniques are complemented by mid-retrieval strategies that involve iterative or adaptive retrieval, where the system decides whether additional information is needed based on the initial results. Post-retrieval processing, however, has gained significant attention as a means to curate the final context provided to the generator. This stage encompasses a variety of contextual compression methods that aim to maximize information density while minimizing noise, thereby presenting the reasoning components with a more tractable and focused evidence set.

Within post-retrieval processing, the most direct form of compression is extractive filtering, which selects the most relevant sentences or passages from the retrieved documents. The challenge lies in determining relevance, particularly under varying query complexities and retrieval quality. A key insight is that the optimal amount of context is query-dependent; complex, multi-hop queries may require more documents, while simple queries may be answered with a single passage. To address this, adaptive compression methods have been developed. For instance, AdaComp introduces a low-cost extractive compression method that trains a predictor to determine the minimum number of top-k documents necessary for the RAG system to answer the current query, based on both query complexity and retrieval quality \cite{ref154}. Similarly, Dynamic Context Selection proposes a context-size classifier that dynamically predicts the optimal number of documents to retrieve, mitigating the introduction of distractors and the positional bias inherent in fixed-k strategies \cite{ref155}. These approaches highlight a shift away from static heuristics toward query-aware, performance-driven compression.

Beyond simple extractive methods, more sophisticated compression techniques leverage the internal signals of LLMs. Attention-guided methods, for example, utilize the model's own attention patterns to identify relevant information. QUITO introduces a query-guided attention compression method that calculates the attention distribution of the context in response to the question, using this signal to filter out useless information \cite{ref156}. Building on this, AttentionRAG reformulates RAG queries into a next-token prediction paradigm to isolate the query's semantic focus, enabling precise attention calculation and achieving significant context compression while outperforming methods like LLMLingua \cite{ref157}. AttnComp further refines this idea by leveraging the attention mechanism to identify relevant information and employing a Top-P compression algorithm to retain the minimal set of documents whose cumulative attention weights exceed a threshold, while also estimating response confidence \cite{ref158}. Sentinel takes a different approach by decoding inference-time contextual utilization behaviors from head-wise attention patterns, training a lightweight probe to perform compression using only a single forward pass without dedicated compression training \cite{ref159}.

Instruction-aware and semantic-based methods offer another dimension of contextual compression. Instruction-Aware Contextual Compression filters out less informative content by considering the user's specific instruction, achieving substantial reductions in context-related costs and inference latency while maintaining performance \cite{ref160}. A more linguistically grounded approach leverages Abstract Meaning Representation (AMR) graphs to preserve semantically essential information. These methods quantify node-level entropy within AMR graphs to estimate conceptual importance, distilling raw retrieved documents into a compact set of crucial concepts that filter out interference information \cite{ref161,ref162}. This concept-based distillation explicitly constrains the LLM to focus on vital information, proving particularly effective as the number of supporting documents increases. Clustering-based methods, such as EDC2-RAG, exploit latent inter-document relationships to simultaneously remove irrelevant information and redundant content, demonstrating robust performance improvements across various QA and hallucination-detection datasets \cite{ref163}.

Soft compression methods represent a paradigm shift by encoding long contexts into compact, continuous representations (e.g., embeddings) rather than pruning discrete tokens. xRAG reinterprets document embeddings from dense retrieval as features from a separate modality and fuses them into the language model's representation space, effectively eliminating the need for textual context and achieving extreme compression rates with a single token \cite{ref164}. COCOM similarly reduces long contexts to a handful of Context Embeddings, significantly speeding up decoding time while maintaining performance \cite{ref165}. However, soft compression methods are not without challenges. A critical analysis has revealed that autoencoder-like full-compression methods are both infeasible and non-necessary, as they conflict with downstream generation behavior and dilute task-relevant information. This has motivated selector-based frameworks like SeleCom, which redefines the encoder's role as a query-conditioned information selector, trained with curriculum learning to achieve competitive performance with reduced computation \cite{ref166}. Hybrid approaches, such as HyCo\(_2\), integrate both global and local perspectives, using a hybrid adapter to refine global semantics and a classification layer for local token retention, balancing essential semantics with critical details \cite{ref167}.

The design of the context itself, irrespective of compression, has a significant impact on performance. Research has shown that seemingly superficial choices, such as delimiters or structural markers, can induce substantial shifts in accuracy and stability. Contextual Normalization, a lightweight strategy that adaptively standardizes context representations before generation, has been shown to improve robustness to order variation and strengthen long-context utilization \cite{ref168}. Principled approaches to context engineering are also emerging. Conformal prediction offers a coverage-controlled filtering framework that removes irrelevant content while preserving the recall of supporting evidence, providing statistical guarantees on the retained context and enabling reliable, model-agnostic context reduction \cite{ref169}.

Overall, the landscape of augmentation strategies and contextual compression is diverse and rapidly evolving, encompassing extractive, generative, and soft compression paradigms. The field is converging on the principle that effective context is not merely a function of retrieval accuracy but also of intelligent curation and presentation. This curated context forms the substrate upon which the reasoning-enhanced generation frameworks and conflict-resolution mechanisms described in the following section must operate. Adaptive methods that query-aware compression rates, attention-guided pruning, semantic distillation, and principled statistical guarantees represent the current frontier, offering promising avenues for enhancing the efficiency, faithfulness, and reliability of RAG systems by equipping the generation and reasoning stages with the highest-quality evidence possible.

\subsection{Reasoning-Enhanced Generation and Knowledge Conflict Resolution}

The integration of retrieved knowledge into the generation process is not a trivial feed-forward operation; it requires sophisticated reasoning frameworks to synthesize disparate information and robust mechanisms to resolve conflicts between the parametric memory of the LLM and the non-parametric evidence provided \cite{ref170}. The effectiveness of a RAG system hinges on its ability to not only retrieve relevant information but also to reason over it correctly and to navigate the complex knowledge landscape where internal priors and external evidence may disagree \cite{ref171}. Building upon the curated and compressed context described in the previous subsection, this section addresses the computational and inferential processes that transform that evidence into reliable answers, while also laying the groundwork for the faithfulness and attribution concerns that will be examined in the next subsection.

A primary challenge is that standard RAG pipelines treat retrieved documents as independent, unstructured chunks, which forces models to implicitly connect fragmented information \cite{ref172}. This limitation becomes critical for multi-hop queries, which require synthesizing facts scattered across multiple sources. To address this, reasoning-enhanced generation paradigms have been developed to explicitly structure the reasoning process. For instance, Structure-Augmented Reasoning Generation (SARG) post-processes retrieved text by extracting relational triples, organizing them into a knowledge graph, and then performing multi-hop traversal to identify relevant reasoning chains \cite{ref172}. Similarly, planning-based frameworks like Recursive Evaluation and Adaptive Planning (REAP) maintain an explicit global perspective through a Sub-task Planner and a Fact Extractor module, which perform fine-grained analysis to iteratively update the task-solving trajectory \cite{ref173}. Other approaches, such as the Contrastive-RAG (C-RAG) framework, elicit critical reasoning by generating explanations that explicitly contrast the relevance of retrieved passages to support the final answer \cite{ref174}. To make this reasoning process more transparent and interpretable, frameworks like ArgRAG replace black-box reasoning with structured inference using a Quantitative Bipolar Argumentation Framework (QBAF), enabling decisions to be explained and contested \cite{ref175}.

Beyond structural reasoning, the core issue of knowledge conflict resolution remains a central focus. Knowledge conflicts arise when the retrieved context contradicts the model's internal parametric knowledge, or when retrieved documents themselves are inconsistent, and these conflicts can severely undermine the reliability of RAG systems \cite{ref176,ref177}. Research has shown that the failure to resolve these conflicts is often a more significant driver of hallucination than retrieval accuracy itself, as evidence can be retrieved but not correctly integrated during generation due to prior-driven overrides \cite{ref178}. This suggests the problem is not merely a retrieval issue but a deep-seated generation or ``integration bottleneck'' \cite{ref179}. A broad spectrum of techniques has been proposed to tackle this, ranging from fine-tuning and training-based methods to inference-time decoding and architectural interventions.

Training and fine-tuning approaches aim to align the LLM's behavior with the desired conflict-resolution policy. For example, a reasoning-trace-augmented RAG framework adds structured, interpretable reasoning across document-level adjudication, conflict analysis, and grounded synthesis, demonstrating significant gains in answer correctness and behavioral adherence after supervised fine-tuning \cite{ref180}. Another approach, Knowledgeable-R1, uses a reinforcement-learning framework to train models to resist contextual interference while still exploiting helpful external evidence \cite{ref181}. Inference-time methods are also popular, including GuarantRAG, which explicitly decouples reasoning from evidence integration by generating an ``Inner-Answer'' from parametric knowledge and a ``Refer-Answer'' using a Contrastive DPO objective, then fusing them with joint decoding \cite{ref179}. The Grounded Decoding framework offers a training-free alternative that fuses the RAG distribution with a retrieval-only distribution using a KL-barycenter objective, with an adaptive weighting scheme based on distributional disagreement \cite{ref182}.

To address the challenge of conflicting evidence within a single retrieved set, some methods focus on identifying conflicts at a structural or representation level. The ConflictRAG framework introduces a two-stage conflict detection module to identify, classify, and resolve conflicts prior to generation \cite{ref183}. Similarly, ArbGraph decomposes retrieved documents into atomic claims and organizes them into a conflict-aware evidence graph to pre-generate evidence arbitration \cite{ref184}. Another class of methods, including ProbeRAG, moves beyond black-box interventions to probe the model's internal latent space, discovering that conflicting and aligned knowledge states are linearly separable, and using this to perform fine-grained knowledge pruning and conflict-aware attention modulation \cite{ref185}. Building on similar mechanistic insights, SHIFT adopts a gated activation steering approach to modulate internal activations, and CoRect contrasts logits from contextualized and non-contextualized forward passes to identify and rectify layers exhibiting high parametric bias \cite{ref186,ref187}. Parametric Knowledge Muting through FFN Suppression (ParamMute) goes further by identifying and suppressing specific feed-forward networks that are disproportionately activated during unfaithful generation \cite{ref188}.

Finally, to convincingly demonstrate the generality of these problems, researchers have developed realistic benchmarks and evaluation protocols that go beyond simple accuracy metrics to measure context compliance and conflict handling \cite{ref189,ref171}. These diagnostic studies, such as Context-Driven Decomposition (CDD), provide inference-time interventions to reveal whether a model relied on retrieved evidence or its parametric prior \cite{ref171}. Moreover, the question of when to rely on which source is central; frameworks like TCR (Transparent Conflict Resolution) and SABER aim to make this decision process observable, controllable, and self-aware, providing mechanisms for the model to trust its parameters, trust the retrieved context, or abstain from answering when both are unreliable \cite{ref190,ref191}. The reasoning and conflict-resolution mechanisms discussed here thus directly condition the faithfulness of the final output, establishing the conceptual bridge to the attribution frameworks and hallucination mitigation strategies explored next.

\subsection{Faithfulness, Citation Attribution, and Hallucination Reduction}

The integration of retrieved knowledge into the generation process is fundamentally challenged by two intertwined problems: ensuring that the generated text remains faithful to the provided evidence, and providing reliable, granular attributions that allow users to verify the model's claims. Trustworthy RAG systems must therefore address not only \emph{what} to generate, but also \emph{why} a particular claim is made and \emph{where} the supporting evidence resides. This subsection explores the mechanisms developed to address these challenges, focusing on citation generation frameworks, metrics for evaluating attribution quality, and a spectrum of hallucination detection and mitigation strategies that operate at various stages of the generation pipeline.

A primary approach to enhancing verifiability is the generation of explicit citations. However, generating effective citations is not a single, monolithic task. Frameworks often decompose this process into distinct stages to improve precision. For instance, the ``Attribute First, then Generate'' paradigm breaks down the conventional end-to-end generation into content selection, sentence planning, and sequential sentence generation, ensuring that selected source segments act as fine-grained attributions, yielding more concise citations and reducing human verification time \cite{ref192}. Building on this idea of structured decomposition, the VeriCite framework introduces a three-stage process of initial answer generation, supporting evidence selection, and final answer refinement, using an NLI model to verify claims and optimize citation quality without sacrificing answer correctness \cite{ref193}. Another framework, FullCite, focuses on structured inline citations linking each claim to both its source document and supporting evidence, revealing that while LLMs are effective at identifying relevant documents, they often struggle with precise evidence span identification \cite{ref194}. These findings suggest that the granularity of citations is a critical design choice; enforcing overly fine-grained, sentence-level citations can disrupt the semantic dependencies needed for synthesis, while excessively coarse citations introduce noise, with performance often peaking at intermediate, paragraph-level granularity \cite{ref195}.

Beyond the architecture of citation generation, a crucial distinction arises between citation correctness and citation faithfulness. Correctness evaluates whether the cited document supports the statement, while faithfulness ensures the model genuinely relied on that document during generation, rather than producing a post-rationalization. Research highlights that prevalent ``post-rationalization'' undermines reliable attribution, as current attributed answers often lack citation faithfulness \cite{ref196}. To address this, some methods condition generation on the context surrounding citation markers, transforming them from generic placeholders into active knowledge pointers \cite{ref197}. The choice between generation-time and post-hoc citation paradigms also presents a trade-off. Generation-Time Citation (G-Cite) produces answers and citations in one pass, prioritizing precision, while Post-hoc Citation (P-Cite) adds citations after drafting, achieving higher coverage with competitive correctness, leading to recommendations for a retrieval-centric, P-Cite-first approach in high-stakes applications \cite{ref198}. To systematically evaluate the quality of attributions, novel protocols such as citation masking and recovery have been proposed, which assess whether a model can correctly recover the evidence for a given claim \cite{ref199}. Furthermore, the usefulness of citations can be measured by their sufficiency, ensuring they contain all essential information, achieved through sub-sentence citation generation guided by credit models \cite{ref200}.

Parallel to advances in attribution, substantial effort has been directed toward the core problem of hallucination detection and mitigation---a problem that is intimately connected to the knowledge conflict resolution mechanisms discussed earlier. Indeed, unresolved conflicts between parametric memory and external evidence are a primary driver of hallucination, as the model may default to prior knowledge or fail to integrate conflicting evidence correctly \cite{ref178}. A nuanced approach to this problem involves moving beyond answer-level evaluation to facet-level diagnostics. By decomposing questions into atomic reasoning facets and analyzing a Facet x Chunk matrix of retrieval relevance and NLI-based faithfulness, this framework reveals that hallucinations are often driven less by retrieval accuracy and more by the integration of evidence during generation, exposing systematic evidence override patterns \cite{ref178}. This aligns with findings that the dominant knowledge source behind an answer---parametric memory or provided context---can be identified with a linear probe, and attribution mismatches can raise error rates by up to 70\% \cite{ref201}. For a more structured analysis, HART formalizes hallucination tracing as a four-stage task of span localization, mechanism attribution, evidence retrieval, and causal tracing, providing a structured dataset for causal-level interpretability \cite{ref202}.

Mitigation strategies largely fall into training-based, decoding-based, and post-hoc categories. Training-time approaches, such as VeriFY, teach models to reason about their own factual uncertainty through consistency-based self-verification, augmenting training with verification traces but using loss masking to avoid reinforcing hallucinated content \cite{ref203}. On the other hand, decoding-time interventions offer a flexible, training-free alternative that is particularly well-suited to RAG settings where the retrieved context must be dynamically integrated. TRACE corrects hallucinations by selecting the corrective layer and operator based on the trajectory of candidate tokens within the model's own forward passes, improving factuality across numerous models without finetuning \cite{ref204}. Similarly, Attribution-Guided Decoding (AGD) steers generation by selecting token candidates that exhibit the highest attribution to a user-defined region of interest, such as contextual sources, reducing hallucinations and improving factual accuracy \cite{ref205}. Another lightweight approach, HAVE, addresses two issues---input-agnostic head importance and inaccurate raw attention weights---through head-adaptive gating and value calibration to construct token-level evidence \cite{ref206}. Monitoring Decoding (MD) dynamically monitors the generation process, revising hallucination-prone tokens identified by a monitor function through tree-based decoding \cite{ref207}. Finally, GRAD constructs a sparse token transition graph from a small retrieved corpus and adaptively fuses graph-derived logits with model logits to favor high-evidence continuations, offering a lightweight and plug-and-play method for grounding \cite{ref208}. These diverse strategies, spanning training, decoding, and architectural interventions, collectively illustrate a robust ecosystem for enhancing the faithfulness and trustworthiness of RAG systems. Together with the attribution frameworks described earlier, they form a comprehensive toolkit for ensuring that RAG outputs are not only accurate but also transparent and verifiable---a prerequisite for deployment in high-stakes applications where users must be able to trust and audit the model's claims.

\subsection{Robustness to Noisy and Adversarial Evidence}

Retrieval-Augmented Generation (RAG) systems are fundamentally dependent on the quality of retrieved evidence, yet this dependency introduces a critical vulnerability: in real-world deployments, the retrieval pipeline is rarely perfect, and the external knowledge base can be compromised or inherently noisy. While the preceding discussion centered on improving attribution and mitigating hallucination under the assumption of reasonably reliable retrieval, this subsection confronts harsher realities. It addresses challenges that arise when retrieved passages are irrelevant, misleading, or conflicting---and extends to malicious scenarios where adversaries actively inject crafted documents designed to manipulate the model's output. The discussion is structured around two broad categories of solutions: those that enhance the generator's intrinsic resilience through fine-grained knowledge processing and training, and those that implement explicit defense mechanisms against misinformation pollution and adversarial content.

A primary source of vulnerability stems from the fact that retrieved contexts are often a mixture of relevant and irrelevant information, which, if left unchecked, can undermine the very faithfulness that the attribution and hallucination techniques from the previous subsection aim to secure. To mitigate this, recent work emphasizes the importance of noise-robust generation techniques that operate at a fine-grained level. For instance, the concept of fine-grained knowledge pruning has been proposed to filter out irrelevant fragments from the retrieved context before it is fed to the generator \cite{ref185}. This approach moves beyond simple document-level filtering to identify and remove noisy tokens or sentences that could otherwise derail the generation process. Furthermore, the internal state of the LLM can be leveraged to assess the reliability of external information. Research has shown that conflicting and aligned knowledge states are linearly separable in the model's latent space, and that the presence of contextual noise increases the entropy of these representations \cite{ref185}. This insight has led to the development of frameworks that probe the model's latent space to detect hard conflicts between parametric and external knowledge, and subsequently modulate attention heads to prioritize faithful context integration---an approach that has demonstrated substantial improvements in contextual faithfulness \cite{ref185}.

Another critical direction for enhancing intrinsic robustness is through specialized fine-tuning, which complements the training-based hallucination mitigation strategies described earlier. Traditional fine-tuning of large language models often suffers from catastrophic forgetting when exposed to adversarial examples, failing to retain the robust linguistic features captured during pre-training \cite{ref209}. In the context of RAG, a dedicated method called Robust Fine-Tuning (RbFT) has been introduced to bolster the resilience of LLMs against retrieval defects by employing two targeted fine-tuning tasks \cite{ref210}. This approach is designed to maintain high inference efficiency while making the model more resistant to noisy, irrelevant, or misleading counterfactual information retrieved from the knowledge base \cite{ref210}. Another fine-tuning-based strategy focuses on teaching models to discern and process information based on its source credibility. The Credibility-aware Generation (CAG) framework generates training data based on information credibility, thereby equipping models with the ability to weigh the influence of documents according to their reliability, which significantly improves performance and resilience against noisy documents \cite{ref211}.

Beyond training-time interventions, inference-time defense mechanisms are crucial for safeguarding RAG systems against misinformation pollution, addressing the other side of the reliability coin introduced by imperfect retrieval. A prominent threat model involves scenarios where semantically relevant passages contain subtle misinformation, framing, or fabrications. To counter this, the MIRAGE framework constructs an NLI-based cross-document claim graph to identify a consistent, multi-source supported subset of evidence \cite{ref212}. By applying a ``Defended-Claims Gate,'' MIRAGE either conditions generation on this verified subset or blocks retrieval entirely to answer parametrically, thereby restoring factuality under polluted evidence conditions \cite{ref212}. Related to this, a ``skeptical prompting'' strategy can activate the LLM's internal reasoning to partially self-defend against adversarial passages, though its effectiveness is highly dependent on the model's inherent reasoning capacity \cite{ref213}. However, a significant finding in the literature highlights a ``monitoring-control gap'': models can readily acknowledge contradictory evidence while failing to let this awareness constrain their final recommendations, suggesting that simple detection mechanisms are insufficient for safe resolution \cite{ref214}.

To address the challenge of unreliable contexts more holistically, sophisticated adaptation mechanisms have been developed that go beyond static defenses, dynamically balancing internal parametric knowledge with external evidence---a theme that resonates with the knowledge conflict resolution discussed in prior sections. The Astute RAG approach addresses imperfect retrieval by adaptively eliciting essential information from the LLM's internal knowledge and iteratively consolidating internal and external knowledge with source-awareness, finalizing answers based on information reliability \cite{ref215}. This method is particularly effective in resolving knowledge conflicts, achieving performance comparable to or even surpassing the conventional use of LLMs even under worst-case retrieval scenarios \cite{ref215}. Similarly, the Stateful Evidence-Driven RAG framework models question answering as a progressive evidence accumulation process, maintaining a persistent evidence pool that captures both supportive and non-supportive information \cite{ref216}. By performing evidence-driven deficiency analysis to identify gaps and conflicts, this framework iteratively refines queries and guides subsequent retrieval, thereby achieving stable evidence aggregation and robustness against substantial retrieval noise \cite{ref216}. Another innovative approach focuses on the retrieval pipeline itself, employing counterfactual risk minimization to align retrieval with decision safety rather than mere semantic similarity. This framework, known as CoRM-RAG, simulates user biases during training and uses a lightweight Evidence Critic to identify documents with sufficient evidential strength, which significantly outperforms dense retrievers in adversarial settings \cite{ref217}.

The empirical evaluation of RAG robustness in high-stakes domains such as healthcare has provided actionable insights that tie these techniques back to the broader trustworthiness goals of this survey. Controlled studies varying the type and composition of retrieved documents (helpful, harmful, and adversarial) reveal that adversarial documents substantially degrade alignment, but robustness can be preserved when helpful evidence is also present in the retrieval pool \cite{ref218}. Extensive factorial studies on poisoning further show that robustness is not attributable to a single component but arises from the interaction of retrieval, generation, and knowledge-base configuration, with retriever architecture and retrieval depth being the strongest factors affecting poisoning exposure \cite{ref219}. In addition, the ``Relevance-Robustness Gap'' highlights a critical failure mode where maximizing semantic relevance paradoxically promotes the retrieval of sycophantic evidence that reinforces hallucinations, a challenge directly addressed by counterfactual risk minimization frameworks \cite{ref217}.

In summary, robustness to noisy and adversarial evidence in RAG requires a multi-faceted approach that complements the attribution and hallucination controls discussed earlier. It involves making the generator intrinsically more robust through fine-grained knowledge pruning and specialized fine-tuning \cite{ref185,ref210,ref211}, while simultaneously implementing explicit inference-time defenses that verify evidence consistency and judge source credibility \cite{ref212,ref211}. Furthermore, dynamic adaptation frameworks that iteratively consolidate internal and external knowledge based on reliability signals offer a pathway toward more trustworthy RAG systems that can gracefully handle the inevitable imperfections of the retrieval process \cite{ref215,ref216}. These robustness mechanisms, in turn, lay the groundwork for the next set of considerations in designing deployable RAG systems, where efficiency and scalability must be balanced against the reliability guarantees established here.

\section{Training, Fine-Tuning, and Optimization Paradigms}

\subsection{Training Paradigms for RAG Systems}

\begin{center}\rule{0.5\linewidth}{0.5pt}\end{center}

The effectiveness of a Retrieval-Augmented Generation (RAG) system hinges not only on its architectural design but also on the paradigms used to train and adapt its constituent modules. Early RAG systems often relied on off-the-shelf, frozen components. However, the research community has increasingly recognized that joint optimization and targeted training strategies are essential for unlocking the full potential of the retriever-generator pipeline. This subsection delves into the core training paradigms for RAG, exploring end-to-end frameworks, supervised fine-tuning, retrieval-augmented pretraining, and advanced data-centric strategies. These paradigms establish the foundational training objectives and data utilization principles upon which the subsequent fine-tuning techniques---such as parameter-efficient adaptation---are built.

A fundamental challenge in training RAG systems is the end-to-end optimization of the retriever and generator, which are often heterogeneous and coupled through discrete latent variables (i.e., the retrieved documents). A major hurdle is that marginalizing over all relevant passages from a knowledge base is computationally intractable. Traditional top-K marginalization and variational approaches like VRAG often suffer from biased or high-variance gradient estimates. To address these issues, joint stochastic approximation (JSA) has been proposed as a powerful alternative. This algorithm, a stochastic extension of the expectation-maximization (EM) method, estimates discrete latent variable models more effectively, leading to significant performance improvements over vanilla RAG and VRAG in both open-domain question answering and knowledge-grounded dialog tasks \cite{ref220}. Building on this, the Direct Retrieval-augmented Optimization (DRO) framework offers another perspective. DRO enables end-to-end training of a generative knowledge selection model and an LLM generator without relying on the simplifying assumption of document independence, which has been criticized for being unrealistic. By treating document permutation as a latent variable and alternating between estimation and maximization phases, DRO progressively improves both components. Theoretically analogous to policy-gradient methods in reinforcement learning, DRO has demonstrated substantial 5\%-15\% improvements in EM and F1 over strong baselines \cite{ref221}. Similarly, Differentiable Data Rewards (DDR) proposes an end-to-end training method that aligns data preferences between different RAG modules. Using a rollout method to sample potential responses as perturbations, DDR evaluates the impact of these perturbations on the overall system and optimizes each agent (retriever and generator) to produce outputs that improve the system's collective performance. This approach has proven particularly effective for smaller-parameter LLMs that rely more heavily on retrieved knowledge, enhancing their ability to extract key information and mitigate conflicts between parametric memory and external evidence \cite{ref20}.

Beyond end-to-end optimization, supervised fine-tuning (SFT) has emerged as a dominant strategy for adapting pretrained LLMs to specific RAG tasks and domains. A pioneering approach in this area is Retrieval Augmented FineTuning (RAFT), which provides a training recipe to improve a model's ability to answer questions in an ``open-book'' in-domain setting. RAFT trains the LLM to identify and ignore distractor documents that do not help answer the question, while learning to cite verbatim the correct sequence from the relevant document. This, combined with chain-of-thought-style responses, significantly improves the model's reasoning ability in domain-specific RAG settings, as demonstrated on datasets like PubMed, HotpotQA, and Gorilla \cite{ref222}. However, full fine-tuning of increasingly large models is resource-intensive. To address this, the CRAFT approach combines RAFT with parameter-efficient fine-tuning (PEFT) techniques, specifically Low-Rank Adaptation (LoRA). CRAFT reduces fine-tuning and storage requirements while achieving faster inference times and maintaining comparable RAG performance. This makes it particularly useful for knowledge-intensive QA tasks in resource-constrained environments where hardware and internet access are limited \cite{ref223}. Further extending this idea, the P-RAG framework proposes a hybrid architecture that integrates parametric knowledge within the LLM and retrieved evidence, guided by Chain-of-Thought prompting and LoRA fine-tuning, showcasing state-of-the-art results on biomedical and multi-hop QA datasets \cite{ref224}. These SFT-based methods demonstrate how the choice of fine-tuning strategy directly impacts the practical feasibility of deploying RAG systems, naturally leading to the deeper investigation of parameter-efficient and adaptive LoRA variants presented in the following subsection.

Another line of research moves beyond simply fine-tuning on static datasets by focusing on the pretraining phase itself. The ``To Memorize or to Retrieve'' study systematically investigates the trade-off between pretraining corpus size and retrieval store size under fixed data budgets. By training models ranging from 30M to 3B parameters, this work introduces a three-dimensional scaling framework that models performance as a function of model size, pretraining tokens, and retrieval corpus size. The findings reveal that retrieval consistently improves performance over parametric-only baselines, and the framework enables the estimation of optimal data budget allocation between pretraining and retrieval, providing quantitative guidance for designing scalable language modeling systems \cite{ref225}.

The quality and structure of training data are paramount to the success of these fine-tuning paradigms. A common issue is that the top-k retrieved documents in a training set can vary greatly in quality, with some being irrelevant or misleading. The DACL-RAG framework tackles this by introducing a multi-stage RAG training framework that combines a multi-level data augmentation strategy with a multi-stage curriculum learning paradigm. The data augmentation strategy constructs comprehensive and diverse training sets with controllable difficulty levels through sample evolution. The curriculum learning paradigm then organizes these examples into progressive training stages, ensuring stable and consistent improvements. This approach demonstrates consistent gains of 2\% to 4\% across four open-domain QA datasets \cite{ref226}. Similarly, systematic knowledge injection through diverse augmentation is the core idea behind another framework that enhances the fine-tuning process by varying the relevance of retrieved information (context augmentation) and by using multiple answers for the same question (knowledge paraphrasing). This teaches the model when to ignore and when to rely on retrieved content, mitigating the risk of memorization without true extraction and improving token-level recall by up to 10\% \cite{ref227}.

Finally, instead of treating RAG modules as separate entities to be optimized in a pipeline, a more holistic approach advocates for joint optimization. The SmartRAG pipeline exemplifies this by incorporating a policy network and a retriever that are jointly optimized using reinforcement learning. The policy network serves as a decision-maker for when to retrieve, a query rewriter, and an answer generator, all trained together with a reward designed to maximize performance while minimizing retrieval cost. When jointly optimized, all modules become aware of how other modules are working, enabling them to find the best way to cooperate as a complete system, achieving better performance than separately optimized counterparts \cite{ref228}. This shift toward coordinated, system-level optimization---whether through end-to-end differentiable rewards, RL-based policy learning, or sophisticated curriculum-based data strategies---is a defining trend in the evolution of RAG training paradigms. Together, these training paradigms and data-centric strategies provide the contextual foundation for understanding why module-specific and resource-efficient fine-tuning methods, such as the LoRA-based approaches discussed next, have become indispensable for advancing RAG systems in practice.

\begin{center}\rule{0.5\linewidth}{0.5pt}\end{center}

\subsection{Parameter-Efficient Fine-Tuning and Adaptation Methods}

Building upon the foundational training paradigms of end-to-end optimization and supervised fine-tuning, the practical deployment of RAG systems introduces a critical bottleneck: the computational cost and memory footprint of adapting increasingly large language models (LLMs). Full fine-tuning of models with billions of parameters is often computationally prohibitive and risks catastrophic forgetting of the parametric knowledge that is essential for generation. To address these constraints, parameter-efficient fine-tuning (PEFT) methods, particularly the Low-Rank Adaptation (LoRA) family, have become the de facto standard for adapting LLMs in RAG systems. This subsection investigates the landscape of LoRA-based techniques, focusing on their architectural innovations, resource-efficient variants, and strategic applications for enhancing RAG performance, thereby providing the practical toolkit that makes the end-to-end and data-centric training strategies discussed earlier viable in real-world deployments.

The foundational principle of LoRA is to constrain the weight update \(\Delta\)W to a low-rank subspace, parameterized by the product of two smaller matrices, A and B \cite{ref229}. While this approach drastically reduces the number of trainable parameters, its fixed-rank nature and linear structure present known limitations. For instance, a fixed rank may not be optimal for all layers of a transformer, as the representational complexity varies significantly across the network \cite{ref230}. Consequently, a substantial body of research has focused on developing dynamic rank allocation and more expressive adaptation mechanisms---innovations that are particularly salient for RAG, where the complexity of integrating retrieved evidence can vary widely by query.

Addressing the issue of static rank allocation, several methods have been proposed to enable dynamic adjustments during training. For example, ALoRA introduces a method to estimate the importance of each rank and then prunes less important ranks, reallocating the parameter budget to modules that require higher capacity \cite{ref231}. Similarly, SoRA employs a gate unit optimized with proximal gradient methods to dynamically control the cardinality of the rank, effectively allowing the adapter to expand and then contract to an optimal, task-specific rank \cite{ref232}. The Learnable Rank LoRA (LR-LoRA) method further questions the fixed-rank inductive bias by making the adapter rank itself a learnable parameter, allowing the optimizer to determine the appropriate rank for each layer automatically \cite{ref230}. These approaches highlight a shift toward more flexible and adaptive PEFT methods, which are crucial for RAG systems where the complexity of integrating retrieved evidence can vary widely by query.

Beyond adaptive rank, researchers have explored architectural modifications to enhance LoRA's expressiveness. One line of work focuses on decomposing the update matrix into finer components. For instance, DoRA decomposes high-rank layers into structured single-rank components to allow for more granular pruning \cite{ref233}. Another approach, GraLoRA, partitions weight matrices into sub-blocks, each with its own low-rank adapter, to overcome the structural bottleneck of standard LoRA and improve gradient propagation \cite{ref234}. The problem of parameter interference, where updates for one task degrade performance on another, is tackled by methods like FlyLoRA, which uses an implicit Mixture-of-Experts (MoE) architecture with frozen sparse random projections to decouple tasks \cite{ref235}. Complementing this, NeuroLoRA introduces a learnable neuromodulation gate that contextually rescales the projection space before expert selection, further enhancing multi-task and continual learning capabilities \cite{ref236}. For large-scale multi-task deployment, frameworks like MoRE align different ranks of a LoRA module with different tasks, using an adaptive rank selector to jointly train low-rank experts \cite{ref237}.

While architectural and algorithmic innovations are critical, the practical deployment of RAG systems in resource-constrained environments has driven the development of highly efficient LoRA variants, which directly complement the resource-efficient fine-tuning strategies (e.g., CRAFT, P-RAG) discussed in the context of supervised fine-tuning. Quantization-aware methods are particularly important. LowRA presents a framework for LoRA fine-tuning below 2 bits per parameter by optimizing fine-grained quantization mappings and precision assignment \cite{ref238}. RoLoRA tackles the challenge of weight-activation quantization by applying rotation to eliminate activation outliers, ensuring that the fine-tuned model remains amenable to low-bit quantization \cite{ref239}. Memory efficiency is another key concern. CARE-LoRA addresses the activation memory bottleneck by exploiting the projection structure of LoRA to reconstruct full activations from compressed counterparts, reducing the overall memory footprint during fine-tuning \cite{ref240}. Similarly, LoRA-FA freezes the projection-down weight and only updates the projection-up weight, a simple yet effective strategy that eliminates the need to store full-rank input activations for backpropagation \cite{ref241}. These memory- and quantization-efficient methods are essential for adapting RAG systems on edge devices or in federated learning settings where communication and memory are at a premium. In federated learning, for example, it has been shown that freezing the randomly initialized non-zero matrix and fine-tuning only the zero-initialized one (FFA-LoRA) can enhance stability when combined with differential privacy and reduce communication costs \cite{ref242}.

The performance of LoRA is also highly sensitive to initialization, scaling factors, and optimization strategies, which determines whether the adapted model can effectively leverage the retrieved knowledge in a RAG pipeline. Standard LoRA's slow convergence and suboptimal performance can be mitigated through careful initialization. LoRA-GA, for instance, introduces an initialization method that aligns the gradients of the low-rank product with those of full fine-tuning at the first optimization step, leading to significantly faster convergence and better final performance \cite{ref243}. SC-LoRA further refines this by constraining the output of the adapters in a subspace that balances new task learning with the preservation of general knowledge, mitigating the phenomenon of knowledge forgetting---a risk that is particularly acute in RAG systems where the model must balance parametric memory with external evidence \cite{ref244}. Another persistent issue is the learning dynamics associated with the rank. The rank-stabilized LoRA (rsLoRA) method provides a theoretical justification for using a scaling factor that divides adapters by the square root of the rank instead of the rank itself, which stabilizes training and allows for higher-rank adaptations to be effective \cite{ref245}. On the optimization side, Flat-LoRA focuses on finding adaptation weights that lie in a flat region of the full parameter space, using a Bayesian expectation loss objective to improve generalization \cite{ref246}. Furthermore, the convergence of LoRA can be accelerated by identifying an optimal low-rank factorization at each step, as proposed by RefLoRA, which promotes a flatter loss landscape and faster, more stable convergence \cite{ref247}.

In the context of RAG, the choice of which modules to adapt is a critical architectural decision that directly influences how effectively the model integrates retrieved information. The ``dominant adaptation module'' perspective, identified using a gradient-based sensitivity probe called PAGE, suggests that a single adapter placed on a specific shallow feed-forward network down-projection can yield surprisingly strong performance, sometimes outperforming vanilla LoRA with a fraction of the parameters \cite{ref248}. This insight challenges the common practice of uniformly distributing adapters across all attention modules. Methods like PLoP provide a lightweight, automatic way to identify the most effective module types for adaptation for a given task \cite{ref249}. These findings are particularly relevant for RAG, where the interaction between the retriever and generator might benefit from targeted adaptations in specific layers rather than broad, unfocused updates. The success of these placement strategies suggests that a hierarchical configuration of adapter experts, where both the number and rank of experts are adjusted per layer, may offer a more efficient path to fine-tuning LLMs for RAG tasks \cite{ref250}.

Finally, for improving reasoning capabilities in RAG, LoRA can be strategically integrated with Chain-of-Thought (CoT) prompting, thereby connecting back to the supervised fine-tuning strategies like RAFT and P-RAG that emphasized reasoning in domain-specific settings. The ability to fine-tune models on CoT examples using LoRA offers a computationally efficient way to enhance multi-step reasoning without full fine-tuning. Furthermore, the modularity of LoRA enables layer-wise model merging, allowing for the combination of adapters trained for different tasks (e.g., retrieval-oriented generation and reasoning) into a single model, as explored in frameworks for continual learning and multi-task model merging \cite{ref251,ref252}. This compositional property of LoRA makes it a powerful tool for building complex RAG systems, where different components of the pipeline might require specialized knowledge or capabilities. More importantly, these parameter-efficient adaptations define \emph{how} the model parameters are updated, but the \emph{objectives} that guide these updates---such as preference alignment and reinforcement learning signals---are the focus of the next subsection. As we will see, the synergy between efficient adaptation mechanisms and sophisticated alignment objectives is what ultimately enables the development of RAG systems that are not only deployable but also precisely tuned to user preferences and task requirements.

\subsection{Preference Optimization and Reinforcement Learning for RAG Alignment}

The alignment of retrieval-augmented generation (RAG) systems with downstream task objectives and user preferences represents a critical frontier in optimizing the synergy between retrieval and generation, extending well beyond the parameter-efficient fine-tuning strategies discussed previously. While LoRA-based approaches provide efficient mechanisms for adapting LLMs within RAG pipelines, they primarily focus on the parametric side of the equation. This subsection delves into the advanced alignment techniques that explicitly optimize the \emph{interplay} between retrieval and generation, focusing on two principal paradigms: Direct Preference Optimization (DPO) variants and reinforcement learning (RL) methods, alongside the reward modeling innovations that underpin them. In doing so, it bridges the architectural concerns of the preceding section with the broader system-level coordination challenges addressed in the next.

\textbf{Direct Preference Optimization Variants for RAG Alignment}

The standard DPO framework, while effective for general instruction following \cite{ref253,ref254}, has been adapted and extended to address the unique challenges of RAG, where the generator must not only be helpful but also correctly utilize and cite retrieved evidence. A primary concern is aligning the generator with retrieval relevance. For instance, a multi-perspective preference alignment approach, PA-RAG \cite{ref255}, constructs high-quality instruction fine-tuning data and multi-perspective preference data by sampling varied quality responses from the generator across different document quality scenarios. It then optimizes the generator using SFT and DPO, significantly enhancing response informativeness, robustness, and citation quality.

However, standard DPO and its derivatives are not without their limitations. A fundamental issue is overfitting to the offline preference data, which can lead to likelihood displacement and degeneracy \cite{ref256,ref257}. To mitigate these problems, a robust alignment method named \textbf{RPO} was introduced to adaptively leverage multi-source knowledge based on retrieval relevance. Unlike other methods, RPO integrates an implicit representation of retrieval relevance directly into the reward model, thereby unifying retrieval evaluation and response generation within a single model and eliminating the need for a separate procedure to assess retrieval quality \cite{ref258}. This highlights a key trend: moving from purely quality-based preference optimization to relevance-aware alignment.

Other DPO variants have been developed to address specific weaknesses of the base algorithm. For example, \textbf{BPO} tackles the ``Degraded Chosen Responses'' issue, where DPO can inadvertently decrease the likelihood of chosen responses, by dynamically balancing the optimization of chosen and rejected responses through a balanced reward margin and gap adaptor, requiring only a single line of code modification \cite{ref259}. Similarly, research like \textbf{MPO} generalizes DPO to set-level contrasts, optimizing over entire sets of responses by extending the Bradley-Terry model to group-wise comparisons, which reduces alignment bias with respect to the number of responses per query \cite{ref260}.

To further enhance DPO, some works propose incorporating additional mechanisms. Integrating a temporal decay factor, for instance, has been shown to mitigate length bias by prioritizing earlier tokens that are more critical for alignment \cite{ref261}. In reasoning-heavy or multi-step retrieval contexts, hierarchical preference optimization (HiPO) decomposes responses into reasoning segments (query clarification, reasoning steps, and answer) and computes loss as a weighted sum of DPO loss for each segment, enabling segment-specific training while maintaining efficiency \cite{ref262}. Furthermore, to overcome the weak learning signal in standard DPO where chosen and rejected responses are similar, \textbf{RPO} also proposed hint-guided reflection, using external models to identify hallucination sources and generate reflective hints, thereby creating on-policy preference pairs with stronger contrast and clearer preference signals \cite{ref263}.

\textbf{Reinforcement Learning for RAG and Reasoning Optimization}

While DPO variants offer computationally efficient alignment, reinforcement learning methods, particularly those based on Group Relative Policy Optimization (GRPO), have become the de facto standard for optimizing complex reasoning and retrieval policies in RAG systems. GRPO's critic-free, group-relative advantage estimation has proven effective for training large language models (LLMs) on agentic tasks \cite{ref264,ref265}. However, this approach is not without challenges, especially concerning credit assignment. A single final reward often conflates planning errors with execution errors in multi-hop retrieval, a phenomenon termed ``hierarchical credit entanglement.'' To refine this, frameworks like \textbf{Search-P1} propose path-centric reward shaping, which evaluates the structural quality of reasoning trajectories through order-agnostic step coverage and soft scoring, thereby extracting learning signals even from failed samples \cite{ref266}.

The design of the reward signal is paramount in RL-based RAG. Many systems now move beyond sparse, outcome-only rewards to process-level or multi-faceted reward functions. For instance, EVO-RAG couples a seven-factor, step-level reward vector covering relevance, redundancy, efficiency, and answer correctness within a curriculum-guided reinforcement learning framework, allowing the agent to learn when to search, backtrack, or answer \cite{ref267}. Similarly, DyKnow-RAG, designed for e-commerce, uses a dual-group rollout strategy (parametric-only vs.~with-context) and a posterior-driven inter-group advantage scaling mechanism to teach LLMs to learn to trust external knowledge, effectively optimizing context utilization without human process labels \cite{ref268}.

To address the inherent instability and variance in GRPO, several algorithmic innovations have been proposed. Some methods focus on improving credit assignment at the token level. \textbf{EP-GRPO}, for example, uses entropy-gated modulation to prioritize high-entropy decision pivots and implicit process signals from policy divergence to provide directional token-level feedback without external reward models \cite{ref269}. Another approach, \textbf{SRPO}, routes correct samples to GRPO's reward-aligned reinforcement and failed samples to Self-Distillation Policy Optimization (SDPO) for targeted logit-level correction, combining the rapid early improvement of the latter with the long-horizon stability of the former \cite{ref270}. Meanwhile, frameworks like \textbf{AMIR-GRPO} augment GRPO with an implicit DPO-style contrastive regularizer constructed from intra-group reward rankings, which amplifies the suppression of low-reward trajectories and attenuates response-level length bias \cite{ref271}.

A further challenge in RL fine-tuning is the multi-reward setting. As models are expected to align with diverse human preferences, RL pipelines often incorporate multiple reward models, each capturing a distinct preference. This creates conflicts where a single rollout can have positive advantages on one reward dimension but negative on another. To resolve this, \textbf{GDPO} (Group reward-Decoupled normalization Policy Optimization) was introduced, which decouples the normalization of individual rewards, thereby preserving their relative differences and enabling more accurate multi-reward optimization \cite{ref272,ref273}. Its successor, \textbf{GD\(^2\)PO}, further employs a conflict-aware filtering mechanism to mask out rollouts suffering from severe reward-wise disagreement, preventing conflicting signals from canceling each other out \cite{ref273}.

Finally, the stability of GRPO is a known concern. Curriculum learning has emerged as a powerful tool to stabilize training. \textbf{CuSearch} introduces a curriculum rollout sampling framework built on Search-Depth Greedy Allocation (SDGA), which reallocates a fixed update budget toward deeper-search trajectories, providing denser supervision for the retrieval sub-policy and yielding significant performance gains \cite{ref274}. Similarly, the ACQO framework employs a Curriculum Reinforcement Learning (CRL) approach to stabilize the training of an adaptive query reformulation module by progressively introducing more challenging queries \cite{ref131}. These examples underscore the importance of carefully designed reward structures, advanced credit assignment, and curriculum-guided training strategies to fully harness the potential of RL for developing robust, reliable, and efficient agentic RAG systems.

Collectively, these DPO and RL-based alignment strategies represent a shift from static, one-time fine-tuning toward dynamic, feedback-driven optimization. They directly complement the parameter-efficient adaptation methods discussed earlier by providing the \emph{learning objectives} and \emph{reward signals} that guide \emph{what} should be learned during fine-tuning. At the same time, they lay the groundwork for the next level of system-level coordination, where the focus expands beyond a single generator to jointly align the retriever, the generator, and their interface, ensuring that the non-parametric knowledge retrieved is not just accessible but is actively and synergistically integrated to maximize the generative and reasoning capabilities of the LLM.

\subsection{Retriever-Generator Alignment and Cooperative Optimization}

The effectiveness of Retrieval-Augmented Generation (RAG) is fundamentally constrained by the alignment---or misalignment---between its two primary components: the retriever, which surfaces potentially relevant documents, and the generator, which consumes this context to produce an answer. While the DPO and RL-based approaches discussed previously focus on optimizing the generator's behavior given a fixed retrieval pipeline, a core challenge in RAG system design is that retrievers and generators are often developed and optimized independently, leading to a disconnect in which the retriever presents information that is semantically relevant but not optimally useful for the specific LLM's reasoning and generation process. This subsection surveys the array of techniques developed to bridge this gap, spanning from bridge models and fine-grained guidance to advanced preference optimization, reward-based alignment, and cooperative training paradigms. In doing so, it moves from the generator-centric alignment of the preceding section toward the holistic, system-level coordination that will be examined in the following discussion.

\textbf{Bridge Models and Utility-Aware Retrieval.} One foundational approach to aligning retrievers and generators is the introduction of a bridge mechanism trained to connect the two components. Rather than simply passing retrieved documents to the LLM, a bridge model is explicitly optimized to transform or select information with the downstream generation objective in mind. For instance, prior work has validated that both the ranking and selection assumptions of retrievers are often misaligned with what an LLM finds useful, and proposed a framework that chains supervised learning with reinforcement learning (RL) to train a bridge model that optimizes the connection between the retriever and the LLM, improving the utility of the context for the generator \cite{ref275}. This approach moves beyond treating the retriever and generator as independent black boxes, recognizing that the interface between them is a critical point for optimization.

More broadly, this insight has led to a shift from optimizing retrieval based on static relevance labels to optimizing it based on downstream generation utility. The R3 framework explicitly adopts an RL training paradigm that enables a retriever to explore and self-improve within a given RAG environment, learning to retrieve context that is directly beneficial for the target task, rather than relying solely on human or synthetic relevance labels that may not reflect the generator's needs \cite{ref276}. This utility-focused paradigm is further exemplified by FiGRet (Fine-grained Guidance for Retrievers), which is inspired by pedagogical theories like Guided Discovery Learning. Instead of relying on abstract preference signals that are hard for dense retrievers to learn from, FiGRet leverages LLMs' reasoning capability to construct fine-grained, information-centric examples from cases where the retriever performs poorly. These examples focus on learning objectives for the retriever that are highly relevant to RAG, specifically \texttt{relevance}, \texttt{comprehensiveness}, and \texttt{purity}, thereby scaffolding the retriever's learning process to better align with the LLM's needs \cite{ref277}.

\textbf{Retrieval Preference Optimization and Pairwise Alignment.} To explicitly train a generator to be a better consumer of retrieved information, a range of preference optimization techniques have been introduced. A prominent example is Retrieval Preference Optimization (RPO), a lightweight alignment method designed to adaptively leverage multi-source knowledge based on retrieval relevance. RPO integrates an implicit representation of retrieval relevance into the reward model, allowing the system to combine retrieval evaluation and response generation into a single model and thereby training the LLM to handle knowledge conflicts between its parametric memory and the retrieved non-parametric knowledge \cite{ref258}. By becoming aware of retrieval relevance during training, the generator learns when to trust the retrieved context and when to rely on its internal knowledge. Notably, this approach shares conceptual ground with the relevance-aware DPO variants discussed in the previous subsection, yet it operates at the interface of retrieval and generation rather than purely on response quality, illustrating the continuum of alignment strategies across the RAG pipeline.

Complementary approaches focus on aligning the entire RAG pipeline with user preferences at multiple stages. DPA-RAG, for instance, proposes a dual preference alignment framework to address the difficulty of aligning a retriever with the diverse knowledge preferences of different LLMs. It performs external preference alignment by jointly integrating pair-wise, point-wise, and contrastive signals into a reranker, ensuring the retrieval component aligns with downstream generation. It also performs internal preference alignment by introducing a pre-aligned stage before vanilla Supervised Fine-tuning (SFT), which allows the LLM to implicitly capture knowledge aligned with its own reasoning preferences \cite{ref278}. This two-pronged approach connects both the quality of what is retrieved and how the generator internalizes it.

\textbf{Differentiable Rewards and End-to-End Cooperative Optimization.} A major frontier in RAG alignment is the end-to-end optimization of the entire system, where retriever and generator are trained jointly towards a unified objective. The Differentiable Data Rewards (DDR) method tackles this by aligning the data preferences between different RAG modules. DDR end-to-end trains the system by collecting rewards via a rollout method where agents sample potential responses as perturbations. By evaluating the impact of these perturbations on the performance of the whole RAG system, the agents are optimized to produce outputs that improve the overall system, fostering a cooperative dynamic and enabling the generation module to better extract key information \cite{ref20}. This approach moves beyond sequential optimization to a more holistic, system-level consideration of mutual benefit.

Another powerful approach to cooperative optimization is the use of multi-agent reinforcement learning (MARL), particularly when RAG pipelines are composed of multiple modules such as query rewriting, document retrieval, filtering, and generation. Treating each component of the RAG pipeline as an RL agent, MMOA-RAG harmonizes all agents' goals toward a unified reward---such as the F1 score of the final answer---using multi-agent reinforcement learning. This addresses the misalignment that occurs when components are optimized separately, imposing an over-arching goal on each sub-module to cooperate in the service of the final answer \cite{ref279}. A similar collaborative philosophy is evident in frameworks that introduce a proxy-centric multi-agent system as an intermediary. C-3PO, for example, implements three specialized agents that assess the need for retrieval, generate effective queries, and select information suitable for the LLM. This cooperative design allows for the optimization of the entire pipeline without altering the core retriever or LLM, utilizing a tree-structured rollout for effective reward credit assignment \cite{ref280}. This plug-and-play nature is critical for maintaining the broad applicability of the alignment technique across different base models.

\textbf{Training-Free, Self-Practicing, and Consistency-Driven Approaches.} The field has also seen a rise in alignment techniques that avoid the high cost of fine-tuning or that leverage the model's own capabilities. RAG-Pref, an online and training-free algorithm, demonstrates the power of conditioning a model on both preferred and dispreferred samples during inference. This contrastive information, retrieved from a small set of pre-aligned examples, shifts the model's sampling distribution toward the high-reward region and can significantly boost its refusal guardrails against harmful or agentic attacks, often outperforming other online alignment methods that require retraining \cite{ref281}. In contrast, SIM-RAG introduces a framework that improves multi-round retrieval capabilities through self-practicing. Rather than relying on expensive human-labeled process supervision, it has the RAG system generate synthetic training data by practicing multi-round retrieval. A lightweight information sufficiency Critic is then trained on this data (including both successful and unsuccessful trajectories) to evaluate whether the system has retrieved sufficient information at each round, guiding retrieval decisions at inference time. This data-efficient and system-efficient framework enhances the system's self-awareness without needing to modify the existing LLM or search engine \cite{ref282}.

The pursuit of alignment for accurate content has also been extended to ensure consistency. In high-stakes settings, RAG systems must produce outputs that are stable across semantically equivalent queries, yet variability in both the retriever and the generator often leads to inconsistent answers. To address this, Con-RAG uses a reinforcement learning approach, specifically Paraphrased Set Group Relative Policy Optimization (PS-GRPO), to train the generator to produce consistent outputs across paraphrased query sets. By assigning group similarity rewards, the generator is trained to maintain core information content in the face of retrieval-induced variability, thereby improving both reliability and trustworthiness of the overall system \cite{ref283}. This focus on distributional robustness across retrieval perturbations complements the reward-shaping and credit-assignment innovations highlighted in the previous subsection, extending them from single-query optimization to multi-query consistency.

In summary, aligning the retriever and generator is no longer a simplistic, one-time pre-processing step but a multifaceted optimization challenge that builds directly upon the generator-centric alignment techniques discussed earlier. Methods have evolved from training explicit bridge models and fine-tuning retrievers with fine-grained, pedagogically inspired guidance, to sophisticated end-to-end training. This modern era is characterized by differentiable reward pathways \cite{ref20}, multi-agent cooperative learning \cite{ref279}, and utility-aware exploration \cite{ref276}. In parallel, there is a growing repertoire of training-free and self-practicing methods that offer cost-effective alternatives \cite{ref281,ref282}, alongside a newer focus on system-level properties like consistency \cite{ref283}. Collectively, these strategies ensure that the non-parametric knowledge retrieved is not just accessible, but is actively and synergistically integrated to maximize the generative and reasoning capabilities of the LLM, setting the stage for the broader system-level coordination challenges that follow.

\section{Advanced Architectures: GraphRAG, Agentic RAG, and Multi-Agent Systems}

\subsection{Graph-Based RAG: Knowledge Graphs, Construction Paradigms, and Topological Retrieval}

Graph-based Retrieval-Augmented Generation (GraphRAG) has emerged as a transformative paradigm that extends the foundational retrieve-then-read pipeline by organizing external corpora into structured knowledge graphs. Unlike flat, chunk-level retrieval, graph-based approaches exploit explicit relational topologies to support multi-hop reasoning, effectively bridging fragmented evidence that standard dense retrieval fails to connect. This subsection examines the foundational techniques underpinning GraphRAG, focusing on three interlocking dimensions: graph construction paradigms, graph representation choices, and topological retrieval strategies. We analyze the trade-offs among granularity, construction cost, scalability, and robustness, providing a systematic overview of the design space for modern graph-augmented retrieval systems.

The construction of knowledge graphs is the first critical decision in any GraphRAG pipeline, and current methods span a spectrum from lightweight statistical approaches to semantically rich LLM-based extraction. At the lightweight extreme, TIGRAG (Token-Induced GraphRAG) builds a token co-occurrence knowledge graph using sliding-window statistics, directly modeling topological relationships between tokens without any LLM involvement, thereby enabling scalable and cost-effective construction \cite{ref284}. Similarly, LinearRAG constructs a relation-free hierarchical graph, termed Tri-Graph, using only lightweight entity extraction and semantic linking, avoiding the unstable relation modeling that plagues LLM-based pipelines and scaling linearly with corpus size at zero extra token consumption \cite{ref285}. These statistical construction methods are particularly appealing for large-scale or cost-sensitive deployments, as they avoid the hallucination and error propagation risks inherent in generative extraction. In contrast, many advanced systems rely on LLM-extracted triples and entity-centric pipelines to capture richer semantics. However, this approach introduces significant indexing costs and potential hallucinations. To mitigate these issues, AGRAG substitutes LLM-based entity extraction with a statistics-based method, addressing the problem of inaccurate graph construction caused by LLM hallucination while still enabling structured reasoning paths \cite{ref286}. Clue-RAG further optimizes this trade-off by introducing a multi-partite graph index that incorporates chunks, knowledge units, and entities at multiple granularities, coupled with a hybrid extraction strategy that reduces LLM token usage while producing accurate and disambiguated knowledge units \cite{ref287}. KET-RAG offers a middle ground by identifying a small set of key text chunks for LLM-based skeleton construction and building a lightweight text-keyword bipartite graph for the remaining corpus, achieving comparable retrieval quality to MS-GraphRAG while reducing indexing costs by over an order of magnitude \cite{ref288}. The choice of construction paradigm fundamentally shapes the trade-off between graph quality and cost, with token co-occurrence and entity-centric statistical methods providing scalable defaults, while LLM-driven extraction offers deeper semantics at a premium.

Beyond how graphs are built, the choice of graph representation significantly influences retrieval effectiveness. Most foundational approaches rely on heterogeneous graphs with entity and chunk nodes, as exemplified by HippoRAG-style architectures and their successors. For instance, the unified framework in \cite{ref289} constructs hybrid graphs with both chunk and entity nodes, then iteratively clusters them and generates LLM-based summaries to enable context-aware and relation-aware retrieval across abstraction levels. However, pairwise entity-relation graphs are limited to low-order associations, which has motivated the development of hypergraph-based representations that connect multiple entities simultaneously. HyperSU formulates hyperedge construction as an entity-aware minimum-description-length optimization problem, inducing source-grounded semantic-unit hyperedges that balance sentence-level semantic coherence and entity compactness \cite{ref290}. Building on this, Cog-RAG introduces a theme-aligned dual-hypergraph architecture, using a theme hypergraph to capture inter-chunk thematic structure and an entity hypergraph to model high-order semantic relations, drawing inspiration from the top-down cognitive process of human reasoning \cite{ref291}. EHRAG further advances this direction by constructing a hypergraph that captures both structural and semantic relationships, employing a hybrid structural-semantic retrieval mechanism that outperforms state-of-the-art baselines while maintaining linear indexing complexity \cite{ref292}. HyperGraphPro goes a step further by integrating a structure-aware hypergraph retrieval mechanism that jointly considers semantic relevance and graph connectivity, promoting coherent traversal along multi-hop reasoning paths \cite{ref293}. An alternative representation paradigm shifts away from entity-centric graphs entirely, modeling knowledge as proposition-level structures. ToPG (Traversal over Proposition Graphs) models the knowledge base as a heterogeneous graph of propositions, entities, and passages, effectively combining the granular fact density of propositions with graph connectivity \cite{ref294}. This proposition-level granularity offers a finer-grained alternative to entity-based graphs, improving performance on fact-oriented single-hop queries while retaining multi-hop capability. Topology-enhanced approaches have also been explored, such as TopoRAG, which lifts textual graphs into cellular complexes to model multi-dimensional topological structures including cycles, capturing higher-dimensional relational dependencies that standard graphs overlook \cite{ref295}. The representation choice involves a fundamental tension: entity-centric graphs offer interpretability and multi-hop connectivity but may lose context, while proposition-level graphs preserve factual granularity at the cost of increased complexity.

The retrieval stage over these structures has evolved from static, query-blind traversal to dynamic, query-aware strategies. Personalized PageRank (PPR) remains a dominant diffusion mechanism, as popularized by HippoRAG and refined in subsequent works. Aethel combines bipartite PPR with coreference-aware teleportation to propagate relevance through entity-passage graphs for multi-hop financial diligence \cite{ref296}. QAFD-RAG introduces query-aware traversal by dynamically weighting edges based on semantic alignment between endpoints and the query embedding, providing the first statistical guarantees for query-aware graph retrieval and recovering relevant subgraphs with high probability \cite{ref44}. Building on this, CatRAG transforms static knowledge graphs into query-adaptive navigation structures through symbolic anchoring, query-aware dynamic edge weighting, and key-fact passage weight enhancement, mitigating the ``static graph fallacy'' where fixed transition probabilities cause semantic drift \cite{ref297}. Spreading activation methods offer an efficient alternative to flow diffusion: \cite{ref298} proposes a spreading-activation approach with a single per-step semantic gate based on cosine similarity, executing the entire retrieval procedure as a single Cypher query in one round-trip to Neo4j. Similarly, SemFlowRAG reconstructs the flat retrieval space into a corpus-adaptive semantic gradient graph, transforming static undirected edges into directed semantic constraints and using an abstractness-guided directed PageRank algorithm to force retrieval along a ``high-to-low semantic abstractness'' gradient \cite{ref299}. DA-RAG leverages attributed community search to extract relevant subgraphs dynamically based on the query, capturing high-order graph structures for self-complementary knowledge retrieval \cite{ref39}. For query routing, EA-GraphRAG introduces a syntax-aware complexity scorer that routes queries to either dense RAG or graph-based retrieval, achieving state-of-the-art performance in mixed scenarios while reducing latency \cite{ref300}. Dynamic and iterative strategies further enhance retrieval quality: Clue-RAG's Q-Iter mechanism performs query-driven iterative retrieval through semantic search and constrained graph traversal \cite{ref287}, while GRAIL frames graph retrieval as an interactive learning problem, using LLM-guided random exploration and path filtering to synthesize training data and learn a policy that dynamically balances retrieval breadth and precision \cite{ref301}. GraphSearch extends search-augmented reasoning to graph learning, combining a graph-aware query planner with a topology-aware retriever for zero-shot node classification and link prediction \cite{ref302}.

The foundational techniques discussed herein establish the structural and algorithmic basis upon which more sophisticated retrieval paradigms are built. Critically, while static graph construction and query-adaptive retrieval strategies---such as PPR, spreading activation, and learning-based policies---provide powerful mechanisms for multi-hop evidence integration, they remain largely reactive: retrieval is triggered by a single query and executed in a predetermined manner. This limitation becomes particularly pronounced when facing complex, multi-hop queries that require iterative refinement, where initial retrievals may be incomplete or ambiguous. As the following subsection will elaborate, agentic RAG addresses this gap by embedding dynamic, multi-round decision-making into the inference process. Rather than treating retrieval as a fixed pre-processing step with static graph traversal, agentic systems empower LLMs to autonomously decide when to retrieve, what to search for, how to reformulate queries based on intermediate reasoning, and when to terminate the loop. The graph-based structures and retrieval strategies surveyed here---ranging from PPR diffusion to policy-learned navigation---serve as foundational action primitives that agentic controllers can invoke selectively and adaptively, transforming GraphRAG from a one-shot pipeline into a sequential decision-making problem. Thus, the design space of graph construction, representation, and static retrieval examined in this subsection directly informs and constrains the agentic architectures that follow, highlighting a continuum from fixed topological retrieval to fully adaptive, reasoning-driven search.

In summary, the foundational techniques in graph-based RAG span a rich design space defined by construction paradigms, representation choices, and retrieval strategies. Lightweight statistical construction methods such as token co-occurrence graphs \cite{ref284} and hierarchical semantic graphs \cite{ref285} enable scalable deployment, while LLM-driven extraction offers richer semantics at higher cost \cite{ref288}. Representation choices range from heterogeneous graphs to hypergraphs \cite{ref303} and proposition-level structures \cite{ref294}, each balancing granularity and connectivity. Finally, retrieval has evolved from static PPR toward query-aware dynamic weighting \cite{ref304}, spreading activation \cite{ref298}, and learning-based interactive policies \cite{ref301}. The inherent trade-offs among indexing cost, granularity, scalability, and robustness to noisy graphs remain central to the design of effective GraphRAG systems, setting the stage for the agentic and multi-agent architectures discussed in subsequent subsections.

\subsection{Agentic and Dynamic Retrieval: Adaptive Planning, Iterative Reasoning, and Self-Reflection}

Agentic retrieval-augmented generation represents a fundamental departure from static, single-pass retrieval pipelines by embedding dynamic, multi-round decision-making directly into the inference process. Rather than treating retrieval as a fixed pre-processing step, agentic RAG systems empower large language models (LLMs) to autonomously decide when to retrieve, what to search for, how to reformulate queries, and when to terminate the retrieval-reasoning loop, effectively transforming RAG into a sequential decision-making problem \cite{ref25}. This paradigm shift is driven by the recognition that complex, multi-hop queries often require iterative refinement: initial retrievals may be incomplete, ambiguous, or noisy, and subsequent reasoning steps must be able to adaptively acquire additional evidence, correct earlier mistakes, and compress context to maintain coherence \cite{ref305,ref170}. Importantly, this agentic capability directly addresses the central limitation identified in the preceding subsection on graph-based retrieval: while graph structures such as knowledge graphs and hypergraphs provide powerful mechanisms for multi-hop evidence integration, static graph traversal remains fundamentally reactive---triggered by a single query and executed in a predetermined manner. Agentic systems overcome this constraint by embedding dynamic, multi-round decision-making into the inference process, transforming GraphRAG from a one-shot pipeline into a sequential decision-making problem where graph-based retrieval strategies serve as action primitives that agents can invoke selectively and adaptively.

A central theoretical foundation for agentic RAG is the formalization of retrieval as a Markov Decision Process (MDP), which provides a principled framework for modeling the sequential nature of retrieval decisions and for applying reinforcement learning (RL) to optimize agent behavior. For instance, Graph-R1 models retrieval as a multi-turn agent-environment interaction within a graph-based knowledge structure, optimizing the entire agentic process through end-to-end reinforcement learning, thereby outperforming traditional GraphRAG and RL-enhanced RAG methods in reasoning accuracy and retrieval efficiency \cite{ref306}. Similarly, EvoGraph-R1 extends this formulation to multimodal settings by reconceptualizing knowledge hypergraphs as dynamic environments shaped through agent interactions: the agent observes the graph state and executes actions to query, expand via web search, refine the graph structure, or terminate with an answer \cite{ref42}. This closed-loop formulation enables the hypergraph to evolve by integrating new evidence and correcting errors, supporting multi-hop reasoning across modalities \cite{ref42}. The SoK systematization further formalizes agentic retrieval-generation loops as finite-horizon partially observable MDPs, explicitly modeling control policies and state transitions to unify the fragmented landscape of agentic architectures \cite{ref25}. This MDP-based formulation also provides the theoretical bridge to the multi-agent architectures discussed in the following subsection: just as a single agent makes sequential decisions about retrieval, multi-agent systems distribute these decisions across specialized agents that must coordinate their actions, extending the sequential decision-making framework into a distributed decision-making paradigm.

Beyond MDP formalization, adaptive controllers have emerged as a key mechanism for dynamically modulating retrieval effort based on query complexity and evidence sufficiency. A2RAG couples an adaptive controller that verifies evidence sufficiency and triggers targeted refinement only when necessary with an agentic retriever that progressively escalates retrieval effort, achieving substantial improvements in recall while cutting token consumption and end-to-end latency by approximately 50\% relative to iterative multi-hop baselines \cite{ref307}. EA-GraphRAG similarly introduces a syntax-aware complexity scorer that routes queries to either dense RAG or graph-based retrieval based on a continuous complexity score, demonstrating that rigid application of GraphRAG to all queries---regardless of complexity---is the root cause of performance degradation and prohibitive latency \cite{ref300}. Stop-RAG casts iterative RAG as a finite-horizon MDP and introduces a value-based controller trained with full-width forward-view Q(\(\lambda\)) targets to adaptively decide when to stop retrieving, consistently outperforming fixed-iteration baselines and prompting-based stopping strategies \cite{ref308}. These adaptive approaches collectively underscore the importance of cost-aware orchestration: agentic enhancements are not universally beneficial and must be applied selectively according to query and domain characteristics \cite{ref31}. This cost-aware design principle, introduced here at the single-agent level, becomes a central organizing theme in the following subsection, where multi-agent orchestration must balance accuracy gains against computational overhead across distributed agents.

Iterative retrieval methods with self-reflection constitute another major strand of agentic RAG research. These systems interleave retrieval and reasoning over multiple rounds, using feedback from intermediate reasoning steps to guide subsequent searches. Self-RAG and IRCoT represent early instantiations of this idea, where the model reflects on whether retrieved information is sufficient and refines its query accordingly \cite{ref170}. DoctorRAG introduces a diagnose-and-repair framework that corrects failures via explicit error localization and prefix reuse: a distilled diagnosis model jointly assesses evidence sufficiency, classifies failure types, and localizes the earliest failure point, enabling targeted, efficient correction of known failed trajectories rather than blind reruns \cite{ref309}. MA-RAG facilitates test-time scaling for complex medical reasoning by iteratively evolving both external evidence and internal reasoning history within an agentic refinement loop, using semantic conflict among candidate responses as a proactive signal for multi-round retrieval \cite{ref310}. SIM-RAG enhances self-awareness by training a lightweight information sufficiency critic through self-practicing multi-round retrieval, where the system explores multiple retrieval paths labeled as successful or unsuccessful, and the critic guides retrieval decisions at inference time without requiring modifications to existing LLMs or search engines \cite{ref282}. These self-reflection mechanisms, while developed primarily for single-agent systems, naturally extend to multi-agent contexts where distinct agents can specialize in retrieval, reflection, and revision---a pattern that the following subsection examines in depth through hierarchical and consensus-driven multi-agent frameworks.

The integration of self-reflection with context compression and token-efficiency mechanisms is critical for practical deployment. TeaRAG addresses the substantial token overhead of agentic RAG by compressing both retrieved content and reasoning steps: it augments chunk-based semantic retrieval with graph retrieval using concise triplets, leverages Personalized PageRank to highlight key knowledge, and introduces Iterative Process-aware Direct Preference Optimization to penalize excessive reasoning steps, achieving significant token reductions while improving accuracy \cite{ref311}. Similarly, GRASP provides agents with semantic search, keyword search, and paragraph-reading actions, enabling them to retrieve sentence-level evidence and expand context only when needed, with the learned policy developing interpretable skimming and scanning behaviors \cite{ref312}. Plan*RAG isolates the reasoning plan as a directed acyclic graph outside the language model's working memory, enabling systematic exploration of reasoning paths, atomic subqueries, and parallel execution while maintaining bounded context window utilization \cite{ref313}. These token-efficiency mechanisms reflect a broader concern that becomes even more pronounced in multi-agent settings, where the coordination overhead and inter-agent communication costs can rapidly accumulate, motivating the efficiency-oriented frameworks examined in the subsequent subsection.

The training of agentic RAG systems has also seen significant advances, particularly through reinforcement learning with process-level supervision. TreePS-RAG models agentic RAG reasoning as a rollout tree where each reasoning step maps to a node, enabling step-wise credit assignment via Monte Carlo estimation over descendant outcomes without requiring intermediate labels \cite{ref314}. DecEx-RAG models RAG as an MDP incorporating decision-making and execution, introducing an efficient pruning strategy to optimize data expansion and achieving significant performance improvements through process-level policy optimization \cite{ref315}. HiPRAG introduces hierarchical process rewards that evaluate the necessity of each search decision on-the-fly, reducing over-search rates to just 2.3\% while concurrently lowering under-search rates \cite{ref316}. Search-P1 introduces path-centric reward shaping that evaluates the structural quality of reasoning trajectories through order-agnostic step coverage and soft scoring, extracting learning signals even from failed samples \cite{ref266}. A comparative study of process versus outcome rewards finds that process-level rewards improve training stability, reduce computational costs, and enhance efficiency, requiring substantially fewer training instances than outcome-based approaches \cite{ref317}. These training methodologies, developed for single-agent policies, are directly adapted in multi-agent systems where credit assignment must be distributed across cooperating agents---a challenge that frameworks like HERA address through role-specific credit assignment and experience accumulation, as discussed in the following subsection.

Finally, diagnostic tools and benchmarks have been developed to enable hop-level failure analysis in agentic RAG systems. AgenticRAGTracer provides a benchmark constructed automatically by LLMs that supports step-by-step validation, revealing that failures are primarily driven by distorted reasoning chains---either collapsing prematurely or wandering into over-extension \cite{ref318}. Query performance prediction (QPP) has also been applied to agentic RAG, showing that QPP estimates of generated queries are positively correlated with final answer quality, suggesting a path toward adaptive retrieval where QPP informs the model whether retrieved results are likely to be useful \cite{ref319}. These diagnostic efforts, combined with controlled ablation studies demonstrating that simpler fixed choices often outperform their adaptive versions on a fixed local-model budget \cite{ref305}, collectively highlight the need for careful, evidence-based design of agentic RAG systems that balance accuracy, efficiency, and interpretability. These cautionary findings directly inform the design of multi-agent systems in the following subsection, where controlled studies similarly reveal that simpler fixed hybrid retrieval strategies can rival or exceed complex adaptive routing on fixed budgets, underscoring the importance of empirical validation and cost-aware design across both single-agent and multi-agent paradigms.

\subsection{Multi-Agent Collaboration and Hybrid Orchestration for Complex Reasoning}

The progression from single-agent pipelines to multi-agent architectures marks a fundamental shift in Retrieval-Augmented Generation (RAG), particularly for complex reasoning tasks that demand coordinated evidence gathering and synthesis across heterogeneous data ecosystems. While single-agent systems are constrained by fixed retrieval and generation loops, multi-agent frameworks decompose complex queries, assign specialized roles, and orchestrate iterative collaboration to achieve deeper reasoning and more robust answers. This paradigm builds directly upon the agentic principles discussed previously---where individual agents autonomously decide when and how to retrieve---by distributing these decisions across multiple specialized agents that must coordinate their actions. This subsection examines the architectural paradigms of multi-agent RAG, focusing on hierarchical decomposition, consensus-driven integration, and hybrid orchestration of graph topology, document semantics, and source reliability, while also considering the efficiency implications that the subsequent section will address in greater depth.

A dominant design pattern is hierarchical multi-agent collaboration, where a coordinating agent decomposes a complex query into tractable sub-tasks, delegates them to specialized retrieval agents, and then integrates their outputs. This pattern extends the adaptive controllers discussed in the agentic RAG context by introducing explicit role specialization and inter-agent communication. HM-RAG exemplifies this approach with a three-tier architecture comprising a Decomposition Agent, Multi-source Retrieval Agents, and a Decision Agent. The Decomposition Agent dissects complex queries into semantically coherent sub-tasks via query rewriting and schema-guided context augmentation, while parallel retrieval agents execute modality-specific searches across vector, graph, and web-based databases. The Decision Agent then employs consistency voting to reconcile discrepancies across sources and refine answers through expert model integration. This hierarchical design yields substantial gains---a 12.95\% improvement in answer accuracy and a 3.56\% boost in question classification accuracy on ScienceQA and CrisisMMD benchmarks \cite{ref320}. Similarly, the MA-RAG framework orchestrates Planner, Step Definer, Extractor, and QA agents that progressively refine retrieval and synthesis through chain-of-thought reasoning, demonstrating that even small models like LLaMA3-8B can surpass larger standalone LLMs on multi-hop benchmarks through collaborative modular reasoning \cite{ref321}. This hierarchical decomposition is particularly effective because it transforms a monolithic reasoning problem into parallelizable sub-problems, each handled by an agent specialized for its particular retrieval modality or reasoning step.

A critical challenge in multi-agent collaboration is the quality of shared knowledge, particularly in graph-based retrieval where fragmented extraction can undermine downstream reasoning. MemGraphRAG addresses this by introducing a memory-based multi-agent society where agents leverage shared memory to maintain global context throughout knowledge graph construction. This memory mechanism enables agents to dynamically resolve logical conflicts and preserve structural connectivity across the corpus, overcoming the fragmentation that plagues isolated, fragment-level extraction approaches \cite{ref322}. Complementary to this, HERA introduces a hierarchical framework that jointly evolves multi-agent orchestration and role-specific prompts. At the global level, HERA optimizes query-specific agent topologies through reward-guided sampling and experience accumulation, while at the local level, Role-Aware Prompt Evolution refines agent behaviors via credit assignment. This dual-level adaptation achieves a 38.69\% average improvement over recent baselines while maintaining token efficiency, demonstrating the value of learning-based orchestration over static workflows \cite{ref323}. These approaches highlight that the effectiveness of multi-agent systems depends not only on individual agent capabilities but also on the mechanisms for knowledge sharing and coordination that bind them together.

Consensus-driven frameworks represent another essential paradigm for handling ambiguous or conflicting retrieval results, addressing a challenge that becomes more pronounced when multiple specialized agents return potentially divergent evidence. AnchorRAG addresses the open-world challenge where anchor entities cannot be reliably pre-identified by deploying a predictor agent that dynamically aligns query terms with knowledge graph nodes and initializes multiple independent retriever agents for parallel multi-hop exploration. A supervisor agent then formulates iterative retrieval strategies and synthesizes knowledge paths, mitigating the impact of erroneous anchors and achieving state-of-the-art results on real-world reasoning tasks \cite{ref324}. ConRAG takes a different consensus-driven approach by systematically optimizing the query and corpus sides and leveraging multi-view evidence---relation, entity, and text signals---to improve retrieval accuracy on multi-hop QA, achieving up to 26.9\% average performance gains over vanilla RAG and enabling a 31B model to achieve new state-of-the-art results on MuSiQue \cite{ref325}. SPD-RAG further applies consensus principles along the document axis: each document is processed by a dedicated agent, and a coordinator agent dispatches tasks and aggregates partial answers through a token-bounded synthesis layer. On the LOONG benchmark, SPD-RAG achieves an Avg Score of 58.1, substantially outperforming both Normal RAG (33.0) and Agentic RAG (32.8) while using only 38\% of the full-context baseline's API cost \cite{ref326}. The token-bounded synthesis in SPD-RAG also foreshadows the efficiency concerns that become central in the following subsection, where cost-aware orchestration across agents and retrieval strategies is examined in detail.

Beyond hierarchical decomposition and consensus voting, hybrid systems that unify graph topology, document semantics, and source reliability represent a frontier in multi-agent orchestration. HydraRAG is a training-free framework that integrates these three dimensions to support deep, faithful reasoning. It employs agent-driven exploration combining structured graph traversal and unstructured document retrieval to handle multi-hop and multi-entity problems, and uses tri-factor cross-source verification---source trustworthiness assessment, cross-source corroboration, and entity-path alignment---to balance topic relevance with cross-modal agreement. HydraRAG achieves state-of-the-art results across seven benchmarks, outperforming the strong hybrid baseline ToG-2 by an average of 20.3\% and up to 30.1\%, while enabling smaller models like Llama-3.1-8B to reach reasoning performance comparable to GPT-4-Turbo \cite{ref327}. Multi-agent designs also prove valuable in distributed and access-restricted settings. SCOUT-RAG employs four cooperative agents to estimate domain relevance, decide when to expand retrieval across domains, adapt traversal depth, and synthesize answers, achieving performance comparable to centralized baselines while substantially reducing cross-domain calls, total tokens, and latency \cite{ref328}. SPLIT-RAG similarly addresses scalability by using question-driven semantic graph partitioning to create specialized subgraphs, then assigning lightweight LLM agents to activated partitions with a hierarchical merging module for consistency resolution \cite{ref329}. These hybrid approaches demonstrate that multi-agent architectures can simultaneously enhance both reasoning quality and operational efficiency---a dual objective that becomes the central focus of the subsequent subsection on scalability and efficiency.

Cooperative designs where agents exchange information or supervise iterative retrieval extend the orchestration paradigm further. The mRAG framework, optimized through self-training with reward-guided trajectory sampling, showcases how inter-agent collaboration can be refined through learning across planning, searching, reasoning, and coordination subtasks \cite{ref330}. TechRAG implements a comprehensive evidence-gated pipeline with Planner, Researcher, Writer, and Critic agents, incorporating evidence sufficiency scoring, drift-guarded retries, and critic-driven revision to ensure citation integrity for technical literature reasoning \cite{ref331}. Meanwhile, the multi-agent GraphRAG framework translates natural language to Cypher queries for property graphs, using iterative content-aware correction and aggregated feedback loops to ensure query validity across diverse data sources \cite{ref332}. These cooperative frameworks draw heavily on the reinforcement learning and process supervision techniques introduced in the agentic RAG subsection, applying them at the inter-agent level rather than solely within a single agent's decision loop.

Despite these advances, empirical studies temper expectations regarding the universal benefit of multi-agent complexity, echoing similar cautions raised for agentic RAG systems in the previous subsection. A controlled ablation on HotpotQA with a local 7B model found that fixed hybrid retrieval via reciprocal rank fusion outperformed rule-based adaptive routing, and that two retrieval iterations captured 95\% of the gains of five, suggesting that simpler fixed choices can rival or exceed adaptive versions on a fixed budget \cite{ref305}. Likewise, an agent-orchestrated adaptive RAG study found that query decomposition yields consistent gains in structured domains but degrades ranking precision on multi-hop benchmarks, while reflection improves citation accuracy at substantial latency cost, arguing for adaptive, cost-aware orchestration rather than uniformly aggressive pipelines \cite{ref31}. Together, these findings indicate that the design of multi-agent RAG systems must be guided by query complexity, domain characteristics, and resource constraints, presenting both significant opportunities and important trade-offs for complex reasoning. Moreover, these trade-offs---between accuracy gains and computational overhead, between orchestration sophistication and latency---establish the central tension that motivates the efficiency-oriented frameworks and cost-aware design principles examined in the following subsection on scalability and deployment considerations.

\subsection{Scalability, Efficiency, and Benchmarking of Advanced RAG Architectures}

The advanced RAG architectures discussed thus far---GraphRAG, agentic, and multi-agent systems---offer substantial gains in reasoning and accuracy, but their practical deployment hinges critically on scalability and efficiency. The sophisticated indexing pipelines, iterative retrieval loops, and multi-agent orchestration that enable complex reasoning can introduce prohibitive computational costs, high token consumption, and significant latency, limiting their viability in production environments. The multi-agent coordination strategies examined previously---particularly the token-bounded synthesis in SPD-RAG and the adaptive retrieval triggers in A2RAG---already hinted at these concerns; this subsection now places efficiency at the center of analysis, examining when and how the complexity of these systems can be justified.

\textbf{Efficiency-Oriented Frameworks and Construction Cost Reduction}

A primary bottleneck for GraphRAG is the high cost and instability of offline knowledge graph construction, which often relies on expensive LLM-based relation extraction---a cost that is amplified in multi-agent settings where multiple specialized agents must coordinate over such structures. To address this, several frameworks propose lightweight alternatives. For instance, \href{LinearRAG\%20Linear\%20Graph\%20Retrieval\%20Augmented\%20Generation\%20on\%20Large-scale\%20Corpora}{285} introduces a relation-free hierarchical graph, termed Tri-Graph, constructed using only lightweight entity extraction and semantic linking. This paradigm scales linearly with corpus size and incurs no extra token consumption, circumventing the noisy and costly relation modeling of prior systems. Similarly, \href{Efficient\%20Retrieval-Augmented\%20Generation\%20via\%20Token\%20Co-occurrence\%20Graphs}{286} constructs a token co-occurrence knowledge graph using sliding-window statistics, enabling scalable graph construction without LLM-based extraction and substantially reducing indexing time, inference latency, and prompt footprint. The \href{EHRAG\%20Bridging\%20Semantic\%20Gaps\%20in\%20Lightweight\%20GraphRAG\%20via\%20Hybrid\%20Hypergraph\%20Construction\%20and\%20Retrieval}{292} framework further exemplifies this trend by employing Named Entity Recognition (NER) and text embedding clustering to build a hybrid structural-semantic hypergraph with zero token consumption during construction, maintaining linear indexing complexity while capturing latent semantic connections between disjoint entities.

This focus on architecture efficiency is reflected in design philosophies like that of \href{NodeRAG\%20Structuring\%20Graph-based\%20RAG\%20with\%20Heterogeneous\%20Nodes}{333}, which emphasizes that inadequately designed graph structures impede the seamless integration of graph algorithms and degrade performance. By proposing a heterogeneous graph structure, NodeRAG demonstrates performance advantages over prior methods like GraphRAG and LightRAG, not only in question-answering quality but also in indexing time, query time, and storage efficiency. Similarly, \href{E\%5E2GraphRAG\%20Streamlining\%20Graph-based\%20RAG\%20for\%20High\%20Efficiency\%20and\%20Effectiveness}{35} streamlines the pipeline by constructing a summary tree with LLMs and an entity graph with \emph{SpaCy}, achieving up to 10 times faster indexing than GraphRAG and a 100 times speedup in retrieval over LightRAG while maintaining competitive performance. The \href{TERAG\%20Token-Efficient\%20Graph-Based\%20Retrieval-Augmented\%20Generation}{286} framework tackles the cost issue head-on by using Personalized PageRank (PPR) for retrieval, achieving at least 80\% of the accuracy of widely used graph-based methods while consuming only 3\%--11\% of the output tokens, making it well-suited for large-scale, cost-sensitive deployments. These frameworks, alongside \href{SlimRAG\%20Retrieval\%20without\%20Graphs\%20via\%20Entity-Aware\%20Context\%20Selection}{334} and \href{LeanRAG\%20Knowledge-Graph-Based\%20Generation\%20with\%20Semantic\%20Aggregation\%20and\%20Hierarchical\%20Retrieval}{335}, collectively demonstrate a paradigm shift toward minimizing the structural overhead and token footprint of graph-based systems. \href{SlimRAG\%20Retrieval\%20without\%20Graphs\%20via\%20Entity-Aware\%20Context\%20Selection}{334} argues that semantic similarity does not imply semantic relevance and replaces structure-heavy components with a compact entity-to-chunk table, while \href{LeanRAG\%20Knowledge-Graph-Based\%20Generation\%20with\%20Semantic\%20Aggregation\%20and\%20Hierarchical\%20Retrieval}{335} addresses the inefficiency of flat search and disconnected ``semantic islands'' by aggregating entities into a navigable semantic network and using a bottom-up, structure-guided retrieval strategy that reduces retrieval redundancy by 46\%.

\textbf{Retrieval Latency and Token Consumption in Agentic Systems}

Agentic RAG systems, while offering dynamic planning and iterative refinement---as described in the preceding subsections on agentic and multi-agent architectures---often incur substantial token overhead from multi-round search and reasoning processes. This overhead is particularly acute in multi-agent settings, where inter-agent communication and consensus-building add layers of coordination cost on top of individual retrieval and reasoning steps. To address this, \href{TeaRAG\%20A\%20Token-Efficient\%20Agentic\%20Retrieval-Augmented\%20Generation\%20Framework}{311} proposes a token-efficient framework that compresses both retrieved content---by augmenting chunk-based retrieval with concise graph triplets and using Personalized PageRank to highlight key knowledge---and reasoning steps via a novel Iterative Process-aware Direct Preference Optimization (IP-DPO). This design reduces output tokens by up to 61\% while improving exact match accuracy. Further pushing the efficiency boundary, \href{LatentRAG\%20Latent\%20Reasoning\%20and\%20Retrieval\%20for\%20Efficient\%20Agentic\%20RAG}{336} shifts both reasoning and retrieval from discrete language space to continuous latent space, generating latent tokens directly from hidden states in a single forward pass. This approach reduces inference latency by approximately 90\% compared to explicit agentic methods, narrowing the gap with traditional single-step RAG. Additionally, frameworks like \href{A2RAG\%20Adaptive\%20Agentic\%20Graph\%20Retrieval\%20for\%20Cost-Aware\%20and\%20Reliable\%20Reasoning}{40} and \href{Use\%20Graph\%20When\%20It\%20Needs\%20Efficiently\%20and\%20Adaptively\%20Integrating\%20Retrieval-Augmented\%20Generation\%20with\%20Graphs}{35} introduce adaptive controllers that dynamically adjust retrieval effort based on query complexity. \href{A2RAG\%20Adaptive\%20Agentic\%20Graph\%20Retrieval\%20for\%20Cost-Aware\%20and\%20Reliable\%20Reasoning}{40} couples an adaptive controller that triggers targeted refinement only when evidence is insufficient, cutting token consumption and latency by about 50\% relative to iterative baselines, while \href{Use\%20Graph\%20When\%20It\%20Needs\%20Efficiently\%20and\%20Adaptively\%20Integrating\%20Retrieval-Augmented\%20Generation\%20with\%20Graphs}{35} uses syntax-aware complexity analysis to route simple queries to dense RAG and complex ones to graph-based retrieval, improving accuracy and reducing latency in mixed scenarios. The \href{GraphRAG-Router\%20Learning\%20Cost-Efficient\%20Routing\%20over\%20GraphRAGs\%20and\%20LLMs\%20with\%20Reinforcement\%20Learning}{337} also addresses cost-efficiency by using a hierarchical routing strategy to coordinate heterogeneous GraphRAGs and generator LLMs, reducing overuse of large LLMs by nearly 30\% through a curriculum cost-aware reward. These adaptive controllers resonate directly with the orchestration strategies discussed in the multi-agent subsection, where cost-aware routing and selective delegation emerged as key design principles for balancing reasoning depth with operational overhead.

\textbf{Benchmarking and Comparative Evaluation}

Given the proliferation of these advanced architectures, a critical question emerges: when are these complex systems truly beneficial, and when does a simpler RAG suffice---particularly in light of the empirical cautions raised at the end of the previous subsection regarding multi-agent complexity? Several benchmarking efforts have been developed to provide data-driven answers. The \href{GraphRAG-Bench\%20Challenging\%20Domain-Specific\%20Reasoning\%20for\%20Evaluating\%20Graph\%20Retrieval-Augmented\%20Generation}{338} benchmark is designed to rigorously evaluate GraphRAG models with college-level, domain-specific questions that demand multi-hop reasoning, providing a holistic evaluation across graph construction, knowledge retrieval, and answer generation. Similarly, \href{When\%20to\%20use\%20Graphs\%20in\%20RAG\%20A\%20Comprehensive\%20Analysis\%20for\%20Graph\%20Retrieval-Augmented\%20Generation}{338} offers a comprehensive dataset covering fact retrieval, complex reasoning, contextual summarization, and creative generation to systematically investigate the conditions under which GraphRAG surpasses traditional RAG.

Comparative studies provide crucial insights into this ``when to use'' paradigm. One systematic evaluation, \href{RAG\%20vs.\%20GraphRAG\%20A\%20Systematic\%20Evaluation\%20and\%20Key\%20Insights}{339}, introduces a unified protocol and finds that RAG and GraphRAG have distinct strengths depending on the task, leading the authors to propose selection and integration strategies that combine both paradigms. Notably, research on agentic search challenges the assumption that explicit graph structure is always necessary. The \href{Do\%20We\%20Still\%20Need\%20GraphRAG\%20Benchmarking\%20RAG\%20and\%20GraphRAG\%20for\%20Agentic\%20Search\%20Systems}{340} benchmark demonstrates that agentic search can substantially improve dense RAG, narrowing the performance gap to GraphRAG, especially in reinforcement learning-based settings, while acknowledging GraphRAG's advantage in complex multi-hop reasoning when its offline cost is amortized. This is echoed by research on routing and orchestration, which shows that agentic enhancements are not universally beneficial and must be applied selectively according to query and domain characteristics \href{Agent-Orchestrated\%20Adaptive\%20RAG\%20A\%20Comparative\%20Study\%20on\%20Structured\%20and\%20Multi-Hop\%20Retrieval}{31}---a finding that parallels the controlled ablations on multi-agent complexity discussed in the previous subsection.

Furthermore, studies like the one presented in \href{Is\%20GraphRAG\%20Needed\%20From\%20Basic\%20RAG\%20to\%20Graph-\%20Agentic\%20Solutions\%20with\%20Context\%20Optimization}{335} provide a framework for evaluating nine standardized RAG scenarios and identify a ``retrieval-generation gap'' where expanded retrieval does not proportionally improve generation quality. This finding cautions against over-relying on retrieval-oriented metrics, suggesting that advanced retrieval benefits can be overstated. A scaling study, \href{BM25\%20Wins\%20at\%20Scale\%20A\%20Scaling\%20Study\%20of\%20Retrieval-Augmented\%20Generation\%20Paradigms}{100}, reveals a scale-dependent crossover in performance, showing that lexical retrieval (BM25) becomes a stronger default as corpus size grows, while agentic reasoning works best after ranked discovery. Finally, the practical deployment of these systems is addressed by system-level engineering efforts and declarative frameworks. For instance, \href{AutoRAGTuner\%20A\%20Declarative\%20Framework\%20for\%20Automatic\%20Optimization\%20of\%20RAG\%20Pipelines}{146} provides a declarative, configuration-driven framework that automates the RAG life cycle and mitigates engineering overhead by enabling architectural adjustments with up to a 95\% reduction in code churn. In resource-constrained serving environments, \href{VectorLiteRAG\%20Latency-Aware\%20and\%20Fine-Grained\%20Resource\%20Partitioning\%20for\%20Efficient\%20RAG}{341} demonstrates that fine-grained GPU resource allocation based on performance modeling can significantly improve SLO-compliant throughput without needing additional hardware. These benchmarking and system-level efforts collectively underscore that the choice of architecture---whether simple RAG, GraphRAG, agentic, or multi-agent---must be driven by a careful analysis of task complexity, corpus scale, cost budgets, and performance requirements, rather than a one-size-fits-all approach. This cost-aware perspective provides the practical lens through which the theoretical advances of the preceding subsections can be translated into sustainable, production-ready deployments.

\section{Multimodal, Multilingual, and Domain-Specific RAG Systems}

\subsection{Foundations and Architectures of Multimodal RAG}

The emergence of Multimodal Retrieval-Augmented Generation (MRAG) marks a fundamental departure from text-only RAG, which is inherently constrained by its reliance on parsing real-world documents into plain text---a process that inevitably discards crucial visual cues such as layouts, tables, and charts that often contain essential information \cite{ref342}. By enabling the retrieval and generation process to natively handle multiple modalities---including images, tables, audio, and video---MRAG grounds responses in more complete and faithful evidence \cite{ref343,ref45}, thereby addressing a core limitation of its predecessor and setting the stage for a new generation of knowledge-intensive multimodal applications. The following subsections build on this foundation to examine the benchmarks, methodologies, and open challenges that have emerged alongside these architectural advances.

The foundational architecture of MRAG can be understood as an extension of the tripartite RAG pipeline---retriever, generator, and augmentation mechanism---where each component is adapted to handle multimodal inputs and outputs. Early work in this area, such as MuRAG, pioneered the concept of a Multimodal Retrieval-Augmented Transformer, which integrates a non-parametric multimodal memory into the generation process. MuRAG is pre-trained using a joint contrastive and generative loss on a mixture of large-scale image-text and text-only corpora, enabling it to retrieve and reason over both images and text to answer queries \cite{ref344}. This model established a critical precedent: that retrieval could operate beyond the textual domain to augment vision-language tasks with external cross-modal knowledge.

A core architectural challenge in MRAG lies in the design of the retrieval backbone. One dominant approach leverages vision-language models (VLMs) as unified representers for both documents and queries. This strategy is exemplified by VisRAG, which proposes a VLM-based pipeline where entire document pages are embedded directly as images, eliminating the need for a brittle text-extraction phase. By treating the document as an image, VisRAG maximizes the retention of original information, including layout and visual elements, and has demonstrated that this approach yields significant end-to-end performance gains of 20--40\% over traditional text-based RAG pipelines \cite{ref342}. Similarly, frameworks like Document Haystacks introduce vision-centric RAG, leveraging a suite of specialized multimodal vision encoders optimized for different strengths to power a question-document relevance module, enabling efficient retrieval over thousands of images \cite{ref345}. This ``image-as-document'' paradigm contrasts with approaches that rely on OCR or text extraction, which, while often effective, are sensitive to layout shifts and lose spatial cues that may contain the answer \cite{ref346}.

Beyond using VLMs as unified backbones, other architectures focus on the fusion of separate modality-specific encoders. The ``Co-Modality'' approach, as implemented in CMRAG, uses a Unified Encoding Model (UEM) to project queries, parsed text, and images into a shared embedding space via triplet-based training. A subsequent Unified Co-Modality-informed Retrieval (UCMR) method statistically normalizes similarity scores to effectively fuse these cross-modal signals, demonstrating that jointly leveraging text and images yields more accurate retrieval than single-modality baselines \cite{ref347}. This focus on alignment between different embedding spaces is a recurring theme, as the modality gap between text and vision is a fundamental obstacle that fine-tuning often cannot easily bridge; lightweight, training-free linear mappings have been proposed to extend across this gap for RAG purposes \cite{ref348}.

A critical design consideration in MRAG is the granularity at which retrieval and generation operate. Early systems often treated entire documents or pages as the atomic retrieval unit. However, this assumption is frequently suboptimal, as a page may contain a large amount of irrelevant noise that dilutes the signal for the generator \cite{ref349}. To address this, recent frameworks have shifted toward finer granularities. RegionRAG, for instance, shifts the retrieval paradigm from document-level to region-level, identifying salient patches within a page and dynamically grouping them into coherent semantic regions to serve as the retrieval unit. This approach not only improves retrieval accuracy but also reduces the visual token consumption for the generator \cite{ref350}. Likewise, FES-RAG reframes MRAG as a fine-grained evidence selection problem, decomposing documents into atomic sentence-level textual and region-level visual fragments, and uses a principled metric to measure each fragment's marginal contribution to generation confidence \cite{ref349}. This push towards finer granularity is also mirrored in layout-oriented frameworks like LFRAG, which moves from page-level to block-level retrieval by performing layout segmentation to construct semantically coherent retrieval units, reducing token consumption by over 70\% in generation \cite{ref351}.

To support research in this direction, several frameworks have been developed to handle the complexity of multimodal documents. For instance, RAG-Anything reconceptualizes multimodal content as interconnected knowledge entities. It introduces a dual-graph construction to capture both cross-modal relationships and textual semantics, enabling a cross-modal hybrid retrieval that combines structural knowledge navigation with semantic matching. This allows for effective reasoning over heterogeneous content where relevant evidence is scattered across modalities \cite{ref44}. Furthermore, sophisticated frameworks like Multi-Granularity Graph (MG\(^2\)-RAG) build hierarchical multimodal knowledge graphs by combining lightweight textual parsing with entity-driven visual grounding, fusing textual entities and visual regions into unified multimodal nodes that preserve atomic evidence, which supports multi-hop structured reasoning while significantly reducing graph construction costs \cite{ref46}. These graph-based methods represent a more structural approach to MRAG, capturing the complex interdependencies between modalities.

The training strategies for multimodal retrievers and generators are diverse and crucial for system performance. Some approaches, like RAVEN, focus on efficient task-specific fine-tuning of a base VLM with retrieval-augmented samples, demonstrating that retrieval properties can be acquired without adding retrieval-specific parameters, leading to significant improvements in tasks such as image captioning and VQA \cite{ref352}. Others explore more specialized training paradigms, such as the two-stage framework proposed for Nyx, which pre-trains a unified mixed-modal retriever on a constructed dataset and then fine-tunes it using feedback from downstream VLMs to align retrieval outputs with generative preferences \cite{ref353}. Additionally, end-to-end frameworks like CoRe-MMRAG introduce a specialized training paradigm to enhance knowledge source discrimination and multimodal integration, directly addressing the inconsistencies that arise between parametric knowledge and retrieved multimodal evidence \cite{ref354}. Overall, the field of MRAG foundations is characterized by a move away from brute-force page-level processing towards more intelligent, fine-grained, and structurally aware systems that leverage VLMs to unify the representational space of different modalities, a trend that is essential for achieving reliable and efficient multimodal knowledge access \cite{ref355}.

\subsection{Benchmarking and Evaluation of Multimodal RAG Systems}

The rapid proliferation of Multimodal Retrieval-Augmented Generation (MRAG) systems---which extend the RAG paradigm by retrieving and integrating evidence across text, images, tables, and video---has created an urgent need for rigorous, task-appropriate evaluation benchmarks and methodologies. Unlike text-only RAG evaluation, which benefits from mature metrics and long-standing datasets, the evaluation of MRAG systems is complicated by heterogeneous input modalities, the challenge of assessing cross-modal relevance, and the difficulty of verifying whether generated responses are truly grounded in both textual and visual evidence. This subsection provides a comprehensive review of the key benchmarks and evaluation methodologies that have emerged to address this gap, analyzing task taxonomies, metric design at multiple granularities, LLM-as-a-judge approaches, and modality-specific challenges.

\textbf{Foundational Benchmarks and Task Taxonomies.} Early efforts to benchmark MRAG systems focused on defining task categories that reflect real-world multimodal reasoning requirements. The M\^{}\{2\}RAG benchmark formalized four core tasks inherent to multimodal contexts: image captioning, multimodal question answering, multimodal fact verification, and image reranking \cite{ref356}. All of these tasks are evaluated in an open-domain setting, requiring models to retrieve query-relevant information from a multimodal collection and utilize that evidence for generation. This taxonomy proved influential, as it established a direct link between retrieval quality and downstream task performance while also highlighting the need for instruction tuning to optimize models for multimodal context utilization \cite{ref356}. In a similar spirit, the MRAG-Bench benchmark introduced a vision-centric perspective by systematically identifying nine distinct scenarios---ranging from varying viewpoints to visual analogies---where visually augmented knowledge is demonstrably more beneficial than pure textual augmentation \cite{ref357}. This benchmark comprised 16,130 images and 1,353 human-annotated multiple-choice questions, providing controlled conditions for evaluating the marginal contribution of retrieved images. Notably, while all evaluated vision-language models (LVLMs) showed greater improvement when augmented with images compared to textual knowledge, even top-performing proprietary models such as GPT-4o achieved only a 5.82\% improvement with ground-truth visual information---a stark contrast to the 33.16\% improvement observed in human participants \cite{ref357}. This finding exposed a critical gap between the availability of relevant visual evidence and the models' ability to effectively exploit it.

\textbf{Document-Centric and Visually Rich Benchmarks.} A separate line of benchmarks emphasizes the evaluation of MRAG systems on real-world documents where information is embedded in complex visual layouts. UniDoc-Bench was introduced as the first large-scale, realistic benchmark for document-centric MM-RAG, constructed from real-world PDF pages across domains and featuring a rigorous pipeline that extracts and links evidence from text, tables, and figures to generate multimodal QA pairs spanning factual retrieval, comparison, summarization, and logical reasoning \cite{ref358}. A key design principle of UniDoc-Bench is its support for apples-to-apples comparison across four paradigms---text-only, image-only, multimodal text-image fusion, and multimodal joint retrieval---under a unified protocol with standardized candidate pools and evaluation metrics \cite{ref358}. The experiments revealed a substantive finding: multimodal text-image fusion consistently outperformed both unimodal and jointly multimodal embedding-based retrieval, indicating that neither text nor images alone suffice for document-centric QA and that current multimodal embedding models remain inadequate for capturing complementary information. Similarly, VisDoMBench introduced the first benchmark specifically designed for multi-document QA with visually rich elements such as tables, charts, and presentation slides, and proposed VisDoMRAG---a framework that runs concurrent textual and visual RAG pipelines and employs a consistency-constrained modality fusion mechanism to align reasoning across modalities \cite{ref359}. This benchmark highlighted scenarios where critical information is distributed across modalities and demonstrated that multi-modal fusion improves accuracy by 12--20\% over unimodal and long-context baselines \cite{ref359}. Addressing a higher level of complexity, the MMDocRAG benchmark introduced a large-scale evaluation for document VQA featuring 4,055 expert-annotated QA pairs with multi-page, cross-modal evidence chains \cite{ref360}. MMDocRAG pioneered metrics for evaluating multimodal quote selection and enabled interleaved text-plus-image answers. Through experiments with 60 VLM/LLM models and 14 retrieval systems, the authors identified persistent challenges in multimodal evidence retrieval and integration, and importantly found that fine-tuned LLMs using detailed image descriptions achieved substantial improvements---suggesting a scalable alternative to end-to-end multimodal models \cite{ref360}. For the specific challenge of large-scale visual retrieval across thousands of documents, the DocHaystack and InfoHaystack benchmarks extended the scope to 1,000+ documents, revealing limitations of LMMs in long-horizon retrieval and leading to the proposal of V-RAG, a vision-centric RAG framework with specialized multimodal encoders that achieved 9\% and 11\% improvements in Recall@1 \cite{ref345}.

\textbf{Generative and Multimodal-Output Evaluation.} A more recent trend addresses the task of Multimodal Retrieval-Augmented Multimodal Generation (MRAMG), where the system output itself must contain a combination of text and images, fully leveraging the multimodal nature of the corpus. MRAMG-Bench was constructed as a meticulously curated, human-annotated benchmark comprising 4,346 documents, 14,190 images, and 4,800 QA pairs across six datasets spanning web, academia, and lifestyle domains \cite{ref361}. This benchmark introduced a comprehensive suite of statistical and LLM-based metrics tailored for evaluating multimodal outputs, acknowledging that traditional text-only metrics are inadequate for assessing generation quality when images are part of the answer \cite{ref361}. Alongside this, RAG-IGBench was proposed to specifically evaluate interleaved image-text generation grounded in retrieved multimodal evidence, drawing on the latest publicly available social platform content and introducing metrics that measure text quality, image quality, and text-image consistency \cite{ref362}. A key methodological contribution of RAG-IGBench was validating that its proposed evaluation metrics correlate highly with human judgments, thereby establishing their practical utility for future system comparison \cite{ref362}. For medical applications, MedGEN-Bench introduced a comprehensive multimodal benchmark for open-ended medical generation with 6,422 expert-validated image-text pairs spanning six imaging modalities and 16 clinical tasks, and proposed a novel three-tier assessment framework combining pixel-level metrics, semantic text analysis, and expert-guided clinical relevance scoring \cite{ref363}.

\textbf{Diagnostic Evaluation and Source Attribution.} The need to move beyond single aggregate scores toward fine-grained, diagnosable evaluation led to several important initiatives. MRAG-Suite was introduced as a diagnostic evaluation platform integrating diverse multimodal benchmarks (WebQA, Chart-RAG, Visual-RAG, MRAG-Bench) and introducing difficulty-based and ambiguity-aware filtering strategies alongside an MM-RAGChecker---a claim-level diagnostic tool \cite{ref364}. The empirical results demonstrated substantial accuracy reductions under difficult and ambiguous queries, and the claim-level diagnostic proved effective in identifying prevalent hallucination patterns, offering guidance for system improvement \cite{ref364}. Recognizing that current evaluations often fail to systematically account for query difficulty, this work highlighted the need to stratify benchmark results by query characteristics rather than reporting only average performance. In a similar diagnostic spirit, MAVIS was introduced as the first benchmark designed to evaluate multimodal source attribution over long-form visual question answering \cite{ref365}. The dataset comprises 157K visual QA instances with fact-level citations referring to multimodal documents and introduced fine-grained automatic metrics spanning informativeness, groundedness, and fluency, demonstrating strong correlation with human judgments \cite{ref365}. Notable findings included that LVLMs with multimodal RAG generate more informative and fluent answers than unimodal RAG, but exhibit weaker groundedness for image documents than text documents---a gap amplified in multimodal settings---revealing modality-specific fragility in grounding \cite{ref365}. These diagnostic efforts directly complement the retrieval-centric architectural advances discussed in the previous subsection, providing the measurement tools necessary to validate whether fine-grained, region-level, and graph-based retrieval mechanisms actually deliver the promised gains in grounding and hallucination reduction.

\textbf{Sector-Specific Benchmarks for Multimodal RAG.} To address the specific demands of chart-based information, CHARGE (CHARt-based document question-answering GEneration) introduced a semi-automated framework for generating evaluation data via structured keypoint extraction, cross-modal verification, and keypoint-based generation, resulting in the Chart-MRAG Bench with 4,738 QA pairs across 8 domains \cite{ref366}. The evaluation of current approaches on this benchmark surfaced three critical limitations: unified multimodal embedding retrieval methods struggle in chart-based scenarios; even with ground-truth retrieval, state-of-the-art MLLMs achieve only 58.19\% Correctness and 73.87\% Coverage; and MLLMs demonstrate a consistent text-over-visual modality bias \cite{ref366}. For the medical domain, the MRAG benchmark addressed the lack of comprehensive evaluation in scientific and clinical QA, covering various tasks in both English and Chinese while building a corpus from Wikipedia and PubMed, and accompanied the benchmark with the MRAG-Toolkit to facilitate systematic exploration of RAG components \cite{ref367}. Experiments revealed that retrieval model size and prompting strategies significantly influence system performance, and highlighted a subtle trade-off: while RAG improves usefulness and reasoning quality, LLM responses may become slightly less readable for long-form questions \cite{ref367}. Finally, M4-RAG stands out for its scale and cross-cultural coverage, encompassing 42 languages, 56 regional dialects, and 189 countries with over 80,000 culturally diverse image-question pairs \cite{ref368}. A controlled retrieval environment---containing millions of curated multilingual documents---balances realism with reproducibility. The systematic evaluation revealed a crucial finding: although RAG consistently benefits smaller VLMs, it often fails to scale to larger models and can even degrade their performance, exposing a mismatch between model capacity and current retrieval effectiveness \cite{ref368}. This finding urges the community to invest in more effective retrieval mechanisms that can complement larger model parametric knowledge. Notably, M4-RAG also serves as a bridge to the multilingual RAG challenges examined in the following subsection, demonstrating that multimodal evaluation and cross-lingual evaluation increasingly overlap in the design of globally deployable RAG systems.

\textbf{Synthesis and Open Challenges.} Taken together, the benchmarks surveyed above reveal a rapidly maturing evaluation ecosystem for MRAG systems. Across diverse benchmarks, the key themes converge on the importance of: (1) task taxonomies that incorporate both retrieval- and generation-stage challenges \cite{ref356,ref357}; (2) document-centric evaluation that reflects the complexity of real-world visually rich documents \cite{ref358,ref359}; (3) diagnostic evaluation that provides root-cause analysis beyond aggregate scores \cite{ref364}; (4) fact-level source attribution metrics that verify the grounding of each generated claim \cite{ref365}; and (5) specialized benchmarks for charts, medical imaging, and cross-lingual cultural contexts \cite{ref366,ref367,ref368}. Standardized protocols comparing four paradigms, as introduced by UniDoc-Bench, are essential for making meaningful progress claims \cite{ref358}. Finally, benchmarks should increasingly incorporate multimodality in the expected output (MRAMG) rather than only in the retrieved context, as the ability to generate interleaved image-text responses is a defining capability of next-generation agents \cite{ref361}.

\subsection{Multilingual and Cross-Lingual RAG Systems}

Multilingual Retrieval-Augmented Generation (mRAG) extends the RAG paradigm to settings where user queries, supporting corpora, and generated responses may span multiple languages. While the core principle of augmenting parametric knowledge with non-parametric evidence remains unchanged, mRAG introduces unique challenges related to cross-lingual evidence retrieval, language preference biases, and support for low-resource languages. These challenges require a rethinking of both the retrieval and generation components of the RAG pipeline, making mRAG a distinct and active area of research. The following sections examine the architectural strategies, bias-mitigation techniques, empirical findings, and evaluation benchmarks that collectively define the current state of the art in this domain.

A fundamental challenge in mRAG is determining where and how to perform retrieval in a multilingual context. Early strategies often adopted English-pivoting approaches, where non-English queries are translated into English before retrieval, and the retrieved English documents are used directly. This approach, termed tRAG (question-translation), can benefit from the superior English-centric capabilities of many retrievers and generators, yet suffers from limited coverage and inefficiency because it discards potentially valuable non-English documents \cite{ref369}. In contrast, MultiRAG (Multilingual RAG) performs retrieval directly across the multilingual corpus, treating all languages uniformly. While this increases efficiency and recall of potentially relevant documents, it introduces inconsistencies in the retrieved content due to cross-lingual variations and potential redundancies \cite{ref369}. To address the issue of inconsistent retrieved content, CrossRAG proposes translating the retrieved documents into a common language, such as English, before generation. This approach standardizes the evidence for the downstream generator and has been shown to enhance performance on knowledge-intensive tasks across both high-resource and low-resource languages \cite{ref369}. Building on this idea, further work has shown that translation quality is a critical bottleneck. The Quality-Aware Translation Tagging in mRAG (QTT-RAG) method explicitly evaluates translation quality along several dimensions---semantic equivalence, grammatical accuracy, and naturalness---and attaches these quality scores as metadata to the translated documents. This technique allows the generator to weigh the evidence based on its reliability, preserving factual integrity and outperforming approaches that assume perfect translation quality, especially for low-resource languages like Korean and Finnish \cite{ref370}.

A particularly deep-seated issue in mRAG is the phenomenon of language preference, where both retrievers and generators exhibit biases toward certain languages. Research on language preference has systematically investigated these biases in the retrieval and generation stages of mRAG \cite{ref371}. Retrievers tend to favor documents in high-resource or query languages, but this preference does not always lead to improved generation performance. Furthermore, generators often display a preference for the query language or for Latin scripts, leading to inconsistent output. In response, frameworks like Dual Knowledge Multilingual RAG (DKM-RAG) aim to mitigate these issues by fusing translated multilingual passages with the model's complementary parametric knowledge, effectively bridging the gap between external evidence and internal knowledge \cite{ref371}. This bias is not merely a performance concern but can result in ``linguistic nepotism,'' where models preferentially cite sources in a specific language, even trading off document relevance for language preference, a behavior amplified for lower-resource languages and mid-context documents \cite{ref372}. Related work has also shown that current rerankers in mRAG pipelines systematically favor English and the query's native language, suppressing ``answer-critical'' documents from other languages. To counter this, methods like LAURA (Language-Agnostic Utility-driven Reranker Alignment) align multilingual evidence ranking with downstream generative utility, mitigating language bias and consistently improving mRAG performance \cite{ref373}. Complementing these system-level approaches, more recent work has proposed debiased metrics, such as DeLP (Debiased Language Preference), to explicitly factor out structural confounds like exposure bias and topic locality. Debiased analysis reveals that the apparent English preference is largely a byproduct of evidence distribution, while retrievers actually favor monolingual alignment between the query and document language. This insight has led to the development of DELTA (DEbiased Language preference-guided Text Augmentation), a lightweight framework that leverages monolingual alignment to optimize cross-lingual retrieval \cite{ref374}. During generation, especially when retrieved evidence is in a different language than the query, models often suffer from language drift, producing responses in an unintended language. This is often attributed to decoder-level collapse rather than a failure of comprehension. To mitigate this, training-free decoding strategies like Soft Constrained Decoding (SCD) gently steer generation toward the target language by penalizing token selection from non-target languages, improving language alignment without architectural changes \cite{ref375}. Similarly, teacher-regularized reinforcement learning methods, such as TR-RAG, have been proposed to couple reward optimization with on-policy distillation to prevent language-consistency collapses (drift) and improve evidence-grounded correctness in English-evidence cross-lingual RAG regimes \cite{ref376}.

Beyond general mRAG techniques, a significant body of work investigates the behavior of LLMs in controlled multilingual settings. Studies that isolate the effect of context language on generation quality have found that while LLMs are surprisingly capable of extracting relevant information from passages written in different languages, they are much weaker at consistently formulating full answers in the correct language, particularly when distractor passages in other languages are included \cite{ref377}. In cross-lingual retrieval for specific domains, the language mismatch between query and document languages has been found to be a critical bottleneck, with substantial performance drops in ranking, suggesting that focusing on retriever optimization may be more beneficial than simply improving generation capabilities \cite{ref378}. While high-quality multilingual retrievers and generators exist, task-specific prompt engineering is still needed to enable generation in user languages, as are evaluation metrics that account for variations like named entities across languages \cite{ref379}. For complex reasoning tasks, frameworks such as DaPT have been designed to handle multi-hop queries by generating sub-question graphs in both the source language and an English translation, merging them in a bilingual retrieval-and-answer sequence. This strategy produces more accurate and concise answers and significantly boosts performance on multilingual multi-hop benchmarks \cite{ref380}. In addition, search-augmented reinforcement learning frameworks like CroSearch-R1 integrate multilingual knowledge into GRPO processes using multi-turn retrieval and cross-lingual knowledge integration to dynamically align and supplement evidence, while language-coupled RL (LcRL) methods regularize anti-consistency penalties to manage knowledge conflicts in large multilingual collections \cite{ref381,ref382}.

The evaluation of mRAG systems has also evolved to address its unique complexity. Early benchmarks like M4-RAG provide a massive-scale evaluation for multilingual multimodal RAG, spanning 42 languages, 56 regional dialects, and 189 countries, revealing that RAG may not scale effectively to larger models and that performance degrades for non-English prompts and contexts \cite{ref368}. The XRAG benchmark specifically targets cross-lingual RAG settings, using recent news articles to ensure external knowledge is required, and has uncovered challenges in response language correctness and reasoning over multilingual information \cite{ref383}. Benchmarks like MEMERAG emphasize native-language evaluation, using expert human annotators to assess faithfulness and relevance, thereby providing a more representative user experience compared to translated data \cite{ref384}. Furthermore, MMed-Bench-IR addresses the specific needs of multilingual medical information retrieval, disentangling cross-lingual alignment, concept discrimination, and evidence retrieval to reveal severe degradation in biomedical encoder performance across languages \cite{ref385}. IndicRAGSuite focuses on Indian languages, providing large-scale datasets and benchmarks to address the lack of resources for multilingual retrieval and generation in a linguistically diverse region \cite{ref386}. Finally, benchmarks like BordIRLines tackle culturally sensitive tasks, showing that retrieving multilingual documents can improve response consistency and decrease geopolitical bias compared to in-language retrieval, highlighting complex interactions between language and culture in RAG \cite{ref387}. Collectively, these benchmark efforts show that tackling cross-lingual RAG challenges---evidenced by work on RAG for tasks like morphological translation and comprehension-based QA---can lead to more robust and equitable systems that serve a broader range of languages and users.

In summary, mRAG represents a distinct and challenging frontier in retrieval-augmented generation, where the interplay between language, retrieval, and generation creates unique obstacles not present in monolingual settings. The architectural choices between English-pivoting, direct multilingual retrieval, and cross-lingual translation are foundational and directly impact downstream performance. The pervasive issue of language preference---manifesting in both retrievers and generators---demands specialized mitigation strategies at the ranking, decoding, and training levels. Empirical findings underscore the need for retriever optimization as a primary lever, while evaluation benchmarks continue to reveal performance gaps that motivate future research. As the field progresses, the insights from mRAG will likely inform broader RAG design, particularly in multi-domain and multi-cultural applications, making it an essential area of study for the next generation of language technologies.

\subsection{Domain-Specific RAG Applications and Specialized Adaptations}

Domain-specific RAG systems represent one of the most consequential and rapidly expanding frontiers of retrieval-augmented generation, where generic architectures are systematically adapted to meet the unique epistemic, regulatory, and operational demands of high-stakes fields. Just as multilingual RAG required rethinking retrieval and generation components to handle language diversity, domain-specific RAG demands a similar reconceptualization of what constitutes relevant evidence, how that evidence should be retrieved, and how it must be integrated to support decision-making that carries real-world consequences. This subsection surveys tailored RAG systems across healthcare, legal, finance, chemistry, and enterprise settings, with a focus on the specialized adaptation strategies that distinguish these systems from their general-purpose counterparts.

The healthcare domain exemplifies how domain-specific constraints fundamentally reshape the RAG pipeline. Diagnostic errors carry direct patient consequences, and the heterogeneity of medical data---spanning clinical notes, imaging studies, structured laboratory results, and biomedical literature---demands sophisticated retrieval architectures that go beyond text-only processing. \cite{ref388} exemplifies the multimodal turn in clinical RAG by integrating cross-modal attention fusion techniques to process both textual and visual data from PubMed articles, targeting Alzheimer's Disease case studies. This system demonstrates that clinical retrieval cannot be confined to text alone; visual information embedded in biomedical literature carries diagnostic significance that text-only pipelines would discard. Similarly, \cite{ref389} pushes this paradigm further by retrieving and reasoning over document page images rather than OCR'd text, leveraging a sharded MapReduce LLM filter and a vision-language model that iteratively refines queries and accumulates evidence in a memory bank. The authors report a +5.8 point gain over a no-retrieval baseline and +1.0 point from page-image versus text-chunk retrieval, empirically validating the value of preserving visual document structure in medical reasoning. For medical vision-language tasks requiring retrieval across heterogeneous sources, \cite{ref390} introduces Modality-specific CLIPs for effective report retrieval and a Multi-corpora Query Generator that tailors queries to diverse text corpora, followed by Heterogeneous Knowledge Preference Tuning to achieve cross-modality and multi-source knowledge alignment. This approach directly addresses the challenge that medical multimodal RAG systems often fail when retrieval sources are conceptually disparate, such as radiology reports versus clinical guidelines.

Beyond architectural innovations in evidence representation, medical RAG also requires strategies that emulate clinical reasoning processes. \cite{ref391} integrates both explicit clinical knowledge and implicit case-based experience, enhancing retrieval precision by allocating conceptual tags to queries and knowledge sources and employing a hybrid retrieval mechanism across relevant knowledge and patient cases. The Med-TextGrad module using multi-agent textual gradients ensures final outputs adhere to retrieved knowledge and patient queries. In a related vein, \cite{ref392} proposes a self-training approach that equips LLMs with joint capabilities of question answering and question generation for domain adaptation, prompting the model to generate diverse domain-relevant questions from unlabeled corpora and filtering for high-quality synthetic examples. Across 11 datasets spanning three domains, SimRAG outperforms baselines by 1.2\%--8.6\%, demonstrating the viability of self-improvement when domain-specific labeled data is scarce---a challenge particularly acute in specialized fields where expert annotation is expensive and time-consuming. The iterative follow-up capability is further explored in \cite{ref393}, where the i-MedRAG system allows LLMs to ask follow-up queries based on prior information-seeking attempts, outperforming all existing prompt engineering and fine-tuning methods on GPT-3.5 with an accuracy of 69.68\% on the MedQA dataset. This system reveals that complex clinical vignettes often require multiple rounds of information-seeking, a capability that single-shot RAG pipelines fundamentally lack.

The legal domain presents a different set of challenges rooted in the requirements of precision, precedential grounding, and multi-turn interaction. Unlike the medical domain where multimodal evidence is paramount, legal RAG prioritizes the faithful alignment of generated responses with specific legal provisions and the ability to sustain coherent reasoning over extended consultations. \cite{ref394} represents the first benchmark specifically designed to evaluate RAG systems for multi-turn legal consultations, comprising 1,013 multi-turn dialogue samples and 17,228 candidate legal articles, with each sample annotated by legal experts across five rounds of progressive questioning. The companion LexiT toolkit provides a comprehensive implementation of RAG components tailored for the legal domain, while an LLM-as-a-judge evaluation pipeline enables detailed assessment. The fine-grained evaluation of legal RAG is further advanced by \cite{ref395}, which introduces ClaimRAG-LAW, a dataset supporting French and English that targets both legal experts and non-experts. This work applies a claim-level evaluation framework revealing limitations in retrieval, generation, and claim-level analysis, highlighting that standard document-level metrics obscure critical failures in legal reasoning. General-purpose legal RAG systems, the authors find, still hallucinate at varying rates even when retrieval succeeds, necessitating fine-grained evaluation that separates retrieval from generation performance---a lesson that parallels the language preference biases observed in mRAG, where system-level metrics often mask underlying deficiencies. The challenge of legal RAG is compounded by the need for multi-perspective views. \cite{ref396} introduces MVRAG, a multi-view RAG framework for knowledge-dense domains that utilizes intention-aware query rewriting from multiple domain viewpoints to enhance retrieval precision. On legal and medical case retrieval tasks, MVRAG demonstrates significant improvements in recall and precision, addressing the limitation that previous multi-view retrieval efforts focused solely on different semantic forms of queries while neglecting the expression of specific domain knowledge perspectives.

In the financial sector, RAG systems must contend with highly heterogeneous document structures---including dense narrative text, structured tables, and complex figures---alongside evolving regulatory standards and proprietary data constraints. These requirements mirror the cross-lingual challenges of mRAG in that evidence must be integrated from disparate sources with varying degrees of reliability and accessibility. \cite{ref397} addresses three critical gaps in financial RAG assessment: a fully modular architecture where components can be dynamically interchanged during runtime, a document-centric evaluation paradigm generating domain-specific QA pairs from newly ingested financial documents, and an intuitive interface bridging research-implementation divides. The evaluation quantifies retrieval efficacy and response quality, revealing significant performance variations across configurations and providing evidence for the necessity of modular benchmarking in a domain where no single configuration dominates. For compliance-critical financial workflows, \cite{ref398} proposes a multi-aspect framework with a multimodal pre-processing pipeline that unifies diverse data formats and generates chunk-level metadata summaries, a multi-path sparse-dense retrieval system augmented with query expansion (HyDE) and metadata-aware semantic search, and a domain-specialized re-ranking module fine-tuned via Direct Preference Optimization (DPO). FinSage achieves 92.51\% recall on 75 expert-curated questions and surpasses the best baseline on FinanceBench by 24.06\% in accuracy, and has been deployed as a financial question-answering agent serving over 1,200 users. The challenge of proprietary data in finance is explicitly addressed by \cite{ref399}, which employs task-aware prompt engineering, hybrid retrieval integrating open-source and proprietary data, and a four-tier task classification (explicit factual, implicit factual, interpretable rationale, and hidden rationale). The framework incorporates multi-layered security mechanisms including differential privacy, data anonymization, and role-based access controls, alongside real-time compliance monitoring through automated regulatory validation systems.

The chemistry domain presents a unique set of requirements due to the highly specialized, dynamic nature of chemical knowledge spanning molecular structures, reactions, and properties. \cite{ref400} introduces ChemRAG-Bench, a comprehensive benchmark integrating heterogeneous knowledge sources including scientific literature, PubChem, PubMed abstracts, textbooks, and Wikipedia entries. The accompanying ChemRAG-Toolkit supports five retrieval algorithms and eight LLMs, demonstrating that RAG yields an average relative improvement of 17.4\% over direct inference methods. Crucially, this work conducts in-depth analyses on retriever architectures, corpus selection, and the number of retrieved passages, culminating in practical recommendations that guide RAG deployment in chemistry. The findings reveal that retrieval-augmented generation is not uniformly beneficial across all chemical tasks; optimal configurations depend on the nature of the target question and the composition of the external corpus---an insight that echoes the mRAG finding that retrieval strategies must be tuned to the specific linguistic and epistemic characteristics of the target setting.

The enterprise and industrial sectors represent another critical deployment frontier for domain-specific RAG, where the lessons from healthcare, legal, and financial applications converge into practical engineering challenges. \cite{ref401} presents five domain-specific RAG applications developed for real-world scenarios across governance, cybersecurity, agriculture, industrial research, and medical diagnostics. Each system incorporates multilingual OCR, semantic retrieval via vector embeddings, and domain-adapted LLMs, and a web-based evaluation with 100 participants assessed systems across six dimensions: ease of use, relevance, transparency, responsiveness, accuracy, and likelihood of recommendation. The resulting twelve key lessons highlight technical, operational, and ethical challenges affecting RAG reliability and usability in practice. \cite{ref402} advances the enterprise search paradigm by enhancing semantic retrievers through a metadata generation pipeline and hybrid query indexes using dense and sparse vectors. By leveraging key concepts, topics, and acronyms, this approach creates metadata-enriched semantic indexes and boosted hybrid queries, achieving 82\% retrieval accuracy and 77\% RAG accuracy on PubMedQA without fine-tuning---surpassing prior zero-shot RAG benchmarks and rivaling fine-tuned models. Domain-specific knowledge organization is further advanced by \cite{ref56}, which leverages domain-specific documents as the primary knowledge source, integrating text, images, and tables to construct a multimodal knowledge graph spanning conceptual and instance layers. The framework introduces semantic pruning and structured subgraph retrieval, combining knowledge graph context with vector retrieval to produce more reliable responses, demonstrating superior performance in domain-specific question answering through Langfuse multidimensional scoring. For biomedical domains requiring reasoning over molecular interactions, \cite{ref55} provides a QA dataset on polypharmacy composed of a multimodal knowledge graph with text and molecular structure, revealing that existing LLMs struggle to answer complex questions unless given necessary background data---signaling the essential role of strong RAG frameworks in scientific reasoning.

Synthesizing across these domains, several patterns emerge in the adaptation of RAG to specialized fields, patterns that also resonate with the architectural and bias-mitigation strategies discussed in the context of multilingual RAG. First, domain-specific RAG consistently requires richer representations of evidence than text-only pipelines can provide, whether through visual document retention in medicine \cite{ref389}, multimodal knowledge graphs in enterprise and molecular domains \cite{ref56}, or metadata-enriched hybrid indexes \cite{ref402}. Second, domain-specific retrieval strategies must move beyond similarity-based ranking toward intention-aware and multi-view approaches that capture the domain's epistemic structure \cite{ref396}, just as mRAG systems must move beyond language-agnostic ranking to account for cross-lingual utility. Third, self-improving and iterative mechanisms are increasingly critical for domains where labeled data is scarce and reasoning is multi-step, paralleling the way mRAG systems employ iterative query refinement and translation quality assessment to handle low-resource languages \cite{ref392}. Fourth, fine-grained evaluation that distinguishes retrieval, generation, and claim-level performance is essential for identifying where domain-specific systems actually fail, mirroring the debiased metrics and comprehensive benchmarks developed for multilingual settings \cite{ref395}. Finally, the integration of domain-specific reasoning frameworks---whether case-based analogy in medicine \cite{ref391}, task-aware prompt engineering in finance \cite{ref399}, or multi-view query rewriting in law and medicine \cite{ref396}---represents the defining characteristic of state-of-the-art domain-specific RAG. As these systems continue to evolve, the boundary between general-purpose RAG and domain-specialized RAG will likely blur, with cross-fertilization of techniques such as multimodal knowledge graphs, self-improving adaptation, and preference-based re-ranking driving advances across all high-stakes application areas, much as insights from multilingual RAG are informing broader RAG design in multi-domain and multi-cultural contexts.

\section{Evaluation, Benchmarking, and Robustness}

\subsection{Evaluation Methodologies and Metrics for RAG Systems}

Evaluating Retrieval-Augmented Generation (RAG) systems is a multifaceted challenge that extends beyond simple accuracy measurements, requiring a structured assessment of both the retrieval component and the generation component, as well as their interaction. A fundamental distinction in this domain is between component-level evaluation, which isolates the performance of the retriever or generator, and end-to-end assessment, which measures the quality of the final system output. Traditional information retrieval metrics are often repurposed for evaluating the retrieval stage, but this practice is not without problems. Classical metrics like nDCG, MAP, and MRR were designed under the assumption that human users sequentially examine documents with diminishing attention to lower ranks. However, in RAG systems, the retrieved documents are consumed by a large language model that processes the entire context as a whole, not sequentially. Furthermore, these metrics do not account for the detrimental effect of related-but-irrelevant documents, which do not merely get ignored but can actively degrade generation quality. To bridge this gap, utility-based annotation schemas and new metrics like UDCG (Utility and Distraction-aware Cumulative Gain) have been proposed, which not only quantify the positive contribution of relevant passages but also the negative impact of distracting ones, using an LLM-oriented positional discount to better correlate with end-to-end answer accuracy \cite{ref403}. Beyond ranking metrics, the concept of retrieval utility has been re-examined through the lens of semantic information gain. For instance, Semantic Perplexity (SePer) measures the reduction in an LLM's internal belief about the correctness of the retrieved information post-retrieval, providing a more direct measure of how useful the retrieval is for the generation task \cite{ref404}.

For evaluating the generation component, a rich ecosystem of metrics has emerged to capture different dimensions of quality, primarily faithfulness, relevance, and correctness. Reference-based metrics, which compare the generated answer to a golden answer, and reference-free metrics, which assess quality without ground truth, both play crucial roles. Given the complexity of natural language, evaluation frameworks have increasingly turned to LLM-as-a-Judge approaches. The RAGAS (Retrieval Augmented Generation Assessment) framework is a foundational example, providing a suite of metrics for a reference-free evaluation of RAG pipelines by considering the ability of the retrieval system to identify relevant context, the LLM's ability to exploit these passages faithfully, and the overall quality of the generation \cite{ref405}. Building on this, frameworks like CCRS (Contextual Coherence and Relevance Score) propose a novel suite of metrics that uses a single, powerful, pretrained LLM as a zero-shot, end-to-end judge, evaluating aspects such as contextual coherence, question relevance, information density, answer correctness, and information recall \cite{ref406}. While these approaches have advanced the field, they have also been critiqued for their lack of diagnostic depth. Frameworks such as RAGChecker and CoFE-RAG address this by moving beyond a single holistic score. RAGChecker provides a fine-grained evaluation framework with a suite of diagnostic metrics for both retrieval and generation modules, showing better correlation with human judgments than other metrics and revealing insightful patterns and trade-offs in RAG architecture choices \cite{ref407}. Similarly, CoFE-RAG (Comprehensive Full-chain Evaluation Framework) facilitates a thorough evaluation across the entire pipeline---including chunking, retrieval, reranking, and generation---using multi-granularity keywords to assess the retrieved context and provide unique insights into how RAG systems handle diverse data scenarios \cite{ref408}.

A significant recent trend in RAG evaluation is the incorporation of uncertainty quantification to assess the trustworthiness and reliability of systems, moving beyond mere correctness. Traditional evaluations concentrate on correctness, which does not fully capture the impact of retrieval on LLM uncertainty. The URAG benchmark is designed to assess the uncertainty of RAG systems across various fields by reformulating open-ended generation tasks into multiple-choice question answering, allowing for principled uncertainty quantification via conformal prediction. It measures performance through both accuracy and prediction-set sizes, revealing that accuracy gains often coincide with reduced uncertainty, but this relationship can break under retrieval noise \cite{ref409}. Another promising direction for uncertainty is the application of conformal prediction to the retrieval process itself. CONFLARE (CONFormal LArge language model REtrieval) introduces a four-step framework that applies conformal prediction to quantify retrieval uncertainty. By calibrating similarity scores on a set of questions answerable from the knowledge base, it establishes a threshold to retrieve all chunks with similarity exceeding this threshold during inference, ensuring the true answer is captured in the context with a (1-alpha) confidence level \cite{ref410}. Similarly, other work has applied conformal prediction to ensure coverage of relevant evidence while pruning irrelevant context, demonstrating reliable, coverage-controlled context reduction in RAG \cite{ref169}.

The evaluation landscape is further enriched by specialized frameworks that focus on specific challenges. For instance, given the growing importance of citations, evaluation frameworks have been developed to assess the quality of generated citations and the grounding of claims. A comparative evaluation framework has been proposed to assess the effectiveness of faithfulness metrics in distinguishing between different levels of citation support---full, partial, and no support---revealing that no single metric consistently excels across all evaluations \cite{ref411}. Furthermore, metrics like Recall Conversion Rate (RCR) have been introduced to diagnose the gap between retrieval performance and downstream reasoning accuracy, quantifying the contribution of retrieval to reasoning and highlighting that improvements in retrieval recall do not always translate to commensurate gains in downstream reasoning quality \cite{ref412}.

Finally, a critical line of research focuses on the validity and reliability of the evaluation metrics themselves. The quality of a RAG system is fundamentally limited by the retrieval component, and evaluating the retrieval quality is often done with heuristically constructed query sets that introduce a hidden intrinsic bias. Frameworks such as semantic stratification ground evaluation in corpus structure by organizing documents into an interpretable global space of entity-based clusters and systematically generating queries for missing strata, yielding formal semantic coverage guarantees and interpretable visibility into retrieval failure modes \cite{ref413}. On a component level, meta-evaluation frameworks are used to assess the calibration and discrimination capabilities of the judge models themselves. For example, GroUSE, a meta-evaluation benchmark of 144 unit tests, reveals that existing automated RAG evaluation frameworks often overlook important generator failure modes, even when using GPT-4 as a judge, and shows that correlation with GPT-4 is an incomplete proxy for the practical performance of judge models \cite{ref414}. Overall, the field is shifting from aggregate score-centric evaluation toward more diagnostic, interpretable, and uncertainty-aware methodologies that can provide actionable guidance for improving RAG systems \cite{ref415}.

\subsection{Benchmarking Frameworks and Diagnostic Analysis}

The proliferation of Retrieval-Augmented Generation systems has necessitated a corresponding evolution in evaluation methodologies, moving beyond the accuracy-centric metrics discussed previously toward interpretable, stage-decomposed diagnostics that can isolate performance bottlenecks and guide actionable improvements. While the evaluation landscape reviewed above has largely focused on measuring \emph{how well} a RAG system performs---through component-level metrics, end-to-end frameworks, and uncertainty quantification---a complementary line of research has emerged that emphasizes \emph{why} a system underperforms and \emph{where} in the pipeline the root cause resides. Rather than treating RAG as a monolithic black-box, these diagnostic frameworks decompose the pipeline into its constituent stages---retrieval, augmentation, and generation---to provide explainable, actionable insights that extend beyond aggregate scores.

A foundational contribution in this direction is the RAGChecker framework, which incorporates a suite of diagnostic metrics for both the retrieval and generation modules \cite{ref407}. Its meta-evaluation demonstrates significantly better correlations with human judgments than other metrics, emphasizing the necessity of fine-grained, module-level assessment over holistic scoring. Similarly, the PRGB benchmark introduces a multi-level fine-grained evaluation centered on document utilization, using an innovative placeholder-based approach to decouple the contributions of the LLM's parametric knowledge from the external retrieved knowledge \cite{ref416}. This methodology allows for a nuanced understanding of an LLM's role and its limitations in error resilience and context faithfulness.

The shift toward interpretable evaluation is further exemplified by RAGVUE, a diagnostic and explainable framework that decomposes RAG behavior into retrieval quality, answer relevance and completeness, strict claim-level faithfulness, and judge calibration, with each metric accompanied by a structured explanation \cite{ref417}. This transparency is critical for identifying whether errors are grounded in retrieval or reasoning. RAGXplain takes this a step further, structuring evaluation around a `Metric Diamond' that connects user input, retrieved context, generated answer, and ground truth via six diagnostic dimensions, and using LLM reasoning to produce natural-language failure-mode explanations and prioritized interventions \cite{ref415}. Its empirical validation across several QA benchmarks shows that following its recommendations can consistently improve pipeline performance, highlighting the value of turning diagnostics into prescriptive guidance---a natural complement to the uncertainty-aware and meta-evaluation approaches discussed in the previous section.

Beyond generic frameworks, significant efforts have focused on domain-specific and scenario-specific benchmarking. RAGEval provides a framework for generating high-quality, scenario-specific datasets through a schema-based pipeline, accompanied by three novel metrics for factual accuracy: Completeness, Hallucination, and Irrelevance \cite{ref418}. This approach enables rigorous evaluation in specialized domains where existing public benchmarks are inadequate. In the industrial sector, FAB-Bench addresses the complexities of evaluating RAG in semiconductor manufacturing, defining six diagnostic metrics that measure factual accuracy, contextual utilization, completeness, retrieval relevance, technical depth, and reasoning consistency \cite{ref419}. Notably, FAB-Bench studies how retrieval precision and generative fidelity co-evolve with context window size, identifying distinct context-scaling behaviors and pinpointing attention dilution as a primary mechanism for performance degradation at extreme lengths. For enterprise multi-turn scenarios, a case-aware LLM-as-a-Judge evaluation framework has been proposed, which evaluates each turn using eight operationally grounded metrics that separate retrieval quality, grounding fidelity, answer utility, precision integrity, and case/workflow alignment, thus capturing failure modes conventional benchmarks miss \cite{ref420}.

A key trend in modern diagnostics is the emphasis on root-cause analysis and the isolation of specific component failures. The XRAG benchmarking suite categorizes RAG components into four phases---pre-retrieval, retrieval, post-retrieval, and generation---to systematically identify potential failure points \cite{ref421}. This granular approach facilitates targeted optimization for known failure modes. The FATHOMS-RAG benchmark differentiates itself by evaluating the entire pipeline, including its ability to ingest and reason over multiple modalities, using a phrase-level recall metric and a nearest-neighbor classifier to identify potential hallucinations, thereby providing a holistic view that includes cross-modal and cross-document synthesis \cite{ref422}.

The need for controlled and reproducible comparisons of RAG design choices has also led to the formulation of RAG as an architecture search problem. The RAISE framework provides a comprehensive benchmark and search engine that standardizes the evaluation of optimization methods for RAG pipelines under defined search spaces and budgets \cite{ref144}. Its experiments across multiple text and multimodal datasets show that optimization performance is highly task-dependent, cautioning against the assumption of universally superior strategies. Similarly, RAGSmith treats RAG design as an end-to-end architecture search over a vast configuration space, finding that a genetic search consistently outperforms naive baselines and that improvement magnitude correlates with question type, providing practical, domain-aware guidance for assembling effective RAG systems \cite{ref145}. Finally, RAGPulse underscores the importance of real-world workload traces in benchmarking, providing an open-source dataset of user queries that reveals significant temporal locality and a highly skewed hot document access pattern \cite{ref63}. This trace provides a high-fidelity foundation for developing and validating optimization strategies, bridging the gap between academic research and real-world deployment. Collectively, these frameworks demonstrate a clear shift towards interpretable, stage-decomposed evaluation that not only measures performance but also explains it, offering a roadmap for more reliable and actionable RAG system development. This diagnostic foundation becomes particularly critical when systems are deployed in adversarial environments, where understanding component-level vulnerabilities---as explored in the following section---is essential for designing robust defenses.

\subsection{Robustness and Adversarial Attack Evaluation}

The robustness of Retrieval-Augmented Generation (RAG) systems against adversarial manipulation is a critical concern, as these systems are exposed to new attack surfaces through their reliance on external knowledge bases. The integration of retrieved evidence introduces vulnerabilities that are not present in standalone LLMs, including corpus poisoning, indirect prompt injection, and various forms of textual perturbation. A systematic understanding of these threats requires a formalized threat model that categorizes adversary capabilities and access levels. For instance, a comprehensive threat model for RAG systems has been proposed that defines a taxonomy of adversary types based on their access to model components and data, along with formal threat vectors such as document-level membership inference and data poisoning \cite{ref79}. This formalization is essential for establishing a rigorous foundation for evaluating and designing robust RAG systems. Building on this, the SLOT taxonomy provides a four-axis framework---attack Surface, defense Layer, Objective (following CIA properties), and Target---to organize the literature on RAG security, mapping attacks and defenses onto a six-stage knowledge-access pipeline \cite{ref423}.

Adversarial attack methodologies against RAG systems have evolved significantly, moving from simplistic text injection to sophisticated, purpose-built optimization frameworks. One of the foundational and most prevalent threat vectors is corpus poisoning, where an attacker injects malicious documents into the knowledge base to manipulate the output for a target query. The PoisonedRAG framework formalized this as an optimization problem and demonstrated that by injecting as few as five poisoned texts into a database with millions of entries, an attacker can achieve a high success rate in forcing the LLM to generate a chosen answer, a finding that highlights the fragility of the retrieval pipeline \cite{ref424}. Subsequent work has introduced more practical attack paradigms that address the constraints of real-world deployments. For example, CorruptRAG was designed to be a practical poisoning attack where the adversary injects only a single poisoned text per query, a more feasible and stealthy approach compared to earlier strategies that assumed the ability to inject a sufficient number of texts to dominate natural results \cite{ref425}. In a similar vein, HijackRAG introduced the ``retrieval prompt hijack attack,'' demonstrating the ability to manipulate the retrieval mechanisms themselves to generate attacker-predetermined answers, with both black-box and white-box strategies available depending on the adversary's knowledge \cite{ref426}.

Attack methodologies are generally categorized by their access to the target system, which directly shapes the design and expected efficacy of the attack. White-box attacks, such as TPARAG, leverage a surrogate LLM to generate and iteratively optimize malicious passages at the token level, jointly considering both the retrieval and generation stages to ensure both retrievability and high attack success, even against more robust black-box RAG systems \cite{ref427}. Another advanced white-box approach, Joint-GCG, takes this further by unifying gradient-based attacks across both the retriever and generator models through innovations like Cross-Vocabulary Projection and Adaptive Weighted Fusion, achieving higher attack success rates and showing unprecedented transferability to unseen models \cite{ref428}. In contrast, black-box attacks operate under more realistic constraints where the adversary has no access to model internals. The RIPRAG attack framework, for instance, treats the target RAG system as a black box and uses Reinforcement Learning from Black-box Feedback (RLBF) to optimize poison generation, effectively improving attack success rates against complex systems without requiring internal knowledge \cite{ref429}. Beyond these access-based distinctions, more specific attacks have been developed to target the reasoning process of agentic RAG systems, which involve multi-step planning and tool use. KidnapRAG is a sequential black-box poisoning attack that hijacks multi-step reasoning by injecting role-specific documents (Bait, Chain-Link, Mal-Ins) to attract initial retrieval, induce query reformulation, and provide attacker-controlled evidence, respectively \cite{ref430}. Furthermore, some attacks are designed to disable the self-correction ability of LLMs, a key defense mechanism in real-world scenarios; DisarmRAG compromises the retriever itself (rather than just poisoning the knowledge base) to inject anti-self-correction instructions into the context, achieving high success rates even under detection defenses \cite{ref431}.

The attack surface extends beyond simple fact manipulation to include opinion manipulation and system availability, reflecting the diverse objectives adversaries may pursue. Topic-FlipRAG and FlippedRAG address the more practical scenario of adversarial opinion manipulation on controversial topics, where LLMs are required to reason and synthesize multiple perspectives, making them particularly susceptible to systematic knowledge poisoning that can alter the opinion polarity of generated content and potentially shift user cognition \cite{ref432,ref433}. The stealthiness of these attacks is a growing concern, with research proposing methods that craft human-imperceptible adversarial examples to retrieve target documents and influence final answers \cite{ref434}. Furthermore, attackers are not limited to manipulating the final answer; they can also target system efficiency. The Retrieval-Augmented Inference Cost Attack (RA-ICA) injects malicious documents into the knowledge corpus to cause an abnormal increase in token consumption during the inference phase, presenting a denial-of-service threat by exhausting computational resources \cite{ref435}. In multimodal settings, adversarial attacks also pose a significant risk, with frameworks like MM-PoisonRAG demonstrating that both localized, query-specific poisoning (LPA) and a single, untargeted global adversarial injection (GPA) can disrupt multimodal RAG systems, with GPA capable of collapsing generation quality to 0\% accuracy, bypassing existing defenses \cite{ref436}.

A crucial aspect of evaluating RAG robustness is understanding how system components interact under attack, as this interaction often determines whether an attack succeeds or fails in practice. A rigorously controlled empirical study on adversarial poisoning in RAG introduced a taxonomy of context types---adversarial, untouched, and guiding---and found that adversarial documents substantially degrade alignment with ground truth, but that robustness can be preserved when helpful evidence is also present in the retrieval pool \cite{ref213}. This finding is echoed by research showing that while ``skeptical prompting'' can activate LLMs' internal reasoning to partially defend against adversarial passages, its effectiveness depends strongly on the model's reasoning capacity \cite{ref213}. A large-scale factorial study analyzing 432 configurations across datasets, retrievers, and generation models confirmed that poisoning vulnerability is not attributable to a single component but arises from the interaction of retrieval, generation, and knowledge-base configuration \cite{ref219}. The study also revealed that retriever architecture and retrieval depth are strong factors affecting poisoning exposure, and that replicating poisoned content across multiple databases amplifies adversarial influence \cite{ref219}. The realism of the evaluation is also critical, as a study on realistic multi-stage retrieval pipelines showed that many existing poisoning attacks degrade substantially after reranking, despite achieving high retrieval-stage relevance, and proposed a Chunk-aware and Rerank-Consistent Poisoning (CRCP) framework to address the granularity mismatch \cite{ref437}. Similarly, end-to-end evaluations using frameworks like Rag 'n Roll show that existing attacks are often optimized to boost ranking but do not translate into reliable attacks, settling around a 40\% success rate across various configurations, and that the effectiveness of unoptimized documents can be increased by deploying multiple of them \cite{ref70}.

Finally, the evaluation of robustness must consider the diversity of attack scenarios, including adversarial perturbations at a character level that simulate real-world noise. The Genetic Attack on RAG (GARAG) was introduced to assess vulnerability to low-level noisy perturbations, revealing that minor textual inaccuracies can devastate the performance of RAG components and their synergy \cite{ref438}. Benchmarking efforts are also vital for standardizing robustness evaluation and enabling fair comparisons across different defenses. A comprehensive benchmark framework for poisoning attacks on RAG covers 5 QA datasets, 13 attack methods, and 7 defense mechanisms, revealing that while attacks perform well on standard datasets, their effectiveness drops on expanded versions, and that advanced RAG architectures remain susceptible \cite{ref439}. Even under multi-adversary conditions, where high-value queries attract multiple attackers with conflicting objectives, a standardized framework called PoisonArena has been developed to evaluate attacks and defenses, revealing that many strategies that succeed in single-attacker scenarios degrade markedly under competition \cite{ref440}. The potential for building defense mechanisms based on these insights is promising; for instance, a network-based defense inspired by robust statistics, which aggregates a robust subset of retrieved texts to generate the final answer, has been shown to reduce the attack success rate to as low as 1\% while maintaining an accuracy of 71\% \cite{ref441}. Collectively, these attack analyses and empirical evaluations lay the groundwork for the defense strategies discussed next, which are organized along the retrieval, generation, and system-level dimensions of the RAG pipeline.

\subsection{Defense Mechanisms and Resilience Enhancement}

The vulnerability of Retrieval-Augmented Generation (RAG) systems to adversarial attacks---particularly corpus poisoning, misinformation pollution, and prompt injection---has motivated a rich body of defense research. These defenses operate across the pipeline, spanning retrieval-side filtering, generation-side interventions, training-based robustness, and system-level orchestration, directly countering the attack methodologies and empirical findings presented previously. A central challenge uniting these approaches is the trade-off between defense efficacy, computational overhead, and task performance, particularly in high-stakes domains where both security and answer quality are paramount.

\textbf{Retrieval-Side Defenses.} A first line of defense targets the retrieval stage, aiming to prevent malicious documents from entering the context before generation. One prominent strategy is adversarial filtering, which screens retrieved documents for statistical or semantic anomalies. For example, \cite{ref442} exploit a newly proposed property that distinguishes adversarial texts from clean ones in the knowledge source, enabling effective filtering against PoisonedRAG-style attacks. Similarly, \cite{ref443} operate directly on the retriever: RAGPart leverages document partitioning to mitigate the effect of poisoned points, while RAGMask identifies suspicious tokens based on significant similarity shifts under targeted token masking, providing lightweight, generation-model-agnostic protection. Another retrieval-side approach, \cite{ref444}, stress-tests query--passage pairs under mild randomized perturbations and extracts probe gradients to derive instability signals---representational consistency and dispersion risk---for reranking, preserving passage content and requiring no retraining.

Cluster-based detection offers an alternative pattern-recognition strategy. \cite{ref445} implements a two-stage defense: a cluster filtering strategy detects potential attack patterns, followed by an LLM-based self-assessment to identify malicious documents and resolve inconsistencies. Importantly, TrustRAG is a plug-and-play, training-free module that works with both open- and closed-source models. On the more theoretical end, \cite{ref446} adopts a graph-theoretic perspective, identifying a ``consistent majority'' among retrieved documents via a Maximum Independent Set (MIS) algorithm that prioritizes higher-reliability documents, providing provable robustness guarantees against bounded adversarial corruption. \cite{ref447} capitalizes on a key discriminative pattern---poisoned documents exhibit significantly stronger alignment between their backward rankings and the query's forward ranking---to build a dual-signal framework that simultaneously achieves efficiency and robustness. Finally, \cite{ref448} proposes the Normalized Passage Attention Score (NPAS) and an Attention-Variance Filter (AV Filter) that flag anomalous passages by analyzing attention weights, yielding up to \textasciitilde20\% higher accuracy than baseline defenses. \cite{ref185} complements these with a fine-grained knowledge pruning stage that filters irrelevant context before generation, demonstrating that retrieval-side filtering can be informed by the model's latent representations.

\textbf{Generation-Side Interventions.} When malicious or noisy documents evade retrieval-side filters, defenses must operate during generation itself. A key insight from \cite{ref213} is that ``skeptical prompting'' can activate LLMs' internal reasoning, enabling partial self-defense against adversarial passages---though its effectiveness depends strongly on the model's reasoning capacity. This finding underscores the utility of generation-side interventions that leverage the model's parametric knowledge as a check on retrieved evidence.

NLI-based claim verification provides a more systematic approach. \cite{ref212} builds an NLI-based cross-document claim graph and applies a Defended-Claims Gate that conditions generation on a consistent, multi-source supported subset, or blocks retrieval entirely and answers parametrically. This training-free, model-agnostic defense consistently restores factuality under mixed and fully polluted evidence across multiple long-form QA benchmarks. Conflict-aware attention represents a more fine-grained intervention. \cite{ref185} discovers that conflicting and aligned knowledge states are linearly separable in the model's latent space and that contextual noise increases the entropy of these representations. Building on this, ProbeRAG operates in three stages---fine-grained knowledge pruning, latent conflict probing, and conflict-aware attention modulation---to steer the model toward faithful context integration. The attention mechanism itself can also be structurally modified: \cite{ref449} introduces Sparse Document Attention RAG (SDAG), a block-sparse attention mechanism that disallows cross-attention between retrieved documents, requiring only a minimal inference-time change to the attention mask with no fine-tuning. SDAG demonstrates clear merits when integrated with state-of-the-art defenses, suggesting that architectural attention modifications can complement other defensive layers.

\textbf{Training-Based Approaches.} Beyond inference-time interventions, training-based defenses aim to make the system intrinsically robust. \cite{ref450} combines adversarial retriever fine-tuning with a Zero-Knowledge Inference Patch (ZKIP)---a label-free, black-box filter that performs leave-one-out decoding and scores documents by the semantic shift and entropy change their removal induces. ZKIP drives attack success rate to 0.000 in every defended configuration while keeping Recall@5 within 0.03 of the clean baseline. \cite{ref451} proposes RAPS-DA, a regime-aware peer specialization framework that divides conflicts into Grounding, Arbitration, and Resistance regimes, training one peer specialist per regime from a shared base model and routing samples to their regime-matched peer. This approach surpasses prompting, decoding, fine-tuning, and RL baselines across five conflict scenarios. Preference optimization has also been adapted for robustness. While not explicitly a defense against adversarial attacks, \cite{ref20} introduces a Differentiable Data Rewards method that end-to-end trains RAG systems by aligning data preferences between modules, enabling the generator to better extract key information and mitigate conflicts between parametric and external knowledge---capabilities that indirectly strengthen resilience. Structured critic learning offers yet another avenue. \cite{ref452} proposes an ensemble framework (EPD) that aggregates outputs from a knowledge-injected LLM, a base LLM, and a dedicated judge model, reducing membership inference attack success by up to 526.3\% for RAG while maintaining answer quality. Importantly, \cite{ref453} cautions that the marginal robustness benefit of sophisticated training strategies decreases substantially as model capacity increases---stronger models naturally exhibit better confidence calibration and attention patterns---suggesting that simplified pipelines may suffice for powerful models.

\textbf{System-Level Frameworks.} Given the diversity of attacks and the need for practical deployment, several frameworks adopt a system-level view. \cite{ref454} operates during the post-retrieval phase, using lightweight machine learning to detect and filter adversarial content without additional model training, reducing attack success rate against the Gemini model from 0.89 to as low as 0.02. \cite{ref455} uses sentence-level processing and bait-guided context diversity detection to identify malicious content without relying on LLM internal knowledge, delivering state-of-the-art security with plug-and-play deployment while improving clean-scenario performance and reducing token consumption by 48--80\%. \cite{ref456} places three independent, formally specified rings across the pipeline---an Ingest Guard screening documents for lexical and statistical poisoning signatures, a Retrieval Scorer re-ranking by provenance and consistency-weighted trust scores, and a Cross-LLM Consensus stage polling three architecturally diverse language models---reducing attack success from roughly 91\% to 13\% while preserving benign accuracy. \cite{ref457} addresses the security-utility trade-off directly: a Sentinel detects anomalous retrieval behavior, after which a Strategist selectively deploys only the defenses warranted by the query context, eliminating membership inference leakage while recovering retrieval utility relative to a static defense stack. The RAG security landscape is further systematized by \cite{ref423} and \cite{ref68}, which organize defense literature along the retrieval, context construction, and generation pipeline, and by \cite{ref75}, which constructs a taxonomy from input-side (dynamic access control, homomorphic encryption, adversarial pre-filtering) and output-side (federated learning isolation, differential privacy, data sanitization) perspectives.

\textbf{Trade-offs and Open Challenges.} Across these defense strategies, consistent trade-offs emerge, directly informing the defense landscape that the following subsection will elaborate upon. First, defense efficacy versus task performance: \cite{ref457} shows that an always-on defense stack can reduce contextual recall by more than 40\%, motivating adaptive orchestration. Second, computational cost: \cite{ref450} costs k+1 generator passes per query (6× for k=5), while \cite{ref454} explicitly targets resource efficiency. Third, the stealth--attack relationship: \cite{ref448} shows that stealthy attacks that control the response must bias inference more than benign passages, inherently compromising stealth, yet \cite{ref433} demonstrates that existing defenses remain insufficient against sophisticated manipulation attacks. The benchmark \cite{ref439} further reveals that current defense techniques fail to provide robust protection against state-of-the-art attacks across diverse RAG architectures, underscoring the pressing need for more resilient and generalizable defenses. Looking forward, \cite{ref69} highlights a novel vulnerability where safety guardrails themselves can be exploited for denial-of-service attacks, suggesting that defense strategies must anticipate second-order effects. The field is converging toward adaptive, multi-layered, and provably robust frameworks that balance security with utility---an essential direction as RAG systems are increasingly deployed in high-stakes domains where both reliability and trustworthiness are non-negotiable.

\section{Security, Privacy, and Trustworthiness}

\subsection{Attack Surface and Threat Modeling for RAG Systems}

Retrieval-Augmented Generation (RAG) systems, while effectively mitigating hallucination and adapting LLMs to domain-specific knowledge, fundamentally alter the security posture of LLM-based applications. By coupling a parametric generator with a non-parametric, often externally sourced, knowledge base, RAG expands the attack surface beyond the model itself to encompass the entire retrieval, context construction, and generation pipeline. A comprehensive understanding of this expanded attack surface is the first prerequisite for designing robust and trustworthy systems. This subsection establishes a formal taxonomy of security threats unique to RAG architectures, organizing the attack surface along the RAG workflow and introducing formal threat models that categorize adversaries based on their access levels.

\textbf{A Unified Threat Taxonomy Along the RAG Pipeline}

The RAG architecture can be decomposed into three primary stages---retrieval, context construction, and generation---each presenting distinct vulnerabilities. The retrieval stage is susceptible to attacks targeting the knowledge base itself, such as data poisoning, and attacks on the retrieval mechanism, such as adversarial query manipulation. The context construction stage, where retrieved documents are combined and formatted, is vulnerable to indirect prompt injection, where malicious instructions are embedded within the retrieved content. Finally, the generation stage, where the LLM synthesizes the final response, faces threats from jailbreak attempts, data leakage, and the propagation of misinformation. This pipeline-level view is critical; as noted in a comprehensive review of RAG security, vulnerabilities are often conflated with inherent LLM flaws, whereas in RAG they are frequently introduced by the integration of external knowledge access \cite{ref423}. A more granular taxonomy, organized along the attack surface, defense layer, objective (following the CIA triad of confidentiality, integrity, and availability), and target, has been proposed to systematically structure the literature \cite{ref423}. Similarly, the security implications of each module---from indexing and retrieval to generation---must be considered holistically, as an attacker can target any point in this multi-module architecture \cite{ref75,ref68}.

\textbf{Adversarial Attacks: Data Poisoning}

Among the most significant threats to RAG systems is data poisoning, a class of attacks where an adversary injects malicious content into the knowledge base to manipulate system behavior. These attacks are particularly effective because knowledge bases are often populated from public or user-contributed sources, providing a low-barrier entry point for attackers. The overarching goal can range from causing a denial-of-service to steering opinions or leaking private data \cite{ref458,ref459}. Within this category, several sophisticated attack methodologies have been developed.

One line of work focuses on backdoor attacks that compromise the retriever or the generator. For instance, a two-stage optimization framework can craft a single malicious document that, when retrieved by a specific trigger sequence, induces a variety of adversarial objectives on the LLM output, including refusal to answer or privacy violations \cite{ref460}. Other backdoor attacks target the retrieval process itself, with a joint backdoor attack framework that constructs elaborate target contexts and trigger sets optimized via contrastive learning to manipulate LLMs in universal attack scenarios \cite{ref461}. More recent work has explored backdooring the fine-tuning process of the dense retriever itself to ensure malicious documents are retrieved for trigger queries \cite{ref462}. Researchers have also demonstrated backdoor attacks that embed triggers to enable the extraction of data from the RAG knowledge base \cite{ref463}.

Another significant line of research centers on the creation of increasingly realistic and robust poisoning attacks. While early attacks assumed favorable conditions, later work highlights the challenges of attacking practical RAG systems. A learning-to-poison framework was introduced to generate robust poisons that can survive data preprocessing and are effective even when the user's exact query is unknown, addressing the limitations of earlier attacks \cite{ref459}. Similarly, the gap between simplified evaluations and real-world retrieval pipelines has been exposed, demonstrating that attacks optimized for retrieval-stage relevance often fail after reranking and chunking \cite{ref437}. Research has also delved into the influence of various system components on poisoning susceptibility, revealing that attack success arises from the complex interaction of retriever, generator, and knowledge-base configuration rather than a single point of failure \cite{ref219}. The ``architecture matters'' finding is particularly crucial: evaluations across different RAG architectures (vanilla, agentic, multi-agent debate) under adversarial poisoning show a wide variance in attack success rates, emphasizing that architectural choices significantly impact robustness \cite{ref464}.

\textbf{Adversarial Attacks: Knowledge Base Extraction}

A distinct and severe class of threats involves knowledge base extraction, where attackers aim to exfiltrate the proprietary or confidential contents of the RAG system's retrieval database. The goal is often to steal intellectual property or replicate the functionality of a commercial application. Attack methodologies have evolved from simple prompt-based queries to sophisticated, adaptive frameworks. A relevance-based black-box attack utilizes an attacker-side open-source LLM to generate effective queries that force the RAG system to progressively leak its hidden knowledge base \cite{ref465}. Another framework provides a configurable environment for benchmarking leakage attacks, revealing that a model's instruction-following capability is a strong predictor of its leakage risk, establishing that query generation and adversarial instructions contribute independently to the risk \cite{ref466}. In black-box settings, agent-based attacks that reason from feedback to generate new adversarial queries can achieve high extraction coverage, even against commercial platforms \cite{ref467}. While some attacks rely on the LLM's instruction-following abilities, a more robust method involves backdooring the LLM during the fine-tuning phase, after which the adversary can trigger the leak and extract documents verbatim \cite{ref463}. The inherent difficulty of preventing this leakage is exacerbated by the ``RAG paradox,'' where improvements in RAG faithfulness can inadvertently increase the risk of revealing the underlying knowledge base \cite{ref466}.

\textbf{User-Query and Context-Construction Attacks}

Beyond text-based knowledge bases, the attack surface extends to other data modalities and system components. Visual document RAG systems, which use page screenshots as their knowledge base, have been shown to be vulnerable to poisoning attacks that require injecting only a single adversarial image, capable of spreading targeted disinformation or causing a denial-of-service \cite{ref468}. The data loading stage itself is a critical, often-overlooked vulnerability. Research has proposed a taxonomy of knowledge-based poisoning attacks and introduced novel threat vectors like content obfuscation and injection that target common document formats (DOCX, HTML, PDF), demonstrating high success rates against a range of data loaders \cite{ref469}. Finally, the logic of hybrid RAG pipelines, which combine vector search with graph expansion, can introduce a ``retrieval pivot'' vulnerability where a vector-retrieved seed chunk can pivot via entity links into sensitive graph neighborhoods, causing cross-tenant data leakage \cite{ref470}.

\textbf{Threat Modeling and Adversary Categorization}

To rigorously analyze these diverse threats, a formal threat model is essential. A foundational effort proposed a structured taxonomy of adversary types based on their access to model components and data, and formally defined key threat vectors such as document-level membership inference and data poisoning \cite{ref79}. Other work has categorized the attack surface from the perspective of risk management, listing primary attack vectors such as prompt injection, data poisoning, and adversarial query manipulation, and proposing a prioritized list of mitigating actions \cite{ref471}. The anatomy of threats within RAG systems is often attributed to the inherent trust placed in retrieved content and the transparency of the system, which can be exploited for both integrity and confidentiality violations, a concept known as the ``confused deputy'' problem where an LLM is tricked into performing actions it was not intended to \cite{ref472}. A key differentiator in threat modeling is the adversary's knowledge level---white-box (full knowledge of the system), gray-box (partial knowledge), and black-box (no knowledge). Many sophisticated attacks, such as a sequential poisoning attack on agentic RAG systems, strictly adhere to the black-box setting, where the attacker can only publish externally retrievable documents \cite{ref430}. Similarly, a two-stage black-box attack for document poisoning on commercial LLMs has been shown to achieve high success rates without any knowledge of the target system's internals \cite{ref473}. This formal categorization of adversaries is crucial for evaluating the practical feasibility and real-world impact of security threats and defenses.

\subsection{Privacy Attacks: Membership Inference, Data Leakage, and Side Channels}

Privacy attacks against Retrieval-Augmented Generation (RAG) systems target the two most valuable assets of the architecture: the external knowledge base and the generated outputs that draw from it. Unlike traditional attacks on parametric models, which focus on extracting training data or memorized information, RAG systems expose a unique and expanded attack surface because the retrieval database is accessed dynamically at inference time, and the retrieved content---rather than model weights or training data alone---becomes a primary source of leakage. This subsection systematically examines the primary privacy violation mechanisms, including membership inference attacks (MIAs), contextual leakage and knowledge extraction, and analyzes the system-level factors that determine their success, thereby bridging the preceding discussion of adversarial threats (e.g., data poisoning, knowledge base extraction) with the subsequent treatment of defense mechanisms.

\textbf{Membership Inference Attacks (MIAs)} aim to determine whether a specific document or data sample is present in the RAG system's retrieval database. This attack class has received significant attention due to the sensitive nature of proprietary or private corpora often stored in RAG knowledge bases, and it represents a more granular privacy violation than the large-scale knowledge base extraction discussed previously. The S\(^2\)MIA framework was among the first to systematically investigate this threat, leveraging the semantic similarity between a candidate sample and the text generated by the RAG system. The core intuition is that if a document is in the external database, the generated response will exhibit high semantic similarity to that document, enabling an attacker to infer membership with high accuracy while remaining robust against several baseline defenses \cite{ref71}. Building on this idea, the MBA (Mask-Based Membership Inference Attack) framework proposed a more explainable approach. Instead of relying solely on response similarity, MBA masks a portion of the target document and prompts the RAG system to predict the masked values. If the complete document exists in the knowledge base, the masked variant retrieves the full context, leading to accurate mask predictions that serve as a reliable membership signal. This masking strategy is more document-specific and reduces interference from unrelated retrieved documents or the LLM's internal parametric knowledge \cite{ref474}.

Recent work has further enhanced MIA efficacy through differential calibration and side-channel analysis. DCMI addresses a key limitation of prior methods: the interference of non-member-retrieved documents on RAG outputs. By generating perturbed queries for calibration, DCMI isolates the contribution of member-retrieved documents while minimizing noise from non-member documents. Experiments show that DCMI achieves 97.42\% AUC against Flan-T5-based RAG systems, significantly outperforming the MBA baseline, and maintains a 10--20\% advantage on real-world RAG platforms like Dify and MaxKB \cite{ref475}. A particularly novel attack, BudgetLeak, identifies an overlooked side channel: the generation budget---the maximum number of tokens allowed in the response. By probing responses under varying budgets, attackers observe distinct behavioral patterns between member and non-member queries, as members gain output quality more rapidly with larger budgets. This approach demonstrates that even seemingly innocuous system configuration parameters can leak membership information, underscoring the breadth of the attack surface \cite{ref476}.

Beyond these, query-efficiency and stealth have emerged as critical design goals, particularly in black-box settings where attackers must minimize their footprint. MEntA leverages natural-language entailment to maximize the information gained per query, achieving up to 0.991 AUC with only five queries across NFCorpus, SCIDOCS, and TREC-COVID, while remaining effective under state-of-the-art defenses and reducing attack cost by up to 65× compared to prior methods \cite{ref477}. The Interrogation Attack (IA) crafts natural-text queries that are answerable only with the target document's presence, enabling successful inference with just 30 queries while remaining stealthy; straightforward detectors identify adversarial prompts from existing methods up to 76× more frequently than those generated by IA, and the attack achieves a 2× improvement in TPR@1\%FPR \cite{ref478}. Additionally, E-MIA converts verifiable hard evidence---such as fine-grained details, proper nouns, and causal relations---into an exam with four objectively gradable question types, using the aggregated score as a membership signal that improves member/non-member separability while preserving natural, stealthy queries \cite{ref479}. The threat extends beyond text: MrM introduces the first black-box MIA framework for multimodal RAG systems, using a multi-object data perturbation framework constrained by counterfactual attacks to concurrently induce retrieval of target data and generate membership-leaking information, demonstrated effectively against GPT-4o and Gemini-2 \cite{ref480}.

\textbf{Contextual leakage and knowledge extraction attacks} constitute the second major category of privacy violations, extending the earlier discussion of knowledge base extraction to more subtle and implicit mechanisms. These attacks exploit the instruction-following capabilities of LLMs to force the RAG system to regurgitate the content of retrieved documents, often verbatim. The Pirates of the RAG framework presents a black-box, adaptive attack that uses a relevance-based mechanism and an attacker-side open-source LLM to generate effective queries that leak most of the hidden knowledge base, outperforming prior approaches across diverse domains \cite{ref465}. SECRET formalizes external data extraction attacks (EDEAs) with a unified framework decomposing attack design into extraction instruction, jailbreak operator, and retrieval trigger. Through adaptive optimization using LLMs as optimizers and cluster-focused triggering, SECRET successfully extracts 35\% of the data from RAG powered by Claude 3.7 Sonnet---where other attacks achieve 0\%---highlighting the severity of this threat on state-of-the-art systems \cite{ref481}.

Importantly, not all knowledge extraction requires explicit malicious prompts. IKEA (Implicit Knowledge Extraction Attack) demonstrates that benign-looking queries can silently extract knowledge by leveraging anchor concepts and iteratively mutating them under similarity constraints to explore the embedding space. This attack surpasses baselines by over 80\% in extraction efficiency and 90\% in attack success rate, evading input- and output-level detection mechanisms \cite{ref482}. Data extraction attacks can also be amplified through backdoors, connecting this subsection to the poisoning threats discussed earlier. By injecting a small portion of poisoned data during the fine-tuning phase, attackers can create a backdoor within the LLM that, when integrated into a RAG system, manipulates the model to leak documents from the retrieval database upon specific triggers. With only 5\% poisoned data, this approach achieves 94.1\% success for verbatim extraction and 63.6\% for paraphrased extraction on Gemma-2B-IT, underscoring the privacy risks associated with the supply chain when deploying RAG systems \cite{ref463}. Furthermore, jailbreaking can escalate the scale and severity of these attacks: attackers can escalate membership and entity extraction to full document extraction (80--99.8\% of the database for a Q\&A chatbot) and even craft self-replicating prompts that spread worms across a GenAI ecosystem \cite{ref483}.

\textbf{The success of these attacks is governed by several system-level factors} that must be understood to design effective defenses, a topic elaborated in the following subsection. First, the retriever type and embedding model significantly influence attack effectiveness. The Benchmarking Knowledge-Extraction Attack and Defense framework systematically evaluates heterogeneous retrieval embeddings, diverse generators, and graph-based indexing, demonstrating that embedding algorithms and retriever architectures create varying degrees of vulnerability \cite{ref484}. Second, the generator's instruction-following capability is a dominant factor. LeakDojo's large-scale study across fourteen LLMs and four datasets reveals that stronger instruction-following capability correlates with higher leakage risk, and that improvements in RAG faithfulness can paradoxically introduce increased leakage vulnerability. Notably, query generation and adversarial instructions contribute independently to leakage, with overall leakage well approximated by their product \cite{ref466}. Third, the RAG paradox emerges clearly: mechanisms designed to improve transparency and faithfulness---such as better context utilization and instruction adherence---simultaneously broaden the attack surface, creating a fundamental tension between utility and privacy \cite{ref466,ref485}. This paradox is compounded by the fact that stronger models exhibit higher vulnerability: GPT-4 and Claude-1.3 achieve up to 99\% leakage in multi-turn prompt leakage scenarios \cite{ref486}.

In synthesis, privacy attacks on RAG systems are multifaceted and evolving rapidly, representing a critical dimension of the expanded attack surface introduced in the previous subsection. MIAs leverage semantic, differential, and side-channel signals, while knowledge extraction attacks exploit instruction following, backdoors, and even implicit queries. The interplay between retriever architecture, embedding choices, and generator capabilities determines overall vulnerability, and the inherent transparency mechanisms of RAG create a paradox where improvements in faithfulness amplify privacy risks. These findings collectively emphasize the urgent need for privacy-aware retrieval and generation mechanisms---including cryptographic, differential privacy, and architectural defenses detailed next---that anticipate the full spectrum of attack vectors \cite{ref477,ref466}.

\subsection{Defense Mechanisms: Cryptographic, Differential Privacy, and Architectural Countermeasures}

The security vulnerabilities inherent to Retrieval-Augmented Generation (RAG) systems---ranging from knowledge base extraction and membership inference to data poisoning and prompt injection---necessitate robust and multi-layered defense mechanisms. These countermeasures can be systematically categorized along the three primary stages of the RAG pipeline: the input/retrieval stage, the processing/context-construction stage, and the output/generation stage. The defenses span three core methodological pillars: cryptographic techniques for data confidentiality, differential privacy (DP) frameworks for formal privacy guarantees, and architectural or system-level countermeasures for dynamic and context-aware protection.

\textbf{Cryptographic Approaches: Securing Data at Rest and in Computation}

At the foundational level, cryptographic methods aim to protect the confidentiality and integrity of the knowledge base itself. A prominent strategy is the use of encryption to secure the data and its corresponding vector embeddings prior to storage. This ensures that the data remains inaccessible to unauthorized parties, even in the event of a security breach. One such approach involves encrypting both textual content and embeddings, restricting access to entities with the appropriate decryption keys and thereby significantly reducing the risk of unintended data exposure. Such frameworks are often supported by formal security proofs that rigorously verify confidentiality and integrity under a computational security model, establishing a theoretical foundation for verifiably secure RAG systems \cite{ref487}.

However, a significant limitation of simple encryption is that it often prevents conventional similarity computations necessary for retrieval. To overcome this, advanced techniques employ Homomorphic Encryption (HE) to enable computations directly on encrypted data. For instance, frameworks have been developed that use HE to select documents based on similarity scores exceeding a pre-defined threshold, reducing computational complexity from quadratic to linear in corpus size, and enabling fully encrypted similarity evaluation without revealing query content or intermediate scores \cite{ref488}. Similarly, other frameworks focus on encrypting embeddings while preserving the ability to compute distance comparisons between encrypted queries and database embeddings, defending against both vector-to-text reconstruction and structural leakage attacks \cite{ref489}. While HE provides strong guarantees, it often incurs high computational latency, which is a critical consideration for practical deployment \cite{ref488}. To address the privacy-utility-efficiency trilemma, some systems propose fully encrypted retrieval with a ``pre-storage full-encryption scheme'' ensuring dual protection of both retrieved content and vector embeddings, balancing security with performance \cite{ref487}.

\textbf{Differential Privacy Frameworks: Providing Formal Privacy Guarantees}

While cryptography protects data in storage and computation, Differential Privacy (DP) offers a formal mathematical guarantee against inference attacks on the system's outputs. The core principle is to add calibrated noise to the output or intermediate computations so that the presence or absence of any single data point does not significantly affect the final result. In the context of RAG, a key challenge is that the retrieved context may directly leak sensitive information from the knowledge base. Early DP approaches to RAG have focused on mechanisms to generate long, accurate answers within a moderate privacy budget by spending the budget only on tokens that require sensitive information and using the non-private LLM for others \cite{ref490}.

However, applying DP to RAG is non-trivial. A significant advancement was the development of algorithms that use the propose-test-release paradigm, which acknowledges that many queries can be sufficiently answered with a few keywords. This involves obtaining an ensemble of relevant contexts, generating responses for each, and then privatively releasing the most frequent keywords to be augmented into the final prompt. This method compresses the semantic space, improving the privacy-utility trade-off \cite{ref491}. Furthermore, while many early DP-RAG algorithms functioned in single-query settings, practical usage often involves multiple queries. To address this, Multi-Query DP-RAG (MURAG) algorithms have been proposed, which leverage an individual privacy filter so that the accumulated privacy loss depends only on how frequently each document is retrieved, not the total number of queries. This allows for scaling to hundreds of queries within a practical DP budget \cite{ref492}. An extension, MURAG-ADA, further improves utility by privately releasing query-specific thresholds for more precise document selection \cite{ref492}.

A critical insight is that the privacy risk of a document is highly dependent on the user's query, meaning that a static, document-level risk assumption is inadequate. This has led to the development of Prompt-Aware Dynamic Hierarchical Differential Privacy frameworks, which perform a prompt-aware risk hierarchy to dynamically assess privacy risks under different queries and apply adaptive sensitive entity replacement and text selection. This approach minimizes unnecessary modifications to the retrieval corpus, achieving a better privacy-utility trade-off than prior methods \cite{ref493}. An alternative and highly scalable approach is to use DP to generate synthetic RAG databases. By instructing LLMs to generate differentially private synthetic text that mimics the database records, one can reuse the synthetic text for multiple queries, thereby avoiding repeated noise injection and additional privacy costs, offering a fixed-budget solution \cite{ref494}.

\textbf{Architectural and System-Level Countermeasures: Dynamic and Context-Aware Defense}

Beyond cryptographic and DP-based techniques, a range of architectural defenses provides dynamic and adaptable protection. A key realization is that enabling all defenses simultaneously incurs substantial utility costs, making adaptive, query-aware defense crucial. The \textbf{Sentinel-Strategist architecture} embodies this principle, employing a Sentinel to detect anomalous retrieval behavior and a Strategist to selectively deploy only the defenses warranted by the query context. This approach has been shown to effectively eliminate membership-inference leakage and substantially recover retrieval utility compared to a fully static defense stack \cite{ref457}.

Another critical architectural defense is \textbf{selective disclosure}, which decouples the enforcement of security constraints from the generation process itself. Instead of relying on prompt-level safeguards, which are vulnerable to prompt injection attacks, SD-RAG applies sanitization and disclosure controls during the retrieval phase, prior to augmenting the language model's input. This model incorporates semantic, human-readable dynamic security and privacy constraints alongside an optimized graph-based data model for fine-grained, policy-aware retrieval \cite{ref495}.

In the absence of robust disclosure control, \textbf{dynamic anonymization} frameworks address privacy by anonymizing the knowledge base. Some approaches, like TRIP-RAG, move beyond the ``one-size-fits-all'' full processing of sensitive entities. They evaluate entities from the perspectives of marginal privacy risk, knowledge divergence, and topical relevance to identify highly sensitive entities while trading off utility, thereby achieving variable-intensity privacy protection \cite{ref496}.

Finally, for cross-institutional settings where data cannot be centralized due to privacy regulations, \textbf{federated learning and confidential computing} offer promising solutions. The fundamental conflict in federated RAG arises from the Transformer's self-attention mechanism, which mandates cross-node access to distributed Key-Value caches. To overcome this, frameworks like FedRAG use a novel Scrambled Distributed Attention protocol that decouples attention execution from data localization without exposing plaintext, robustly defending against intermediate state inversion attacks \cite{ref497}. Furthermore, Confidential Computing (CC) techniques, such as those implemented in C-FedRAG, enable secure connection and scaling of RAG workflows across a decentralized network of data providers by ensuring context confidentiality, often leveraging hardware-based trusted execution environments to protect data in use \cite{ref498}.

In summary, a comprehensive defense strategy for RAG systems requires a layered approach. Cryptographic methods provide the foundation for data confidentiality, differential privacy offers formal guarantees against inference, and adaptive, context-aware architectural controls balance security with utility. The future of secure RAG lies in the intelligent and dynamic orchestration of these diverse defenses to navigate the complex security-utility-efficiency trilemma.

\subsection{Fairness, Ethical Trade-offs, and Trustworthiness in RAG Deployment}

The deployment of Retrieval-Augmented Generation (RAG) systems in real-world applications introduces a complex array of ethical and trustworthiness challenges that extend far beyond technical accuracy. Building upon the security vulnerabilities and defense mechanisms discussed in the previous subsection---which established cryptographic, differential privacy, and architectural countermeasures---this subsection shifts focus to the broader socio-technical implications of RAG deployment. As these systems increasingly mediate critical decisions in domains like healthcare, finance, and law, the interplay between fairness, privacy, security, and utility becomes central to responsible system design. This subsection examines these intersections, focusing on fairness vulnerabilities, the privacy-utility-security trilemma, transparency versus security trade-offs, access control in enterprise settings, and the evaluation frameworks needed to navigate this intricate landscape.

A growing body of research reveals that RAG systems introduce unique fairness vulnerabilities that can undermine their ethical deployment---vulnerabilities that persist even when the security mechanisms from the previous subsection are properly implemented. Unlike conventional backdoors that rely on direct trigger-to-target mappings, fairness-driven attacks exploit the interaction between retrieval and generation models to manipulate semantic relationships between target groups and social biases. The BiasRAG framework formalizes this threat through a two-phase backdoor attack: during the pre-training phase, the query encoder is compromised to align target groups with intended social biases, ensuring long-term persistence; in the post-deployment phase, adversarial documents are injected into knowledge bases to reinforce the backdoor, maintaining a subtle and evolving influence on content generation while remaining undetectable under standard fairness evaluations \cite{ref499}. Additional empirical studies using metamorphic testing on small language models integrated into RAG pipelines show that minor demographic variations---particularly those involving racial cues---can break up to one third of metamorphic relations, revealing a consistent bias hierarchy and underscoring the need for careful curation of retrieval components to prevent bias amplification \cite{ref500}. Even more concerning is the finding that RAG can undermine fairness alignment without any fine-tuning or retraining. Experiments using uncensored, partially censored, and fully censored external datasets demonstrate that even with fully censored and supposedly unbiased external datasets, RAG can still lead to biased outputs, highlighting that the effectiveness and cost-efficiency of RAG come with hidden fairness costs \cite{ref72}. A comprehensive empirical evaluation of fairness in RAG methods, using scenario-based questions and analyzing disparities across demographic attributes such as gender and geographic location, confirms that fairness issues persist in both the retrieval and generation stages, suggesting that fairness considerations have been largely overlooked as models prioritize utility like exact match accuracy \cite{ref501}.

Beyond fairness, RAG deployment inherently involves navigating the privacy-utility-security trilemma---a tension that directly extends the trade-offs introduced in the previous subsection's discussion of cryptographic and differential privacy defenses. While those technical mechanisms provide formal guarantees, their practical implementation reveals that optimizing for one dimension often comes at the expense of others. A key finding in this context is that enabling all relevant defenses simultaneously---an ``always-on'' defense stack---incurs substantial utility costs, with experiments showing more than 40\% reduction in contextual recall, indicating that retrieval degradation is the primary failure mode. Adaptive, context-aware defense orchestration, such as the Sentinel-Strategist architecture introduced earlier, addresses this by detecting anomalous retrieval behavior and selectively deploying only the defenses warranted by the query context, eliminating membership inference leakage while substantially recovering retrieval utility relative to static defense stacks \cite{ref457}. In the specific context of healthcare, the MedPriv-Bench benchmark demonstrates a pervasive privacy-utility trade-off by jointly evaluating privacy preservation and clinical utility in medical open-ended question answering, using a multi-agent, human-in-the-loop pipeline to synthesize sensitive medical contexts that create realistic privacy pressure \cite{ref502}. This trade-off is further complicated by the fact that different stakeholder groups have different needs and access rights, suggesting that a one-size-fits-all approach to privacy and security may be neither optimal nor ethical \cite{ref503}. Indeed, the architectural defenses described previously---from homomorphic encryption to dynamic anonymization---each occupy different points along this trilemma frontier, and the choice of which to deploy must be guided by the specific ethical priorities of the application domain.

The tension between transparency and security represents another critical challenge in RAG deployment, one that interacts with the system-level defenses discussed earlier in a complex and often counterintuitive manner. While the previous subsection highlighted how architectural countermeasures like selective disclosure enhance security by controlling information flow during retrieval, the very act of making RAG systems transparent to users---a design choice that builds trust---can inadvertently create new attack surfaces. A structural vulnerability known as the ``RAG paradox'' arises from the system's effort to enhance trust by revealing both retrieved documents and their sources to users. This transparency enables attackers to observe which sources are used and how information is phrased, allowing them to craft poisoned documents that are more likely to be retrieved and upload them to the identified sources. Unlike prior work that focuses solely on improving document retrievability, this attack method explicitly considers both retrievability and user trust in the retrieved content, demonstrating that systems designed to be transparent can inadvertently create attack vectors \cite{ref504}. This transparency-security tension is also reflected in the broader challenge of access control in multi-user enterprise deployments. RAG systems in enterprise settings face the dual challenge of supporting diverse user needs while preventing unauthorized access to proprietary information---a challenge that complements the cryptographic and federated approaches described in the previous subsection, which protect data at rest and in cross-institutional settings but do not fully address the dynamic, user-specific nature of access control in shared deployments. Frameworks that retrieve from pre-generated expert knowledge and on-demand external LLM-generated knowledge, while adopting local open-source generators and selectively using external LLMs only when prompts are deemed safe by a filtering mechanism, demonstrate that it is possible to enhance completeness and prevent data leakage simultaneously \cite{ref505}. However, effective access control requires going beyond simple filtering to account for the heterogeneous needs of different user groups, with guardrails that dynamically moderate access to sensitive content based on user trust metrics, incorporating both direct interaction trust and authority-verified trust \cite{ref503}. This dynamic approach aligns with the recognition that ethical deployment of RAG systems requires careful consideration of who has access to what information and under what conditions, particularly in domains where data sensitivity varies considerably \cite{ref506}. Thus, while the cryptographic and architectural defenses from the previous subsection provide the foundational security layer, enterprise access control demands an additional, policy-driven layer that operates at the intersection of security, trust, and user intent.

To systematically evaluate and ensure trustworthiness in RAG deployment, several comprehensive frameworks have been proposed, building upon and complementing the defense-oriented evaluations of the previous subsection. While those security analyses focus on attack resistance and privacy guarantees, the frameworks discussed here take a broader view that encompasses fairness, transparency, and accountability alongside security and privacy. The Trust-RAG Compass provides a unified framework that assesses trustworthiness across six key dimensions: factuality, robustness, fairness, transparency, accountability, and privacy. Companion to this is the TRC Bench, which conducts comprehensive evaluations of proprietary and open-source models across these dimensions, revealing performance gaps between different types of LLMs and shedding light on varying levels of trustworthiness \cite{ref77}. Similarly, a broader roadmap for developing trustworthy RAG systems organizes the discussion around five key perspectives---reliability, privacy, safety, fairness, and explainability, and accountability---presenting a general framework and taxonomy for each to provide structured understanding of current challenges and identify promising future research directions \cite{ref507}. These evaluation frameworks are complemented by more specialized tools. For instance, benchmarks like SafeRAG evaluate security by classifying attack tasks into categories such as silver noise, inter-context conflict, soft ad, and white Denial-of-Service, demonstrating that RAG exhibits significant vulnerability to all attack tasks and that even the most apparent attack task can easily bypass existing retrievers, filters, or advanced LLMs \cite{ref508}. Additionally, the RAGChecker framework and similar diagnostic tools are evolving to provide fine-grained, stage-decomposed evaluation that isolates performance bottlenecks across the pipeline, moving beyond accuracy-centric assessments toward holistic evaluation that captures the multi-faceted nature of trustworthiness \cite{ref77}. The evaluation of trustworthiness in RAG systems remains a significant challenge, requiring assessment across multiple ethical dimensions that are often in tension with each other, and the development of standardized, comprehensive evaluation frameworks continues to be an active area of research \cite{ref507}.

In summary, the fair, ethical, and trustworthy deployment of RAG systems requires acknowledging and actively managing inherent tensions between competing objectives, extending the defense-in-depth perspective of the previous subsection into the socio-technical domain. While the cryptographic, differential privacy, and architectural mechanisms described earlier provide essential protections against specific attack vectors, they cannot by themselves guarantee ethical deployment. Fairness vulnerabilities can be exploited through backdoor attacks \cite{ref499}, and demographic biases persist even with careful dataset curation \cite{ref72}. The privacy-utility-security trilemma demands adaptive, context-aware solutions rather than static defense stacks \cite{ref457}, and the transparency that builds user trust can simultaneously create attack surfaces \cite{ref504}. Moving forward, the research community and practitioners must prioritize the development of integrated frameworks that address all dimensions of trustworthiness---factuality, robustness, fairness, transparency, accountability, and privacy---in a unified manner, building upon the technical defense mechanisms established in the previous subsection while extending them with ethical considerations that ensure RAG systems can be deployed in high-stakes environments with confidence in their reliability and ethical integrity.

\section{Efficiency, Scalability, and System Engineering}

\subsection{Performance Bottlenecks and Trade-offs in RAG Inference}

Retrieval-Augmented Generation (RAG) systems enhance the factual accuracy and relevance of large language model (LLM) outputs by integrating external knowledge bases into the generation process. However, this augmentation introduces a set of fundamental performance bottlenecks that stem from the very nature of the pipeline: the retrieval and subsequent processing of long, external contexts. The dominant challenge is the severe computational and memory overhead imposed by the large number of input tokens, which leads to a substantial increase in prefill-stage computation and the memory footprint of the key-value (KV) cache. This overhead directly translates into increased latency, reduced throughput, and higher operational costs, creating a complex trade-off landscape that system designers must navigate.

The prefill stage, where the LLM processes the concatenated user query and retrieved documents, emerges as the primary performance bottleneck in RAG inference. Unlike the token-by-token autoregressive generation phase, the prefill stage must process all input tokens in parallel, which is computationally intensive. The time-to-first-token (TTFT) is a critical user-perceived latency metric that is directly and severely impacted by the length of these inputs. The computational cost of prefill grows linearly with the number of input tokens, and as RAG systems are designed to retrieve more documents for better context, the TTFT can become prohibitively high. This problem is exacerbated in high-throughput serving scenarios where multiple requests compete for the same computational resources. For instance, while KV-cache reuse offers a promising solution by storing previously computed KV states, its effectiveness is often limited by low cache hit rates and high CPU-GPU data transfer overhead, as highlighted in the design of systems like PCR, which aims to maximize cache reuse efficiency through intelligent prefetching \cite{ref509}. The fundamental challenge is that the long sequence generation inherent to RAG introduces high computation and memory costs, a bottleneck specifically pinpointed in the analysis of current RAG systems that motivates the need for specialized caching systems \cite{ref510}. Techniques that pre-compute and store KV caches for documents offline can eliminate online computation of KV caches during inference, drastically reducing TTFT, as demonstrated by TurboRAG, which reports up to a 9.4x reduction in TTFT across a suite of RAG benchmarks \cite{ref511}.

The KV cache, which stores the intermediate key and value states for each token to avoid redundant computation, is the second major bottleneck, primarily due to its immense and rapidly growing memory footprint. The memory required by the KV cache scales linearly with both the context length and the batch size, placing significant pressure on the limited high-bandwidth memory (HBM) of GPUs. The memory footprint is so substantial that it often becomes the primary factor limiting the batch size that can be processed concurrently on a single GPU, directly impacting throughput and cost-efficiency. As context lengths in production systems push from thousands to millions of tokens, the KV cache can easily reach several gigabytes, creating critical bottlenecks on GPU memory capacity, memory bandwidth, and inference throughput \cite{ref512}. This memory pressure forces system designers to make a critical trade-off between the number of concurrent requests (batch size) and the context length that can be accommodated. The issue is further complicated by the fact that KV cache must be reserved upon request arrival while the output token length is unknown, leading to a central trade-off between memory efficiency and preemption risk \cite{ref513}. Consequently, the growth of KV caches poses significant system-level challenges, particularly as model sizes increase and concurrent requests compete for limited memory resources \cite{ref514}.

The interplay between these prefill-time and memory bottlenecks creates a complex web of operational trade-offs among latency (TTFT), throughput, memory usage, and cost. The most fundamental trade-off is between generation quality and efficiency. While retrieving more documents generally improves the factual grounding and quality of the answer, it also leads to longer input sequences, which in turn increases prefill computation and KV cache memory. The computational overhead of processing these long contexts during the prefill stage often makes it a dominant serving cost. For example, in systems serving RAG applications, the long sequences generated by knowledge injection are identified as the performance bottleneck, and the need to cache intermediate states becomes a key optimization opportunity to mitigate these costs \cite{ref510}. The cost associated with this stage is not only in terms of latency but also in computational resources, as processing the same retrieved documents repeatedly across different user queries represents a significant amount of redundant computation on expensive GPUs \cite{ref515}. This redundancy creates a direct trade-off between minimizing latency by pre-computing and caching KV states versus the cost of the storage and the high latency of I/O operations required to retrieve them. When KV caches are offloaded to CPU DRAM or SSDs, the limited bandwidth of the PCIe interconnect can become a new bottleneck, potentially erasing the benefits of caching; it is a key finding that in offloaded serving, the GPU may consume only a fraction of its rated power while spending the vast majority of time waiting for cache transfers \cite{ref516}. This highlights a critical trade-off between using expensive HBM for speed or cheaper but slower storage for capacity, a decision that fundamentally shapes system performance and cost.

Real-world workload characterization is essential for understanding and addressing these bottlenecks. While many studies rely on synthetic or generic workloads, the distinctive characteristics of real-world RAG systems, such as significant temporal locality and highly skewed hot document access patterns, have been captured in open-source trace datasets like RAGPulse \cite{ref63}. This kind of analysis reveals that the performance of RAG systems is not uniform but is highly dependent on the workload's specific characteristics, such as the distribution of query types and the reuse patterns of external documents. The performance variability across the wide range of emerging RAG variants, which may have different numbers of pipeline stages (retrieval, reranking, generation) and different interaction patterns, makes the challenge of systematic performance optimization even more complex. The lack of a unified abstraction to describe these different RAG algorithms and their workloads makes it difficult to design general-purpose optimizations, as systems are often tuned for a specific architecture \cite{ref517}. Furthermore, the comparative analysis of different KV cache management frameworks, such as vLLM, InfiniGen, and H2O, reveals that there is no universally optimal strategy; the best choice depends heavily on the specific parameter settings like request rates, model sizes, and the operational environment \cite{ref514}. This finding underscores the need for adaptive and configurable system designs.

These profound inefficiencies---the computational overhead of long contexts, the prefill bottleneck, and the immense memory footprint of KV caches---constitute the fundamental operational constraints that motivate the broad spectrum of optimization techniques discussed in the following subsections. The high latency and cost associated with the prefill stage have driven research into advanced KV-cache reuse mechanisms and selective recomputation techniques. The memory pressure has spurred the development of offloading strategies, quantization, and eviction policies. Ultimately, the goal is to find a balance between the quality benefits of rich context and the efficiency demands of real-world deployment, a balance that requires co-designing algorithms, systems, and hardware. The challenge lies in developing holistic solutions that can adapt to workload variability and navigate the complex trade-off space, an endeavor that sits at the heart of system engineering for RAG.

\subsection{Retrieval-Side Optimization and Caching Strategies}

The retrieval stage is a critical performance bottleneck in Retrieval-Augmented Generation (RAG) systems, as it directly impacts end-to-end latency and throughput. Building on the operational constraints outlined earlier---where the prefill bottleneck and KV cache memory pressure were identified as fundamental challenges---this subsection explores the optimization landscape for the retrieval stage, focusing on two complementary strategies: hierarchical cache-aware systems for managing intermediate computational states and semantic caching mechanisms that exploit query similarities. These techniques are essential for mitigating the computational overhead of vector database lookups and the prefill-stage processing of long retrieved contexts.

A significant body of work has focused on caching the intermediate Key-Value (KV) states of retrieved documents to avoid redundant computation. This directly addresses the prefill-time bottleneck described in the previous subsection, where repeated processing of the same documents across queries was identified as a source of substantial computational waste. The foundational system in this area, \cite{ref510}, proposes a multilevel dynamic caching system that organizes the intermediate states of retrieved knowledge in a knowledge tree, caching them across the GPU and host memory hierarchy. Its replacement policy is aware of both LLM inference characteristics and RAG retrieval patterns, and it dynamically overlaps retrieval and inference steps to minimize latency. This system demonstrated that caching knowledge's intermediate states can reduce time-to-first-token (TTFT) by up to 4x and improve throughput by up to 2.1x. Building on this concept, \cite{ref511} redesigns the RAG inference paradigm by pre-computing and storing KV caches for document chunks offline. During online inference, it directly retrieves the precomputed KV cache, thereby eliminating the online prefill computation for retrieved documents. By addressing critical insights into mask matrices and positional embeddings, TurboRAG achieves up to a 9.4x reduction in TTFT while maintaining comparable performance, and it requires no modification to the underlying models or inference systems.

The challenge of reusing precomputed chunk-level KV caches in arbitrary contexts is addressed by \cite{ref515}. This system introduces the concept of ``chunk-caches'' and presents methods to identify reusable caches, efficiently perform a small fraction of recomputation to maintain output quality, and manage cache storage and eviction in hardware. By masking the overheads associated with cache management, Cache-Craft reduces redundant computation by 51\% over state-of-the-art prefix-caching and 75\% over full recomputation, leading to a 1.6x throughput improvement and 2x reduction in end-to-end response latency on production workloads. Furthering the granularity of cache fusion, \cite{ref518} introduces a compressed-view query-aware selector to decide which tokens to recompute during cache fusion. This approach uses chunk-anchor query probing and critical-layer profiling to identify important tokens without requiring all-layer inspection, achieving an average prefill-time speedup of 1.7x over full prefill while maintaining full-prefill-level quality. Similarly, \cite{ref519} addresses the TTFT bottleneck by using a small auxiliary LLM to identify critical tokens for restoring inter-chunk attention. This guidance enables selective KV cache recomputation, and combined with shared prefixes and a sliding-window grouping strategy, CacheClip retains up to 91.1\% of full-attention performance on LongBench while accelerating prefill time by up to 3.33x.

Beyond chunk-level and document-level caches, another line of research explores query-level semantic caching, which addresses the memory footprint concerns raised in the previous subsection by reducing the frequency of expensive retrievals and the subsequent prefill of long contexts. For instance, \cite{ref520} introduces Proximity, an approximate key-value cache that reuses previously retrieved documents for similar queries. By leveraging query similarities and a locality-sensitive hashing (LSH) scheme for fast lookups, it significantly reduces the reliance on expensive vector database calls, achieving a 77.2\% reduction in database calls while maintaining recall and accuracy on MMLU and MedRAG benchmarks. This semantic, query-level caching approach is complementary to KV cache reuse. In a similar vein, \cite{ref521} enhances semantic caching for LLMs by integrating Multi-Vector Retrieval (MVR). It uses a learnable segmentation model to split prompts for fine-grained similarity comparisons via a MaxSim operator, with its training objective derived from a theoretical analysis to maximize cache hits under correctness constraints. This method improves cache hit rates by up to 37\% compared to state-of-the-art methods.

The effectiveness of these caching systems relies heavily on efficient cache management strategies, including prefetching, eviction, and data placement. These strategies directly address the trade-offs between fast but expensive HBM and slower but cheaper storage that were identified as a central challenge in the previous subsection. For instance, \cite{ref509} proposes a system to maximize KV-cache reuse efficiency in high-throughput scenarios. It introduces a prefix-tree caching structure with a look-ahead LRU replacement policy that uses pending requests in the scheduler queue to improve cache hit ratios, along with layer-wise overlapping and queue-based prefetching to hide communication latency and proactively load KV caches from SSD into DRAM. This system achieves up to a 2.47x speedup in average TTFT. Addressing the unique challenges of mobile platforms, \cite{ref522} proposes a hierarchical cache solution that predictively populates cache with likely-future queries and progressively matches similar queries and QKV cache to maximize the reuse of intermediate results. This adaptive configuration aims to maximize caching utility with minimal resource consumption, resulting in a 34.4\% latency reduction across various applications on commodity mobile phones. The concept of hotness-aware data management is explored by \cite{ref523}, which leverages the frequency of KV chunk access to inform a mixed-precision compression strategy and a data placement policy that stores hot chunks in faster memory. This approach achieves an average 2.10x speedup in TTFT over TurboRAG. Finally, to manage the growing KV cache footprint across multiple storage tiers, \cite{ref524} introduces Kareto, an optimizer that identifies the Pareto frontier of cost, throughput, and latency within a storage configuration space. It uses a diminishing-return-guided pruning method to navigate the large configuration space and implements fine-grained adaptive tuning with eviction policies based on KV block access patterns, improving cache efficiency and allowing for better cost-effectiveness and adaptability to workload changes.

In summary, retrieval-side optimization in RAG systems is evolving from simple vector search to a complex, cache-centric engineering discipline. The integration of document-level, chunk-level, and query-level caching, combined with sophisticated prefetching, eviction, and data placement policies, is crucial for overcoming the latency and throughput bottlenecks of production RAG deployments. These techniques provide the efficient and compact context representations that the generation-stage optimizations---discussed next---build upon to further accelerate decoding and maintain output quality.

\subsection{Generation-Side Acceleration and Context Management}

The generation stage of a RAG pipeline---where the LLM synthesizes an answer from the retrieved context---presents a distinct set of performance bottlenecks that complement the retrieval-side optimizations discussed previously. While retrieval-side systems focus on efficiently sourcing and caching relevant information, generation-side acceleration is primarily concerned with reducing the computational and memory overhead of processing long, concatenated contexts and generating tokens. The primary challenges are threefold: the latency and compute associated with the prefill stage, the growing memory footprint of the key-value (KV) cache during decoding, and the semantic noise introduced by irrelevant retrieved documents. Research in this area has produced a variety of context management and decoding optimization techniques to address these issues, often balancing the trade-off between efficiency and the faithfulness of the final output. Notably, many of these techniques build directly upon the caching infrastructure described in the previous subsection---for instance, KV-cache reduction methods complement the cache-aware retrieval systems by making the cached representations themselves more compact and efficient.

A primary source of inefficiency in RAG inference is the long sequence of retrieved documents that must be processed in the prefill stage. This process generates a large KV cache, which not only consumes significant memory but also slows down the time-to-first-token (TTFT). This challenge is exacerbated by the retrieval-stage optimizations that may return extensive contextual chunks, making the prefill bottleneck even more pronounced. A large body of work addresses this through \textbf{context pruning and compression}, which aims to reduce the volume of information the LLM must process. These methods operate on the principle that not all retrieved tokens are equally important for answering the query. For instance, \cite{ref157} introduces an attention-guided context pruning method that reformulates RAG queries into a next-token prediction paradigm to isolate the semantic focus, enabling precise attention calculation between queries and contexts. This approach achieves significant context compression while outperforming other methods. Similarly, \cite{ref6} provides a comprehensive overview of this evolving field, detailing various paradigms for reducing the overhead of extensive contextual data.

More advanced methods move beyond simple token pruning to \textbf{compressive and adaptive context management}. Instead of heuristically filtering tokens, these systems learn to compress the retrieved information into more compact representations. \cite{ref525} proposes a novel approach where retrieved contexts are compressed into compact embeddings before being encoded by the LLM. This reduces computational cost and also optimizes the compressed embeddings for downstream task performance. Building on this idea, \cite{ref526} integrates a lightweight module into language model layers to help the model better utilize highly compressed context representations, using an adaptive gating system to manage complexity and maintain generation speed. The degree of compression can also be dynamically adjusted based on query complexity. For example, \cite{ref527} introduces a framework that adjusts compression rates based on input complexity, preventing both over-compression of simple queries and under-compression of complex ones. Another framework, \cite{ref528}, presents a hierarchical multi-scale approach that compresses and partitions documents into coarse-to-fine granularities, adaptively merging relevant contexts to construct query-specific long contexts that balance precision and coverage. These methods highlight a shift from static, fixed-ratio compression to more intelligent, content-aware, and adaptive systems.

Another significant area of optimization is \textbf{KV-cache reduction and efficient decoding}, which directly extends the caching mechanisms introduced in the retrieval stage. While the retrieval stage focused on caching and reusing KV states across requests, the generation stage must also manage the ongoing growth of the KV cache during decoding. The KV cache grows linearly with context length and is a major bottleneck for memory and latency, especially during the decoding phase. A key observation is that much of the computation over the RAG context during decoding is unnecessary, as only a small subset of tokens is directly relevant to generation. Based on this insight, \cite{ref529} proposes an efficient decoding framework that compresses, senses, and expands to exploit the sparsity structure in attention patterns, leading to substantial acceleration in time-to-first-token without loss in accuracy. A different approach to the KV-cache bottleneck involves being selective about what gets stored in the first place. The \cite{ref530} paper proposes a mechanism to learn to predict token utility before cache entry, filtering out low-utility states early to maintain a compact cache and deliver significant prefill and decode speedups. To further enhance efficiency, \textbf{speculative decoding} can also be integrated with RAG, where a smaller draft model generates candidate tokens that are then verified by the larger target model. \cite{ref531} extends this by using a RAG drafter operating on shortened retrieval contexts to speculate on the generation of long-context target LLMs, achieving performance improvements and more than 2x speedups.

At the system level, \textbf{cache-augmented generation (CAG)} presents an alternative paradigm to real-time RAG, effectively bypassing the online retrieval stage altogether and thereby merging retrieval and generation into a single, pre-computed inference phase. In contrast to the ``retrieve-always'' assumption of standard RAG, which incurs significant computational overhead and inference latency, \cite{ref532} proposes preloading all relevant resources into the LLM's extended context and caching its runtime parameters. During inference, the model answers queries using these preloaded parameters without additional retrieval steps. This method eliminates retrieval latency and minimizes retrieval errors but is most effective when the knowledge base is of a manageable size. To scale CAG for larger and more dynamic knowledge bases, \cite{ref533} introduces Adaptive Contextual Compression to dynamically manage context inputs and proposes a Hybrid CAG-RAG Framework that integrates selective retrieval to augment preloaded contexts when necessary.

Complementing these approaches is \textbf{entropy-based lazy retrieval}, which offers a smarter way to decide when retrieval is even necessary, thus bridging the retrieval and generation stages more intelligently. Instead of always querying a vector database, \cite{ref534} introduces an adaptive framework that first processes a query with a compact document summary and only triggers expensive chunk retrieval when the model's predictive entropy exceeds a calibrated threshold, signaling genuine uncertainty. This provides a tunable accuracy-efficiency trade-off, reducing retrieval calls without significant performance degradation. Finally, the choice between compression, caching, and retrieval can be informed by the nature of the task itself. As \cite{ref535} points out, while RAG is suitable for tasks where sparse evidence suffices, task-aware compression of external knowledge into a compacted representation can be superior for broad knowledge tasks, enabling the model to reason efficiently over all relevant information at once.

In summary, generation-side acceleration in RAG is a multi-faceted problem that both builds upon and complements the retrieval-stage optimizations. Techniques range from micro-level optimizations like KV-cache pruning and speculative decoding to macro-level strategic decisions about when to retrieve, compress, or pre-cache information. These techniques also set the stage for the system-level orchestration challenges discussed next, where end-to-end RAG serving frameworks must integrate retrieval and generation optimizations into cohesive pipelines with heterogeneous resource allocation. Ultimately, the challenge is to integrate these techniques cohesively to not only make the system faster and more cost-effective but also to ensure that the integrity and faithfulness of the ground knowledge are maintained for high-quality generation across diverse deployment scenarios.

\subsection{System-Level Frameworks, Scheduling, and Edge Deployment}

As RAG transitions from research prototypes to production-scale deployments, the focus of system engineering shifts from algorithmic innovation to the efficient orchestration of heterogeneous computational and memory resources. The generation-side optimizations discussed previously---such as KV-cache reduction, context compression, and speculative decoding---provide powerful building blocks, but their effective deployment in real-world systems requires careful coordination across the entire RAG pipeline. End-to-end RAG serving involves a complex interplay between vector databases, embedding models, rerankers, and the LLM itself, each with distinct compute and memory footprints. This subsection examines the system-level frameworks, scheduling strategies, and edge deployment paradigms that have emerged to address these challenges, spanning from datacenter-scale optimization to resource-constrained mobile environments.

A foundational challenge in RAG serving is the heterogeneity of its constituent stages. Modern RAG pipelines are multi-stage workflows that combine retrieval, reranking, prompt construction, and generation, each exhibiting different latency profiles and resource demands. To navigate this complexity, \cite{ref536} introduces a co-design simulation framework that models diverse request stages---including RAG, KV retrieval, reasoning, prefill, and decode---across complex hardware hierarchies. This simulator enables system designers to explore critical trade-offs such as memory bandwidth contention, inter-cluster communication latency, and batching efficiency in hybrid CPU-accelerator deployments, providing an alternative to costly empirical benchmarking. Building on this analytical foundation, \cite{ref517} proposes RAGSchema, a structured abstraction that captures the wide range of RAG algorithm variants, and RAGO, an optimization framework that leverages this schema to achieve up to a 2× increase in queries per second per chip and a 55\% reduction in time-to-first-token latency compared to systems built on general-purpose LLM serving extensions.

The realization that retrieval and generation impose fundamentally different computational demands has motivated a wave of systems that decouple these stages to enable parallel execution and specialized resource allocation. This decoupling directly complements the generation-side techniques discussed in the previous subsection---by isolating the retrieval stage, these systems can more effectively apply compression, caching, and speculative decoding strategies without interference. \cite{ref537} presents an end-to-end serving framework built on three pillars: a flexible pipeline specification interface, heterogeneity-aware deployment that provisions components as a distributed inference system, and a closed-loop runtime controller that reduces SLO violations through request prioritization and auto-scaling, achieving a 2.04× throughput improvement over commercial alternatives. Similarly, \cite{ref538} exploits the insight that serialized RAG pipelines leave substantial idle GPU time. By decoupling retrieval and generation into parallel pipelines with joint memory placement and dynamic batch scheduling, RAGDoll achieves up to a 3.6× speedup in average latency on consumer-grade hardware, making RAG deployment feasible on a single GPU. \cite{ref539} further advances this direction by introducing a runtime system built on a graph-based abstraction that exposes optimization opportunities in stage-level parallelism and intra-request similarity, achieving speedups of up to 5× across diverse RAG workflows through dynamic graph transformations such as node splitting and dependency rewiring.

A significant bottleneck in RAG inference is the prefill stage, where long retrieved contexts necessitate extensive computation of key-value (KV) caches---a challenge directly inherited from the generation-side optimizations that aim to reduce KV-cache growth and memory pressure. TeleRAG addresses this by introducing lookahead retrieval, a prefetching mechanism that predicts required data and transfers them from CPU to GPU in parallel with LLM generation \cite{ref540}. Combined with a prefetching scheduler and cache-aware scheduler for multi-GPU support, TeleRAG achieves up to a 1.53× average end-to-end latency reduction for single queries and a 1.83× higher throughput for batched requests. Building on the concept of KV cache sharing across instances, \cite{ref541} introduces a disk-based shared cache management system that exploits the locality of documents related to user queries. By proactively generating and storing disk KV caches and sharing them across multiple LLM instances, this system achieves a 15--71\% increase in throughput and up to a 12--65\% reduction in latency. At a finer granularity, \cite{ref519} leverages the insight that small auxiliary LLMs exhibit similar attention distributions to primary LLMs, enabling efficient identification of critical tokens for selective KV cache recomputation, and achieving up to a 3.33× prefill acceleration while retaining most of the full-attention quality. For resource-constrained environments, \cite{ref509} introduces a prefix-tree caching structure with look-ahead LRU replacement, layer-wise overlapping for pipelined KV-cache loading, and queue-based prefetching from SSD, achieving up to a 2.47× speedup in average time-to-first-token.

At the system level, layered routing and hierarchical caching strategies have proven effective in managing the trade-off between freshness, latency, and cost. These architectures extend the cache-augmented generation (CAG) paradigm introduced in the previous subsection, moving beyond simple preloading to sophisticated multi-tier caching that adapts to query patterns and knowledge dynamics. \cite{ref542} presents a five-layer architecture that routes queries through fixed KV caches, semantic caches, a memory-recall mode exploiting the LLM's own weights, adaptive session memory, and a conventional retrieval layer. This design reduces average GPU time to 0.248 seconds per query---roughly half that of a naive RAG baseline---while sustaining a throughput of approximately 100,000 queries per second. In multi-GPU and disaggregated settings, \cite{ref543} addresses the challenge of massive KV-cache migrations, which dominate the execution critical path in agentic RAG. TENT's declarative orchestration decouples transfer intent from physical execution, dynamically ``spraying'' fine-grained slices across rails based on real-time telemetry, achieving up to a 1.36× higher throughput and 26\% lower P90 time-to-first-token compared to state-of-the-art transfer engines.

The deployment of RAG on edge and mobile devices presents a distinct set of challenges, including limited memory, heterogeneous SoC architectures, energy constraints, and intermittent connectivity. These constraints motivate a different class of optimizations that complement the datacenter-scale strategies discussed earlier---here, the emphasis shifts from maximizing throughput to minimizing resource footprint while maintaining acceptable latency. \cite{ref544} addresses the memory constraint by pruning embeddings within clusters and generating embeddings on-demand during retrieval, while adaptively caching embeddings for large tail clusters, achieving significant latency reduction over baseline IVF indexes while allowing datasets to fit into device memory. Extending this to collaborative settings, \cite{ref545} proposes a hierarchical scheduling framework with an online query identification mechanism using proximal policy optimization, dynamic inter-node scheduling, and an intra-node scheduler based on online convex optimization, achieving performance gains of 4.23\% to 91.39\% over baseline methods. Similarly, \cite{ref546} employs a hierarchical collaborative gating mechanism to dynamically select among local, edge-assisted, and cloud-based strategies, matching the accuracy of cloud-based knowledge graph RAG while reducing costs by up to 84.6\%.

For heterogeneous mobile SoCs, \cite{ref547} builds profiling-based performance models for each sub-stage and model-processing-unit configuration, capturing latency, workload shape, and contention-induced slowdown, and leverages these in an online scheduler combining shape-aware partitioning, criticality-based accelerator mapping, and bandwidth-aware concurrency control, reducing end-to-end latency by up to 10.94×. \cite{ref548} presents the first end-to-end RAG pipeline running all neural stages on a mobile NPU, achieving 18.1× faster LLM prefilling and 4.0× lower end-to-end latency compared to CPU baselines, with answer quality on par with CPU and GPU backends. \cite{ref522} introduces a hierarchical cache architecture that progressively matches similar queries and QKV caches at different computing stages, combined with predictive population and adaptive configuration, surpassing the best-performing baseline by 34.4\% latency reduction. For extreme resource constraints, \cite{ref549} presents a monolithic architecture consolidating ingestion, extraction, and hybrid retrieval into a single SQLite container, reducing disk footprint by approximately 99.5\% while achieving 100\% Recall@1 for entity retrieval. On the network side, \cite{ref550} decomposes multi-document RAG into parallel token generation processes across cloud and device, employing a dual-side speculative aggregation algorithm to avoid frequent synchronization, achieving up to 1.9× greater gains over standalone small language models compared to centralized RAG.

In multi-tenant and federated settings, systems must balance privacy, efficiency, and collaboration. \cite{ref551} introduces an edge-cloud collaborative framework supporting querying of structured SQL, knowledge graphs, and unstructured documents, with a three-tier caching strategy that optimizes response latency. \cite{ref552} compresses and shares knowledge graphs across heterogeneous edge devices, using hybrid retrieval to expand knowledge coverage and an optimizer that adaptively selects small language models and RAG parameters per query, reducing costs by up to 98.4\% while increasing peak throughput by 97.8\%. These collaborative edge frameworks, along with hierarchical orchestration schemes such as \cite{ref553} and \cite{ref554}, which reduces memory usage by \textasciitilde70\% for MoE models, collectively demonstrate the viability of decentralized RAG deployments. Ultimately, the diversity of system architectures surveyed---from datacenter-scale simulators like \cite{ref536} to single-file edge containers like \cite{ref549}---underscores the need for unified abstraction frameworks and systematic performance engineering to navigate the expanding design space of production RAG systems.

\section{Applications, Case Studies, and Industrial Deployment}

\subsection{Industry-Specific Applications and Case Studies}

Having established the fundamental architectures, advanced techniques, and evaluation methodologies of RAG in the preceding sections, we now turn our attention to the crucible where these systems are truly tested: real-world deployment in domain-specific contexts. This subsection examines how RAG systems are adapted to meet the distinct constraints and technical requirements of high-stakes industries, while also framing the engineering and operational challenges elaborated in the subsequent discussion on production deployment. Across medicine, law, finance, and industrial/enterprise settings, we observe that the transition from generic frameworks to domain-grounded systems reveals a set of recurring patterns---domain-specific retrieval strategies, privacy and security-driven architectural choices, and the need for fine-grained, domain-aware evaluation---that collectively inform the engineering practices presented in the next section.

\textbf{Medical and Healthcare Applications}

The medical domain is arguably the most demanding environment for RAG systems, where diagnostic accuracy, patient safety, and strict privacy regulations (e.g., HIPAA, GDPR) dictate architectures. Self-BioRAG exemplifies a domain-specific framework that trains a 7B-parameter model on 84k filtered biomedical instruction sets, enabling it to generate explanations, retrieve relevant documents, and self-reflect on its outputs using customized reflective tokens \cite{ref555}. This approach demonstrates that domain-specific components---a tailored retriever, a domain corpus, and instruction sets---are essential for adhering to medical reasoning standards. Similarly, MedRAG enhances retrieval with a knowledge graph-elicited reasoning framework, constructing a four-tier hierarchical diagnostic KG that captures critical diagnostic differences across diseases \cite{ref556}. By dynamically integrating these KG insights with similar Electronic Health Records (EHRs) retrieved from databases, MedRAG improves diagnostic specificity, particularly for diseases with similar manifestations, and proactively generates follow-up questions for personalized decision-making. The system was validated on both a public dataset (DDXPlus) and a private chronic pain diagnostic dataset from Tan Tock Seng Hospital, outperforming existing RAG baselines in reducing misdiagnosis rates.

Beyond diagnostic support, medical RAG systems must operate under strict data localization constraints. One study developed a RAG system for recommending research collaborators based on PubMed publications, using PubMedBERT for domain-specific embeddings and a locally deployed LLaMA3 model for generation---demonstrating feasibility under fully local infrastructure requirements common in hospital settings \cite{ref557}. The broader medical evaluation landscape is captured by benchmarks like MedRGB, which tests models across sufficiency, integration, and robustness scenarios, revealing that current models exhibit limited ability to handle noise and misinformation in retrieved medical documents \cite{ref558}. Hybrid self-reflective approaches, such as Self-MedRAG, integrate sparse (BM25) and dense (Contriever) retrievers with Reciprocal Rank Fusion (RRF) and an iterative NLI-based self-reflection loop, achieving substantial accuracy gains on MedQA (80.00\% to 83.33\%) and PubMedQA (69.10\% to 79.82\%), confirming that iterative evidence verification enhances clinical reliability \cite{ref559}. The integration of standardized health data exchange formats is also emerging: FHIR-RAG-MEDS combines HL7 FHIR resources with RAG to enhance personalized medical decision support grounded in evidence-based clinical guidelines \cite{ref560}.

\textbf{Legal Domain Applications}

Legal RAG systems face distinct challenges: the structured, hierarchical nature of legal provisions, the need for clause-level precision, and the requirement for citation-grounded outputs that comply with legal norms. Legal-DC addresses the gap in Chinese legal RAG evaluation by constructing a benchmark of 480 legal documents and 2,475 question-answer pairs with clause-level annotations \cite{ref561}. The accompanying LegRAG framework integrates legal adaptive indexing (clause-boundary segmentation) with a dual-path self-reflection mechanism, outperforming state-of-the-art methods by 1.3\% to 5.6\% across key metrics. ClaimRAG-LAW extends legal RAG evaluation to French and English, targeting both expert and non-expert users with diverse question types, and applies a fine-grained claim-level evaluation framework that reveals limitations in retrieval, generation, and claim-level analysis \cite{ref395}. For litigation support, a Japanese medical litigation RAG system explicitly addresses the principle prohibiting private knowledge use, requiring the retrieval module to reference external knowledge with appropriate timestamps and to generate responses strictly faithful to the retrieved context \cite{ref562}. The LegalBench-RAG benchmark focuses on precise retrieval of minimal, highly relevant text segments rather than document-level chunks, producing 6,858 query-answer pairs over a 79M-character corpus, with human-annotated ground truth to support citation generation for end users \cite{ref563}. A reasoning-focused legal retrieval benchmark introduces Bar Exam QA and Housing Statute QA, constructed through annotation processes resembling real legal research, confirming that legal RAG remains a challenging application \cite{ref564}. Additionally, MVRAG introduces intention-aware query rewriting from multiple domain viewpoints for knowledge-dense domains like law, demonstrating significant improvements in recall and precision for legal and medical case retrieval \cite{ref396}.

\textbf{Financial Services Applications}

The financial sector demands RAG systems capable of handling heterogeneous data (text, tables, diagrams), evolving regulatory standards, and strict compliance requirements. SMARTFinRAG addresses the evaluation gap in financial RAG by introducing a fully modular architecture where components can be dynamically interchanged at runtime, combined with a document-centric evaluation paradigm that generates domain-specific QA pairs from newly ingested financial documents \cite{ref397}. FinSage proposes a multi-aspect RAG framework for regulatory compliance analysis in multi-modal financial filings, incorporating a multi-modal pre-processing pipeline that generates chunk-level metadata summaries, a multi-path sparse-dense retrieval with HyDE query expansion, and a domain-specialized re-ranking module fine-tuned via Direct Preference Optimization (DPO) \cite{ref398}. Deployed as a financial QA agent in online meetings serving over 1,200 people, FinSage achieves 92.51\% recall on expert-curated questions and surpasses FinanceBench baselines by 24.06\% in accuracy. AstuteRAG-FQA tackles proprietary data challenges in financial question answering through a task-aware prompt engineering framework that classifies queries into four tiers (explicit factual, implicit factual, interpretable rationale, hidden rationale), with a hybrid retrieval strategy combining open-source and proprietary data while maintaining strict security protocols including differential privacy and role-based access controls \cite{ref399}.

Peer-aware comparative reasoning represents an emerging paradigm in financial RAG. Rather than retrieving isolated facts, a contrastive insight layer on top of RAG retrieves comparable cases and highlights nuances between peer companies, generating context-specific risk insights that outperform baseline RAG on ROUGE and BERTScore against human-generated equity research \cite{ref565}. For fintech applications specifically, an agentic RAG architecture with modular specialized agents---including query reformulation, iterative sub-query decomposition guided by keyphrase extraction, contextual acronym resolution, and cross-encoder re-ranking---demonstrates improved retrieval precision over standard RAG, albeit with increased latency \cite{ref566}.

\textbf{Industrial and Enterprise Applications}

Industrial settings impose additional constraints: multilingual OCR requirements, heterogeneous documentation formats, and the need for production-grade performance metrics. An empirical evaluation of RAG for automotive requirements engineering using authentic manufacturing documentation (669 requirements across four specification standards) achieved 98.2\% extraction accuracy with complete traceability, outperforming BERT-based and ungrounded LLM baselines by 24.4\% and 19.6\%, respectively \cite{ref567}. The system achieved 83\% reduction in manual analysis time and 47\% cost savings through multi-provider LLM orchestration, with hybrid semantic-lexical retrieval reaching an MRR of 0.847. For enterprise report generation tasks, Secure Multifaceted-RAG (SecMulti-RAG) augments internal document retrieval with pre-generated expert knowledge and external LLM-generated knowledge, but uses a local open-source generator and filtering mechanisms to prevent data leakage---achieving 79.3 to 91.9\% win rates over traditional RAG in LLM-based evaluation \cite{ref505}. EASI-RAG provides a structured agile method for deploying RAG in industrial SMEs, validated in an environmental testing laboratory where the system was deployed in under a month by teams with no prior RAG experience \cite{ref568}.

Cybersecurity applications benefit from RAG-enhanced multi-agent incident response. AutoBnB-RAG extends the Backdoors \& Breaches tabletop game with retrieval-augmented generation, enabling agents to issue retrieval queries and incorporate external evidence during collaborative investigations \cite{ref569}. Evaluated across eight team structures, retrieval augmentation improved decision quality and success rates, with demonstrated ability to reconstruct complex multi-stage attacks from public breach reports. In the e-commerce domain, BSharedRAG proposes a backbone-shared framework where a domain-specific backbone model is continually pre-trained, then two plug-and-play LoRA modules minimize retrieval and generation losses respectively, outperforming baselines by up to 13\% in Hit@3 for retrieval and 23\% in BLEU-3 for generation \cite{ref570}.

Across these sectors, several consistent themes emerge: domain-specific retrieval strategies are non-negotiable, privacy and security constraints frequently dictate on-premises or federated architectures, and evaluation must move beyond generic accuracy metrics toward domain-grounded, fine-grained assessment \cite{ref401,ref571}. The ``seven failure points'' identified in engineering RAG systems---including validation only feasible during operation and robustness evolving rather than being designed in---underscore that industrial deployments require continuous monitoring and iterative refinement \cite{ref572}. These findings directly motivate the systematic engineering methodologies, deployment frameworks, and operational practices that constitute the focus of the following subsection, where we examine how enterprises operationalize RAG systems to meet production-grade requirements.

\subsection{Deployment Frameworks and Engineering Methodologies}

The successful deployment of Retrieval-Augmented Generation (RAG) systems in production environments hinges not only on algorithmic innovation but equally on robust engineering methodologies, structured deployment frameworks, and systematic operational practices. While much of the academic literature focuses on improving retrieval accuracy or generation quality in isolation, the transition from research prototype to enterprise-grade system introduces a distinct set of technical and organizational challenges. These include data preprocessing at scale, integration with legacy IT infrastructure, continuous monitoring, and the management of evolving knowledge bases. As the preceding subsection demonstrated, domain-specific deployments across medicine, law, finance, and industry consistently surface recurring engineering patterns---privacy-driven architectural choices, fine-grained evaluation needs, and iterative refinement requirements. This subsection synthesizes the current landscape of deployment frameworks and engineering methodologies that operationalize these patterns, drawing on agile methods tailored for resource-constrained settings, comprehensive blueprints for on-premises deployment, modular toolkits designed for research and prototyping, and empirical lessons learned from multi-domain industrial implementations. In doing so, it bridges the gap between the domain-grounded systems described earlier and the evaluation paradigms examined in the subsequent subsection, where continuous monitoring and diagnostic frameworks become essential to sustaining production-grade performance.

One prominent challenge in industrial adoption is the resource disparity between large technology companies and Small and Medium Enterprises (SMEs). SMEs often lack the specialized NLP expertise and computational infrastructure required for sophisticated RAG deployments. To address this, the EASI-RAG method (Enterprise Application Support for Industrial RAG) was proposed as a structured, agile methodology explicitly designed for industrial SME contexts \cite{ref568}. EASI-RAG is grounded in method engineering principles and comprises well-defined roles, activities, and techniques. Its practical viability was demonstrated through a real-world case study in an environmental testing laboratory, where a RAG tool was implemented to answer operator queries using data extracted from operational procedures. Significantly, the system was deployed in under a month by a team with no prior RAG experience, and it was subsequently refined iteratively based on user feedback. The case study reported fast implementation, high user adoption, and accurate answers, highlighting the potential of structured agile methods to lower the barrier to entry for RAG adoption in industrial settings \cite{ref568}. This agile approach directly addresses the ``seven failure points'' identified in the previous subsection, particularly the need for continuous validation and iterative refinement that only become feasible during live operation.

While agile methods address the ``how'' of incremental development, organizations with stringent data governance requirements also need comprehensive architectural guidance for ``where'' and ``with what'' to deploy. This is particularly acute for on-premises deployments, where data protection regulations---such as those governing the medical and legal domains discussed earlier---prohibit the use of cloud-based services. Addressing this gap, an AI engineering blueprint for scalable on-premises enterprise RAG solutions was introduced \cite{ref573}. This blueprint provides: (1) an end-to-end reference architecture described using the 4+1 view model, which facilitates communication among stakeholders with different concerns; (2) a reference application for on-premises deployment; and (3) best practices for tooling, development, and CI/CD pipelines. By making these artifacts publicly available, the blueprint offers a concrete starting point for organizations navigating the complexities of self-hosted RAG systems, particularly in regulated industries. The emphasis on the 4+1 view model is noteworthy, as it aligns with established software architecture documentation practices, thereby bridging the gap between AI research and conventional software engineering \cite{ref573}. Complementary to this, the broader operational lifecycle of such systems is addressed by the concept of RAGOps, which extends LLMOps by incorporating a strong focus on data management to handle the continuous changes in external data sources \cite{ref574}. RAGOps outlines the integrated management lifecycles of both the LLM and the data, defining key design considerations across different stages and highlighting the need for automated evaluation and testing of data operations to sustain retrieval relevance and generation quality---a theme that directly anticipates the evaluation frameworks discussed in the following subsection.

Beyond high-level methodologies and blueprints, the practical engineering of RAG systems is greatly facilitated by modular toolkits that abstract away implementation details and support rapid experimentation. For instance, FlexRAG is an open-source framework designed specifically for research and prototyping, supporting text-based, multimodal, and network-based RAG with comprehensive lifecycle support, efficient asynchronous processing, and persistent caching capabilities \cite{ref575}. By enabling researchers to rapidly develop, deploy, and share advanced RAG systems, FlexRAG addresses common difficulties such as algorithm reproduction, lack of new techniques, and high system overhead. Similarly, UltraRAG provides a modular and automated toolkit that automates knowledge adaptation throughout the entire workflow---from data construction and training to evaluation---while featuring a user-friendly WebUI that allows users to build and optimize systems without coding expertise \cite{ref576}. UltraRAG supports multimodal input and offers comprehensive tools for managing the knowledge base, positioning itself as an end-to-end development solution for diverse user scenarios. These toolkits are instrumental in democratizing RAG development, allowing practitioners to focus on domain-specific challenges rather than re-implementing core components, and they provide the component-level flexibility that the modular evaluation benchmarks in the following subsection assume as a prerequisite.

Nevertheless, even with structured methods and powerful toolkits, real-world deployments consistently reveal a set of recurring challenges that transcend individual domains. Empirical studies, such as the interview study with 13 industry practitioners, have identified that current RAG applications are often limited to domain-specific QA tasks and remain in prototype stages \cite{ref571}. The study highlighted that industry requirements focus heavily on data protection, security, and quality---priorities that mirror the privacy and security-driven architectural choices observed across the medical, legal, and financial sectors in the preceding subsection---while issues such as ethics, bias, and scalability receive comparatively less attention. Data preprocessing emerged as a key challenge, and system evaluation was found to be predominantly conducted by humans rather than automated methods \cite{ref571}. These findings are echoed in multi-domain development experiences, where five domain-specific RAG applications were built for governance, cybersecurity, agriculture, industrial research, and medical diagnostics \cite{ref401}. This work underscored the importance of integrating multilingual OCR, semantic retrieval via vector embeddings, and domain-adapted LLMs, and it documented twelve key lessons learned, highlighting technical, operational, and ethical challenges affecting reliability and usability. Furthermore, comparative studies of ``Enhanced'' and ``Agentic'' RAG paradigms provide empirical insights into the trade-offs between performance and cost, guiding practitioners in selecting the most effective design for real-world applications \cite{ref577}. For highly specialized industrial needs, such as requirements engineering in automotive manufacturing, comprehensive empirical evaluations have demonstrated that RAG frameworks can achieve significant gains in extraction accuracy and cost savings, provided that the system is meticulously designed and evaluated against production-grade metrics \cite{ref567}.

A critical dimension of production deployment is the establishment of a sound evaluation and testing strategy---an area that serves as the direct conceptual bridge to the following subsection on real-world evaluation and benchmarking. Traditional binary notions of correctness are insufficient for probabilistic AI systems. A comprehensive assurance strategy proposes that AI testing should focus on continuous risk reduction rather than strict correctness verification, introducing a structured AI Failure Taxonomy and a revised five-layer AI Assurance Pyramid \cite{ref578}. This philosophy is operationalized in frameworks like the LLM Readiness Harness, which combines automated benchmarks, OpenTelemetry observability, and CI quality gates to produce scenario-weighted readiness scores, enabling teams to block risky releases based on groundedness, retrieval hit rate, cost, and latency \cite{ref579}. Additionally, systematic methodologies for RAG evaluation, which include the crucial role of choosing appropriate baselines and metrics, the necessity for systematic refinements derived from qualitative failure analysis, and the reporting of key design decisions, are essential for sound and reliable results \cite{ref580}. These engineering-side evaluation practices complement the benchmark-driven and LLM-as-a-Judge approaches detailed in the subsequent subsection, together forming a comprehensive evaluation ecosystem that spans both development-time and operation-time assessment.

Finally, the engineering of production RAG systems must also consider architectural evolution, particularly as systems mature and operational demands grow. The journey from naive RAG to deep agentic retrieval, as documented in a production pipeline for regulatory compliance at Ontario Power Generation, illustrates a common trajectory \cite{ref24}. This case study formalizes a mature architecture as Progressive Evidence Acquisition with Cost-Aware Escalation (PEA-CAE), where the system begins with low-cost, high-precision retrieval and escalates to more expensive operations only when justified by expected evidence gain. This evolution underscores that context engineering---the iterative design of how information is retrieved, structured, and presented to the LLM---is often a more tractable and economically viable path than domain-specific fine-tuning for large, evolving corpora. The lessons from such deployments, combined with the structured methods, architectural blueprints, and toolkits discussed above, form a foundation for building reliable, maintainable, and effective RAG systems in real-world settings. They also establish the operational context that motivates the evaluation paradigms in the following subsection, where domain-specific benchmarks, black-box assessment tools, and LLM-as-a-Judge methodologies provide the diagnostic capabilities necessary to sustain and improve these deployed systems over time.

\subsection{Evaluation and Benchmarking in Real-World Settings}

The preceding subsections have established that deploying RAG in production is as much an exercise in continuous operational management as it is in algorithmic design. As systems transition from prototypes to live environments, the limitations of offline, accuracy-focused evaluation become increasingly apparent. The industrial experience reports discussed earlier reveal that validation is only truly feasible during live operation and that robustness must be continuously monitored. This subsection addresses that need by examining the evaluation paradigms and benchmarking frameworks specifically designed to assess RAG systems under the complexities of real-world deployment---where domain-specific nuances, multi-turn interactions, operational constraints, and the demand for interpretable diagnostics fundamentally reshape what constitutes meaningful evaluation.

\textbf{Domain-Specific Benchmarks for Industrial Vertical Markets.} The proliferation of RAG applications in specialized domains has driven the creation of dedicated benchmarks that reflect the unique linguistic, structural, and reasoning challenges of each field. In biomedicine, the Medical Retrieval-Augmented Generation (MRAG) benchmark addresses the paucity of comprehensive evaluation resources in the medical domain by covering diverse tasks in both English and Chinese, leveraging a corpus constructed from Wikipedia and PubMed, and providing the MRAG-Toolkit for systematic component exploration \cite{ref367}. Similarly, the chemistry domain has seen the introduction of ChemRAG-Bench, a comprehensive benchmark with an accompanying heterogeneous corpus integrating scientific literature, PubChem, PubMed abstracts, textbooks, and Wikipedia entries, systematically demonstrating substantial performance gains from RAG over direct inference \cite{ref400}.

In the legal domain, where the precision requirements discussed in the previous subsection are especially acute, several benchmarks have emerged to address the critical gap in evaluating retrieval components specifically. LegalBench-RAG is the first benchmark designed to evaluate the retrieval step of RAG pipelines within legal spaces, emphasizing precise retrieval of minimal, highly relevant text segments---a crucial requirement in legal applications where long context windows are costly and induce hallucinations \cite{ref563}. Building on this, Legal RAG Bench offers an end-to-end benchmark and evaluation methodology for legal RAG, employing a full factorial design and a novel hierarchical error decomposition framework to enable apples-to-apples comparisons of retrieval and reasoning model contributions \cite{ref581}. Extending this to a multilingual and non-expert context, ClaimRAG-LAW provides a comprehensive dataset for legal RAG in French and English, targeting both experts and non-experts with diverse question types and enabling claim-level analysis \cite{ref395}.

The finance sector has also produced dedicated evaluation resources. OmniEval presents an omnidirectional and automatic RAG benchmark in the financial domain, featuring a matrix-based scenario system categorizing queries into five task classes and sixteen financial topics, combined with a multi-stage evaluation system that assesses both retrieval and generation performance \cite{ref582}. For a more interactive and modular approach, SMARTFinRAG introduces a financial RAG benchmark with a fully modular architecture allowing dynamic component interchange at runtime, a document-centric evaluation paradigm for generating QA pairs from newly ingested financial documents, and an intuitive interface bridging research and implementation divides \cite{ref397}. In the semiconductor manufacturing vertical, FAB-Bench provides an end-to-end framework for adaptive benchmarking of RAG systems, defining six diagnostic metrics and revealing context-scaling behaviors such as logarithmic growth and attention dilution at extreme context lengths \cite{ref419}.

Perhaps the most industrially relevant benchmark is RAGBench, which stands out as the first comprehensive, large-scale RAG benchmark dataset comprising 100,000 examples across five unique industry-specific domains and various RAG task types, sourced from corpora like user manuals \cite{ref583}. This benchmark is paired with the TRACe evaluation framework, a set of explainable and actionable metrics that enable holistic evaluation and continuous improvement of production applications---a critical capability for industry deployment \cite{ref583}. Together, these domain-specific benchmarks provide the foundation for rigorous, context-aware evaluation that complements the operational failure-point analysis emphasized in the prior subsection.

\textbf{Black-Box Assessment Tools for Deployed Systems.} Beyond domain benchmarks, industrial practitioners require tools to evaluate operational RAG systems without intrusive instrumentation---a need directly aligned with the continuous monitoring requirements highlighted in the deployment frameworks discussed earlier. SCARF (System for Comprehensive Assessment of RAG Frameworks) addresses this need as a modular and flexible evaluation framework designed to benchmark deployed RAG applications systematically \cite{ref584}. It provides an end-to-end, black-box evaluation methodology enabling limited-effort comparison across diverse RAG frameworks, supporting multiple deployment configurations and automated testing across vector databases and LLM serving strategies \cite{ref584}. Similarly, RAGe offers a modular framework for benchmarking and guiding the efficient development of RAG applications, focusing on resource telemetry and component recommendation to evaluate trade-offs among accuracy, efficiency, and scalability, even on consumer-grade hardware \cite{ref585}. These tools are essential for facilitating rapid prototyping and systematic performance engineering in production environments \cite{ref585}. Complementing these, RAGPerf provides an end-to-end benchmarking framework that decouples the RAG workflow into modular components---embedding, indexing, retrieval, reranking, and generation---allowing detailed profiling and fine-grained performance analysis with negligible overhead \cite{ref586}.

\textbf{LLM-as-a-Judge Methodologies for Enterprise Multi-Turn Systems.} The complexity of enterprise RAG assistants operating in multi-turn, case-based workflows---such as technical support or IT operations---demands evaluation frameworks that capture operational constraints, structured identifiers, and resolution workflows, moving beyond the static benchmarks described above. A case-aware LLM-as-a-Judge evaluation framework has been proposed for enterprise multi-turn RAG systems, evaluating each turn using eight operationally grounded metrics that separate retrieval quality, grounding fidelity, answer utility, precision integrity, and case/workflow alignment \cite{ref420}. This framework uses deterministic prompting with strict JSON outputs, enabling scalable batch evaluation, regression testing, and production monitoring, while a severity-aware scoring protocol reduces score inflation and improves diagnostic clarity \cite{ref420}. For multi-turn conversational benchmarks, MTRAG provides an end-to-end human-generated multi-turn RAG benchmark reflecting real-world properties, demonstrating that even state-of-the-art LLM RAG systems struggle with later turns, unanswerable questions, non-standalone questions, and multiple domains \cite{ref587}.

General-purpose LLM-as-a-Judge frameworks have also been adapted for comprehensive RAG evaluation, offering scalable alternatives to human annotation that align with the automated testing strategies advocated in the deployment methodologies. CCRS (Contextual Coherence and Relevance Score) offers a zero-shot LLM-as-a-Judge suite of five metrics---Contextual Coherence, Question Relevance, Information Density, Answer Correctness, and Information Recall---using a single pretrained LLM as an end-to-end judge \cite{ref406}. Compared to complex frameworks like RAGChecker, CCRS provides comparable or superior discriminative power for key aspects such as recall and faithfulness while being significantly more computationally efficient \cite{ref406}. For domain-specific evaluation with human alignment, RAGalyst presents an automated, human-aligned agentic framework that generates high-quality synthetic QA datasets and refines LLM-as-a-Judge metrics (Answer Correctness and Answerability) through prompt optimization to achieve strong correlation with human annotations \cite{ref588}.

\textbf{Role of Human Evaluation and Trade-offs.} Despite advances in automation, the role of human evaluation remains critical in real-world settings, particularly given the human-AI interaction concerns and trust issues identified in the following subsection. Human pairwise judgments have been shown to provide reliable and cost-effective results compared to LLM-based pairwise judgments or automated comparisons with reference responses \cite{ref589}. However, human and LLM ratings often lack agreement, and humans struggle to focus strictly on metric descriptions \cite{ref78}. Metrics must therefore balance interpretability, efficiency, and fidelity, with frameworks like RAGXplain translating performance metrics into actionable guidance through a `Metric Diamond' structure \cite{ref415}. Ultimately, the trade-offs between performance, efficiency, and interpretability must be carefully managed, with evaluation frameworks providing diagnostic clarity without imposing prohibitive computational costs, a balance exemplified by tools like RAGVUE and InspectorRAGet \cite{ref417,ref590}. These evaluation strategies provide the empirical foundation for identifying the failure modes and success factors that the subsequent subsection will explore in detail, thereby completing the bridge from deployment engineering to operational learning.

\subsection{Lessons Learned, Failure Points, and Emerging Industrial Practices}

Building upon the foundational evaluation frameworks and benchmarking methodologies established in the preceding subsection, this subsection transitions from the question of \emph{how} to measure RAG systems to the question of \emph{what} those measurements reveal about real-world industrial deployment. The shift from controlled experimental settings to live production environments consistently exposes a substantial gap between theoretical promise and operational reality. Synthesizing findings from case studies, user studies, and engineering experience reports, this subsection distills the recurring failure modes that undermine deployed systems, identifies the critical success factors that separate robust implementations from brittle ones, and surveys the emerging engineering practices that are shaping the next generation of industrial RAG architectures.

A central theme emerging from industrial experience is that RAG systems are fundamentally brittle in production environments. A pivotal experience report, based on three case studies across research, education, and biomedical domains, codified this fragility into seven distinct failure points that software engineers must consider when designing RAG systems \cite{ref572}. These failure points span the entire pipeline, from document ingestion and chunking strategies to retrieval quality, prompt construction, and the final generation stage. Critically, this work concluded that validation of a RAG system is only truly feasible during live operation, and that robustness is an evolving property that must be continuously monitored and adapted, rather than a static characteristic that can be designed in at the outset \cite{ref572}. This finding directly challenges the assumption---implicit in many of the offline benchmarks discussed in the previous subsection---that static evaluation on curated datasets is sufficient to guarantee production readiness.

Echoing these findings, empirical studies have underscored the highly selective nature of RAG deployment. A systematic investigation across multiple LLMs and datasets revealed that deploying RAG is not universally beneficial; in some cases, it can harm performance for up to 12.6\% of samples even when the retrieved documents are perfectly relevant \cite{ref62}. This suggests that RAG can introduce a form of ``noise'' or misalignment that conflicts with a model's parametric knowledge, leading to a performance degradation that is more pronounced than simply generating without retrieval. Furthermore, the optimal amount of retrieved context is task-dependent, with question-answering tasks showing a universal sweet spot of 5--10 documents, while code generation tasks require more nuanced, scenario-specific optimization \cite{ref62}. These findings argue for context-aware design decisions rather than a one-size-fits-all RAG architecture, and they underscore the necessity of a ``selective deployment'' mentality where the decision to use RAG is as critical as the choice of its components---a nuance that the domain-specific benchmarks in the previous subsection only partially capture.

Recurring failure modes in real-world deployments extend beyond general inaccuracies to specific operational vulnerabilities. Hallucination remains the paramount concern, manifesting not only as fabrication but also as a failure to adequately ground responses in the provided context. Evaluations of production systems have highlighted issues such as off-topic responses, failed citations, and inappropriate abstentions, necessitating the robust evaluation protocols combining human annotations with LLM-based judges that were surveyed in the prior subsection \cite{ref591}. In high-stakes environments like industrial troubleshooting for complex cyber-physical systems, RAG has shown promise in assisting operators by retrieving relevant procedures, yet the need for cross-validation of recommendations before they are acted upon was identified as a critical safety measure \cite{ref592}. In the legal domain, where the precision requirements are most acute, RAG systems face unique challenges related to lexical redundancy in corpora and the generation of hallucinated clauses or precedents, with small, locally deployed models exacerbating these issues due to their limited capacity \cite{ref593}.

Retrieval noise and the quality of the knowledge base itself constitute another major class of industrial failure points. The assumption that all retrieved information is beneficial is demonstrably false; a comprehensive benchmark introduced seven distinct types of noise from a linguistic perspective, showing that while some noise can be paradoxically beneficial, other types are categorically harmful to LLM performance \cite{ref64}. This finding complicates the design of retrieval filters and highlights the need for noise-aware generation strategies. Furthermore, the structure of the knowledge base is paramount. In software bug diagnosis, a RAG system that indexed resolved issues as structured symptom-root cause-resolution triples significantly outperformed a system relying on unstructured document chunks, demonstrating that the way knowledge is organized is as important as the retrieval algorithm itself \cite{ref594}. This experience highlights that operational failure history, often an underutilized asset, can serve as a superior knowledge source compared to generic system documentation---a lesson that directly informs the knowledge graph approaches discussed later.

Privacy and security are emerging as non-negotiable constraints in industrial RAG deployments. The integration of external, often sensitive, knowledge bases introduces new attack surfaces that extend well beyond the traditional concerns of LLM deployments. The ``ConfusedPilot'' class of vulnerabilities demonstrated how injected malicious text can corrupt responses and how the caching mechanisms used for efficiency can be exploited to leak confidential enterprise data \cite{ref472}. Beyond intentional attacks, the very nature of RAG makes it susceptible to membership inference attacks, where an adversary can determine if a specific document is present in the knowledge base, threatening data producers' rights \cite{ref474,ref71}. These findings underscore that security must be a foundational design consideration from the outset, as retrofitting defenses onto a vulnerable architecture is both complex and often cost-prohibitive.

From these challenges, a set of key success factors and emerging engineering practices has crystallized. User studies have shown that even when RAG systems are accurate, they can introduce cognitive friction. For instance, while tailored warnings about potential hallucinations can improve user accuracy, they can also lead to confusion and diminished trust in the system, highlighting the delicate balance between system transparency and user experience \cite{ref595}. Therefore, the design of the human-AI interaction loop is a critical success factor. For Small and Medium Enterprises (SMEs) with limited resources, agile and structured deployment methodologies are essential. The EASI-RAG method, validated in an environmental testing laboratory, demonstrated that a team with no prior RAG experience could deploy a functional tool in under a month, achieving high user adoption and accurate answers, using a structured method-engineering approach \cite{ref568}. This suggests that with the right scaffolding, the deployment barrier can be significantly lowered.

Emerging industrial practices are moving toward more sophisticated, context-aware architectures---a trajectory that directly connects to the evaluation metrics and diagnostic tools described in the previous subsection. One key trend is the integration of application-aware reasoning, often termed ``RAG+,'' which moves beyond simple retrieval-then-read to incorporate task-specific logic and multi-step workflows. Another is the development of case-based reasoning, where retrieval is structured around prior resolved cases rather than raw text chunks, as demonstrated by the operational memory RAG system in production for over six months \cite{ref594}. This approach aligns with the move towards ``knowledge graphs for next-generation industrial RAG solutions,'' where structured representations of entities and their relationships enable more precise, multi-hop reasoning and mitigate the noise inherent in unstructured text retrieval. For example, enhancing the non-parametric data store with a knowledge graph has been shown to improve the analytical and semantic question-answering capabilities in structured analyses like Failure Mode and Effects Analysis (FMEA) \cite{ref596}. Finally, the increasing awareness of security has led to frameworks that restrict data exposure, such as dynamically integrating external LLM-generated knowledge only when queries are deemed safe by a filtering mechanism \cite{ref505}.

In conclusion, the collective lesson from industrial RAG deployments is that a successful system is not merely a sum of powerful components, but a carefully calibrated integration of retrieval, generation, human factors, and robust security that aligns with the evaluation-driven, continuously monitored operational paradigm established in the prior subsection. These operational insights---the failure modes, the success factors, and the emerging architectural patterns---form the empirical foundation for the subsequent discussion of user trust, human-AI interaction, and the broader implications of deploying RAG systems in sociotechnical contexts where their impact extends far beyond mere technical performance.

\section{Open Challenges, Ethical Considerations, and Future Directions}

\subsection{Persistent Technical Challenges in RAG Systems}

Despite the rapid advancement of Retrieval-Augmented Generation (RAG) systems across architectural paradigms, application domains, and optimization frameworks, a set of persistent and deeply intertwined technical challenges continues to constrain their effectiveness, reliability, and practical deployment. These bottlenecks are not merely engineering inconveniences but fundamental limitations that arise from the interaction between parametric memory, non-parametric retrieval, and complex reasoning requirements. In this subsection, we synthesize the most critical unresolved technical challenges that define the current frontier of RAG research, thereby motivating the architectural innovations and training paradigms discussed in the following subsection.

A first fundamental challenge lies in the degradation of retrieval quality as both corpora and query complexity scale. Conventional wisdom often assumes that more data and more sophisticated retrieval mechanisms lead to better performance. However, controlled scaling studies have revealed a more nuanced and scale-dependent reality. A systematic comparison of RAG paradigms across 28 nested corpus tiers spanning a 450-fold size range showed that the relative effectiveness of retrieval strategies is not fixed; a scale-dependent crossover occurs where lexical methods like BM25 overtake agentic and graph-based approaches as the corpus grows, with agentic search becoming less effective and significantly more token-expensive as the search space expands \cite{ref100}. This finding challenges the assumption that more sophisticated retrieval automatically translates to better performance at scale and points to a pressing need for retrieval strategies that maintain precision without incurring prohibitive computational costs. Furthermore, retrieval effectiveness is not merely a function of the retriever itself. Systematic analyses of RAG pipelines have revealed that retrieval-oriented metrics can improve under broader retrieval settings while downstream generation metrics behave non-monotonically, indicating that optimizing for retrieval quality alone does not guarantee improvements in end-to-end task performance \cite{ref597}. This disconnect between retrieval objectives and generation outcomes remains a core unresolved issue.

Second, the way retrieved information is packed into the model's context introduces severe and persistent biases. The ``lost-in-the-middle'' phenomenon, where models disproportionately attend to information at the beginning and end of the input context while ignoring the middle, remains a salient challenge even in modern LLMs. While reproducibility studies show that the magnitude of this positional bias depends heavily on evaluation choices such as topic sampling and retrieval quality, the effect is real and interacts strongly with retrieval order and context size under imperfect retrieval conditions \cite{ref598}. The problem is further exacerbated by the presence of distractors---semantically irrelevant passages that slip through the retrieval filter. Research has demonstrated that a static top-k retrieval strategy either misses relevant information or introduces such distractors, which can actively degrade output quality, and that the position of relevant passages within the context window significantly modulates their influence on generation \cite{ref155}. Moreover, the effectiveness of context utilization is not uniform across model scales. Empirical studies with oracle retrieval---where the retrieved passage is guaranteed to contain the answer---reveal that models of 7B parameters or smaller fail to extract the correct answer 85\% to 100\% of the time on questions they cannot answer alone, and that adding retrieval context can destroy 42\% to 100\% of answers the model previously knew, a distraction effect driven by the presence of context rather than its quality \cite{ref599}. This points to a fundamental utilization bottleneck, where the principal limitation of RAG for smaller models is not retrieval quality but the ability to effectively process and extract information from the provided context.

Third, RAG systems must navigate the delicate and often conflicting relationship between parametric memory (knowledge stored in model weights during pretraining) and non-parametric knowledge (information retrieved at inference). Dynamic knowledge conflicts arise when the retrieved evidence contradicts or is inconsistent with the model's internal knowledge. Under realistic conditions, imperfect retrieval augmentation is inevitable, common, and harmful, and knowledge conflicts between LLM-internal and external knowledge are identified as a key bottleneck in the post-retrieval stage of RAG \cite{ref215}. Models must decide which source to trust, and naive merging strategies fail to handle the full spectrum of real-world scenarios where one source is unreliable or the two conflict \cite{ref600}. The problem is compounded by the earlier finding that even when models are prompted to use confidence scores provided with retrieved documents, they comply with the instruction but produce worse answers---a gap between utilization and accuracy that remains unmeasured and poorly understood in prior work \cite{ref601}. Moreover, the optimal balance between parametric and non-parametric knowledge is not fixed. Scaling law studies exploring the trade-off between pretraining corpus size and retrieval store size show that the marginal utility of retrieval depends strongly on model scale, task type, and the degree of pretraining saturation, though the acquisition of a fully generalizable theory of optimal allocation remains an open question \cite{ref225}.

Fourth, there is a persistent and unresolved tension between retrieval precision and generation flexibility. On one hand, retrieving too few passages risks missing critical evidence, while retrieving too many overwhelms the prompt window, diluting relevance and increasing cost \cite{ref602}. On the other hand, the static retrieve-then-generate pipeline is suboptimal for complex tasks requiring multi-hop reasoning and adaptive information access \cite{ref19}. The decision of how much information to retrieve is task-dependent and difficult to universalize: QA tasks show universal patterns where 5-10 documents are optimal, while code generation requires scenario-specific optimization, and universal RAG strategies are demonstrably inadequate \cite{ref62}. Even the definition of a ``good'' retrieval metric is challenged under budget constraints. Standard document recall is the wrong quantity to optimize when only a fraction of retrieved evidence can be shown to the reader under a token budget; a novel diagnostic, ``answer-in-context,'' which measures whether the gold answer survives as a contiguous span in the packed context, is a better predictor of answer F1 and carries information beyond retrieval metrics \cite{ref603}. The interaction between retrieval precision and the reader's capacity to use the retrieved evidence further complicates this trade-off. Controlled experiments show that a packing change that raises coverage but not answer-in-context yields no accuracy gain, highlighting that evidence density and reader capacity, not just retrieval quality, are the true bottlenecks \cite{ref603}.

Fifth, scalability constraints extend beyond retrieval accuracy and the parametric/non-parametric trade-off to fundamental computational and memory costs. RAG introduces long-context processing requirements, significantly increasing memory consumption and inference latency, particularly in the prefill stage \cite{ref523,ref510}. Systematic characterization studies identify long sequence generation due to knowledge injection as the primary performance bottleneck in RAG, and the memory footprint of KV caches presents a substantial deployment challenge \cite{ref510}. The cost-latency trade-off is not monotonic or predictable: gating based on calibrated confidence can increase latency by 27\% on one model size while saving 8\% on another, demonstrating that system-level effects are model-dependent and often counter-intuitive \cite{ref604}. Static retrieval configurations cannot resolve the three-way tension between deeper retrieval for factual grounding, inflated token costs, and end-to-end latency across heterogeneous workloads \cite{ref605}.

Finally, the brittleness of adaptive retrieval mechanisms themselves constitutes a significant and under-appreciated challenge. Adaptive RAG methods, which dynamically trigger retrieval only when needed, are highly vulnerable to query variations that preserve identical semantics. Small surface-level changes in query formulation can dramatically alter retrieval behavior and accuracy, and while larger models show better overall performance, their robustness does not improve accordingly \cite{ref66}. Similarly, retrievers degrade significantly under even minor query perturbations, and the sensitivity of each module in the RAG pipeline varies \cite{ref606}. Retrieval fusion techniques, such as multi-query retrieval and reciprocal rank fusion, often fail to deliver end-to-end gains in production settings; recall-oriented fusion techniques exhibit diminishing returns once realistic re-ranking limits and context budgets are applied, and retrieval-level improvements do not reliably translate into downstream impact \cite{ref607}. Furthermore, adaptive retrieval mechanisms that rely on uncertainty signals can be miscalibrated. Calibration of sequence log-probability and prefix-logit uncertainty signals into probabilities of correctness significantly improves decision quality, yet a significant gap remains between ideal calibration and real-world performance, particularly in production environments with heterogeneous queries \cite{ref604}. Even when sub-query rewriting is used, increasing the number of documents used for reranking and generation beyond a certain point reduces effectiveness, and there are hard limits on the benefits of rewrites \cite{ref608}.

The persistent technical challenges outlined above---retrieval quality degradation at scale, positional and distraction biases, dynamic knowledge conflicts, the precision-flexibility trade-off, scalability and cost-latency constraints, and the brittleness of adaptive mechanisms---form an interconnected web of limitations. No single-stage optimization can resolve these issues in isolation. The field must move toward holistic, systems-level solutions that jointly consider retrieval, context construction, generation, and efficiency. It is precisely these limitations that motivate the research frontiers explored next: the development of next-generation adaptive retrieval mechanisms capable of self-evaluation, the integration of causal structures for deeper reasoning, the application of process-level reinforcement learning for training agentic systems, and the pursuit of self-evolving knowledge representations. Each of these directions directly responds to the bottlenecks identified here---from the need for cost-aware and selective retrieval that addresses scalability and brittleness, to the demand for structured reasoning that mitigates knowledge conflicts and the precision-flexibility trade-off, to the vision of knowledge systems that continuously refine themselves in response to the very challenges of scale and dynamic information.

\subsection{Emerging Research Frontiers and Next-Generation Architectures}

The evolution of Retrieval-Augmented Generation is accelerating toward a new generation of architectures that transcend the static retrieve-then-read paradigm. As the limitations of fixed pipelines become increasingly apparent, several interconnected research frontiers are emerging that promise to redefine how RAG systems perceive, interact with, and reason over external knowledge. These frontiers include next-generation adaptive retrieval mechanisms, the deep integration of causal structures, reinforcement learning with process-level rewards, the confluence of RAG with long-context and agentic AI systems, and the development of self-evolving knowledge representations. Critically, each of these frontiers directly responds to the bottlenecks identified in the previous subsection: the need for cost-aware and selective retrieval that addresses scalability and brittleness, the demand for structured reasoning that mitigates knowledge conflicts and the precision-flexibility trade-off, and the vision of knowledge systems that continuously refine themselves in response to the very challenges of scale and dynamic information.

A fundamental shift in current research is the movement away from uniform, static retrieval toward next-generation adaptive mechanisms capable of self-evaluation. Traditional RAG pipelines, as well as many early agentic systems, operate with predetermined workflows or heuristic routing rules that fail to account for query complexity or the model's own epistemic state. Research has demonstrated that such rigid designs often produce suboptimal outcomes; for instance, query decomposition, a common adaptive enhancement, yields consistent gains in structured domains but can degrade ranking precision in multi-hop benchmarks, while reflection mechanisms improve citation accuracy only at substantial latency costs \cite{ref31}. Similarly, controlled ablations reveal that fixed hybrid retrieval via reciprocal rank fusion can outperform rule-based adaptive routing, and that most of the performance gains are realized through short retrieval loops rather than deep iterative expansion \cite{ref305}. These findings highlight the need for truly adaptive systems that are cost-aware, selective about when to retrieve, and capable of calibrating their search behavior based on the difficulty of the query.

Central to this new adaptive paradigm is the notion of self-awareness---the ability of a RAG system to recognize its own knowledge boundaries and make informed decisions about whether to retrieve, how much to retrieve, and when to stop, as well as to verify and calibrate the knowledge it has already acquired. Recent work has begun to formalize this capability through the lens of query performance prediction, demonstrating that the predicted quality of generated queries is positively correlated with the quality of the final answer, offering a pathway toward informed retrieval decisions within agentic workflows \cite{ref319}. Frameworks such as SIM-RAG propose lightweight critics trained through self-practicing to explicitly evaluate whether sufficient information has been retrieved at each round, enabling systems to avoid unnecessary searches while ensuring they do not stop prematurely \cite{ref282}. This pursuit of self-skepticism---the capacity to question one's own confidence---is emerging as a cornerstone for reliable adaptive retrieval, allowing models to operate with a clearer sense of when their parametric knowledge is sufficient and when external evidence is indispensable \cite{ref609}.

The integration of causal structures and knowledge graphs represents another promising frontier, moving beyond associative correlations toward a deeper understanding of the mechanisms underlying the retrieved evidence. Current RAG systems predominantly rely on semantic similarity, which often fails to distinguish spurious associations from true causal relationships, yielding responses that are factually grounded yet logically incomplete. Frameworks like CDF-RAG introduce causal dynamic feedback loops that iteratively refine queries, retrieve structured causal graphs, and validate responses against causal pathways, thereby improving both causal consistency and explainability \cite{ref610}. This direction is complemented by a growing body of work on Graph-based RAG, which leverages entity-relation structures to support multi-hop reasoning. However, conventional graph approaches are often constrained by static graph construction and one-time retrieval. Emerging systems such as Graph-R1 reconceptualize knowledge graphs as dynamic environments, modeling retrieval as a Markov Decision Process where an agent observes the graph state and executes actions to query, expand, refine, or terminate, thereby enabling the graph itself to evolve through interaction \cite{ref306}. Similarly, EvoGraph-R1 extends this self-evolving paradigm to multimodal knowledge hypergraphs, demonstrating that graphs which adapt through iterative interaction can support more robust reasoning across modalities \cite{ref42}. The key insight across these works is that knowledge graphs should not be treated as static artifacts but as living structures that grow, correct, and refine themselves through continued use.

Reinforcement learning is rapidly becoming the dominant paradigm for training these next-generation agentic retrieval systems, and a critical frontier lies in the design of reward signals. Early agentic RAG systems, following the success of outcome-based RL approaches, rely on the correctness of the final answer as the sole reward. While effective, this sparse outcome-based supervision suffers from low sample efficiency, gradient conflict, and weak guidance for intermediate reasoning actions \cite{ref317,ref315}. The field is therefore converging on process-level rewards that provide dense, step-wise feedback on the quality of intermediate reasoning and retrieval decisions. Hierarchical process rewards, for instance, evaluate the necessity of each search decision by decomposing the reasoning trajectory into discrete steps, rewarding optimal search and non-search behavior to reduce both over- and under-searching \cite{ref316}. Path-centric reward shaping assesses the structural quality of reasoning trajectories through order-agnostic step coverage, extracting learning signals even from failed samples and improving training stability \cite{ref266}. Tree-based process supervision offers an online framework for step-wise credit assignment without intermediate labels by modeling reasoning as a rollout tree, where step utility is estimated via Monte Carlo estimation over descendant outcomes \cite{ref314}. Fine-grained atomic thought rewards decompose reasoning into functional units, providing supervision anchors that enable more interpretable, human-like reasoning patterns \cite{ref611}. Curriculum-guided RL frameworks further refine these approaches by evolving the agent from broad exploration to concise refinement through time-varying reward schedules \cite{ref267}. The synthesis of these efforts points toward a future where agentic retrieval policies are trained with rich, process-level supervision that aligns every intermediate action with the overarching goal of generating accurate, grounded responses.

A further frontier explores the synergy between RAG, long-context models, and agentic AI systems. As context windows grow, a critical question emerges regarding how RAG should evolve in tandem. Research comparing iterative RAG against idealized static retrieval indicates that the process of staged retrieval can often be more influential than simply providing all oracle evidence at once, demonstrating that the manner in which evidence is presented and reasoned over matters as much as its presence \cite{ref170}. This suggests that long-context capabilities do not eliminate the need for RAG but instead redefine its role: from providing extensive context to supplying carefully curated, reasoning-compatible evidence. The integration of RAG with agentic AI systems is likewise transforming both paradigms. Surveys of this convergence highlight how advanced reasoning optimizes each stage of RAG, while retrieved knowledge supplies missing premises for complex inference, creating a bidirectional synergy that pushes the boundaries of what is achievable on knowledge-intensive benchmarks \cite{ref21,ref53}. This unification of retrieval and reasoning is moving naturally toward multi-agent frameworks where specialized agents decompose complex queries, retrieve heterogeneous evidence, resolve conflicts, and synthesize answers, coordinating through hierarchical or consensus-driven mechanisms \cite{ref320,ref310}. The ability of these agentic systems to justify their retrieval decisions and produce traceable reasoning pathways marks a significant leap toward interpretable and trustworthy AI.

Finally, the concept of self-evolving knowledge representations stands out as a transformative research direction. Beyond static corpora and fixed graphs, the emerging vision is of knowledge bases that continuously update themselves through interaction with the agent and the environment. EvoGraph-R1 exemplifies this by treating knowledge hypergraphs as environments shaped through agent interactions, where actions such as expansion and editing produce feedback that guides the structure's evolution \cite{ref42}. Experience-driven orchestration further advances this principle by encapsulating retrieval strategies as reusable skills, where an agent analyzes the current scene, consults an experience memory, selects an appropriate strategy, and returns structured evidence \cite{ref612}. Hierarchical frameworks that jointly evolve multi-agent orchestration and role-specific prompts demonstrate the power of experience accumulation at both the system and component levels \cite{ref323}. The convergence of these directions promises RAG systems that not only retrieve and reason but also learn---adapting their knowledge structures, retrieval policies, and reasoning strategies in response to new evidence and evolving task demands, ultimately charting a path toward more robust, autonomous, and intelligent retrieval-augmented language models.

\subsection{Trustworthiness, Security, and Ethical Governance}

The deployment of Retrieval-Augmented Generation (RAG) systems in high-stakes environments necessitates a rigorous examination of their trustworthiness, security posture, and ethical governance. While the previous sections have charted the evolution of RAG from static pipelines toward adaptive, self-evolving architectures, this progress simultaneously broadens the attack surface and introduces novel failure modes that must be systematically addressed. As RAG systems become more autonomous---retrieving, reasoning, and acting through iterative loops---the risks of adversarial manipulation, data leakage, and biased amplification grow in tandem. This subsection therefore examines the security threats, privacy vulnerabilities, fairness concerns, and governance challenges that accompany these advances, framing them as first-class constraints that must shape the design of next-generation systems.

From a security perspective, RAG architectures are vulnerable to a distinct set of adversarial threats. The reliance on external knowledge bases makes the system susceptible to corpus poisoning attacks, where adversaries inject malicious or misleading documents into the retrieval source. This vulnerability is particularly acute in healthcare, where attackers with access to a relatively small number of samples can compromise AI systems regardless of dataset size \cite{ref613}. Such attacks can be launched with limited technical skill and often go undetected for extended periods, with detection sometimes taking six to twelve months or not occurring at all. Beyond poisoning, RAG systems face indirect prompt injection attacks, where malicious instructions embedded in retrieved documents can override the system's original directives, and knowledge extraction attacks, where adversaries exploit the system's outputs to infer sensitive information present in the private knowledge base. The concept of a ``confused-deputy'' vulnerability emerges when an LLM agent, operating with elevated privileges, is manipulated via its retrieved context to perform unauthorized actions. Standardized threat taxonomies, such as the nine-domain ontology encompassing Misuse, Poisoning, Privacy, Adversarial, and Biases threats, have been proposed to bridge the communication gap between technical security teams and legal compliance professionals, enabling quantitative risk assessment and a clearer translation of technical vulnerabilities into financial liability \cite{ref614}. This fragmented landscape, where reactive and siloed efforts fail to address the interdependence of technical integrity and societal trust, underscores the critical need for a comprehensive AI governance framework that integrates intrinsic security, derivative security, and social ethics \cite{ref615}.

Privacy presents another critical dimension of trustworthiness for RAG systems. The knowledge base itself can become a vector for data leakage, as retrieval mechanisms may surface sensitive or personally identifiable information embedded in source documents. Membership inference attacks, which aim to determine whether a specific data record was part of the training set, and more generally, the knowledge base, represent a persistent threat. Research has shown that privacy-preserving techniques like differential privacy can inadvertently exacerbate fairness disparities across demographic groups, as the privacy cost is not borne equally; groups experiencing greater overfitting gaps are more vulnerable to inference attacks \cite{ref616}. This necessitates auditing fairness-privacy trade-offs at the subpopulation level, as aggregate evaluations can obscure significant disparities in privacy risk \cite{ref617}. Defensive strategies span a wide spectrum, from input-stage cryptographic approaches like homomorphic encryption for retrieval operations and output-stage differential privacy frameworks, to architectural countermeasures. For instance, adaptive orchestration systems, where a strategist agent controls a separate generator to mitigate prompt-injection risks, represent a promising class of defenses \cite{ref618}. The challenge lies in the privacy-utility-security trilemma, where achieving high levels of all three simultaneously is often constrained by fundamental trade-offs. For example, generating synthetic data with formal privacy guarantees shows significant performance degradation under realistic privacy constraints, highlighting the gap between theoretical promises and practical utility \cite{ref619}.

Fairness and bias amplification within the RAG pipeline constitute a profound ethical concern, one that becomes increasingly pronounced as self-evolving systems adapt their retrieval strategies over time. Bias can be introduced at multiple stages: in the source corpora used for indexing, in the retrieval algorithms that may preferentially surface certain perspectives, and in the generator's own parametric biases. Retrieval of evidence can be skewed by demographic or topical biases present in the underlying document collection, leading to systematically unfair outputs for certain groups. This problem is exacerbated by the fact that fairness interventions often rely on sensitive data that privacy restrictions limit, creating a conflict between fairness and privacy objectives \cite{ref620}. Furthermore, the use of opaque, black-box models in high-stakes applications raises questions about their suitability and the feasibility of verifying their fairness and safety, with some research advocating for a shift toward more interpretable systems with verifiable safety guarantees \cite{ref613}. The ethical integrity of the data pipeline itself, from dataset design to inference and dissemination, is ethically constitutive, and well-intentioned projects can inadvertently slide into renewed extraction or exposure of vulnerable populations \cite{ref621}. To guide the responsible development of RAG systems, a comprehensive framework integrating domain definition, trustworthy AI design, auditability, accountability, and governance is necessary, emphasizing the inter-dependencies and iterative feedback loops across these dimensions \cite{ref622}.

Transparency, accountability, and attribution fidelity are fundamental pillars of trustworthy RAG deployment, and they take on renewed significance in the context of agentic and self-evolving systems described earlier. The ability to trace a generated claim back to specific retrieved evidence is crucial for building user trust and enabling error analysis. However, there is an inherent tension between the goal of transparency and the security requirement for confidentiality, as detailed disclosure of pipeline components may inadvertently reveal vulnerabilities or proprietary data. The open-box fallacy, where authorization to deploy is predicated on mechanistic interpretability, is often misguided; instead, a calibrated verification regime based on domain-scoped, independently checkable, and monitored authorization is more practical \cite{ref623}. This is supported by evidence showing a significant gap between understanding a model's internal representations and being able to correct its outputs. Furthermore, the practice of attribution itself can create a paradoxical vulnerability: by making the system more transparent about its knowledge sources, we may inadvertently create side channels for extraction attacks \cite{ref624}. For instance, a RAG system that confidently cites a specific phrasing from its knowledge base signals to an attacker that this exact string is present, potentially confirming membership of that document. Therefore, governance frameworks must carefully balance the societal need for accountability with the organizational requirement to protect proprietary data and maintain system security, navigating the complex ethical trade-offs between system openness, proprietary data protection, and regulatory compliance. This includes implementing robust access control and selective disclosure mechanisms to ensure that sensitive information is only accessible to authorized users, dynamically adjusting the degree of disclosure based on the user's trust level \cite{ref618}. Ultimately, a holistic perspective that considers the entire AI lifecycle, from data acquisition and model development to deployment and continuous monitoring, is essential for building RAG systems that are not only capable but also worthy of the trust placed in them by users and society at large \cite{ref625}.

\subsection{Roadmap Toward Reliable, Human-Centric RAG Systems}

Building directly upon the security, privacy, fairness, and governance challenges articulated in the preceding subsection, this concluding roadmap confronts a central tension: the very mechanisms that render RAG systems more autonomous and capable---iterative retrieval loops, adaptive orchestration, and self-evolving knowledge integration---simultaneously intensify the risks of adversarial manipulation, cascading errors, and inequitable outcomes. The path forward therefore cannot be purely technical optimization; it demands a deliberate synthesis of technical rigor, ethical foresight, and human-centered design. This subsection articulates an actionable agenda for next-generation retrieval-augmented language models, organizing the roadmap around three imperatives: establishing a unified evaluation paradigm, codifying design principles for human-centric systems, and articulating the open research questions that will define the field's trajectory.

The first pillar of this roadmap is a fundamental reconceptualization of how RAG systems are evaluated. The prevailing emphasis on accuracy-centric metrics, such as exact match or F1 score, has yielded an incomplete and sometimes misleading picture of system capability. Research has repeatedly demonstrated that high factual accuracy does not imply high-quality reasoning---leading models achieve strong precision while struggling with logical coherence and essentiality \cite{ref626}. Moreover, evaluations focusing solely on correctness fail to capture uncertainty dynamics; accuracy gains often coincide with reduced uncertainty, but this relationship breaks under retrieval noise, and no single RAG approach is universally reliable across domains \cite{ref409}. Critically, the security vulnerabilities detailed earlier---corpus poisoning, indirect prompt injection, and knowledge extraction---are rarely reflected in conventional evaluation suites, meaning a system can score highly on accuracy while remaining dangerously susceptible to adversarial exploitation. The field needs a convergence toward a holistic assessment framework that systematically evaluates factuality, robustness, fairness, transparency, and privacy as interdependent dimensions. The Trust-RAG Compass and its accompanying benchmark represent foundational steps in this direction, providing a structured approach to assessing these six dimensions for both proprietary and open-source models \cite{ref77}. However, such frameworks must be extended beyond isolated model assessment to encompass the full RAG pipeline, including retrieval quality, evidence grounding, and citation fidelity---and must integrate the threat models and defensive strategies catalogued in the security literature.

A unified evaluation paradigm must also address the critical gap between automated metrics and human judgment. The research community has made progress with LLM-as-a-judge approaches; the human-centered evaluation work demonstrates that structured questionnaires assessing outputs across utility dimensions can capture nuances that numeric scores miss, but also reveals that LLMs and humans exhibit systematic disagreements in their assessments \cite{ref78}. The proposed evaluation frameworks, such as VERA, which employs cross-encoder mechanisms and bootstrap statistics to establish confidence bounds on LLM-based metrics, offer a path toward more rigorous and reproducible assessment \cite{ref627}. Similarly, the framework developed in the RAGalyst work demonstrates the value of human-aligned agentic evaluation in specialized domains, where performance is highly context-dependent and no single configuration proves universally optimal \cite{ref588}. The roadmap calls for standardizing stage-decomposed diagnostics that can isolate failures at each pipeline component, building on existing diagnostic tools that provide root-cause analysis and actionable guidance. This shift toward interpretable, granular evaluation reflects a deeper insight: the process of staged retrieval is often more influential than the mere presence of ideal evidence, as demonstrated in scientific multi-hop question answering where iterative RAG outperforms even optimal static retrieval \cite{ref170}. Crucially, evaluation frameworks must also incorporate privacy and security metrics as first-class citizens---measuring not merely whether a system performs well, but whether it does so without leaking sensitive information or falling prey to adversarial manipulation.

The second pillar---human-centric design---demands that RAG systems be engineered not merely for optimal performance but for meaningful collaboration with human users. This imperative becomes particularly salient given the governance challenges outlined earlier: the open-box fallacy, the privacy-utility-security trilemma, and the tension between transparency and confidentiality all argue against purely autonomous, black-box deployments. The emerging consensus is that LLMs are moving beyond transactional interactions to act as companions, coaches, mediators, and curators that scaffold human capabilities, with design principles including transparency, personalization, guardrails, memory with privacy, and a balance of empathy and reliability \cite{ref628}. Applying these principles to RAG architecture requires a deliberate balancing of autonomy and oversight. Agentic RAG systems that autonomously coordinate multi-step reasoning and iterative retrieval offer significant advantages in flexibility and adaptability \cite{ref22}. Yet, relinquishing too much control to autonomous loops introduces systemic risks: hallucination propagation, memory poisoning, retrieval misalignment, and cascading tool-execution vulnerabilities can compound over iterative cycles \cite{ref25}. Human-centric design therefore argues for adaptive orchestration mechanisms that calibrate system autonomy to task complexity and user expertise, incorporating control mechanisms that allow users to navigate the accuracy-cost trade-off \cite{ref629}. This aligns with findings from user-controlled frameworks that enable dynamic adjustment between minimal-cost and high-accuracy retrieval based on specific requirements. Critically, such adaptive control must also serve as a security mechanism: by giving users visibility into retrieval decisions and the ability to intervene, systems can mitigate the ``confused-deputy'' vulnerability and other manipulation risks identified in threat analyses.

A human-centric roadmap also mandates attention to fairness and equitable access---a concern that the preceding subsection revealed to be in tension with privacy protections. Empirical studies reveal that fairness issues persist in both the retrieval and generation stages of RAG pipelines, with disparities emerging across demographic attributes \cite{ref501}. The roadmap therefore calls for fairness-aware design embedded at every stage, from corpus curation and retrieval ranking to generation constraints, coupled with systematic auditing frameworks that assess demographic disparities in deployment. Crucially, the fairness-privacy trade-off identified earlier---where differential privacy can exacerbate disparities, and where fairness interventions require sensitive data that privacy restricts---must be addressed through subpopulation-level auditing and the development of techniques that reconcile these competing objectives. This dimension connects directly to the trustworthiness imperative: a RAG system that performs well on average but fails systematically for specific populations cannot be considered reliable, and a system that protects privacy at the cost of fairness cannot be considered just.

The third pillar articulates the open research questions that will define the next generation of RAG systems. First, stable adaptive retrieval remains an unsolved challenge. Current adaptive mechanisms are brittle under query variations, and retrieval failure states are often treated as signals to retry rather than to diagnose \cite{ref630}. Research demonstrates that query-evidence misalignment is a typed, structured phenomenon where different failure modes require distinct corrective actions \cite{ref630}. A promising direction is reinforcement learning with process-level rewards to train agents that can learn stable retrieval policies, following the lead of frameworks like Search-P1 with path-centric reward shaping for agentic RAG training \cite{ref266}, and deep reasoning frameworks that explicitly train retrieval-augmented critics to identify evidence-aligned reasoning \cite{ref631}. Second, formal trajectory evaluation for agentic systems is needed. Current evaluation practices are largely static, failing to capture the dynamic, multi-round interactions that characterize agentic retrieval \cite{ref25}. The roadmap calls for rigorous, formal methods to assess not just final outputs but the quality of the reasoning trajectories that produce them---including their resilience to adversarial perturbations and their privacy properties over extended interactions. Third, the long-term convergence of RAG with reasoning, multimodal understanding, and lifelong learning requires sustained investigation. The synergy between RAG and advanced reasoning capabilities, including Chain-of-Thought and agentic search interleaving, has produced state-of-the-art results across knowledge-intensive benchmarks \cite{ref21}. As models evolve toward long-context processing, the boundary between retrieval and parametric knowledge becomes increasingly fluid. The emergence of streaming multimodal data creates additional pressure for lifelong learning architectures that maintain index freshness and cross-modal consistency without prohibitive re-indexing costs, as exemplified by lifelong multimodal agent designs \cite{ref632}. The integration of causal structures and knowledge graphs into retrieval also promises more principled approaches to multi-hop reasoning. Within these ambitions, the security and governance challenges identified earlier must be addressed preemptively: adaptive systems that continuously learn from retrieved content need robust defenses against poisoning over time, and lifelong learning architectures require privacy-preserving mechanisms for updating knowledge bases without exposing sensitive history.

In conclusion, the roadmap toward reliable, human-centric RAG systems is a call for disciplinary integration. It requires the information retrieval community to contribute robust retrieval methods; the NLP community to advance reasoning and faithful generation; the security community to codify threat models and defenses; and the human-computer interaction community to ensure that systems genuinely serve human needs. The open questions are substantial, but the trajectory is clear: moving from a narrow focus on accuracy toward a holistic engineering discipline that treats factuality, robustness, fairness, transparency, and privacy as first-class design constraints. The preceding analysis of security, privacy, and governance challenges makes evident that these are not peripheral concerns but foundational requirements that must shape system architecture from the outset. The synthesis offered in this survey underscores that RAG is not a static endpoint but a dynamic foundation for the next generation of grounded, accountable, and human-centric artificial intelligence systems---systems that are not only capable but also secure, equitable, and worthy of the trust placed in them by users and society at large.


\begin{thebibliography}{632}
\addcontentsline{toc}{section}{References}
\bibitem{ref1} Retrieval-Augmented Generation for Large Language Models A Survey. arXiv:2312.10997v5, 2023. \url{https://arxiv.org/abs/2312.10997v5}
\bibitem{ref2} RAGLAB A Modular and Research-Oriented Unified Framework for Retrieval-Augmented Generation. arXiv:2408.11381v2, 2024. \url{https://arxiv.org/abs/2408.11381v2}
\bibitem{ref3} \$\textbackslash\{\}text\{Memory\}\textasciicircum{}3\$ Language Modeling with Explicit Memory. arXiv:2407.01178v1, 2024. \url{https://arxiv.org/abs/2407.01178v1}
\bibitem{ref4} MemOS A Memory OS for AI System. arXiv:2507.03724v4, 2025. \url{https://arxiv.org/abs/2507.03724v4}
\bibitem{ref5} Parametric Retrieval Augmented Generation. arXiv:2501.15915v1, 2025. \url{https://arxiv.org/abs/2501.15915v1}
\bibitem{ref6} Contextual Compression in Retrieval-Augmented Generation for Large Language Models A Survey. arXiv:2409.13385v2, 2024. \url{https://arxiv.org/abs/2409.13385v2}
\bibitem{ref7} ReDeEP Detecting Hallucination in Retrieval-Augmented Generation via Mechanistic Interpretability. arXiv:2410.11414v2, 2024. \url{https://arxiv.org/abs/2410.11414v2}
\bibitem{ref8} Retrieval-Augmented Generation for Natural Language Processing A Survey. arXiv:2407.13193v4, 2024. \url{https://arxiv.org/abs/2407.13193v4}
\bibitem{ref9} Mitigating Hallucination in Large Language Models (LLMs) An Application-Oriented Survey on RAG, Reasoning, and Agentic Systems. arXiv:2510.24476v1, 2025. \url{https://arxiv.org/abs/2510.24476v1}
\bibitem{ref10} LUMINA Detecting Hallucinations in RAG System with Context-Knowledge Signals. arXiv:2509.21875v3, 2025. \url{https://arxiv.org/abs/2509.21875v3}
\bibitem{ref11} Bridging External and Parametric Knowledge Mitigating Hallucination of LLMs with Shared-Private Semantic Synergy in Dual-Stream Knowledge. arXiv:2506.06240v2, 2025. \url{https://arxiv.org/abs/2506.06240v2}
\bibitem{ref12} When Retrieval Succeeds and Fails Rethinking Retrieval-Augmented Generation for LLMs. arXiv:2510.09106v1, 2025. \url{https://arxiv.org/abs/2510.09106v1}
\bibitem{ref13} From RAG to Memory Non-Parametric Continual Learning for Large Language Models. arXiv:2502.14802v2, 2025. \url{https://arxiv.org/abs/2502.14802v2}
\bibitem{ref14} MemLLM Finetuning LLMs to Use An Explicit Read-Write Memory. arXiv:2404.11672v1, 2024. \url{https://arxiv.org/abs/2404.11672v1}
\bibitem{ref15} Knowledge Capsules Structured Nonparametric Memory Units for LLMs. arXiv:2604.20487v2, 2026. \url{https://arxiv.org/abs/2604.20487v2}
\bibitem{ref16} Attribution Techniques for Mitigating Hallucinated Information in RAG Systems A Survey. arXiv:2601.19927v1, 2026. \url{https://arxiv.org/abs/2601.19927v1}
\bibitem{ref17} Augmenting LLMs with Knowledge A survey on hallucination prevention. arXiv:2309.16459v1, 2023. \url{https://arxiv.org/abs/2309.16459v1}
\bibitem{ref18} RAG+ Enhancing Retrieval-Augmented Generation with Application-Aware Reasoning. arXiv:2506.11555v4, 2025. \url{https://arxiv.org/abs/2506.11555v4}
\bibitem{ref19} Dynamic and Parametric Retrieval-Augmented Generation. arXiv:2506.06704v1, 2025. \url{https://arxiv.org/abs/2506.06704v1}
\bibitem{ref20} RAG-DDR Optimizing Retrieval-Augmented Generation Using Differentiable Data Rewards. arXiv:2410.13509v2, 2024. \url{https://arxiv.org/abs/2410.13509v2}
\bibitem{ref21} Towards Agentic RAG with Deep Reasoning A Survey of RAG-Reasoning Systems in LLMs. arXiv:2507.09477v2, 2025. \url{https://arxiv.org/abs/2507.09477v2}
\bibitem{ref22} Agentic Retrieval-Augmented Generation A Survey on Agentic RAG. arXiv:2501.09136v4, 2025. \url{https://arxiv.org/abs/2501.09136v4}
\bibitem{ref23} Modular RAG Transforming RAG Systems into LEGO-like Reconfigurable Frameworks. arXiv:2407.21059v1, 2024. \url{https://arxiv.org/abs/2407.21059v1}
\bibitem{ref24} From Naive RAG to Deep Agentic Retrieval An Evolving Context Engineering Pipeline for Regulatory Compliance. arXiv:2607.24791v1, 2026. \url{https://arxiv.org/abs/2607.24791v1}
\bibitem{ref25} SoK Agentic Retrieval-Augmented Generation (RAG) Taxonomy, Architectures, Evaluation, and Research Directions. arXiv:2603.07379v1, 2026. \url{https://arxiv.org/abs/2603.07379v1}
\bibitem{ref26} Reasoning RAG via System 1 or System 2 A Survey on Reasoning Agentic Retrieval-Augmented Generation for Industry Challenges. arXiv:2506.10408v1, 2025. \url{https://arxiv.org/abs/2506.10408v1}
\bibitem{ref27} A-RAG Scaling Agentic Retrieval-Augmented Generation via Hierarchical Retrieval Interfaces. arXiv:2602.03442v1, 2026. \url{https://arxiv.org/abs/2602.03442v1}
\bibitem{ref28} Is GraphRAG Needed From Basic RAG to Graph- Agentic Solutions with Context Optimization. arXiv:2606.25656v1, 2026. \url{https://arxiv.org/abs/2606.25656v1}
\bibitem{ref29} Interact-RAG Reason and Interact with the Corpus, Beyond Black-Box Retrieval. arXiv:2510.27566v3, 2025. \url{https://arxiv.org/abs/2510.27566v3}
\bibitem{ref30} JADE Bridging the Strategic-Operational Gap in Dynamic Agentic RAG. arXiv:2601.21916v1, 2026. \url{https://arxiv.org/abs/2601.21916v1}
\bibitem{ref31} Agent-Orchestrated Adaptive RAG A Comparative Study on Structured and Multi-Hop Retrieval. arXiv:2606.05658v1, 2026. \url{https://arxiv.org/abs/2606.05658v1}
\bibitem{ref32} Creating a Taxonomy for Retrieval Augmented Generation Applications. arXiv:2408.02854v4, 2024. \url{https://arxiv.org/abs/2408.02854v4}
\bibitem{ref33} Graph-Guided Concept Selection for Efficient Retrieval-Augmented Generation. arXiv:2510.24120v1, 2025. \url{https://arxiv.org/abs/2510.24120v1}
\bibitem{ref34} Retrieval-Augmented Generation with Hierarchical Knowledge. arXiv:2503.10150v3, 2025. \url{https://arxiv.org/abs/2503.10150v3}
\bibitem{ref35} Retrieval-Augmented Generation with Graphs (GraphRAG). arXiv:2501.00309v2, 2024. \url{https://arxiv.org/abs/2501.00309v2}
\bibitem{ref36} Graph Retrieval-Augmented Generation A Survey. arXiv:2408.08921v2, 2024. \url{https://arxiv.org/abs/2408.08921v2}
\bibitem{ref37} A Survey of Graph Retrieval-Augmented Generation for Customized Large Language Models. arXiv:2501.13958v3, 2025. \url{https://arxiv.org/abs/2501.13958v3}
\bibitem{ref38} HG-RAG Hierarchy-Guided Retrieval-Augmented Generation for Structured Knowledge Graphs. arXiv:2607.14095v1, 2026. \url{https://arxiv.org/abs/2607.14095v1}
\bibitem{ref39} DA-RAG Dynamic Attributed Community Search for Retrieval-Augmented Generation. arXiv:2602.08545v1, 2026. \url{https://arxiv.org/abs/2602.08545v1}
\bibitem{ref40} LightRAG Simple and Fast Retrieval-Augmented Generation. arXiv:2410.05779v3, 2024. \url{https://arxiv.org/abs/2410.05779v3}
\bibitem{ref41} Simple Is Effective The Roles of Graphs and Large Language Models in Knowledge-Graph-Based Retrieval-Augmented Generation. arXiv:2410.20724v4, 2024. \url{https://arxiv.org/abs/2410.20724v4}
\bibitem{ref42} EvoGraph-R1 Self-Evolving Multimodal Knowledge Hypergraphs for Agentic Retrieval. arXiv:2607.12764v1, 2026. \url{https://arxiv.org/abs/2607.12764v1}
\bibitem{ref43} Graph-based Approaches and Functionalities in Retrieval-Augmented Generation A Comprehensive Survey. arXiv:2504.10499v2, 2025. \url{https://arxiv.org/abs/2504.10499v2}
\bibitem{ref44} RAG-Anything All-in-One RAG Framework. arXiv:2510.12323v1, 2025. \url{https://arxiv.org/abs/2510.12323v1}
\bibitem{ref45} Ask in Any Modality A Comprehensive Survey on Multimodal Retrieval-Augmented Generation. arXiv:2502.08826v3, 2025. \url{https://arxiv.org/abs/2502.08826v3}
\bibitem{ref46} MG\$\textasciicircum{}2\$-RAG Multi-Granularity Graph for Multimodal Retrieval-Augmented Generation. arXiv:2604.04969v2, 2026. \url{https://arxiv.org/abs/2604.04969v2}
\bibitem{ref47} mKG-RAG Leveraging Multimodal Knowledge Graphs in Retrieval-Augmented Generation for Knowledge-intensive VQA. arXiv:2508.05318v2, 2025. \url{https://arxiv.org/abs/2508.05318v2}
\bibitem{ref48} Scaling Beyond Context A Survey of Multimodal Retrieval-Augmented Generation for Document Understanding. arXiv:2510.15253v3, 2025. \url{https://arxiv.org/abs/2510.15253v3}
\bibitem{ref49} MKG-RAG-Bench Benchmarking Retrieval in Multimodal Knowledge Graph-Augmented Generation. arXiv:2606.26458v1, 2026. \url{https://arxiv.org/abs/2606.26458v1}
\bibitem{ref50} Visual-RAG Benchmarking Text-to-Image Retrieval Augmented Generation for Visual Knowledge Intensive Queries. arXiv:2502.16636v2, 2025. \url{https://arxiv.org/abs/2502.16636v2}
\bibitem{ref51} MegaRAG Multimodal Knowledge Graph-Based Retrieval Augmented Generation. arXiv:2512.20626v2, 2025. \url{https://arxiv.org/abs/2512.20626v2}
\bibitem{ref52} A Study on the Implementation Method of an Agent-Based Advanced RAG System Using Graph. arXiv:2407.19994v3, 2024. \url{https://arxiv.org/abs/2407.19994v3}
\bibitem{ref53} Synergizing RAG and Reasoning A Systematic Review. arXiv:2504.15909v2, 2025. \url{https://arxiv.org/abs/2504.15909v2}
\bibitem{ref54} MIXRAG Mixture-of-Experts Retrieval-Augmented Generation for Textual Graph Understanding and Question Answering. arXiv:2509.21391v2, 2025. \url{https://arxiv.org/abs/2509.21391v2}
\bibitem{ref55} BioMol-MQA A Multi-Modal Question Answering Dataset For LLM Reasoning Over Bio-Molecular Interactions. arXiv:2506.05766v1, 2025. \url{https://arxiv.org/abs/2506.05766v1}
\bibitem{ref56} DSRAG A Domain-Specific Retrieval Framework Based on Document-derived Multimodal Knowledge Graph. arXiv:2509.10467v1, 2025. \url{https://arxiv.org/abs/2509.10467v1}
\bibitem{ref57} Comparative Analysis of Neural Retriever-Reranker Pipelines for Retrieval-Augmented Generation over Knowledge Graphs in E-commerce Applications. arXiv:2602.22219v1, 2025. \url{https://arxiv.org/abs/2602.22219v1}
\bibitem{ref58} A Hybrid RAG System with Comprehensive Enhancement on Complex Reasoning. arXiv:2408.05141v3, 2024. \url{https://arxiv.org/abs/2408.05141v3}
\bibitem{ref59} Pursuing Best Industrial Practices for Retrieval-Augmented Generation in the Medical Domain. arXiv:2602.03368v2, 2026. \url{https://arxiv.org/abs/2602.03368v2}
\bibitem{ref60} ModernBERT + ColBERT Enhancing biomedical RAG through an advanced re-ranking retriever. arXiv:2510.04757v1, 2025. \url{https://arxiv.org/abs/2510.04757v1}
\bibitem{ref61} Towards Reliable Retrieval in RAG Systems for Large Legal Datasets. arXiv:2510.06999v1, 2025. \url{https://arxiv.org/abs/2510.06999v1}
\bibitem{ref62} Understanding the Fundamental Design Decisions of Retrieval-Augmented Generation Systems. arXiv:2411.19463v3, 2024. \url{https://arxiv.org/abs/2411.19463v3}
\bibitem{ref63} RAGPulse An Open-Source RAG Workload Trace to Optimize RAG Serving Systems. arXiv:2511.12979v1, 2025. \url{https://arxiv.org/abs/2511.12979v1}
\bibitem{ref64} Pandora's Box or Aladdin's Lamp A Comprehensive Analysis Revealing the Role of RAG Noise in Large Language Models. arXiv:2408.13533v4, 2024. \url{https://arxiv.org/abs/2408.13533v4}
\bibitem{ref65} Out of Style RAG's Fragility to Linguistic Variation. arXiv:2504.08231v2, 2025. \url{https://arxiv.org/abs/2504.08231v2}
\bibitem{ref66} How You Ask Matters! Adaptive RAG Robustness to Query Variations. arXiv:2604.10745v1, 2026. \url{https://arxiv.org/abs/2604.10745v1}
\bibitem{ref67} Securing RAG A Risk Assessment and Mitigation Framework. arXiv:2505.08728v2, 2025. \url{https://arxiv.org/abs/2505.08728v2}
\bibitem{ref68} Security and Privacy in Retrieval-Augmented Generation Architectures, Threats, Defenses, and Future Directions for Building Trustworthy Systems. arXiv:2606.25533v1, 2026. \url{https://arxiv.org/abs/2606.25533v1}
\bibitem{ref69} Hoist with His Own Petard Inducing Guardrails to Facilitate Denial-of-Service Attacks on Retrieval-Augmented Generation of LLMs. arXiv:2504.21680v1, 2025. \url{https://arxiv.org/abs/2504.21680v1}
\bibitem{ref70} Rag and Roll An End-to-End Evaluation of Indirect Prompt Manipulations in LLM-based Application Frameworks. arXiv:2408.05025v2, 2024. \url{https://arxiv.org/abs/2408.05025v2}
\bibitem{ref71} Generating Is Believing Membership Inference Attacks against Retrieval-Augmented Generation. arXiv:2406.19234v2, 2024. \url{https://arxiv.org/abs/2406.19234v2}
\bibitem{ref72} No Free Lunch Retrieval-Augmented Generation Undermines Fairness in LLMs, Even for Vigilant Users. arXiv:2410.07589v1, 2024. \url{https://arxiv.org/abs/2410.07589v1}
\bibitem{ref73} Retrieval-Augmented Generation with Estimation of Source Reliability. arXiv:2410.22954v5, 2024. \url{https://arxiv.org/abs/2410.22954v5}
\bibitem{ref74} ARAGOG Advanced RAG Output Grading. arXiv:2404.01037v1, 2024. \url{https://arxiv.org/abs/2404.01037v1}
\bibitem{ref75} Towards Secure Retrieval-Augmented Generation A Comprehensive Review of Threats, Defenses and Benchmarks. arXiv:2603.21654v1, 2026. \url{https://arxiv.org/abs/2603.21654v1}
\bibitem{ref76} MIRAGE A Metric-Intensive Benchmark for Retrieval-Augmented Generation Evaluation. arXiv:2504.17137v1, 2025. \url{https://arxiv.org/abs/2504.17137v1}
\bibitem{ref77} Trustworthiness in Retrieval-Augmented Generation Systems A Survey. arXiv:2409.10102v2, 2024. \url{https://arxiv.org/abs/2409.10102v2}
\bibitem{ref78} Human-Centered Evaluation of RAG outputs a framework and questionnaire for human-AI collaboration. arXiv:2509.26205v1, 2025. \url{https://arxiv.org/abs/2509.26205v1}
\bibitem{ref79} RAG Security and Privacy Formalizing the Threat Model and Attack Surface. arXiv:2509.20324v2, 2025. \url{https://arxiv.org/abs/2509.20324v2}
\bibitem{ref80} A Systematic Investigation of Document Chunking Strategies and Embedding Sensitivity. arXiv:2603.06976v1, 2026. \url{https://arxiv.org/abs/2603.06976v1}
\bibitem{ref81} Beyond Chunk-Then-Embed A Comprehensive Taxonomy and Evaluation of Document Chunking Strategies for Information Retrieval. arXiv:2602.16974v1, 2026. \url{https://arxiv.org/abs/2602.16974v1}
\bibitem{ref82} Rethinking Chunk Size For Long-Document Retrieval A Multi-Dataset Analysis. arXiv:2505.21700v2, 2025. \url{https://arxiv.org/abs/2505.21700v2}
\bibitem{ref83} A Systematic Analysis of Chunking Strategies for Reliable Question Answering. arXiv:2601.14123v1, 2026. \url{https://arxiv.org/abs/2601.14123v1}
\bibitem{ref84} Query-Adaptive Semantic Chunking for Retrieval-Augmented Generation A Dynamic Strategy with Contextual Window Expansion. arXiv:2605.22834v2, 2026. \url{https://arxiv.org/abs/2605.22834v2}
\bibitem{ref85} Chunking German Legal Code. arXiv:2605.19806v2, 2026. \url{https://arxiv.org/abs/2605.19806v2}
\bibitem{ref86} Structure-Aware Chunking for Tabular Data in Retrieval-Augmented Generation. arXiv:2605.00318v1, 2026. \url{https://arxiv.org/abs/2605.00318v1}
\bibitem{ref87} Financial Report Chunking for Effective Retrieval Augmented Generation. arXiv:2402.05131v3, 2024. \url{https://arxiv.org/abs/2402.05131v3}
\bibitem{ref88} SproutRAG Attention-Guided Tree Search with Progressive Embeddings for Long-Document RAG. arXiv:2606.18381v1, 2026. \url{https://arxiv.org/abs/2606.18381v1}
\bibitem{ref89} Self-Conditioned Positional HNSW for Overlap-Aware Retrieval in Chunked-Document RAG Systems Method and Industrial Evidence-Quality Audit. arXiv:2606.01542v1, 2026. \url{https://arxiv.org/abs/2606.01542v1}
\bibitem{ref90} A Systematic Framework for Enterprise Knowledge Retrieval Leveraging LLM-Generated Metadata to Enhance RAG Systems. arXiv:2512.05411v2, 2025. \url{https://arxiv.org/abs/2512.05411v2}
\bibitem{ref91} Utilizing Metadata for Better Retrieval-Augmented Generation. arXiv:2601.11863v1, 2026. \url{https://arxiv.org/abs/2601.11863v1}
\bibitem{ref92} Cross-Document Topic-Aligned Chunking for Retrieval-Augmented Generation. arXiv:2601.05265v1, 2025. \url{https://arxiv.org/abs/2601.05265v1}
\bibitem{ref93} Multi-Vector Index Compression in Any Modality. arXiv:2602.21202v1, 2026. \url{https://arxiv.org/abs/2602.21202v1}
\bibitem{ref94} Lost in a Single Vector Improving Long-Document Retrieval with Chunk Evidence Aggregation. arXiv:2606.18781v1, 2026. \url{https://arxiv.org/abs/2606.18781v1}
\bibitem{ref95} SlimRAG Retrieval without Graphs via Entity-Aware Context Selection. arXiv:2506.17288v1, 2025. \url{https://arxiv.org/abs/2506.17288v1}
\bibitem{ref96} HiKEY Hierarchical Multimodal Retrieval for Open-Domain Document Question Answering. arXiv:2605.29606v1, 2026. \url{https://arxiv.org/abs/2605.29606v1}
\bibitem{ref97} Reducing Redundancy in Retrieval-Augmented Generation through Chunk Filtering. arXiv:2604.24334v1, 2026. \url{https://arxiv.org/abs/2604.24334v1}
\bibitem{ref98} M-RAG Semantic Key-Value Indexing for Retrieval-Augmented Generation. arXiv:2603.26667v2, 2026. \url{https://arxiv.org/abs/2603.26667v2}
\bibitem{ref99} From BM25 to Corrective RAG Benchmarking Retrieval Strategies for Text-and-Table Documents. arXiv:2604.01733v1, 2026. \url{https://arxiv.org/abs/2604.01733v1}
\bibitem{ref100} BM25 Wins at Scale A Scaling Study of Retrieval-Augmented Generation Paradigms. arXiv:2607.26497v3, 2026. \url{https://arxiv.org/abs/2607.26497v3}
\bibitem{ref101} An Empirical Study of Position Bias in Modern Information Retrieval. arXiv:2505.13950v5, 2025. \url{https://arxiv.org/abs/2505.13950v5}
\bibitem{ref102} CeQe Grounding Lexical Retrieval in Semantic Evidence. arXiv:2608.00452v1, 2026. \url{https://arxiv.org/abs/2608.00452v1}
\bibitem{ref103} From Retrieval to Generation Comparing Different Approaches. arXiv:2502.20245v1, 2025. \url{https://arxiv.org/abs/2502.20245v1}
\bibitem{ref104} ColBERT Efficient and Effective Passage Search via Contextualized Late Interaction over BERT. arXiv:2004.12832v2, 2020. \url{https://arxiv.org/abs/2004.12832v2}
\bibitem{ref105} On Single and Multiple Representations in Dense Passage Retrieval. arXiv:2108.06279v2, 2021. \url{https://arxiv.org/abs/2108.06279v2}
\bibitem{ref106} ColBERTv2 Effective and Efficient Retrieval via Lightweight Late Interaction. arXiv:2112.01488v3, 2021. \url{https://arxiv.org/abs/2112.01488v3}
\bibitem{ref107} SLIM Sparsified Late Interaction for Multi-Vector Retrieval with Inverted Indexes. arXiv:2302.06587v2, 2023. \url{https://arxiv.org/abs/2302.06587v2}
\bibitem{ref108} SPLATE Sparse Late Interaction Retrieval. arXiv:2404.13950v1, 2024. \url{https://arxiv.org/abs/2404.13950v1}
\bibitem{ref109} No More K-means Single-Stage Sparse Coding for Efficient Multi-Vector Retrieval. arXiv:2605.30120v3, 2026. \url{https://arxiv.org/abs/2605.30120v3}
\bibitem{ref110} Query Embedding Pruning for Dense Retrieval. arXiv:2108.10341v1, 2021. \url{https://arxiv.org/abs/2108.10341v1}
\bibitem{ref111} FastLane Efficient Routed Systems for Late-Interaction Retrieval. arXiv:2601.06389v2, 2026. \url{https://arxiv.org/abs/2601.06389v2}
\bibitem{ref112} Mamba Retriever Utilizing Mamba for Effective and Efficient Dense Retrieval. arXiv:2408.08066v2, 2024. \url{https://arxiv.org/abs/2408.08066v2}
\bibitem{ref113} FlexStructRAG Flexible Structure-Aware Multi-Granular Relational Retrieval for RAG. arXiv:2604.16312v1, 2026. \url{https://arxiv.org/abs/2604.16312v1}
\bibitem{ref114} Efficient and Interpretable Information Retrieval for Product Question Answering with Heterogeneous Data. arXiv:2405.13173v1, 2024. \url{https://arxiv.org/abs/2405.13173v1}
\bibitem{ref115} Evidence-Unit Fairness and the Limits of Query-Adaptive Sparse-Dense Fusion in Financial Document Retrieval. arXiv:2608.00183v1, 2026. \url{https://arxiv.org/abs/2608.00183v1}
\bibitem{ref116} DAT Dynamic Alpha Tuning for Hybrid Retrieval in Retrieval-Augmented Generation. arXiv:2503.23013v1, 2025. \url{https://arxiv.org/abs/2503.23013v1}
\bibitem{ref117} Balancing the Blend An Experimental Analysis of Trade-offs in Hybrid Search. arXiv:2508.01405v2, 2025. \url{https://arxiv.org/abs/2508.01405v2}
\bibitem{ref118} Dense X Retrieval What Retrieval Granularity Should We Use. arXiv:2312.06648v2, 2023. \url{https://arxiv.org/abs/2312.06648v2}
\bibitem{ref119} Uncovering the Limitations of Query Performance Prediction Failures, Insights, and Implications for Selective Query Processing. arXiv:2504.01101v1, 2025. \url{https://arxiv.org/abs/2504.01101v1}
\bibitem{ref120} On Coherence-based Predictors for Dense Query Performance Prediction. arXiv:2310.11405v1, 2023. \url{https://arxiv.org/abs/2310.11405v1}
\bibitem{ref121} Predicting Efficiency Effectiveness Trade-offs for Dense vs. Sparse Retrieval Strategy Selection. arXiv:2109.10739v1, 2021. \url{https://arxiv.org/abs/2109.10739v1}
\bibitem{ref122} Faster Learned Sparse Retrieval with Guided Traversal. arXiv:2204.11314v1, 2022. \url{https://arxiv.org/abs/2204.11314v1}
\bibitem{ref123} Query Expansion in the Age of Pre-trained and Large Language Models A Comprehensive Survey. arXiv:2509.07794v3, 2025. \url{https://arxiv.org/abs/2509.07794v3}
\bibitem{ref124} Generative Query Expansion with Multilingual LLMs for Cross-Lingual Information Retrieval. arXiv:2511.19325v1, 2025. \url{https://arxiv.org/abs/2511.19325v1}
\bibitem{ref125} Can Query Expansion Improve Generalization of Strong Cross-Encoder Rankers. arXiv:2311.09175v2, 2023. \url{https://arxiv.org/abs/2311.09175v2}
\bibitem{ref126} LLM-Assisted Pseudo-Relevance Feedback. arXiv:2601.11238v1, 2026. \url{https://arxiv.org/abs/2601.11238v1}
\bibitem{ref127} RAG without Forgetting Continual Query-Infused Key Memory. arXiv:2602.05152v1, 2026. \url{https://arxiv.org/abs/2602.05152v1}
\bibitem{ref128} Question Decomposition for Retrieval-Augmented Generation. arXiv:2507.00355v1, 2025. \url{https://arxiv.org/abs/2507.00355v1}
\bibitem{ref129} When Should Queries Be Decomposed A Stage-Aware Study of Query Decomposition for Multi-Condition Retrieval. arXiv:2606.08577v1, 2026. \url{https://arxiv.org/abs/2606.08577v1}
\bibitem{ref130} Query, Decompose, Compress Structured Query Expansion for Efficient Multi-Hop Retrieval. arXiv:2603.21024v1, 2026. \url{https://arxiv.org/abs/2603.21024v1}
\bibitem{ref131} When should I search more Adaptive Complex Query Optimization with Reinforcement Learning. arXiv:2601.21208v1, 2026. \url{https://arxiv.org/abs/2601.21208v1}
\bibitem{ref132} SAGE Strategy-Adaptive Generation Engine for Query Rewriting. arXiv:2506.19783v2, 2025. \url{https://arxiv.org/abs/2506.19783v2}
\bibitem{ref133} Hypothesis-Conditioned Query Rewriting for Decision-Useful Retrieval. arXiv:2603.19008v1, 2026. \url{https://arxiv.org/abs/2603.19008v1}
\bibitem{ref134} Beyond the Reranker Do RAG Retrieval Enhancements Help Once a Strong Reranker Is Present. arXiv:2606.28367v1, 2026. \url{https://arxiv.org/abs/2606.28367v1}
\bibitem{ref135} Retrieving a Set, Not Independent Passages Set-Level Compatibility Learning for Efficient Set Exploration. arXiv:2607.05712v1, 2026. \url{https://arxiv.org/abs/2607.05712v1}
\bibitem{ref136} Shifting from Ranking to Set Selection for Retrieval Augmented Generation. arXiv:2507.06838v2, 2025. \url{https://arxiv.org/abs/2507.06838v2}
\bibitem{ref137} QUBO-Optimized Evidence Selection for Retrieval-Augmented Question Answering with Unconventional Solvers. arXiv:2607.12334v1, 2026. \url{https://arxiv.org/abs/2607.12334v1}
\bibitem{ref138} OptiSet Unified Optimizing Set Selection and Ranking for Retrieval-Augmented Generation. arXiv:2601.05027v1, 2026. \url{https://arxiv.org/abs/2601.05027v1}
\bibitem{ref139} DualView Adaptive Local-Global Fusion for Multi-Hop Document Reranking. arXiv:2605.18767v1, 2026. \url{https://arxiv.org/abs/2605.18767v1}
\bibitem{ref140} Rerank Before You Reason Analyzing Reranking Tradeoffs through Effective Token Cost in Deep Search Agents. arXiv:2601.14224v2, 2026. \url{https://arxiv.org/abs/2601.14224v2}
\bibitem{ref141} Rethinking Retrieval-Augmentation as Synthesis A Query-Aware Context Merging Approach. arXiv:2603.20286v1, 2026. \url{https://arxiv.org/abs/2603.20286v1}
\bibitem{ref142} BRIEF Bridging Retrieval and Inference for Multi-hop Reasoning via Compression. arXiv:2410.15277v2, 2024. \url{https://arxiv.org/abs/2410.15277v2}
\bibitem{ref143} Self-Correcting RAG Enhancing Faithfulness via MMKP Context Selection and NLI-Guided MCTS. arXiv:2604.10734v1, 2026. \url{https://arxiv.org/abs/2604.10734v1}
\bibitem{ref144} RAISE RAG Design as an Architecture Search Problem. arXiv:2605.30029v1, 2026. \url{https://arxiv.org/abs/2605.30029v1}
\bibitem{ref145} RAGSmith A Framework for Finding the Optimal Composition of Retrieval-Augmented Generation Methods Across Datasets. arXiv:2511.01386v1, 2025. \url{https://arxiv.org/abs/2511.01386v1}
\bibitem{ref146} AutoRAGTuner A Declarative Framework for Automatic Optimization of RAG Pipelines. arXiv:2605.02967v1, 2026. \url{https://arxiv.org/abs/2605.02967v1}
\bibitem{ref147} Benchmarking Retrieval Strategies for Biomedical Retrieval-Augmented Generation A Controlled Empirical Study. arXiv:2605.02520v1, 2026. \url{https://arxiv.org/abs/2605.02520v1}
\bibitem{ref148} Optimizing Retrieval-Augmented Generation Analysis of Hyperparameter Impact on Performance and Efficiency. arXiv:2505.08445v1, 2025. \url{https://arxiv.org/abs/2505.08445v1}
\bibitem{ref149} An Analysis of Hyper-Parameter Optimization Methods for Retrieval Augmented Generation. arXiv:2505.03452v3, 2025. \url{https://arxiv.org/abs/2505.03452v3}
\bibitem{ref150} Faster, Cheaper, Better Multi-Objective Hyperparameter Optimization for LLM and RAG Systems. arXiv:2502.18635v2, 2025. \url{https://arxiv.org/abs/2502.18635v2}
\bibitem{ref151} CAMI Cost-Aware Agent-Guided Multi-Indexing for Semantic Retrieval. arXiv:2606.28365v1, 2026. \url{https://arxiv.org/abs/2606.28365v1}
\bibitem{ref152} Metadata-Driven Retrieval-Augmented Generation for Financial Question Answering. arXiv:2510.24402v1, 2025. \url{https://arxiv.org/abs/2510.24402v1}
\bibitem{ref153} Meta Knowledge for Retrieval Augmented Large Language Models. arXiv:2408.09017v1, 2024. \url{https://arxiv.org/abs/2408.09017v1}
\bibitem{ref154} AdaComp Extractive Context Compression with Adaptive Predictor for Retrieval-Augmented Large Language Models. arXiv:2409.01579v2, 2024. \url{https://arxiv.org/abs/2409.01579v2}
\bibitem{ref155} Dynamic Context Selection for Retrieval-Augmented Generation Mitigating Distractors and Positional Bias. arXiv:2512.14313v1, 2025. \url{https://arxiv.org/abs/2512.14313v1}
\bibitem{ref156} QUITO Accelerating Long-Context Reasoning through Query-Guided Context Compression. arXiv:2408.00274v1, 2024. \url{https://arxiv.org/abs/2408.00274v1}
\bibitem{ref157} AttentionRAG Attention-Guided Context Pruning in Retrieval-Augmented Generation. arXiv:2503.10720v2, 2025. \url{https://arxiv.org/abs/2503.10720v2}
\bibitem{ref158} AttnComp Attention-Guided Adaptive Context Compression for Retrieval-Augmented Generation. arXiv:2509.17486v1, 2025. \url{https://arxiv.org/abs/2509.17486v1}
\bibitem{ref159} Sentinel Decoding Context Utilization via Attention Probing for Efficient LLM Context Compression. arXiv:2505.23277v3, 2025. \url{https://arxiv.org/abs/2505.23277v3}
\bibitem{ref160} Enhancing and Accelerating Large Language Models via Instruction-Aware Contextual Compression. arXiv:2408.15491v1, 2024. \url{https://arxiv.org/abs/2408.15491v1}
\bibitem{ref161} Concept than Document Context Compression via AMR-based Conceptual Entropy. arXiv:2511.18832v1, 2025. \url{https://arxiv.org/abs/2511.18832v1}
\bibitem{ref162} Compressing Long Context for Enhancing RAG with AMR-based Concept Distillation. arXiv:2405.03085v1, 2024. \url{https://arxiv.org/abs/2405.03085v1}
\bibitem{ref163} Efficient Dynamic Clustering-Based Document Compression for Retrieval-Augmented-Generation. arXiv:2504.03165v3, 2025. \url{https://arxiv.org/abs/2504.03165v3}
\bibitem{ref164} xRAG Extreme Context Compression for Retrieval-augmented Generation with One Token. arXiv:2405.13792v2, 2024. \url{https://arxiv.org/abs/2405.13792v2}
\bibitem{ref165} Context Embeddings for Efficient Answer Generation in RAG. arXiv:2407.09252v3, 2024. \url{https://arxiv.org/abs/2407.09252v3}
\bibitem{ref166} Rethinking Soft Compression in Retrieval-Augmented Generation A Query-Conditioned Selector Perspective. arXiv:2602.15856v2, 2026. \url{https://arxiv.org/abs/2602.15856v2}
\bibitem{ref167} Beyond Hard and Soft Hybrid Context Compression for Balancing Local and Global Information Retention. arXiv:2505.15774v1, 2025. \url{https://arxiv.org/abs/2505.15774v1}
\bibitem{ref168} Grounding Long-Context Reasoning with Contextual Normalization for Retrieval-Augmented Generation. arXiv:2510.13191v2, 2025. \url{https://arxiv.org/abs/2510.13191v2}
\bibitem{ref169} Principled Context Engineering for RAG Statistical Guarantees via Conformal Prediction. arXiv:2511.17908v2, 2025. \url{https://arxiv.org/abs/2511.17908v2}
\bibitem{ref170} When Iterative RAG Beats Ideal Evidence A Diagnostic Study in Scientific Multi-hop Question Answering. arXiv:2601.19827v5, 2026. \url{https://arxiv.org/abs/2601.19827v5}
\bibitem{ref171} Does RAG Know When Retrieval Is Wrong Diagnosing Context Compliance under Knowledge Conflict. arXiv:2605.14473v4, 2026. \url{https://arxiv.org/abs/2605.14473v4}
\bibitem{ref172} Structure-Augmented Reasoning Generation. arXiv:2506.08364v4, 2025. \url{https://arxiv.org/abs/2506.08364v4}
\bibitem{ref173} REAP Enhancing RAG with Recursive Evaluation and Adaptive Planning for Multi-Hop Question Answering. arXiv:2511.09966v1, 2025. \url{https://arxiv.org/abs/2511.09966v1}
\bibitem{ref174} Eliciting Critical Reasoning in Retrieval-Augmented Language Models via Contrastive Explanations. arXiv:2410.22874v1, 2024. \url{https://arxiv.org/abs/2410.22874v1}
\bibitem{ref175} ArgRAG Explainable Retrieval Augmented Generation using Quantitative Bipolar Argumentation. arXiv:2508.20131v1, 2025. \url{https://arxiv.org/abs/2508.20131v1}
\bibitem{ref176} CARE-RAG - Clinical Assessment and Reasoning in RAG. arXiv:2511.15994v1, 2025. \url{https://arxiv.org/abs/2511.15994v1}
\bibitem{ref177} Resolving Conflicting Evidence in Automated Fact-Checking A Study on Retrieval-Augmented LLMs. arXiv:2505.17762v1, 2025. \url{https://arxiv.org/abs/2505.17762v1}
\bibitem{ref178} Facet-Level Tracing of Evidence Uncertainty and Hallucination in RAG. arXiv:2604.09174v2, 2026. \url{https://arxiv.org/abs/2604.09174v2}
\bibitem{ref179} Guaranteeing Knowledge Integration with Joint Decoding for Retrieval-Augmented Generation. arXiv:2604.08046v2, 2026. \url{https://arxiv.org/abs/2604.08046v2}
\bibitem{ref180} From Facts to Conclusions Integrating Deductive Reasoning in Retrieval-Augmented LLMs. arXiv:2512.16795v1, 2025. \url{https://arxiv.org/abs/2512.16795v1}
\bibitem{ref181} Resisting Contextual Interference in RAG via Parametric-Knowledge Reinforcement. arXiv:2506.05154v4, 2025. \url{https://arxiv.org/abs/2506.05154v4}
\bibitem{ref182} Grounded Decoding Retrieval-Anchored Probability Fusion for Faithful RAG. arXiv:2606.00432v1, 2026. \url{https://arxiv.org/abs/2606.00432v1}
\bibitem{ref183} ConflictRAG Detecting and Resolving Knowledge Conflicts in Retrieval Augmented Generation. arXiv:2605.17301v2, 2026. \url{https://arxiv.org/abs/2605.17301v2}
\bibitem{ref184} ArbGraph Conflict-Aware Evidence Arbitration for Reliable Long-Form Retrieval-Augmented Generation. arXiv:2604.18362v1, 2026. \url{https://arxiv.org/abs/2604.18362v1}
\bibitem{ref185} Beyond Black-Box Interventions Latent Probing for Faithful Retrieval-Augmented Generation. arXiv:2510.12460v3, 2025. \url{https://arxiv.org/abs/2510.12460v3}
\bibitem{ref186} SHIFT Gate-Modulated Activation Steering for Knowledge Conflict Mitigation in Retrieval-Augmented Generation. arXiv:2606.27786v1, 2026. \url{https://arxiv.org/abs/2606.27786v1}
\bibitem{ref187} CoRect Context-Aware Logit Contrast for Hidden State Rectification to Resolve Knowledge Conflicts. arXiv:2602.08221v1, 2026. \url{https://arxiv.org/abs/2602.08221v1}
\bibitem{ref188} ParamMute Suppressing Knowledge-Critical FFNs for Faithful Retrieval-Augmented Generation. arXiv:2502.15543v4, 2025. \url{https://arxiv.org/abs/2502.15543v4}
\bibitem{ref189} Studying Large Language Model Behaviors Under Realistic Knowledge Conflicts. arXiv:2404.16032v1, 2024. \url{https://arxiv.org/abs/2404.16032v1}
\bibitem{ref190} Seeing through the Conflict Transparent Knowledge Conflict Handling in Retrieval-Augmented Generation. arXiv:2601.06842v1, 2026. \url{https://arxiv.org/abs/2601.06842v1}
\bibitem{ref191} Trust or Abstain A Self-Aware RAG Approach. arXiv:2605.18792v1, 2026. \url{https://arxiv.org/abs/2605.18792v1}
\bibitem{ref192} Attribute First, then Generate Locally-attributable Grounded Text Generation. arXiv:2403.17104v2, 2024. \url{https://arxiv.org/abs/2403.17104v2}
\bibitem{ref193} VeriCite Towards Reliable Citations in Retrieval-Augmented Generation via Rigorous Verification. arXiv:2510.11394v1, 2025. \url{https://arxiv.org/abs/2510.11394v1}
\bibitem{ref194} Explicit Evidence Grounding via Structured Inline Citation Generation. arXiv:2606.07130v1, 2026. \url{https://arxiv.org/abs/2606.07130v1}
\bibitem{ref195} Are Finer Citations Always Better Rethinking Granularity for Attributed Generation. arXiv:2604.01432v2, 2026. \url{https://arxiv.org/abs/2604.01432v2}
\bibitem{ref196} Correctness is not Faithfulness in RAG Attributions. arXiv:2412.18004v1, 2024. \url{https://arxiv.org/abs/2412.18004v1}
\bibitem{ref197} C\$\textasciicircum{}2\$-Cite Contextual-Aware Citation Generation for Attributed Large Language Models. arXiv:2602.00004v2, 2025. \url{https://arxiv.org/abs/2602.00004v2}
\bibitem{ref198} Generation-Time vs. Post-hoc Citation A Holistic Evaluation of LLM Attribution. arXiv:2509.21557v2, 2025. \url{https://arxiv.org/abs/2509.21557v2}
\bibitem{ref199} Evaluating Evidence Attribution in Generated Fact Checking Explanations. arXiv:2406.12645v3, 2024. \url{https://arxiv.org/abs/2406.12645v3}
\bibitem{ref200} Concise and Sufficient Sub-Sentence Citations for Retrieval-Augmented Generation. arXiv:2509.20859v2, 2025. \url{https://arxiv.org/abs/2509.20859v2}
\bibitem{ref201} Probing for Knowledge Attribution in Large Language Models. arXiv:2602.22787v2, 2026. \url{https://arxiv.org/abs/2602.22787v2}
\bibitem{ref202} HART Data-Driven Hallucination Attribution and Evidence-Based Tracing for Large Language Models. arXiv:2603.05828v1, 2026. \url{https://arxiv.org/abs/2603.05828v1}
\bibitem{ref203} Do I Really Know Learning Factual Self-Verification for Hallucination Reduction. arXiv:2602.02018v1, 2026. \url{https://arxiv.org/abs/2602.02018v1}
\bibitem{ref204} TRACE Trajectory Correction from Cross-layer Evidence for Hallucination Reduction. arXiv:2605.18163v1, 2026. \url{https://arxiv.org/abs/2605.18163v1}
\bibitem{ref205} Attribution-Guided Decoding. arXiv:2509.26307v2, 2025. \url{https://arxiv.org/abs/2509.26307v2}
\bibitem{ref206} HAVE Head-Adaptive Gating and ValuE Calibration for Hallucination Mitigation in Large Language Models. arXiv:2509.06596v1, 2025. \url{https://arxiv.org/abs/2509.06596v1}
\bibitem{ref207} Monitoring Decoding Mitigating Hallucination via Evaluating the Factuality of Partial Response during Generation. arXiv:2503.03106v2, 2025. \url{https://arxiv.org/abs/2503.03106v2}
\bibitem{ref208} GRAD Graph-Retrieved Adaptive Decoding for Hallucination Mitigation. arXiv:2511.03900v1, 2025. \url{https://arxiv.org/abs/2511.03900v1}
\bibitem{ref209} How Should Pre-Trained Language Models Be Fine-Tuned Towards Adversarial Robustness. arXiv:2112.11668v1, 2021. \url{https://arxiv.org/abs/2112.11668v1}
\bibitem{ref210} RbFT Robust Fine-tuning for Retrieval-Augmented Generation against Retrieval Defects. arXiv:2501.18365v1, 2025. \url{https://arxiv.org/abs/2501.18365v1}
\bibitem{ref211} Not All Contexts Are Equal Teaching LLMs Credibility-aware Generation. arXiv:2404.06809v1, 2024. \url{https://arxiv.org/abs/2404.06809v1}
\bibitem{ref212} MIRAGE Defending Long-Form RAG Against Misinformation Pollution. arXiv:2607.05069v1, 2026. \url{https://arxiv.org/abs/2607.05069v1}
\bibitem{ref213} Towards More Robust Retrieval-Augmented Generation Evaluating RAG Under Adversarial Poisoning Attacks. arXiv:2412.16708v2, 2024. \url{https://arxiv.org/abs/2412.16708v2}
\bibitem{ref214} Detecting Is Not Resolving The Monitoring Control Gap in Retrieval Augmented LLMs. arXiv:2605.27157v1, 2026. \url{https://arxiv.org/abs/2605.27157v1}
\bibitem{ref215} Astute RAG Overcoming Imperfect Retrieval Augmentation and Knowledge Conflicts for Large Language Models. arXiv:2410.07176v2, 2024. \url{https://arxiv.org/abs/2410.07176v2}
\bibitem{ref216} Stateful Evidence-Driven Retrieval-Augmented Generation with Iterative Reasoning. arXiv:2604.14170v1, 2026. \url{https://arxiv.org/abs/2604.14170v1}
\bibitem{ref217} Beyond Semantic Relevance Counterfactual Risk Minimization for Robust Retrieval-Augmented Generation. arXiv:2605.01302v1, 2026. \url{https://arxiv.org/abs/2605.01302v1}
\bibitem{ref218} Evaluating the Robustness of Retrieval-Augmented Generation to Adversarial Evidence in the Health Domain. arXiv:2509.03787v1, 2025. \url{https://arxiv.org/abs/2509.03787v1}
\bibitem{ref219} Influence Factors on RAG Poisoning. arXiv:2606.12469v1, 2026. \url{https://arxiv.org/abs/2606.12469v1}
\bibitem{ref220} Improving End-to-End Training of Retrieval-Augmented Generation Models via Joint Stochastic Approximation. arXiv:2508.18168v3, 2025. \url{https://arxiv.org/abs/2508.18168v3}
\bibitem{ref221} Direct Retrieval-augmented Optimization Synergizing Knowledge Selection and Language Models. arXiv:2505.03075v1, 2025. \url{https://arxiv.org/abs/2505.03075v1}
\bibitem{ref222} RAFT Adapting Language Model to Domain Specific RAG. arXiv:2403.10131v2, 2024. \url{https://arxiv.org/abs/2403.10131v2}
\bibitem{ref223} Efficient In-Domain Question Answering for Resource-Constrained Environments. arXiv:2409.17648v3, 2024. \url{https://arxiv.org/abs/2409.17648v3}
\bibitem{ref224} P-RAG Prompt-Enhanced Parametric RAG with LoRA and Selective CoT for Biomedical and Multi-Hop QA. arXiv:2602.15874v1, 2026. \url{https://arxiv.org/abs/2602.15874v1}
\bibitem{ref225} To Memorize or to Retrieve Scaling Laws for RAG-Considerate Pretraining. arXiv:2604.00715v1, 2026. \url{https://arxiv.org/abs/2604.00715v1}
\bibitem{ref226} DACL-RAG Data Augmentation Strategy with Curriculum Learning for Retrieval-Augmented Generation. arXiv:2505.10493v2, 2025. \url{https://arxiv.org/abs/2505.10493v2}
\bibitem{ref227} Systematic Knowledge Injection into Large Language Models via Diverse Augmentation for Domain-Specific RAG. arXiv:2502.08356v3, 2025. \url{https://arxiv.org/abs/2502.08356v3}
\bibitem{ref228} SmartRAG Jointly Learn RAG-Related Tasks From the Environment Feedback. arXiv:2410.18141v2, 2024. \url{https://arxiv.org/abs/2410.18141v2}
\bibitem{ref229} Low-Rank Adaptation Redux for Large Models. arXiv:2604.21905v1, 2026. \url{https://arxiv.org/abs/2604.21905v1}
\bibitem{ref230} Parameter-Efficient Fine-Tuning with Learnable Rank. arXiv:2606.04325v1, 2026. \url{https://arxiv.org/abs/2606.04325v1}
\bibitem{ref231} ALoRA Allocating Low-Rank Adaptation for Fine-tuning Large Language Models. arXiv:2403.16187v2, 2024. \url{https://arxiv.org/abs/2403.16187v2}
\bibitem{ref232} Sparse Low-rank Adaptation of Pre-trained Language Models. arXiv:2311.11696v1, 2023. \url{https://arxiv.org/abs/2311.11696v1}
\bibitem{ref233} DoRA Enhancing Parameter-Efficient Fine-Tuning with Dynamic Rank Distribution. arXiv:2405.17357v3, 2024. \url{https://arxiv.org/abs/2405.17357v3}
\bibitem{ref234} GraLoRA Granular Low-Rank Adaptation for Parameter-Efficient Fine-Tuning. arXiv:2505.20355v2, 2025. \url{https://arxiv.org/abs/2505.20355v2}
\bibitem{ref235} FlyLoRA Boosting Task Decoupling and Parameter Efficiency via Implicit Rank-Wise Mixture-of-Experts. arXiv:2510.08396v2, 2025. \url{https://arxiv.org/abs/2510.08396v2}
\bibitem{ref236} NeuroLoRA Context-Aware Neuromodulation for Parameter-Efficient Multi-Task Adaptation. arXiv:2603.12378v1, 2026. \url{https://arxiv.org/abs/2603.12378v1}
\bibitem{ref237} MoRE A Mixture of Low-Rank Experts for Adaptive Multi-Task Learning. arXiv:2505.22694v1, 2025. \url{https://arxiv.org/abs/2505.22694v1}
\bibitem{ref238} LowRA Accurate and Efficient LoRA Fine-Tuning of LLMs under 2 Bits. arXiv:2502.08141v1, 2025. \url{https://arxiv.org/abs/2502.08141v1}
\bibitem{ref239} RoLoRA Fine-tuning Rotated Outlier-free LLMs for Effective Weight-Activation Quantization. arXiv:2407.08044v2, 2024. \url{https://arxiv.org/abs/2407.08044v2}
\bibitem{ref240} CARE-LoRA Compressed Activation REconstruction for Memory-Efficient LoRA. arXiv:2607.11940v1, 2026. \url{https://arxiv.org/abs/2607.11940v1}
\bibitem{ref241} LoRA-FA Memory-efficient Low-rank Adaptation for Large Language Models Fine-tuning. arXiv:2308.03303v1, 2023. \url{https://arxiv.org/abs/2308.03303v1}
\bibitem{ref242} Improving LoRA in Privacy-preserving Federated Learning. arXiv:2403.12313v1, 2024. \url{https://arxiv.org/abs/2403.12313v1}
\bibitem{ref243} LoRA-GA Low-Rank Adaptation with Gradient Approximation. arXiv:2407.05000v2, 2024. \url{https://arxiv.org/abs/2407.05000v2}
\bibitem{ref244} SC-LoRA Balancing Efficient Fine-tuning and Knowledge Preservation via Subspace-Constrained LoRA. arXiv:2505.23724v3, 2025. \url{https://arxiv.org/abs/2505.23724v3}
\bibitem{ref245} A Rank Stabilization Scaling Factor for Fine-Tuning with LoRA. arXiv:2312.03732v1, 2023. \url{https://arxiv.org/abs/2312.03732v1}
\bibitem{ref246} Flat-LoRA Low-Rank Adaptation over a Flat Loss Landscape. arXiv:2409.14396v2, 2024. \url{https://arxiv.org/abs/2409.14396v2}
\bibitem{ref247} RefLoRA Refactored Low-Rank Adaptation for Efficient Fine-Tuning of Large Models. arXiv:2505.18877v4, 2025. \url{https://arxiv.org/abs/2505.18877v4}
\bibitem{ref248} Rethinking Adapter Placement A Dominant Adaptation Module Perspective. arXiv:2605.06183v1, 2026. \url{https://arxiv.org/abs/2605.06183v1}
\bibitem{ref249} PLoP Precise LoRA Placement for Efficient Finetuning of Large Models. arXiv:2506.20629v1, 2025. \url{https://arxiv.org/abs/2506.20629v1}
\bibitem{ref250} Rank Also Matters Hierarchical Configuration for Mixture of Adapter Experts in LLM Fine-Tuning. arXiv:2502.03884v1, 2025. \url{https://arxiv.org/abs/2502.03884v1}
\bibitem{ref251} Continual Fine-Tuning of Large Language Models via Program Memory. arXiv:2605.13162v1, 2026. \url{https://arxiv.org/abs/2605.13162v1}
\bibitem{ref252} Efficient Modular Learning through Naive LoRA Summation Leveraging Orthogonality in High-Dimensional Models. arXiv:2508.11985v1, 2025. \url{https://arxiv.org/abs/2508.11985v1}
\bibitem{ref253} A Survey of Direct Preference Optimization. arXiv:2503.11701v1, 2025. \url{https://arxiv.org/abs/2503.11701v1}
\bibitem{ref254} A Comprehensive Survey of Direct Preference Optimization Datasets, Theories, Variants, and Applications. arXiv:2410.15595v4, 2024. \url{https://arxiv.org/abs/2410.15595v4}
\bibitem{ref255} PA-RAG RAG Alignment via Multi-Perspective Preference Optimization. arXiv:2412.14510v1, 2024. \url{https://arxiv.org/abs/2412.14510v1}
\bibitem{ref256} Robust Preference Optimization through Reward Model Distillation. arXiv:2405.19316v2, 2024. \url{https://arxiv.org/abs/2405.19316v2}
\bibitem{ref257} Understanding Likelihood Over-optimisation in Direct Alignment Algorithms. arXiv:2410.11677v2, 2024. \url{https://arxiv.org/abs/2410.11677v2}
\bibitem{ref258} RPO Retrieval Preference Optimization for Robust Retrieval-Augmented Generation. arXiv:2501.13726v2, 2025. \url{https://arxiv.org/abs/2501.13726v2}
\bibitem{ref259} BPO Revisiting Preference Modeling in Direct Preference Optimization. arXiv:2506.03557v1, 2025. \url{https://arxiv.org/abs/2506.03557v1}
\bibitem{ref260} Multi-Preference Optimization Generalizing DPO via Set-Level Contrasts. arXiv:2412.04628v4, 2024. \url{https://arxiv.org/abs/2412.04628v4}
\bibitem{ref261} Earlier Tokens Contribute More Learning Direct Preference Optimization From Temporal Decay Perspective. arXiv:2502.14340v1, 2025. \url{https://arxiv.org/abs/2502.14340v1}
\bibitem{ref262} HiPO Hierarchical Preference Optimization for Adaptive Reasoning in LLMs. arXiv:2604.20140v2, 2026. \url{https://arxiv.org/abs/2604.20140v2}
\bibitem{ref263} Reflective Preference Optimization (RPO) Enhancing On-Policy Alignment via Hint-Guided Reflection. arXiv:2512.13240v1, 2025. \url{https://arxiv.org/abs/2512.13240v1}
\bibitem{ref264} A Technical Survey of Reinforcement Learning Techniques for Large Language Models. arXiv:2507.04136v2, 2025. \url{https://arxiv.org/abs/2507.04136v2}
\bibitem{ref265} Optimizing Safe and Aligned Language Generation A Multi-Objective GRPO Approach. arXiv:2503.21819v1, 2025. \url{https://arxiv.org/abs/2503.21819v1}
\bibitem{ref266} Search-P1 Path-Centric Reward Shaping for Stable and Efficient Agentic RAG Training. arXiv:2602.22576v1, 2026. \url{https://arxiv.org/abs/2602.22576v1}
\bibitem{ref267} Curriculum Guided Reinforcement Learning for Efficient Multi Hop Retrieval Augmented Generation. arXiv:2505.17391v1, 2025. \url{https://arxiv.org/abs/2505.17391v1}
\bibitem{ref268} Learning to Trust Dynamic Utilization of Retrieval-Augmented Generation for E-commerce Search Relevance. arXiv:2510.11122v2, 2025. \url{https://arxiv.org/abs/2510.11122v2}
\bibitem{ref269} EP-GRPO Entropy-Progress Aligned Group Relative Policy Optimization with Implicit Process Guidance. arXiv:2605.04960v1, 2026. \url{https://arxiv.org/abs/2605.04960v1}
\bibitem{ref270} Unifying Group-Relative and Self-Distillation Policy Optimization via Sample Routing. arXiv:2604.02288v1, 2026. \url{https://arxiv.org/abs/2604.02288v1}
\bibitem{ref271} AMIR-GRPO Inducing Implicit Preference Signals into GRPO. arXiv:2601.03661v1, 2026. \url{https://arxiv.org/abs/2601.03661v1}
\bibitem{ref272} GDPO Group reward-Decoupled Normalization Policy Optimization for Multi-reward RL Optimization. arXiv:2601.05242v1, 2026. \url{https://arxiv.org/abs/2601.05242v1}
\bibitem{ref273} GD\$\textasciicircum{}2\$PO Mitigating Multi-Reward Conflicts via Group-Dynamic reward-Decoupled Policy Optimization. arXiv:2606.16771v1, 2026. \url{https://arxiv.org/abs/2606.16771v1}
\bibitem{ref274} CuSearch Curriculum Rollout Sampling via Search Depth for Agentic RAG. arXiv:2605.11611v2, 2026. \url{https://arxiv.org/abs/2605.11611v2}
\bibitem{ref275} Bridging the Preference Gap between Retrievers and LLMs. arXiv:2401.06954v2, 2024. \url{https://arxiv.org/abs/2401.06954v2}
\bibitem{ref276} Optimizing Retrieval for RAG via Reinforcement Learning. arXiv:2510.24652v2, 2025. \url{https://arxiv.org/abs/2510.24652v2}
\bibitem{ref277} Fine-Grained Guidance for Retrievers Leveraging LLMs' Feedback in Retrieval-Augmented Generation. arXiv:2411.03957v1, 2024. \url{https://arxiv.org/abs/2411.03957v1}
\bibitem{ref278} Understand What LLM Needs Dual Preference Alignment for Retrieval-Augmented Generation. arXiv:2406.18676v2, 2024. \url{https://arxiv.org/abs/2406.18676v2}
\bibitem{ref279} Improving Retrieval-Augmented Generation through Multi-Agent Reinforcement Learning. arXiv:2501.15228v2, 2025. \url{https://arxiv.org/abs/2501.15228v2}
\bibitem{ref280} C-3PO Compact Plug-and-Play Proxy Optimization to Achieve Human-like Retrieval-Augmented Generation. arXiv:2502.06205v2, 2025. \url{https://arxiv.org/abs/2502.06205v2}
\bibitem{ref281} Leveraging RAG for Training-Free Alignment of LLMs. arXiv:2605.11217v1, 2026. \url{https://arxiv.org/abs/2605.11217v1}
\bibitem{ref282} Knowing You Don't Know Learning When to Continue Search in Multi-round RAG through Self-Practicing. arXiv:2505.02811v2, 2025. \url{https://arxiv.org/abs/2505.02811v2}
\bibitem{ref283} Improving Consistency in Retrieval-Augmented Systems with Group Similarity Rewards. arXiv:2510.04392v1, 2025. \url{https://arxiv.org/abs/2510.04392v1}
\bibitem{ref284} Efficient Retrieval-Augmented Generation via Token Co-occurrence Graphs. arXiv:2606.30093v1, 2026. \url{https://arxiv.org/abs/2606.30093v1}
\bibitem{ref285} Linear Constraints. arXiv:2604.21467v1, 2026. \url{https://arxiv.org/abs/2604.21467v1}
\bibitem{ref286} Algorand. arXiv:1607.01341v9, 2016. \url{https://arxiv.org/abs/1607.01341v9}
\bibitem{ref287} Clue-RAG Towards Accurate and Cost-Efficient Graph-based RAG via Multi-Partite Graph and Query-Driven Iterative Retrieval. arXiv:2507.08445v3, 2025. \url{https://arxiv.org/abs/2507.08445v3}
\bibitem{ref288} What are kets. arXiv:2405.10055v1, 2024. \url{https://arxiv.org/abs/2405.10055v1}
\bibitem{ref289} A Unified Framework for Context-Aware and Relation-Aware Graph Retrieval-Augmented Generation. arXiv:2606.18075v1, 2026. \url{https://arxiv.org/abs/2606.18075v1}
\bibitem{ref290} Hyper-Connections. arXiv:2409.19606v3, 2024. \url{https://arxiv.org/abs/2409.19606v3}
\bibitem{ref291} CRAG -- Comprehensive RAG Benchmark. arXiv:2406.04744v2, 2024. \url{https://arxiv.org/abs/2406.04744v2}
\bibitem{ref292} PersianRAG A Retrieval-Augmented Generation System for Persian Language. arXiv:2411.02832v2, 2024. \url{https://arxiv.org/abs/2411.02832v2}
\bibitem{ref293} Hypergraph Categories. arXiv:1806.08304v3, 2018. \url{https://arxiv.org/abs/1806.08304v3}
\bibitem{ref294} A Navigational Approach for Comprehensive RAG via Traversal over Proposition Graphs. arXiv:2601.04859v1, 2026. \url{https://arxiv.org/abs/2601.04859v1}
\bibitem{ref295} Topology of Reasoning Understanding Large Reasoning Models through Reasoning Graph Properties. arXiv:2506.05744v3, 2025. \url{https://arxiv.org/abs/2506.05744v3}
\bibitem{ref296} ÆTHEL Automatically Extracted Typelogical Derivations for Dutch. arXiv:1912.12635v2, 2019. \url{https://arxiv.org/abs/1912.12635v2}
\bibitem{ref297} Petri Automata. arXiv:1702.01804v3, 2017. \url{https://arxiv.org/abs/1702.01804v3}
\bibitem{ref298} Query-Aware Spreading Activation for Multi-Hop Retrieval over Knowledge Graphs. arXiv:2606.30133v1, 2026. \url{https://arxiv.org/abs/2606.30133v1}
\bibitem{ref299} SemFlowRAG Directed Semantic Flow from Abstraction to Evidence for Complex Reasoning. arXiv:2606.28447v1, 2026. \url{https://arxiv.org/abs/2606.28447v1}
\bibitem{ref300} Use Graph When It Needs Efficiently and Adaptively Integrating Retrieval-Augmented Generation with Graphs. arXiv:2602.03578v1, 2026. \url{https://arxiv.org/abs/2602.03578v1}
\bibitem{ref301} Magistral. arXiv:2506.10910v1, 2025. \url{https://arxiv.org/abs/2506.10910v1}
\bibitem{ref302} Graph Learning. arXiv:2507.05636v2, 2025. \url{https://arxiv.org/abs/2507.05636v2}
\bibitem{ref303} Cog-RAG Cognitive-Inspired Dual-Hypergraph with Theme Alignment Retrieval-Augmented Generation. arXiv:2511.13201v2, 2025. \url{https://arxiv.org/abs/2511.13201v2}
\bibitem{ref304} PreQRAG -- Classify and Rewrite for Enhanced RAG. arXiv:2506.17493v1, 2025. \url{https://arxiv.org/abs/2506.17493v1}
\bibitem{ref305} Dissecting Agentic RAG A Component Ablation for Multi-Hop QA with a Local 7B Model. arXiv:2606.21553v1, 2026. \url{https://arxiv.org/abs/2606.21553v1}
\bibitem{ref306} Graph-R1 Towards Agentic GraphRAG Framework via End-to-end Reinforcement Learning. arXiv:2507.21892v2, 2025. \url{https://arxiv.org/abs/2507.21892v2}
\bibitem{ref307} A2RAG Adaptive Agentic Graph Retrieval for Cost-Aware and Reliable Reasoning. arXiv:2601.21162v2, 2026. \url{https://arxiv.org/abs/2601.21162v2}
\bibitem{ref308} Stop-RAG Value-Based Retrieval Control for Iterative RAG. arXiv:2510.14337v1, 2025. \url{https://arxiv.org/abs/2510.14337v1}
\bibitem{ref309} Doctor-RAG A Failure-Aware Repair Framework for Agentic Retrieval-Augmented Generation. arXiv:2604.00865v2, 2026. \url{https://arxiv.org/abs/2604.00865v2}
\bibitem{ref310} From Conflict to Consensus Boosting Medical Reasoning via Multi-Round Agentic RAG. arXiv:2603.03292v3, 2026. \url{https://arxiv.org/abs/2603.03292v3}
\bibitem{ref311} TeaRAG A Token-Efficient Agentic Retrieval-Augmented Generation Framework. arXiv:2511.05385v2, 2025. \url{https://arxiv.org/abs/2511.05385v2}
\bibitem{ref312} GRASP GRanularity-Aware Search Policy for Agentic RAG. arXiv:2607.10463v1, 2026. \url{https://arxiv.org/abs/2607.10463v1}
\bibitem{ref313} Plan RAG Efficient Test-Time Planning for Retrieval Augmented Generation. arXiv:2410.20753v2, 2024. \url{https://arxiv.org/abs/2410.20753v2}
\bibitem{ref314} TreePS-RAG Tree-based Process Supervision for Reinforcement Learning in Agentic RAG. arXiv:2601.06922v1, 2026. \url{https://arxiv.org/abs/2601.06922v1}
\bibitem{ref315} DecEx-RAG Boosting Agentic Retrieval-Augmented Generation with Decision and Execution Optimization via Process Supervision. arXiv:2510.05691v1, 2025. \url{https://arxiv.org/abs/2510.05691v1}
\bibitem{ref316} HiPRAG Hierarchical Process Rewards for Efficient Agentic Retrieval Augmented Generation. arXiv:2510.07794v2, 2025. \url{https://arxiv.org/abs/2510.07794v2}
\bibitem{ref317} Process vs. Outcome Reward Which is Better for Agentic RAG Reinforcement Learning. arXiv:2505.14069v3, 2025. \url{https://arxiv.org/abs/2505.14069v3}
\bibitem{ref318} AgenticRAGTracer A Hop-Aware Benchmark for Diagnosing Multi-Step Retrieval Reasoning in Agentic RAG. arXiv:2602.19127v2, 2026. \url{https://arxiv.org/abs/2602.19127v2}
\bibitem{ref319} Am I on the Right Track What Can Predicted Query Performance Tell Us about the Search Behaviour of Agentic RAG. arXiv:2507.10411v1, 2025. \url{https://arxiv.org/abs/2507.10411v1}
\bibitem{ref320} HM-RAG Hierarchical Multi-Agent Multimodal Retrieval Augmented Generation. arXiv:2504.12330v1, 2025. \url{https://arxiv.org/abs/2504.12330v1}
\bibitem{ref321} MA-RAG Multi-Agent Retrieval-Augmented Generation via Collaborative Chain-of-Thought Reasoning. arXiv:2505.20096v2, 2025. \url{https://arxiv.org/abs/2505.20096v2}
\bibitem{ref322} MemGraphRAG Memory-based Multi-Agent System for Graph Retrieval-Augmented Generation. arXiv:2606.00610v1, 2026. \url{https://arxiv.org/abs/2606.00610v1}
\bibitem{ref323} Experience as a Compass Multi-agent RAG with Evolving Orchestration and Agent Prompts. arXiv:2604.00901v2, 2026. \url{https://arxiv.org/abs/2604.00901v2}
\bibitem{ref324} Towards Open-World Retrieval-Augmented Generation on Knowledge Graph A Multi-Agent Collaboration Framework. arXiv:2509.01238v2, 2025. \url{https://arxiv.org/abs/2509.01238v2}
\bibitem{ref325} ConRAG Consensus-Driven Multi-View Retrieval for Multi-Hop Question Answering. arXiv:2605.28093v2, 2026. \url{https://arxiv.org/abs/2605.28093v2}
\bibitem{ref326} SPD-RAG Sub-Agent Per Document Retrieval-Augmented Generation. arXiv:2603.08329v1, 2026. \url{https://arxiv.org/abs/2603.08329v1}
\bibitem{ref327} HydraRAG Structured Cross-Source Enhanced Large Language Model Reasoning. arXiv:2505.17464v4, 2025. \url{https://arxiv.org/abs/2505.17464v4}
\bibitem{ref328} SCOUT-RAG Scalable and Cost-Efficient Unifying Traversal for Agentic Graph-RAG over Distributed Domains. arXiv:2602.08400v1, 2026. \url{https://arxiv.org/abs/2602.08400v1}
\bibitem{ref329} Divide by Question, Conquer by Agent SPLIT-RAG with Question-Driven Graph Partitioning. arXiv:2505.13994v2, 2025. \url{https://arxiv.org/abs/2505.13994v2}
\bibitem{ref330} CIIR@LiveRAG 2025 Optimizing Multi-Agent Retrieval Augmented Generation through Self-Training. arXiv:2506.10844v1, 2025. \url{https://arxiv.org/abs/2506.10844v1}
\bibitem{ref331} TechRAG Evidence-Gated Multimodal Agentic RAG for Technical Literature Reasoning. arXiv:2606.01613v2, 2026. \url{https://arxiv.org/abs/2606.01613v2}
\bibitem{ref332} Multi-Agent GraphRAG A Text-to-Cypher Framework for Labeled Property Graphs. arXiv:2511.08274v1, 2025. \url{https://arxiv.org/abs/2511.08274v1}
\bibitem{ref333} NodeRAG Structuring Graph-based RAG with Heterogeneous Nodes. arXiv:2504.11544v1, 2025. \url{https://arxiv.org/abs/2504.11544v1}
\bibitem{ref334} SLIM LSTMs. arXiv:1812.11391v1, 2018. \url{https://arxiv.org/abs/1812.11391v1}
\bibitem{ref335} GraphRAG under Fire. arXiv:2501.14050v4, 2025. \url{https://arxiv.org/abs/2501.14050v4}
\bibitem{ref336} LatentRAG Latent Reasoning and Retrieval for Efficient Agentic RAG. arXiv:2605.06285v1, 2026. \url{https://arxiv.org/abs/2605.06285v1}
\bibitem{ref337} GraphRAG-Router Learning Cost-Efficient Routing over GraphRAGs and LLMs with Reinforcement Learning. arXiv:2604.16401v1, 2026. \url{https://arxiv.org/abs/2604.16401v1}
\bibitem{ref338} GraphRAG-Bench Challenging Domain-Specific Reasoning for Evaluating Graph Retrieval-Augmented Generation. arXiv:2506.02404v3, 2025. \url{https://arxiv.org/abs/2506.02404v3}
\bibitem{ref339} RAG vs. GraphRAG A Systematic Evaluation and Key Insights. arXiv:2502.11371v3, 2025. \url{https://arxiv.org/abs/2502.11371v3}
\bibitem{ref340} Crowdsearch. arXiv:2311.08532v1, 2023. \url{https://arxiv.org/abs/2311.08532v1}
\bibitem{ref341} VectorLiteRAG Latency-Aware and Fine-Grained Resource Partitioning for Efficient RAG. arXiv:2504.08930v3, 2025. \url{https://arxiv.org/abs/2504.08930v3}
\bibitem{ref342} VisRAG Vision-based Retrieval-augmented Generation on Multi-modality Documents. arXiv:2410.10594v2, 2024. \url{https://arxiv.org/abs/2410.10594v2}
\bibitem{ref343} A Survey of Multimodal Retrieval-Augmented Generation. arXiv:2504.08748v1, 2025. \url{https://arxiv.org/abs/2504.08748v1}
\bibitem{ref344} MuRAG Multimodal Retrieval-Augmented Generator for Open Question Answering over Images and Text. arXiv:2210.02928v2, 2022. \url{https://arxiv.org/abs/2210.02928v2}
\bibitem{ref345} Document Haystacks Vision-Language Reasoning Over Piles of 1000+ Documents. arXiv:2411.16740v3, 2024. \url{https://arxiv.org/abs/2411.16740v3}
\bibitem{ref346} Beyond Patch Aggregation 3-Pass Pyramid Indexing for Vision-Enhanced Document Retrieval. arXiv:2511.21121v2, 2025. \url{https://arxiv.org/abs/2511.21121v2}
\bibitem{ref347} CMRAG Co-modality-based visual document retrieval and question answering. arXiv:2509.02123v3, 2025. \url{https://arxiv.org/abs/2509.02123v3}
\bibitem{ref348} Multimodal RAG Enhanced Visual Description. arXiv:2508.09170v1, 2025. \url{https://arxiv.org/abs/2508.09170v1}
\bibitem{ref349} Purifying Multimodal Retrieval Fragment-Level Evidence Selection for RAG. arXiv:2604.27600v1, 2026. \url{https://arxiv.org/abs/2604.27600v1}
\bibitem{ref350} RegionRAG Region-level Retrieval-Augmented Generation for Visual Document Understanding. arXiv:2510.27261v3, 2025. \url{https://arxiv.org/abs/2510.27261v3}
\bibitem{ref351} LFRAG Layout-oriented Fine-grained Retrieval-Augmented Generation on Multimodal Document Understanding. arXiv:2605.22829v1, 2026. \url{https://arxiv.org/abs/2605.22829v1}
\bibitem{ref352} RAVEN Multitask Retrieval Augmented Vision-Language Learning. arXiv:2406.19150v1, 2024. \url{https://arxiv.org/abs/2406.19150v1}
\bibitem{ref353} Towards Mixed-Modal Retrieval for Universal Retrieval-Augmented Generation. arXiv:2510.17354v1, 2025. \url{https://arxiv.org/abs/2510.17354v1}
\bibitem{ref354} CoRe-MMRAG Cross-Source Knowledge Reconciliation for Multimodal RAG. arXiv:2506.02544v2, 2025. \url{https://arxiv.org/abs/2506.02544v2}
\bibitem{ref355} Roles of MLLMs in Visually Rich Document Retrieval for RAG A Survey. arXiv:2601.03262v1, 2025. \url{https://arxiv.org/abs/2601.03262v1}
\bibitem{ref356} Benchmarking Retrieval-Augmented Generation in Multi-Modal Contexts. arXiv:2502.17297v2, 2025. \url{https://arxiv.org/abs/2502.17297v2}
\bibitem{ref357} MRAG-Bench Vision-Centric Evaluation for Retrieval-Augmented Multimodal Models. arXiv:2410.08182v2, 2024. \url{https://arxiv.org/abs/2410.08182v2}
\bibitem{ref358} UNIDOC-BENCH A Unified Benchmark for Document-Centric Multimodal RAG. arXiv:2510.03663v3, 2025. \url{https://arxiv.org/abs/2510.03663v3}
\bibitem{ref359} VisDoM Multi-Document QA with Visually Rich Elements Using Multimodal Retrieval-Augmented Generation. arXiv:2412.10704v2, 2024. \url{https://arxiv.org/abs/2412.10704v2}
\bibitem{ref360} Benchmarking Retrieval-Augmented Multimodal Generation for Document Question Answering. arXiv:2505.16470v2, 2025. \url{https://arxiv.org/abs/2505.16470v2}
\bibitem{ref361} MRAMG-Bench A Comprehensive Benchmark for Advancing Multimodal Retrieval-Augmented Multimodal Generation. arXiv:2502.04176v2, 2025. \url{https://arxiv.org/abs/2502.04176v2}
\bibitem{ref362} RAG-IGBench Innovative Evaluation for RAG-based Interleaved Generation in Open-domain Question Answering. arXiv:2512.05119v1, 2025. \url{https://arxiv.org/abs/2512.05119v1}
\bibitem{ref363} MedGEN-Bench Contextually entangled benchmark for open-ended multimodal medical generation. arXiv:2511.13135v2, 2025. \url{https://arxiv.org/abs/2511.13135v2}
\bibitem{ref364} MRAG-Suite A Diagnostic Evaluation Platform for Visual Retrieval-Augmented Generation. arXiv:2509.24253v3, 2025. \url{https://arxiv.org/abs/2509.24253v3}
\bibitem{ref365} MAVIS A Benchmark for Multimodal Source Attribution in Long-form Visual Question Answering. arXiv:2511.12142v2, 2025. \url{https://arxiv.org/abs/2511.12142v2}
\bibitem{ref366} Benchmarking Multimodal RAG through a Chart-based Document Question-Answering Generation Framework. arXiv:2502.14864v1, 2025. \url{https://arxiv.org/abs/2502.14864v1}
\bibitem{ref367} MRAG Benchmarking Retrieval-Augmented Generation for Bio-medicine. arXiv:2601.16503v2, 2026. \url{https://arxiv.org/abs/2601.16503v2}
\bibitem{ref368} M4-RAG A Massive-Scale Multilingual Multi-Cultural Multimodal RAG. arXiv:2512.05959v2, 2025. \url{https://arxiv.org/abs/2512.05959v2}
\bibitem{ref369} Multilingual Retrieval-Augmented Generation for Knowledge-Intensive Task. arXiv:2504.03616v2, 2025. \url{https://arxiv.org/abs/2504.03616v2}
\bibitem{ref370} Quality-Aware Translation Tagging in Multilingual RAG system. arXiv:2510.23070v1, 2025. \url{https://arxiv.org/abs/2510.23070v1}
\bibitem{ref371} Investigating Language Preference of Multilingual RAG Systems. arXiv:2502.11175v4, 2025. \url{https://arxiv.org/abs/2502.11175v4}
\bibitem{ref372} Linguistic Nepotism Trading-off Quality for Language Preference in Multilingual RAG. arXiv:2509.13930v3, 2025. \url{https://arxiv.org/abs/2509.13930v3}
\bibitem{ref373} All Languages Matter Understanding and Mitigating Language Bias in Multilingual RAG. arXiv:2604.20199v1, 2026. \url{https://arxiv.org/abs/2604.20199v1}
\bibitem{ref374} Enhancing Multilingual RAG Systems with Debiased Language Preference-Guided Query Fusion. arXiv:2601.02956v3, 2026. \url{https://arxiv.org/abs/2601.02956v3}
\bibitem{ref375} Language Drift in Multilingual Retrieval-Augmented Generation Characterization and Decoding-Time Mitigation. arXiv:2511.09984v1, 2025. \url{https://arxiv.org/abs/2511.09984v1}
\bibitem{ref376} Distill Where the Student Goes Teacher-Regularized RL for English-Evidence Cross-Lingual RAG. arXiv:2607.02966v2, 2026. \url{https://arxiv.org/abs/2607.02966v2}
\bibitem{ref377} On the Consistency of Multilingual Context Utilization in Retrieval-Augmented Generation. arXiv:2504.00597v4, 2025. \url{https://arxiv.org/abs/2504.00597v4}
\bibitem{ref378} The Cross-Lingual Cost Retrieval Biases in RAG over Arabic-English Corpora. arXiv:2507.07543v2, 2025. \url{https://arxiv.org/abs/2507.07543v2}
\bibitem{ref379} Retrieval-augmented generation in multilingual settings. arXiv:2407.01463v1, 2024. \url{https://arxiv.org/abs/2407.01463v1}
\bibitem{ref380} DaPT A Dual-Path Framework for Multilingual Multi-hop Question Answering. arXiv:2603.19097v1, 2026. \url{https://arxiv.org/abs/2603.19097v1}
\bibitem{ref381} CroSearch-R1 Better Leveraging Cross-lingual Knowledge for Retrieval-Augmented Generation. arXiv:2604.25182v1, 2026. \url{https://arxiv.org/abs/2604.25182v1}
\bibitem{ref382} Language-Coupled Reinforcement Learning for Multilingual Retrieval-Augmented Generation. arXiv:2601.14896v3, 2026. \url{https://arxiv.org/abs/2601.14896v3}
\bibitem{ref383} XRAG Cross-lingual Retrieval-Augmented Generation. arXiv:2505.10089v1, 2025. \url{https://arxiv.org/abs/2505.10089v1}
\bibitem{ref384} MEMERAG A Multilingual End-to-End Meta-Evaluation Benchmark for Retrieval Augmented Generation. arXiv:2502.17163v4, 2025. \url{https://arxiv.org/abs/2502.17163v4}
\bibitem{ref385} MMed-Bench-IR A Heterogeneous Benchmark for Multilingual Medical Information Retrieval. arXiv:2606.24200v1, 2026. \url{https://arxiv.org/abs/2606.24200v1}
\bibitem{ref386} IndicRAGSuite Large-Scale Datasets and a Benchmark for Indian Language RAG Systems. arXiv:2506.01615v2, 2025. \url{https://arxiv.org/abs/2506.01615v2}
\bibitem{ref387} Multilingual Retrieval Augmented Generation for Culturally-Sensitive Tasks A Benchmark for Cross-lingual Robustness. arXiv:2410.01171v3, 2024. \url{https://arxiv.org/abs/2410.01171v3}
\bibitem{ref388} AlzheimerRAG Multimodal Retrieval Augmented Generation for Clinical Use Cases using PubMed articles. arXiv:2412.16701v3, 2024. \url{https://arxiv.org/abs/2412.16701v3}
\bibitem{ref389} Iterative Multimodal Retrieval-Augmented Generation for Medical Question Answering. arXiv:2604.27724v1, 2026. \url{https://arxiv.org/abs/2604.27724v1}
\bibitem{ref390} HeteroRAG A Heterogeneous Retrieval-Augmented Generation Framework for Medical Vision Language Tasks. arXiv:2508.12778v2, 2025. \url{https://arxiv.org/abs/2508.12778v2}
\bibitem{ref391} DoctorRAG Medical RAG Fusing Knowledge with Patient Analogy through Textual Gradients. arXiv:2505.19538v1, 2025. \url{https://arxiv.org/abs/2505.19538v1}
\bibitem{ref392} SimRAG Self-Improving Retrieval-Augmented Generation for Adapting Large Language Models to Specialized Domains. arXiv:2410.17952v2, 2024. \url{https://arxiv.org/abs/2410.17952v2}
\bibitem{ref393} Improving Retrieval-Augmented Generation in Medicine with Iterative Follow-up Questions. arXiv:2408.00727v3, 2024. \url{https://arxiv.org/abs/2408.00727v3}
\bibitem{ref394} LexRAG Benchmarking Retrieval-Augmented Generation in Multi-Turn Legal Consultation Conversation. arXiv:2502.20640v1, 2025. \url{https://arxiv.org/abs/2502.20640v1}
\bibitem{ref395} Fine-grained Claim-level RAG Benchmark for Law. arXiv:2605.21071v3, 2026. \url{https://arxiv.org/abs/2605.21071v3}
\bibitem{ref396} Unlocking Multi-View Insights in Knowledge-Dense Retrieval-Augmented Generation. arXiv:2404.12879v1, 2024. \url{https://arxiv.org/abs/2404.12879v1}
\bibitem{ref397} SMARTFinRAG Interactive Modularized Financial RAG Benchmark. arXiv:2504.18024v1, 2025. \url{https://arxiv.org/abs/2504.18024v1}
\bibitem{ref398} FinSage A Multi-aspect RAG System for Financial Filings Question Answering. arXiv:2504.14493v4, 2025. \url{https://arxiv.org/abs/2504.14493v4}
\bibitem{ref399} AstuteRAG-FQA Task-Aware Retrieval-Augmented Generation Framework for Proprietary Data Challenges in Financial Question Answering. arXiv:2510.27537v1, 2025. \url{https://arxiv.org/abs/2510.27537v1}
\bibitem{ref400} Benchmarking Retrieval-Augmented Generation for Chemistry. arXiv:2505.07671v2, 2025. \url{https://arxiv.org/abs/2505.07671v2}
\bibitem{ref401} Engineering RAG Systems for Real-World Applications Design, Development, and Evaluation. arXiv:2506.20869v3, 2025. \url{https://arxiv.org/abs/2506.20869v3}
\bibitem{ref402} MetaGen Blended RAG Unlocking Zero-Shot Precision for Specialized Domain Question-Answering. arXiv:2505.18247v3, 2025. \url{https://arxiv.org/abs/2505.18247v3}
\bibitem{ref403} Redefining Retrieval Evaluation in the Era of LLMs. arXiv:2510.21440v1, 2025. \url{https://arxiv.org/abs/2510.21440v1}
\bibitem{ref404} SePer Measure Retrieval Utility Through The Lens Of Semantic Perplexity Reduction. arXiv:2503.01478v5, 2025. \url{https://arxiv.org/abs/2503.01478v5}
\bibitem{ref405} RAGAS Automated Evaluation of Retrieval Augmented Generation. arXiv:2309.15217v1, 2023. \url{https://arxiv.org/abs/2309.15217v1}
\bibitem{ref406} CCRS A Zero-Shot LLM-as-a-Judge Framework for Comprehensive RAG Evaluation. arXiv:2506.20128v1, 2025. \url{https://arxiv.org/abs/2506.20128v1}
\bibitem{ref407} RAGChecker A Fine-grained Framework for Diagnosing Retrieval-Augmented Generation. arXiv:2408.08067v2, 2024. \url{https://arxiv.org/abs/2408.08067v2}
\bibitem{ref408} CoFE-RAG A Comprehensive Full-chain Evaluation Framework for Retrieval-Augmented Generation with Enhanced Data Diversity. arXiv:2410.12248v1, 2024. \url{https://arxiv.org/abs/2410.12248v1}
\bibitem{ref409} URAG A Benchmark for Uncertainty Quantification in Retrieval-Augmented Large Language Models. arXiv:2603.19281v1, 2026. \url{https://arxiv.org/abs/2603.19281v1}
\bibitem{ref410} CONFLARE CONFormal LArge language model REtrieval. arXiv:2404.04287v1, 2024. \url{https://arxiv.org/abs/2404.04287v1}
\bibitem{ref411} Towards Fine-Grained Citation Evaluation in Generated Text A Comparative Analysis of Faithfulness Metrics. arXiv:2406.15264v2, 2024. \url{https://arxiv.org/abs/2406.15264v2}
\bibitem{ref412} NeocorRAG Less Irrelevant Information, More Explicit Evidence, and More Effective Recall via Evidence Chains. arXiv:2604.27852v1, 2026. \url{https://arxiv.org/abs/2604.27852v1}
\bibitem{ref413} Coverage, Not Averages Semantic Stratification for Trustworthy Retrieval Evaluation. arXiv:2604.20763v1, 2026. \url{https://arxiv.org/abs/2604.20763v1}
\bibitem{ref414} GroUSE A Benchmark to Evaluate Evaluators in Grounded Question Answering. arXiv:2409.06595v3, 2024. \url{https://arxiv.org/abs/2409.06595v3}
\bibitem{ref415} RAGXplain From Explainable Evaluation to Actionable Guidance of RAG Pipelines. arXiv:2505.13538v2, 2025. \url{https://arxiv.org/abs/2505.13538v2}
\bibitem{ref416} PRGB Benchmark A Robust Placeholder-Assisted Algorithm for Benchmarking Retrieval-Augmented Generation. arXiv:2507.22927v1, 2025. \url{https://arxiv.org/abs/2507.22927v1}
\bibitem{ref417} RAGVUE A Diagnostic View for Explainable and Automated Evaluation of Retrieval-Augmented Generation. arXiv:2601.04196v1, 2025. \url{https://arxiv.org/abs/2601.04196v1}
\bibitem{ref418} RAGEval Scenario Specific RAG Evaluation Dataset Generation Framework. arXiv:2408.01262v5, 2024. \url{https://arxiv.org/abs/2408.01262v5}
\bibitem{ref419} FAB-Bench A Framework for Adaptive RAG Benchmarking in Semiconductor Manufacturing. arXiv:2605.26476v1, 2026. \url{https://arxiv.org/abs/2605.26476v1}
\bibitem{ref420} Case-Aware LLM-as-a-Judge Evaluation for Enterprise-Scale RAG Systems. arXiv:2602.20379v1, 2026. \url{https://arxiv.org/abs/2602.20379v1}
\bibitem{ref421} XRAG eXamining the Core -- Benchmarking Foundational Components in Advanced Retrieval-Augmented Generation. arXiv:2412.15529v4, 2024. \url{https://arxiv.org/abs/2412.15529v4}
\bibitem{ref422} FATHOMS-RAG A Framework for the Assessment of Thinking and Observation in Multimodal Systems that use Retrieval Augmented Generation. arXiv:2510.08945v3, 2025. \url{https://arxiv.org/abs/2510.08945v3}
\bibitem{ref423} Securing Retrieval-Augmented Generation A Taxonomy of Attacks, Defenses, and Future Directions. arXiv:2604.08304v3, 2026. \url{https://arxiv.org/abs/2604.08304v3}
\bibitem{ref424} PoisonedRAG Knowledge Poisoning Attacks to Retrieval-Augmented Generation of Large Language Models. arXiv:2402.07867v1, 2024. \url{https://arxiv.org/abs/2402.07867v1}
\bibitem{ref425} Practical Poisoning Attacks against Retrieval-Augmented Generation. arXiv:2504.03957v2, 2025. \url{https://arxiv.org/abs/2504.03957v2}
\bibitem{ref426} HijackRAG Hijacking Attacks against Retrieval-Augmented Large Language Models. arXiv:2410.22832v1, 2024. \url{https://arxiv.org/abs/2410.22832v1}
\bibitem{ref427} Token-Level Precise Attack on RAG Searching for the Best Alternatives to Mislead Generation. arXiv:2508.03110v1, 2025. \url{https://arxiv.org/abs/2508.03110v1}
\bibitem{ref428} Joint-GCG Unified Gradient-Based Poisoning Attacks on Retrieval-Augmented Generation Systems. arXiv:2506.06151v2, 2025. \url{https://arxiv.org/abs/2506.06151v2}
\bibitem{ref429} RIPRAG Hack a Black-box Retrieval-Augmented Generation Question-Answering System with Reinforcement Learning. arXiv:2510.10008v2, 2025. \url{https://arxiv.org/abs/2510.10008v2}
\bibitem{ref430} KidnapRAG A Black-Box Attack for Hijacking Reasoning in Agentic Retrieval-Augmented Generation Systems. arXiv:2607.00422v1, 2026. \url{https://arxiv.org/abs/2607.00422v1}
\bibitem{ref431} DisarmRAG Stealthy Retriever-Centric Poisoning to Disable Self-Correction in Retrieval-Augmented Generation (Extended Version). arXiv:2508.20083v2, 2025. \url{https://arxiv.org/abs/2508.20083v2}
\bibitem{ref432} Topic-FlipRAG Topic-Orientated Adversarial Opinion Manipulation Attacks to Retrieval-Augmented Generation Models. arXiv:2502.01386v3, 2025. \url{https://arxiv.org/abs/2502.01386v3}
\bibitem{ref433} FlippedRAG Black-Box Opinion Manipulation Adversarial Attacks to Retrieval-Augmented Generation Models. arXiv:2501.02968v6, 2025. \url{https://arxiv.org/abs/2501.02968v6}
\bibitem{ref434} The Silent Saboteur Imperceptible Adversarial Attacks against Black-Box Retrieval-Augmented Generation Systems. arXiv:2505.18583v2, 2025. \url{https://arxiv.org/abs/2505.18583v2}
\bibitem{ref435} Inference Cost Attacks for Retrieval-Augmented Large Language Models. arXiv:2606.02643v1, 2026. \url{https://arxiv.org/abs/2606.02643v1}
\bibitem{ref436} MM-PoisonRAG Disrupting Multimodal RAG with Local and Global Poisoning Attacks. arXiv:2502.17832v4, 2025. \url{https://arxiv.org/abs/2502.17832v4}
\bibitem{ref437} When Poison Fails After Retrieval Revisiting Corpus Poisoning under Chunking and Reranking Pipelines. arXiv:2606.11265v1, 2026. \url{https://arxiv.org/abs/2606.11265v1}
\bibitem{ref438} Typos that Broke the RAG's Back Genetic Attack on RAG Pipeline by Simulating Documents in the Wild via Low-level Perturbations. arXiv:2404.13948v1, 2024. \url{https://arxiv.org/abs/2404.13948v1}
\bibitem{ref439} Benchmarking Poisoning Attacks against Retrieval-Augmented Generation. arXiv:2505.18543v1, 2025. \url{https://arxiv.org/abs/2505.18543v1}
\bibitem{ref440} Uncovering Competing Poisoning Attacks in Retrieval-Augmented Generation. arXiv:2505.12574v5, 2025. \url{https://arxiv.org/abs/2505.12574v5}
\bibitem{ref441} PRA-RAG Provably Robust Aggregation in Retrieval-Augmented Generation against Retrieval Corruption. arXiv:2607.00012v1, 2026. \url{https://arxiv.org/abs/2607.00012v1}
\bibitem{ref442} Defending Against Knowledge Poisoning Attacks During Retrieval-Augmented Generation. arXiv:2508.02835v2, 2025. \url{https://arxiv.org/abs/2508.02835v2}
\bibitem{ref443} RAGPart \& RAGMask Retrieval-Stage Defenses Against Corpus Poisoning in Retrieval-Augmented Generation. arXiv:2512.24268v1, 2025. \url{https://arxiv.org/abs/2512.24268v1}
\bibitem{ref444} ProGRank Probe-Gradient Reranking to Defend Dense-Retriever RAG from Corpus Poisoning. arXiv:2603.22934v3, 2026. \url{https://arxiv.org/abs/2603.22934v3}
\bibitem{ref445} TrustRAG Enhancing Robustness and Trustworthiness in Retrieval-Augmented Generation. arXiv:2501.00879v3, 2025. \url{https://arxiv.org/abs/2501.00879v3}
\bibitem{ref446} ReliabilityRAG Effective and Provably Robust Defense for RAG-based Web-Search. arXiv:2509.23519v2, 2025. \url{https://arxiv.org/abs/2509.23519v2}
\bibitem{ref447} BiRD A Bidirectional Ranking Defense Mechanism for Retrieval Augmented Generation. arXiv:2605.20123v1, 2026. \url{https://arxiv.org/abs/2605.20123v1}
\bibitem{ref448} Through the Stealth Lens Attention-Aware Defenses Against Poisoning in RAG. arXiv:2506.04390v2, 2025. \url{https://arxiv.org/abs/2506.04390v2}
\bibitem{ref449} Addressing Corpus Knowledge Poisoning Attacks on RAG Using Sparse Attention. arXiv:2602.04711v2, 2026. \url{https://arxiv.org/abs/2602.04711v2}
\bibitem{ref450} RAGuard A Layered Defense Framework for Retrieval-Augmented Generation Systems Against Data Poisoning. arXiv:2607.26339v1, 2026. \url{https://arxiv.org/abs/2607.26339v1}
\bibitem{ref451} Regime-Aware Peer Specialization for Robust RAG under Heterogeneous Knowledge Conflicts. arXiv:2606.30518v1, 2026. \url{https://arxiv.org/abs/2606.30518v1}
\bibitem{ref452} Ensemble Privacy Defense for Knowledge-Intensive LLMs against Membership Inference Attacks. arXiv:2512.03100v1, 2025. \url{https://arxiv.org/abs/2512.03100v1}
\bibitem{ref453} On the Diminishing Returns of Complex Robust RAG Training in the Era of Powerful LLMs. arXiv:2502.11400v2, 2025. \url{https://arxiv.org/abs/2502.11400v2}
\bibitem{ref454} Rescuing the Unpoisoned Efficient Defense against Knowledge Corruption Attacks on RAG Systems. arXiv:2511.01268v1, 2025. \url{https://arxiv.org/abs/2511.01268v1}
\bibitem{ref455} EcoSafeRAG Efficient Security through Context Analysis in Retrieval-Augmented Generation. arXiv:2505.13506v1, 2025. \url{https://arxiv.org/abs/2505.13506v1}
\bibitem{ref456} TriShieldRAG A Three-Ring Defense-in-Depth Framework Against Knowledge Corruption in Retrieval-Augmented Generation. arXiv:2607.23838v1, 2026. \url{https://arxiv.org/abs/2607.23838v1}
\bibitem{ref457} Adaptive Defense Orchestration for RAG A Sentinel-Strategist Architecture against Multi-Vector Attacks. arXiv:2604.20932v1, 2026. \url{https://arxiv.org/abs/2604.20932v1}
\bibitem{ref458} RAG Safety Exploring Knowledge Poisoning Attacks to Retrieval-Augmented Generation. arXiv:2507.08862v1, 2025. \url{https://arxiv.org/abs/2507.08862v1}
\bibitem{ref459} Confundo Learning to Generate Robust Poison for Practical RAG Systems. arXiv:2602.06616v1, 2026. \url{https://arxiv.org/abs/2602.06616v1}
\bibitem{ref460} Phantom General Backdoor Attacks on Retrieval Augmented Language Generation. arXiv:2405.20485v3, 2024. \url{https://arxiv.org/abs/2405.20485v3}
\bibitem{ref461} TrojanRAG Retrieval-Augmented Generation Can Be Backdoor Driver in Large Language Models. arXiv:2405.13401v4, 2024. \url{https://arxiv.org/abs/2405.13401v4}
\bibitem{ref462} Backdoored Retrievers for Prompt Injection Attacks on Retrieval Augmented Generation of Large Language Models. arXiv:2410.14479v1, 2024. \url{https://arxiv.org/abs/2410.14479v1}
\bibitem{ref463} Data Extraction Attacks in Retrieval-Augmented Generation via Backdoors. arXiv:2411.01705v2, 2024. \url{https://arxiv.org/abs/2411.01705v2}
\bibitem{ref464} Architecture Matters Comparing RAG Systems under Knowledge Base Poisoning. arXiv:2605.05632v1, 2026. \url{https://arxiv.org/abs/2605.05632v1}
\bibitem{ref465} Pirates of the RAG Adaptively Attacking LLMs to Leak Knowledge Bases. arXiv:2412.18295v2, 2024. \url{https://arxiv.org/abs/2412.18295v2}
\bibitem{ref466} LeakDojo Decoding the Leakage Threats of RAG Systems. arXiv:2605.05818v1, 2026. \url{https://arxiv.org/abs/2605.05818v1}
\bibitem{ref467} Feedback-Guided Extraction of Knowledge Base from Retrieval-Augmented LLM Applications. arXiv:2411.14110v2, 2024. \url{https://arxiv.org/abs/2411.14110v2}
\bibitem{ref468} One Pic is All it Takes Poisoning Visual Document Retrieval Augmented Generation with a Single Image. arXiv:2504.02132v4, 2025. \url{https://arxiv.org/abs/2504.02132v4}
\bibitem{ref469} The Hidden Threat in Plain Text Attacking RAG Data Loaders. arXiv:2507.05093v1, 2025. \url{https://arxiv.org/abs/2507.05093v1}
\bibitem{ref470} Retrieval Pivot Attacks in Hybrid RAG Measuring and Mitigating Amplified Leakage from Vector Seeds to Graph Expansion. arXiv:2602.08668v3, 2026. \url{https://arxiv.org/abs/2602.08668v3}
\bibitem{ref471} Adversarial Threat Vectors and Risk Mitigation for Retrieval-Augmented Generation Systems. arXiv:2506.00281v1, 2025. \url{https://arxiv.org/abs/2506.00281v1}
\bibitem{ref472} ConfusedPilot Confused Deputy Risks in RAG-based LLMs. arXiv:2408.04870v5, 2024. \url{https://arxiv.org/abs/2408.04870v5}
\bibitem{ref473} CtrlRAG Black-box Document Poisoning Attacks for Retrieval-Augmented Generation of Large Language Models. arXiv:2503.06950v2, 2025. \url{https://arxiv.org/abs/2503.06950v2}
\bibitem{ref474} Mask-based Membership Inference Attacks for Retrieval-Augmented Generation. arXiv:2410.20142v2, 2024. \url{https://arxiv.org/abs/2410.20142v2}
\bibitem{ref475} DCMI A Differential Calibration Membership Inference Attack Against Retrieval-Augmented Generation. arXiv:2509.06026v1, 2025. \url{https://arxiv.org/abs/2509.06026v1}
\bibitem{ref476} BudgetLeak Membership Inference Attacks on RAG Systems via the Generation Budget Side Channel. arXiv:2511.12043v1, 2025. \url{https://arxiv.org/abs/2511.12043v1}
\bibitem{ref477} Five Queries Are Enough Query-Efficient and Surrogate-Free Membership Inference Attacks on RAG via Entailment. arXiv:2605.24312v3, 2026. \url{https://arxiv.org/abs/2605.24312v3}
\bibitem{ref478} Riddle Me This! Stealthy Membership Inference for Retrieval-Augmented Generation. arXiv:2502.00306v2, 2025. \url{https://arxiv.org/abs/2502.00306v2}
\bibitem{ref479} E-MIA Exam-Style Black-Box Membership Inference Attacks against RAG Systems. arXiv:2605.00955v1, 2026. \url{https://arxiv.org/abs/2605.00955v1}
\bibitem{ref480} MrM Black-Box Membership Inference Attacks against Multimodal RAG Systems. arXiv:2506.07399v1, 2025. \url{https://arxiv.org/abs/2506.07399v1}
\bibitem{ref481} External Data Extraction Attacks against Retrieval-Augmented Large Language Models. arXiv:2510.02964v2, 2025. \url{https://arxiv.org/abs/2510.02964v2}
\bibitem{ref482} Silent Leaks Implicit Knowledge Extraction Attack on RAG Systems through Benign Queries. arXiv:2505.15420v2, 2025. \url{https://arxiv.org/abs/2505.15420v2}
\bibitem{ref483} Unleashing Worms and Extracting Data Escalating the Outcome of Attacks against RAG-based Inference in Scale and Severity Using Jailbreaking. arXiv:2409.08045v1, 2024. \url{https://arxiv.org/abs/2409.08045v1}
\bibitem{ref484} Benchmarking Knowledge-Extraction Attack and Defense on Retrieval-Augmented Generation. arXiv:2602.09319v3, 2026. \url{https://arxiv.org/abs/2602.09319v3}
\bibitem{ref485} On the Privacy of LLMs An Ablation Study. arXiv:2605.02255v1, 2026. \url{https://arxiv.org/abs/2605.02255v1}
\bibitem{ref486} Investigating the prompt leakage effect and black-box defenses for multi-turn LLM interactions. arXiv:2404.16251v2, 2024. \url{https://arxiv.org/abs/2404.16251v2}
\bibitem{ref487} Provably Secure Retrieval-Augmented Generation. arXiv:2508.01084v1, 2025. \url{https://arxiv.org/abs/2508.01084v1}
\bibitem{ref488} GoldenRetriever Non-Interactive Homomorphic Encrypted Retrieval for Privacy-Preserving RAG. arXiv:2607.29019v1, 2026. \url{https://arxiv.org/abs/2607.29019v1}
\bibitem{ref489} Efficient Privacy-Preserving Retrieval Augmented Generation with Distance-Preserving Encryption. arXiv:2601.12331v1, 2026. \url{https://arxiv.org/abs/2601.12331v1}
\bibitem{ref490} Privacy-Preserving Retrieval-Augmented Generation with Differential Privacy. arXiv:2412.04697v3, 2024. \url{https://arxiv.org/abs/2412.04697v3}
\bibitem{ref491} Differentially Private Retrieval-Augmented Generation. arXiv:2602.14374v1, 2026. \url{https://arxiv.org/abs/2602.14374v1}
\bibitem{ref492} Private-RAG Answering Multiple Queries with LLMs while Keeping Your Data Private. arXiv:2511.07637v1, 2025. \url{https://arxiv.org/abs/2511.07637v1}
\bibitem{ref493} Is External Database Protection Static in Retrieval-Augmented Generation Rethinking Privacy Preservation under Dynamic Queries. arXiv:2607.14811v1, 2026. \url{https://arxiv.org/abs/2607.14811v1}
\bibitem{ref494} Differentially Private Synthetic Text Generation for Retrieval-Augmented Generation (RAG). arXiv:2510.06719v2, 2025. \url{https://arxiv.org/abs/2510.06719v2}
\bibitem{ref495} SD-RAG A Prompt-Injection-Resilient Framework for Selective Disclosure in Retrieval-Augmented Generation. arXiv:2601.11199v1, 2026. \url{https://arxiv.org/abs/2601.11199v1}
\bibitem{ref496} Not All Entities are Created Equal A Dynamic Anonymization Framework for Privacy-Preserving RAG. arXiv:2603.26074v3, 2026. \url{https://arxiv.org/abs/2603.26074v3}
\bibitem{ref497} An Efficient and Privacy-Preserving Architecture for Cross-Institutional Collaborative RAG. arXiv:2605.25716v1, 2026. \url{https://arxiv.org/abs/2605.25716v1}
\bibitem{ref498} C-FedRAG A Confidential Federated Retrieval-Augmented Generation System. arXiv:2412.13163v2, 2024. \url{https://arxiv.org/abs/2412.13163v2}
\bibitem{ref499} Your RAG is Unfair Exposing Fairness Vulnerabilities in Retrieval-Augmented Generation via Backdoor Attacks. arXiv:2509.22486v1, 2025. \url{https://arxiv.org/abs/2509.22486v1}
\bibitem{ref500} Fairness Testing in Retrieval-Augmented Generation How Small Perturbations Reveal Bias in Small Language Models. arXiv:2509.26584v1, 2025. \url{https://arxiv.org/abs/2509.26584v1}
\bibitem{ref501} Does RAG Introduce Unfairness in LLMs Evaluating Fairness in Retrieval-Augmented Generation Systems. arXiv:2409.19804v2, 2024. \url{https://arxiv.org/abs/2409.19804v2}
\bibitem{ref502} MedPriv-Bench Benchmarking the Privacy-Utility Trade-off of Large Language Models in Medical Open-End Question Answering. arXiv:2603.14265v1, 2026. \url{https://arxiv.org/abs/2603.14265v1}
\bibitem{ref503} Trust-Oriented Adaptive Guardrails for Large Language Models. arXiv:2408.08959v3, 2024. \url{https://arxiv.org/abs/2408.08959v3}
\bibitem{ref504} The RAG Paradox A Black-Box Attack Exploiting Unintentional Vulnerabilities in Retrieval-Augmented Generation Systems. arXiv:2502.20995v3, 2025. \url{https://arxiv.org/abs/2502.20995v3}
\bibitem{ref505} Secure Multifaceted-RAG for Enterprise Hybrid Knowledge Retrieval with Security Filtering. arXiv:2504.13425v2, 2025. \url{https://arxiv.org/abs/2504.13425v2}
\bibitem{ref506} Privacy-Aware RAG Secure and Isolated Knowledge Retrieval. arXiv:2503.15548v1, 2025. \url{https://arxiv.org/abs/2503.15548v1}
\bibitem{ref507} Towards Trustworthy Retrieval Augmented Generation for Large Language Models A Survey. arXiv:2502.06872v1, 2025. \url{https://arxiv.org/abs/2502.06872v1}
\bibitem{ref508} SafeRAG Benchmarking Security in Retrieval-Augmented Generation of Large Language Model. arXiv:2501.18636v2, 2025. \url{https://arxiv.org/abs/2501.18636v2}
\bibitem{ref509} PCR A Prefetch-Enhanced Cache Reuse System for Low-Latency RAG Serving. arXiv:2603.23049v1, 2026. \url{https://arxiv.org/abs/2603.23049v1}
\bibitem{ref510} RAGCache Efficient Knowledge Caching for Retrieval-Augmented Generation. arXiv:2404.12457v2, 2024. \url{https://arxiv.org/abs/2404.12457v2}
\bibitem{ref511} TurboRAG Accelerating Retrieval-Augmented Generation with Precomputed KV Caches for Chunked Text. arXiv:2410.07590v1, 2024. \url{https://arxiv.org/abs/2410.07590v1}
\bibitem{ref512} KV Cache Optimization Strategies for Scalable and Efficient LLM Inference. arXiv:2603.20397v1, 2026. \url{https://arxiv.org/abs/2603.20397v1}
\bibitem{ref513} Robust KV Cache Management for LLM Serving under Output Token Length Uncertainty. arXiv:2607.16892v1, 2026. \url{https://arxiv.org/abs/2607.16892v1}
\bibitem{ref514} Comparative Characterization of KV Cache Management Strategies for LLM Inference. arXiv:2604.05012v1, 2026. \url{https://arxiv.org/abs/2604.05012v1}
\bibitem{ref515} Cache-Craft Managing Chunk-Caches for Efficient Retrieval-Augmented Generation. arXiv:2502.15734v1, 2025. \url{https://arxiv.org/abs/2502.15734v1}
\bibitem{ref516} Understanding Bottlenecks for Efficiently Serving LLM Inference With KV Offloading. arXiv:2601.19910v1, 2025. \url{https://arxiv.org/abs/2601.19910v1}
\bibitem{ref517} RAGO Systematic Performance Optimization for Retrieval-Augmented Generation Serving. arXiv:2503.14649v2, 2025. \url{https://arxiv.org/abs/2503.14649v2}
\bibitem{ref518} QCFuse Query-Aware Cache Fusion via Compressed View for Efficient RAG Serving. arXiv:2606.05875v1, 2026. \url{https://arxiv.org/abs/2606.05875v1}
\bibitem{ref519} CacheClip Accelerating RAG with Effective KV Cache Reuse. arXiv:2510.10129v2, 2025. \url{https://arxiv.org/abs/2510.10129v2}
\bibitem{ref520} Leveraging Approximate Caching for Faster Retrieval-Augmented Generation. arXiv:2503.05530v3, 2025. \url{https://arxiv.org/abs/2503.05530v3}
\bibitem{ref521} MVR-cache Optimizing Semantic Caching via Multi-Vector Retrieval and Learned Prompt Segmentation. arXiv:2605.24914v1, 2026. \url{https://arxiv.org/abs/2605.24914v1}
\bibitem{ref522} PerCache Predictive Hierarchical Cache for RAG Applications on Mobile Devices. arXiv:2601.11553v1, 2025. \url{https://arxiv.org/abs/2601.11553v1}
\bibitem{ref523} HA-RAG Hotness-Aware RAG Acceleration via Mixed Precision and Data Placement. arXiv:2510.20878v1, 2025. \url{https://arxiv.org/abs/2510.20878v1}
\bibitem{ref524} Adaptive Multi-Objective Tiered Storage Configuration for KV Cache in LLM Service. arXiv:2603.08739v1, 2026. \url{https://arxiv.org/abs/2603.08739v1}
\bibitem{ref525} Lighter And Better Towards Flexible Context Adaptation For Retrieval Augmented Generation. arXiv:2409.15699v1, 2024. \url{https://arxiv.org/abs/2409.15699v1}
\bibitem{ref526} ArcAligner Adaptive Recursive Aligner for Compressed Context Embeddings in RAG. arXiv:2601.05038v1, 2026. \url{https://arxiv.org/abs/2601.05038v1}
\bibitem{ref527} Enhancing RAG Efficiency with Adaptive Context Compression. arXiv:2507.22931v3, 2025. \url{https://arxiv.org/abs/2507.22931v3}
\bibitem{ref528} MacRAG Compress, Slice, and Scale-up for Multi-Scale Adaptive Context RAG. arXiv:2505.06569v2, 2025. \url{https://arxiv.org/abs/2505.06569v2}
\bibitem{ref529} REFRAG Rethinking RAG based Decoding. arXiv:2509.01092v2, 2025. \url{https://arxiv.org/abs/2509.01092v2}
\bibitem{ref530} KV Admission Learning What to Write for Efficient Long-Context Inference. arXiv:2512.17452v3, 2025. \url{https://arxiv.org/abs/2512.17452v3}
\bibitem{ref531} RAPID Long-Context Inference with Retrieval-Augmented Speculative Decoding. arXiv:2502.20330v2, 2025. \url{https://arxiv.org/abs/2502.20330v2}
\bibitem{ref532} Don't Do RAG When Cache-Augmented Generation is All You Need for Knowledge Tasks. arXiv:2412.15605v2, 2024. \url{https://arxiv.org/abs/2412.15605v2}
\bibitem{ref533} Enhancing Cache-Augmented Generation (CAG) with Adaptive Contextual Compression for Scalable Knowledge Integration. arXiv:2505.08261v1, 2025. \url{https://arxiv.org/abs/2505.08261v1}
\bibitem{ref534} L-RAG Balancing Context and Retrieval with Entropy-Based Lazy Loading. arXiv:2601.06551v1, 2026. \url{https://arxiv.org/abs/2601.06551v1}
\bibitem{ref535} Beyond RAG Task-Aware KV Cache Compression for Comprehensive Knowledge Reasoning. arXiv:2503.04973v1, 2025. \url{https://arxiv.org/abs/2503.04973v1}
\bibitem{ref536} MIST A Co-Design Framework for Heterogeneous, Multi-Stage LLM Inference. arXiv:2504.09775v6, 2025. \url{https://arxiv.org/abs/2504.09775v6}
\bibitem{ref537} Harmonia End-to-End RAG Serving Optimization. arXiv:2505.07833v2, 2025. \url{https://arxiv.org/abs/2505.07833v2}
\bibitem{ref538} RAGDoll Efficient Offloading-based Online RAG System on a Single GPU. arXiv:2504.15302v1, 2025. \url{https://arxiv.org/abs/2504.15302v1}
\bibitem{ref539} HedraRAG Coordinating LLM Generation and Database Retrieval in Heterogeneous RAG Serving. arXiv:2507.09138v1, 2025. \url{https://arxiv.org/abs/2507.09138v1}
\bibitem{ref540} TeleRAG Efficient Retrieval-Augmented Generation Inference with Lookahead Retrieval. arXiv:2502.20969v4, 2025. \url{https://arxiv.org/abs/2502.20969v4}
\bibitem{ref541} Shared Disk KV Cache Management for Efficient Multi-Instance Inference in RAG-Powered LLMs. arXiv:2504.11765v1, 2025. \url{https://arxiv.org/abs/2504.11765v1}
\bibitem{ref542} PentaRAG Large-Scale Intelligent Knowledge Retrieval for Enterprise LLM Applications. arXiv:2506.21593v1, 2025. \url{https://arxiv.org/abs/2506.21593v1}
\bibitem{ref543} TENT A Declarative Slice Spraying Engine for Performant and Resilient Data Movement in Disaggregated LLM Serving. arXiv:2604.00368v2, 2026. \url{https://arxiv.org/abs/2604.00368v2}
\bibitem{ref544} EdgeRAG Online-Indexed RAG for Edge Devices. arXiv:2412.21023v2, 2024. \url{https://arxiv.org/abs/2412.21023v2}
\bibitem{ref545} CoEdge-RAG Optimizing Hierarchical Scheduling for Retrieval-Augmented LLMs in Collaborative Edge Computing. arXiv:2511.05915v1, 2025. \url{https://arxiv.org/abs/2511.05915v1}
\bibitem{ref546} EACO-RAG Towards Distributed Tiered LLM Deployment using Edge-Assisted and Collaborative RAG with Adaptive Knowledge Update. arXiv:2410.20299v2, 2024. \url{https://arxiv.org/abs/2410.20299v2}
\bibitem{ref547} HeRo Adaptive Orchestration of Agentic RAG on Heterogeneous Mobile SoC. arXiv:2603.01661v2, 2026. \url{https://arxiv.org/abs/2603.01661v2}
\bibitem{ref548} Energy-Efficient On-Device RAG on a Mobile NPU System Design and Benchmark on Snapdragon X Elite. arXiv:2606.11257v1, 2026. \url{https://arxiv.org/abs/2606.11257v1}
\bibitem{ref549} RAGdb A Zero-Dependency, Embeddable Architecture for Multimodal Retrieval-Augmented Generation on the Edge. arXiv:2602.22217v1, 2025. \url{https://arxiv.org/abs/2602.22217v1}
\bibitem{ref550} Efficient Distributed Retrieval-Augmented Generation for Enhancing Language Model Performance. arXiv:2504.11197v2, 2025. \url{https://arxiv.org/abs/2504.11197v2}
\bibitem{ref551} HyFedRAG A Federated Retrieval-Augmented Generation Framework for Heterogeneous and Privacy-Sensitive Data. arXiv:2509.06444v1, 2025. \url{https://arxiv.org/abs/2509.06444v1}
\bibitem{ref552} From Cloud to Crowd Democratizing LLM Service with Decentralized Edge Collaboration for RAG. arXiv:2608.00922v1, 2026. \url{https://arxiv.org/abs/2608.00922v1}
\bibitem{ref553} CORE Toward Ubiquitous 6G Intelligence Through Collaborative Orchestration of Large Language Model Agents Over Hierarchical Edge. arXiv:2601.21822v1, 2026. \url{https://arxiv.org/abs/2601.21822v1}
\bibitem{ref554} CoMoE Collaborative Optimization of Expert Aggregation and Offloading for MoE-based LLMs at Edge. arXiv:2508.09208v1, 2025. \url{https://arxiv.org/abs/2508.09208v1}
\bibitem{ref555} Improving Medical Reasoning through Retrieval and Self-Reflection with Retrieval-Augmented Large Language Models. arXiv:2401.15269v2, 2024. \url{https://arxiv.org/abs/2401.15269v2}
\bibitem{ref556} MedRAG Enhancing Retrieval-augmented Generation with Knowledge Graph-Elicited Reasoning for Healthcare Copilot. arXiv:2502.04413v2, 2025. \url{https://arxiv.org/abs/2502.04413v2}
\bibitem{ref557} Leveraging Language Models and RAG for Efficient Knowledge Discovery in Clinical Environments. arXiv:2601.04209v1, 2025. \url{https://arxiv.org/abs/2601.04209v1}
\bibitem{ref558} Comprehensive and Practical Evaluation of Retrieval-Augmented Generation Systems for Medical Question Answering. arXiv:2411.09213v1, 2024. \url{https://arxiv.org/abs/2411.09213v1}
\bibitem{ref559} Self-MedRAG a Self-Reflective Hybrid Retrieval-Augmented Generation Framework for Reliable Medical Question Answering. arXiv:2601.04531v1, 2026. \url{https://arxiv.org/abs/2601.04531v1}
\bibitem{ref560} FHIR-RAG-MEDS Integrating HL7 FHIR with Retrieval-Augmented Large Language Models for Enhanced Medical Decision Support. arXiv:2509.07706v2, 2025. \url{https://arxiv.org/abs/2509.07706v2}
\bibitem{ref561} Legal-DC Benchmarking Retrieval-Augmented Generation for Legal Documents. arXiv:2603.11772v1, 2026. \url{https://arxiv.org/abs/2603.11772v1}
\bibitem{ref562} RAG System for Supporting Japanese Litigation Procedures Faithful Response Generation Complying with Legal Norms. arXiv:2511.22858v1, 2025. \url{https://arxiv.org/abs/2511.22858v1}
\bibitem{ref563} LegalBench-RAG A Benchmark for Retrieval-Augmented Generation in the Legal Domain. arXiv:2408.10343v1, 2024. \url{https://arxiv.org/abs/2408.10343v1}
\bibitem{ref564} A Reasoning-Focused Legal Retrieval Benchmark. arXiv:2505.03970v1, 2025. \url{https://arxiv.org/abs/2505.03970v1}
\bibitem{ref565} Identifying Financial Risk Information Using RAG with a Contrastive Insight. arXiv:2510.03521v1, 2025. \url{https://arxiv.org/abs/2510.03521v1}
\bibitem{ref566} Retrieval Augmented Generation (RAG) for Fintech Agentic Design and Evaluation. arXiv:2510.25518v1, 2025. \url{https://arxiv.org/abs/2510.25518v1}
\bibitem{ref567} An Industrial-Scale Retrieval-Augmented Generation Framework for Requirements Engineering Empirical Evaluation with Automotive Manufacturing Data. arXiv:2603.20534v2, 2026. \url{https://arxiv.org/abs/2603.20534v2}
\bibitem{ref568} An Agile Method for Implementing Retrieval Augmented Generation Tools in Industrial SMEs. arXiv:2508.21024v1, 2025. \url{https://arxiv.org/abs/2508.21024v1}
\bibitem{ref569} AutoBnB-RAG Enhancing Multi-Agent Incident Response with Retrieval-Augmented Generation. arXiv:2508.13118v2, 2025. \url{https://arxiv.org/abs/2508.13118v2}
\bibitem{ref570} BSharedRAG Backbone Shared Retrieval-Augmented Generation for the E-commerce Domain. arXiv:2409.20075v1, 2024. \url{https://arxiv.org/abs/2409.20075v1}
\bibitem{ref571} Retrieval-Augmented Generation in Industry An Interview Study on Use Cases, Requirements, Challenges, and Evaluation. arXiv:2508.14066v1, 2025. \url{https://arxiv.org/abs/2508.14066v1}
\bibitem{ref572} Seven Failure Points When Engineering a Retrieval Augmented Generation System. arXiv:2401.05856v1, 2024. \url{https://arxiv.org/abs/2401.05856v1}
\bibitem{ref573} AI Engineering Blueprint for On-Premises Retrieval-Augmented Generation Systems. arXiv:2604.01395v1, 2026. \url{https://arxiv.org/abs/2604.01395v1}
\bibitem{ref574} RAGOps Operating and Managing Retrieval-Augmented Generation Pipelines. arXiv:2506.03401v1, 2025. \url{https://arxiv.org/abs/2506.03401v1}
\bibitem{ref575} FlexRAG A Flexible and Comprehensive Framework for Retrieval-Augmented Generation. arXiv:2506.12494v2, 2025. \url{https://arxiv.org/abs/2506.12494v2}
\bibitem{ref576} UltraRAG A Modular and Automated Toolkit for Adaptive Retrieval-Augmented Generation. arXiv:2504.08761v1, 2025. \url{https://arxiv.org/abs/2504.08761v1}
\bibitem{ref577} Is Agentic RAG worth it An experimental comparison of RAG approaches. arXiv:2601.07711v2, 2026. \url{https://arxiv.org/abs/2601.07711v2}
\bibitem{ref578} AI Assurance A Comprehensive Testing Strategy for Enterprise AI Systems. arXiv:2605.23459v1, 2026. \url{https://arxiv.org/abs/2605.23459v1}
\bibitem{ref579} LLM Readiness Harness Evaluation, Observability, and CI Gates for LLM RAG Applications. arXiv:2603.27355v2, 2026. \url{https://arxiv.org/abs/2603.27355v2}
\bibitem{ref580} A Methodology for Evaluating RAG Systems A Case Study On Configuration Dependency Validation. arXiv:2410.08801v1, 2024. \url{https://arxiv.org/abs/2410.08801v1}
\bibitem{ref581} Legal RAG Bench an end-to-end benchmark for legal RAG. arXiv:2603.01710v1, 2026. \url{https://arxiv.org/abs/2603.01710v1}
\bibitem{ref582} OmniEval An Omnidirectional and Automatic RAG Evaluation Benchmark in Financial Domain. arXiv:2412.13018v2, 2024. \url{https://arxiv.org/abs/2412.13018v2}
\bibitem{ref583} RAGBench Explainable Benchmark for Retrieval-Augmented Generation Systems. arXiv:2407.11005v2, 2024. \url{https://arxiv.org/abs/2407.11005v2}
\bibitem{ref584} A System for Comprehensive Assessment of RAG Frameworks. arXiv:2504.07803v1, 2025. \url{https://arxiv.org/abs/2504.07803v1}
\bibitem{ref585} RAGe A Retrieval-Augmented Generation Evaluation Framework. arXiv:2605.27445v1, 2026. \url{https://arxiv.org/abs/2605.27445v1}
\bibitem{ref586} RAGPerf An End-to-End Benchmarking Framework for Retrieval-Augmented Generation Systems. arXiv:2603.10765v1, 2026. \url{https://arxiv.org/abs/2603.10765v1}
\bibitem{ref587} MTRAG A Multi-Turn Conversational Benchmark for Evaluating Retrieval-Augmented Generation Systems. arXiv:2501.03468v1, 2025. \url{https://arxiv.org/abs/2501.03468v1}
\bibitem{ref588} RAGalyst Automated Human-Aligned Agentic Evaluation for Domain-Specific RAG. arXiv:2511.04502v1, 2025. \url{https://arxiv.org/abs/2511.04502v1}
\bibitem{ref589} The Viability of Crowdsourcing for RAG Evaluation. arXiv:2504.15689v1, 2025. \url{https://arxiv.org/abs/2504.15689v1}
\bibitem{ref590} InspectorRAGet An Introspection Platform for RAG Evaluation. arXiv:2404.17347v1, 2024. \url{https://arxiv.org/abs/2404.17347v1}
\bibitem{ref591} Towards a rigorous evaluation of RAG systems the challenge of due diligence. arXiv:2507.21753v1, 2025. \url{https://arxiv.org/abs/2507.21753v1}
\bibitem{ref592} Where is My Troubleshooting Procedure Studying the Potential of RAG in Assisting Failure Resolution of Large Cyber-Physical System. arXiv:2601.08706v2, 2026. \url{https://arxiv.org/abs/2601.08706v2}
\bibitem{ref593} Enhancing Legal LLMs through Metadata-Enriched RAG Pipelines and Direct Preference Optimization. arXiv:2603.19251v1, 2026. \url{https://arxiv.org/abs/2603.19251v1}
\bibitem{ref594} Leveraging Resolved Incident History for LLM-Assisted Software Bug Diagnosis. arXiv:2607.21911v1, 2026. \url{https://arxiv.org/abs/2607.21911v1}
\bibitem{ref595} Enhancing Critical Thinking with AI A Tailored Warning System for RAG Models. arXiv:2504.16883v1, 2025. \url{https://arxiv.org/abs/2504.16883v1}
\bibitem{ref596} Knowledge graph enhanced retrieval-augmented generation for failure mode and effects analysis. arXiv:2406.18114v3, 2024. \url{https://arxiv.org/abs/2406.18114v3}
\bibitem{ref597} A Systems-Level Analysis of Sensitivity, Robustness, and Stability in Retrieval-Augmented Generation. arXiv:2606.28337v1, 2026. \url{https://arxiv.org/abs/2606.28337v1}
\bibitem{ref598} Lost in the Evidence Reproducing Document Position and Context Size Effects in RAG. arXiv:2605.27105v2, 2026. \url{https://arxiv.org/abs/2605.27105v2}
\bibitem{ref599} Can Small Language Models Use What They Retrieve An Empirical Study of Retrieval Utilization Across Model Scale. arXiv:2603.11513v1, 2026. \url{https://arxiv.org/abs/2603.11513v1}
\bibitem{ref600} After Retrieval, Before Generation Enhancing the Trustworthiness of Large Language Models in Retrieval-Augmented Generation. arXiv:2505.17118v2, 2025. \url{https://arxiv.org/abs/2505.17118v2}
\bibitem{ref601} Metadata, Structure, or Strategy A Decomposition of RAG Context Enrichment. arXiv:2606.29645v1, 2026. \url{https://arxiv.org/abs/2606.29645v1}
\bibitem{ref602} Enhancing Retrieval-Augmented Generation with Two-Stage Retrieval FlashRank Reranking and Query Expansion. arXiv:2601.03258v1, 2025. \url{https://arxiv.org/abs/2601.03258v1}
\bibitem{ref603} What Survives Into Context A Diagnostic for Budget-Constrained Multi-Hop RAG and When Submodular Evidence Packing Improves It. arXiv:2607.00725v1, 2026. \url{https://arxiv.org/abs/2607.00725v1}
\bibitem{ref604} Know Before You Fetch Calibrated Retrieval-Budget Allocation for Retrieval-Augmented Generation. arXiv:2606.29959v1, 2026. \url{https://arxiv.org/abs/2606.29959v1}
\bibitem{ref605} Cost-Aware Query Routing in RAG Empirical Analysis of Retrieval Depth Tradeoffs. arXiv:2606.02581v1, 2026. \url{https://arxiv.org/abs/2606.02581v1}
\bibitem{ref606} Investigating the Robustness of Retrieval-Augmented Generation at the Query Level. arXiv:2507.06956v1, 2025. \url{https://arxiv.org/abs/2507.06956v1}
\bibitem{ref607} Scaling Retrieval Augmented Generation with RAG Fusion Lessons from an Industry Deployment. arXiv:2603.02153v1, 2026. \url{https://arxiv.org/abs/2603.02153v1}
\bibitem{ref608} UiS-IAI@LiveRAG Retrieval-Augmented Information Nugget-Based Generation of Responses. arXiv:2506.22210v1, 2025. \url{https://arxiv.org/abs/2506.22210v1}
\bibitem{ref609} Reinforced Internal-External Knowledge Synergistic Reasoning for Efficient Adaptive Search Agent. arXiv:2505.07596v1, 2025. \url{https://arxiv.org/abs/2505.07596v1}
\bibitem{ref610} CDF-RAG Causal Dynamic Feedback for Adaptive Retrieval-Augmented Generation. arXiv:2504.12560v1, 2025. \url{https://arxiv.org/abs/2504.12560v1}
\bibitem{ref611} Atom-Searcher Enhancing Agentic Deep Research via Fine-Grained Atomic Thought Reward. arXiv:2508.12800v3, 2025. \url{https://arxiv.org/abs/2508.12800v3}
\bibitem{ref612} An Agent-Oriented Pluggable Experience-RAG Skill for Experience-Driven Retrieval Strategy Orchestration. arXiv:2605.03989v2, 2026. \url{https://arxiv.org/abs/2605.03989v2}
\bibitem{ref613} Data Poisoning Vulnerabilities Across Healthcare AI Architectures A Security Threat Analysis. arXiv:2511.11020v1, 2025. \url{https://arxiv.org/abs/2511.11020v1}
\bibitem{ref614} Standardized Threat Taxonomy for AI Security, Governance, and Regulatory Compliance. arXiv:2511.21901v1, 2025. \url{https://arxiv.org/abs/2511.21901v1}
\bibitem{ref615} Never compromise with vulnerabilities a comprehensive survey on AI governance. arXiv:2508.08789v5, 2025. \url{https://arxiv.org/abs/2508.08789v5}
\bibitem{ref616} Privacy Cost as Equity Input A Group Fairness Criterion for Differentially Private Machine Learning. arXiv:2607.16620v1, 2026. \url{https://arxiv.org/abs/2607.16620v1}
\bibitem{ref617} Auditing Fairness-Privacy Trade-offs Subpopulation-Level Effects of Fairness-Enhancing Algorithms. arXiv:2607.14607v1, 2026. \url{https://arxiv.org/abs/2607.14607v1}
\bibitem{ref618} Trustworthy AI Securing Sensitive Data in Large Language Models. arXiv:2409.18222v1, 2024. \url{https://arxiv.org/abs/2409.18222v1}
\bibitem{ref619} Generating Synthetic Data with Formal Privacy Guarantees State of the Art and the Road Ahead. arXiv:2503.20846v1, 2025. \url{https://arxiv.org/abs/2503.20846v1}
\bibitem{ref620} FAIRPLAI A Human-in-the-Loop Approach to Fair and Private Machine Learning. arXiv:2511.08702v1, 2025. \url{https://arxiv.org/abs/2511.08702v1}
\bibitem{ref621} From Vulnerable Data Subjects to Vulnerabilizing Data Practices Navigating the Protection Paradox in AI-Based Analyses of Platformized Lives. arXiv:2604.15990v1, 2026. \url{https://arxiv.org/abs/2604.15990v1}
\bibitem{ref622} A Framework for Responsible AI Systems Building Societal Trust through Domain Definition, Trustworthy AI Design, Auditability, Accountability, and Governance. arXiv:2503.04739v2, 2025. \url{https://arxiv.org/abs/2503.04739v2}
\bibitem{ref623} The Open-Box Fallacy Why AI Deployment Needs a Calibrated Verification Regime. arXiv:2605.10601v1, 2026. \url{https://arxiv.org/abs/2605.10601v1}
\bibitem{ref624} Private, Verifiable, and Auditable AI Systems. arXiv:2509.00085v1, 2025. \url{https://arxiv.org/abs/2509.00085v1}
\bibitem{ref625} Trustworthy AI From Principles to Practices. arXiv:2110.01167v2, 2021. \url{https://arxiv.org/abs/2110.01167v2}
\bibitem{ref626} Beyond Factual Accuracy Evaluating Global Reasoning Integrity in RAG Systems with LogicScore. arXiv:2601.15050v5, 2026. \url{https://arxiv.org/abs/2601.15050v5}
\bibitem{ref627} VERA Validation and Evaluation of Retrieval-Augmented Systems. arXiv:2409.03759v1, 2024. \url{https://arxiv.org/abs/2409.03759v1}
\bibitem{ref628} Beyond Chat a Framework for LLMs as Human-Centered Support Systems. arXiv:2511.03729v1, 2025. \url{https://arxiv.org/abs/2511.03729v1}
\bibitem{ref629} Fast or Better Balancing Accuracy and Cost in Retrieval-Augmented Generation with Flexible User Control. arXiv:2502.12145v2, 2025. \url{https://arxiv.org/abs/2502.12145v2}
\bibitem{ref630} Skill-RAG Failure-State-Aware Retrieval Augmentation via Hidden-State Probing and Skill Routing. arXiv:2604.15771v3, 2026. \url{https://arxiv.org/abs/2604.15771v3}
\bibitem{ref631} Retrieval is Not Enough Enhancing RAG Reasoning through Test-Time Critique and Optimization. arXiv:2504.14858v4, 2025. \url{https://arxiv.org/abs/2504.14858v4}
\bibitem{ref632} LUMA-RAG Lifelong Multimodal Agents with Provably Stable Streaming Alignment. arXiv:2511.02371v1, 2025. \url{https://arxiv.org/abs/2511.02371v1}
\end{thebibliography}

\end{document}
